%%%% ICLR 2027 submission -- Universality and Convergence of Generative Flows
%%%% Ported from gfnbounds_aaai.tex (AAAI 2026 template) on 2026-08-21.
%%%% The AAAI draft was dropped on 2026-09-13; gfnbounds_aaai.tex survives in git at b735e87.
\documentclass{article}
\usepackage{iclr2027_conference,times}

\usepackage{amsmath,amssymb,amsthm}
\usepackage{graphicx}
\usepackage{algorithm}
\usepackage{algorithmic}
\usepackage[all]{xy}          % xymatrix is used in cv_divergence, introduction, universality
\usepackage{hyperref}
\usepackage{url}
\usepackage{listings}
\lstdefinestyle{lean}{basicstyle=\ttfamily\footnotesize,columns=fullflexible,keepspaces=true,
  showstringspaces=false,breaklines=true,breakatwhitespace=true,breakindent=2em,numbers=none,
  breakautoindent=true,postbreak=\mbox{\ensuremath{\hookrightarrow}\space},
  literate={¬}{{\ensuremath{\neg}}}1 {ν}{{\ensuremath{\nu}}}1 {ϑ}{{\ensuremath{\vartheta}}}1 {ᗮ}{{\ensuremath{^{\perp}}}}1 {ᵐ}{{\ensuremath{^{m}}}}1 {‖}{{\ensuremath{\|}}}1 {₀}{{\ensuremath{_0}}}1 {₁}{{\ensuremath{_1}}}1 {₂}{{\ensuremath{_2}}}1 {₆}{{\ensuremath{_6}}}1 {₈}{{\ensuremath{_8}}}1 {ₗ}{{\ensuremath{_l}}}1 {ℕ}{{\ensuremath{\mathbb{N}}}}1 {ℝ}{{\ensuremath{\mathbb{R}}}}1 {→}{{\ensuremath{\rightarrow}}}1 {∀}{{\ensuremath{\forall}}}1 {∂}{{\ensuremath{\partial}}}1 {∃}{{\ensuremath{\exists}}}1 {∈}{{\ensuremath{\in}}}1 {∑}{{\ensuremath{\sum}}}1 {∙}{{\ensuremath{\bullet}}}1 {∧}{{\ensuremath{\wedge}}}1 {≠}{{\ensuremath{\neq}}}1 {≤}{{\ensuremath{\leq}}}1 {≥}{{\ensuremath{\geq}}}1 {⟨}{{\ensuremath{\langle}}}1 {⟩}{{\ensuremath{\rangle}}}1 {⦃}{{\ensuremath{\{\!\{}}}1 {⦄}{{\ensuremath{\}\!\}}}}1 {⨆}{{\ensuremath{\bigsqcup}}}1 {𝓝}{{\ensuremath{\mathcal{N}}}}1 {λ}{{\ensuremath{\lambda}}}1 {ε}{{\ensuremath{\varepsilon}}}1 {Π}{{\ensuremath{\Pi}}}1 {σ}{{\ensuremath{\sigma}}}1 {⋆}{{\ensuremath{\star}}}1 {−}{{\ensuremath{-}}}1 {∞}{{\ensuremath{\infty}}}1 {∫}{{\ensuremath{\int}}}1 {·}{{\ensuremath{\cdot}}}1 {…}{{\ensuremath{\dots}}}1 {≥}{{\ensuremath{\geq}}}1}

%%%% ---------------------------------------------------------------------
%%%% AUTHOR BLOCK -- DECISION POINT, see README-editorial-pass.txt
%%%% The AAAI draft carried five authors (Brunswic, Clemente, Yang, Rasouli,
%%%% Zhou -- Huawei).  The ICLR 2027 plan records "sole author, independent
%%%% affiliation".  The sole-author block is used below per that decision;
%%%% the original five-author block is in gfnbounds_aaai.tex at git b735e87.
%%%% NOTE: the author list LOCKS at the ICLR abstract deadline (18 Sep 2026) --
%%%% "new authors cannot be added" afterwards.  Settle this before then.
%%%% ---------------------------------------------------------------------
\title{Universality and Convergence of Generative Flows}

\author{Leo Maxime Brunswic \\
Independent Researcher \\
\texttt{leo@brunswic.fr}}

\newcommand{\fix}{\marginpar{FIX}}
\newcommand{\new}{\marginpar{NEW}}

%\iclrfinalcopy % Uncomment for camera-ready, NOT for submission (double blind).

\newtheorem{theorem}{Theorem}
\newtheorem{proposition}[theorem]{Proposition}
\newtheorem{lemma}[theorem]{Lemma}
\newtheorem{corollary}[theorem]{Corollary}
\theoremstyle{definition}
\newtheorem{definition}[theorem]{Definition}
\newtheorem{assumption}[theorem]{Assumption}
\theoremstyle{remark}
\newtheorem{remark}[theorem]{Remark}

%%%% ---------------------------------------------------------------------
%%%% notation.tex --- the paper's notation, defined once.
%%%%
%%%% This file is the single source of the macros below.  It is \input by
%%%% gfnbounds_iclr2027.tex (canonical) and notes_finit_replay.tex (the lab
%%%% record).  Until 2026-09-03 each mainfile, and the since-dropped AAAI draft,
%%%% carried its own copy of the block; the 31 Aug notation pass therefore had to run
%%%% three times, and the copy that went wrong produced \def\lamreg{\lamreg} and a
%%%% build that never terminated.  A macro defined here is never hand-expanded at a
%%%% call site: the notation passes are mechanical substitutions, and a spelled-out
%%%% expansion is what the next one misses.
%%%%
%%%% What deliberately stays in the mainfiles: \appendixAtitle..\appendixHtitle; the
%%%% \newtheorem blocks; and ICLR's editorial \fix / \new marks.
%%%% ---------------------------------------------------------------------

%%%% The primitive \lambda, captured before anything can redefine it.
%%%% notes_finit_replay.tex reads \lambda as its empirical background measure and
%%%% redefines it accordingly.  Because \def bodies expand lazily, \lamreg --- which
%%%% is not a measure --- inherited that redefinition, and the notes rendered
%%%% Morozov et al.'s regularization weight as $\widehat\nu_{\mathrm B}$-reg where
%%%% the paper rendered $\lambda_{\mathrm{reg}}$.  Binding \lamreg to the captured
%%%% primitive settles that without disturbing the notes' own convention.
\let\GFNlambda\lambda

%%%% Measure convention (2026-08-31, extended 2026-09-04).  \lambda is the INVARIANT
%%%% measure of the chain in play; \nuB is the BACKGROUND measure, the one densities and
%%%% kernels are parameterized against; \nutrain is the TRAINING measure; \nuG is the
%%%% GEOMETRY measure of the first-variation theorem, the reference the gradient is read
%%%% against, which fixes the preconditioner and not the objective.  The letter \eta is
%%%% not a measure: it is the scalar lifting level of \foutstareta.  Two further symbols
%%%% resemble a measure and are not: \lamreg is Morozov et al.'s flow-regularization
%%%% weight, and \Lambda in Appendix G is a local smoothness constant.

%%%% --- state space -----------------------------------------------------
\def\Statespace{\mathcal S}

%%%% --- policies --------------------------------------------------------
\def\ForwardPolicy{\pi_{\rightarrow}}
\def\ForwardPolicyStar{\pi_{\rightarrow}^*}
\def\BackwardPolicy{\pi_{\leftarrow}}

%%%% --- flows, and the densities against \nuB ---------------------------
\def\Foutstar{F^*_{\rightarrow}}
\def\Foutmin{F^{*,\min}_{\rightarrow}}
\def\Fin{F_{\leftarrow}}
\def\Finstar{F^*_{\leftarrow}}
\def\Finit{F_{\mathrm{init}}}
\def\Fterm{F_{\mathrm{term}}}
\def\foutstar{f^*_{\rightarrow}}
\def\foutstareta{f^{*,\eta}_{\rightarrow}}
\def\fout{f_{\rightarrow}}
\def\finstar{f^*_{\leftarrow}}
\def\finit{f_{\mathrm{init}}}
\def\fterm{f_{\mathrm{term}}}

%%%% --- the twisted flow that forces flow matching, and its defects -----
\def\Finithat{\widehat {F}_{\mathrm{init}}}
\def\Ftermhat{\widehat{F}_{\mathrm{term}}}
\def\ftermhat{\widehat{f}_{\mathrm{term}}}
\def\deltafinit{\delta f_{\mathrm{init}}}
\def\deltafterm{\delta f_{\mathrm{term}}}

%%%% --- measures --------------------------------------------------------
\def\nuB{\nu_{\mathrm{B}}}
\def\nutrain{\nu_{\mathrm{T}}}
\def\nuG{\nu_{\mathrm{G}}}
\def\target{\kappa}
\def\dln{\frac{d\nuB}{d\nutrain}}
\def\Dq{D_{q}}

%%%% --- losses and divergences ------------------------------------------
\def\LFMstable{\mathcal L_{\mathrm{FM},2}^{\mathrm{stable}}}
\def\LFMstableq{\mathcal L_{\mathrm{FM},q}^{\mathrm{stable}}}
\def\LossWKLFM{\mathcal L_{\mathrm{KL-wFM}}}
\def\TV{\mathrm{TV}}
\def\KL{\mathrm{KL}}

%%%% --- replay mixture weights ------------------------------------------
\def\forwardalpha{\alpha_{\rightarrow} }

%%%% --- constants that are not measures ---------------------------------
\def\lamreg{\GFNlambda_{\mathrm{reg}}}


\newcommand{\appendixAtitle}{Proofs}
\newcommand{\appendixBtitle}{Comparison with the Convergence Bounds of Silva et al.}
\newcommand{\appendixCtitle}{The Initial Flow: Mass, Shape, and the Constant of the Stable Bound}
\newcommand{\appendixDtitle}{The Replay Buffer and the Training Measure}
\newcommand{\appendixEtitle}{Trained Backward Policies: Flow Regularization in Spectral Variables}
\newcommand{\appendixFtitle}{Stochastic Gradients: Noise, Floors, and the Buffer}
\newcommand{\appendixGtitle}{Experimental Program}
\newcommand{\appendixHtitle}{The Doubling Graph: an Unbounded Diffusion Operator at Finite Backward Length}
\newcommand{\appendixItitle}{The Doubling Graph, Certified: a Machine-Checked Theorem and the Ideas Behind It}

\begin{document}

\maketitle

\begin{abstract}
Generative Flows sample from a target distribution by training a flow to satisfy a
flow-matching constraint, yet little is known about how the resulting sampling error relates
to the value of the training loss, or how fast gradient descent reduces it. We introduce a
notion of $L^p$-universality for parameterized generative flows, establish it for ergodic
parameterizations with summable mixing, and prove convergence bounds controlling the sampling
error in total variation by the stable-FM and KL-weak-FM losses. For divergence-based losses
(FM, DB, TB) we prove a minimax theorem: no such bound exists on general state graphs, and the
bounds available in the finite acyclic case necessarily carry constants that degenerate with
the state space. With a frozen backward policy we then compute the gradients of the FM, DB and
TB losses in a unified operator form and show that gradient descent \emph{diffuses} the
current flow through the backward policy, so that the mixing of that policy governs both
universality and the speed of training. Convergence is all one gets: freezing does not restore
the $0$-flow monotonicity whose failure makes these losses unstable, and generators of
vanishing mobility can freeze the diffusion outright --- an obstruction that strictly unimodal
generators such as $(\log x)^2$ escape. Finally, we test the bounds numerically on state spaces
small enough that the sampling error is computed rather than estimated: the stable bound holds
across $21{,}168$ flows and its measured tightness is essentially flat in the state space at
$q=1$, the constant hidden in the KL-weak-FM bound is below $1$ on every instance where the
bound is non-vacuous, and the training rate predicted from the exact coercivity constant
matches the observed one to within $10\%$.
\end{abstract}


\section{Introduction}

Generative Flows (GFlowNets), introduced in \citet{bengio2021flow, bengio2023gflownet}, were initially defined on Directed Acyclic Graphs to sample vertices proportionally to a reward.
The approach was then motivated by comparison to Monte-Carlo Markov Chains (MCMC) and Reinforcement Learning (RL) methods.
The former shares the same objective but suffers from long mixing hence sampling time, while the latter is prone to mode collapse and lacks diversity.
Theoretical improvement of the framework has been very active with introductions of new training losses \cite{malkin2022trajectory,pan2023better,brunswic2024theory, silva2024divergence,hu2024beyond},
extensions beyond graph environments \cite{li2023cflownets, lahlou2023theory} and removing acyclicity constraint \cite{brunswic2024theory,morozov2025revisiting}.

So far, however, universality and convergence of generative flows are an incompletely addressed corner of the Theory of Generative Flows. More precisely, we lack control of the form
\begin{equation}
    \mathrm{dist}(\mu_{\mathrm{target}}, \mu_{\mathrm{sampler}}) \leq F(\mathcal L), \quad F(\min \mathcal L)=0
\end{equation}
where $\mathrm{dist}$ is distance-like in the space of probability distributions and $F$ is an increasing continuous function.
Previous quantitative attempts are restricted to finite acyclic state spaces: \citet{silvageneralization} bound the total variation by an idealized trajectory-balance residual, and \citet{silva2025gflownets} bound the total variation under flow perturbations of an exactly balanced model (see Appendix B for a detailed comparison).
In fact, the instability result of \citet{brunswic2024theory}  implies that \textbf{no convergence bound exists in general for FM, DB and TB}, and more generally for divergence-like losses with a generator continuous at $1$.
Indeed, the instability states that any loop in the state graph implies these losses may be minimized without satisfying the flow-matching objective; we sharpen this observation into a minimax theorem (Theorem \ref{theo:no_bound_divergence}).

The present work aims at filling this gap in the literature by combining three recent advances. First,  \citet{brunswicEGF} proved a general quantitative
control of the total variation between the target and sampling distributions of the generative flows (i.e.\ the $L^1$ sampling error) by the flow-matching $L^1$-error.
The so-called stable FM \cite{brunswic2024theory} and KL-weak-FM \cite{brunswicEGF} losses being essentially estimators of some $L^p$ flow-matching error, they are natural candidates to provide straightforward convergence bounds.
Second, recent results on backward-policy control \cite{gritsaev2024optimizing,morozov2025revisiting} showed that some divergence-like losses are stable in a weak sense if applied only to the stateflow and forward policy
while the backward policy is frozen or trained using a different loss. This suggests that convergence bounds of the type $\mathrm{dist}(\mu_{\mathrm{target}}, \mu_{\mathrm{sampler}}) \leq F(\mathcal L, \BackwardPolicy)$
may exist for such losses. Third, the dependency on the backward policy in such convergence bounds is likely to depend on its exploration properties that ensure the flow does not get stuck in loops.
The Markov Chain perspective \cite{deleu2023generative} together with the universality theorem of \citet{brunswicEGF} allow us to concretize this idea.


\subsection{Positioning}\label{sub:positioning}

A convergence bound is an operational object: it converts a number the practitioner already
logs --- the training loss --- into a certificate on the object they actually care about, the
sampling error, and it names the knobs that move that certificate. That is the standard we
place ourselves against. Compared with other samplers the comparison is one of guarantees
rather than of benchmarks: MCMC pays for exactness with a mixing time that is rarely known,
while annealed importance sampling and sequential Monte-Carlo buy unbiasedness and variance
control but say nothing about the error of a \emph{learned} proposal. Within generative flows,
quantitative guarantees have so far been confined to finite acyclic state spaces
\cite{silvageneralization,silva2025gflownets}; what we ask is what survives when both
restrictions are lifted --- an $L^p$ bound in the measurable setting for the stable losses,
and a proof that \emph{nothing} survives for the divergence losses.

\textbf{Main Contributions}
\begin{enumerate}
    \item an ergodicity-based theory of $L^p$-universality for generative flows, valid on
    Polish state spaces and specializing to graphs;
    \item convergence bounds controlling the sampling error in total variation by the
    stable-FM and KL-weak-FM losses, with explicit constants;
    \item a minimax theorem showing that no such bound exists for divergence-based losses (FM,
    DB, TB) on general state graphs, with an explicit-constant analysis of what is available in
    the finite acyclic case (Appendix \ref{sec:silva_comparison});
    \item a unified first variation of the FM, DB and TB losses under a frozen backward
    policy, exhibiting gradient descent as a diffusion of the flow by that policy;
    \item a training-speed theorem in which the same ergodicity governs both universality and
    convergence rate, with two negative results delimiting it: freezing does \emph{not} restore
    stability in the sense of \citet{brunswic2024theory} (Proposition
    \ref{prop:frozen_unstable}), and global convergence can fail for admissible generators, an
    obstruction that strictly unimodal generators such as $(\log x)^2$ provably escape;
    \item an empirical study on enumerable state spaces where the sampling error is computed
    rather than estimated --- the stable bound on $21{,}168$ flows, the KL-weak-FM bound on
    $3{,}472$, rate constants on $20$ chains --- including a family on which the mixing
    hypothesis provably holds \emph{uniformly} in the state space, via property (T) for
    $\mathrm{Aut}(F_k)$ (Section \ref{sec:experiments});
    \item four appendices carrying the theory into the two objects a practitioner chooses ---
    the initial flow (Appendix \ref{app:finit}) and the training measure (Appendix
    \ref{app:replay}) --- and into two regimes this paper's dynamics do not cover: a trained
    backward policy (Appendix \ref{app:pbtrain}) and sampled gradients (Appendix
    \ref{app:sgd}).
\end{enumerate}
\subsection{Training Objectives of Generative Flows}
    
    
    
    We work in the measurable setting: the state space $\Statespace$ is a Polish space --- a
    vector space or a Riemannian manifold \cite{lee2018introduction, gemici2016normalizing}
    --- carrying a finite background measure $\nuB$. Measure-theoretic notation
    ($\ll$, Radon--Nikodym derivatives, Markov kernels, image measures and $L^p$ spaces) is
    standard and collected in Appendix \ref{app:notation}.

    We adopt the formulation of generative flows used in \cite{brunswicEGF}.  Let $(\Statespace,\nuB)$ be a Polish space with a finite background measure. A generative flow on $\Statespace$ is the data of a star forward policy, i.e. a Markov kernel $\ForwardPolicyStar:\Statespace\rightarrow \Statespace$, and a star outflow, i.e. a finite non-negative measure $\Foutstar$ on $\Statespace$.

    With unnormalized initial and terminal distributions $\Finit$ and $\Fterm$,  it induces a Markov chain $(\underline s_t)_{t\geq 1}$ defined by $\underline s_{1}\sim \Finit$ and $\underline s_{t+1} = \ForwardPolicyStar(\underline s_t)$.
    Define the sampling time $\tau \in \mathbb N_{\geq 1}$ as the random variable given by\footnote{Clearly $\Fterm \ll \Fterm+\Foutstar$ so the Radon-Nikodym derivative is well-defined.} $\mathbb P(\tau\geq 1)=1$ and $\mathbb P(\tau = t|\tau\geq t) = \frac{d\Fterm}{d(\Fterm+\Foutstar)}(\underline s_t)$.
    The sampler associated with a generative flow is obtained by emulating the Markov chain $(\underline s_t)_{t\geq1}$ and sampling $\underline s_\tau$.
    By extending the sampling Markov chain with a source $s_0$ and sink $s_f$ state, we can summarize the flow as follows:
    \begin{equation}
        \xymatrix{
        s_0 \ar[r]^\Finit&\Statespace \ar@(ur,ul)_{\Foutstar\otimes\ForwardPolicyStar}\ar[r]^\Fterm& s_f
        }.
    \end{equation}
    The interest of the formalization is to separate the ``internal'' flow $\Statespace\rightarrow \Statespace$ from the initial and terminal flows which may be twisted for the purpose of a theoretical reasoning and because the terminal flow differs from the target which we assumed to be inaccessible during inference.





    % The fundamental result of generative flows is the sampling Theorem:
  Two results from prior work anchor everything below. First, if the flow-matching constraint
\begin{equation}
\label{equ:FM_const}
    \Finit + \Foutstar\ForwardPolicyStar = \Fterm + \Foutstar
\end{equation}
holds, then $\Finit(\Statespace)=\Fterm(\Statespace)$, the expected sampling time is at most
$\Foutstar(\Statespace)/\Finit(\Statespace)+1$, and the sampler is exact:
$\underline s_\tau\sim\Fterm$ \cite{bengio2021flow, brunswic2024theory}. Second, a flow that
does \emph{not} satisfy \eqref{equ:FM_const} can be made to, by twisting the initial and
terminal distributions: writing $\Finstar := \Foutstar\ForwardPolicyStar$ and
$\delta\Fterm := (\Finit+\Finstar-\Foutstar-\Fterm)^+$,
$\delta\Finit := (\Finit+\Finstar-\Foutstar-\Fterm)^-$, the flow is flow-matching from
$\Finithat := \Finit+\delta\Finit$ to $\Ftermhat := \Fterm+\delta\Fterm$, and its sampling
error is controlled by the initial defect alone,
$\TV\bigl(s_\tau\,\|\,\Ftermhat/\Ftermhat(\Statespace)\bigr)\leq\delta\Finit(\Statespace)/\Ftermhat(\Statespace)$
\cite{brunswicEGF}. Both are restated formally as Theorems \ref{theo:sampling_theorem} and
\ref{theo:negative_control} in Appendix \ref{app:notation}.

The pair $(\Foutstar, \ForwardPolicyStar)$ is trained using gradient descent to minimize a loss function that enforces the flow-matching constraint in equation \eqref{equ:FM_const}. Once convergence is achieved, the generative flow sampler yields samples of the random variable $s_\tau$. Hence, by Theorem \ref{theo:sampling_theorem},  the generative flow sampler samples from $\Fterm$. For example, we can use the most straightforward stable $L^q$ FM-loss:
    \begin{equation}\label{equ:FMloss}
        \LFMstableq := \left(\mathbb E_{s\sim \nutrain}\left[ \left|\finit + \finstar - \fterm - \foutstar\right|^q(s) \right]\right)^{1/q}
    \end{equation}
    where $\finit = \frac{d\Finit}{d\nuB}$, $\fterm=\frac{d\Fterm}{d\nuB}$ and $\finstar=\frac{d\Finstar}{d\nuB} = \frac{d(\Foutstar\ForwardPolicyStar)}{d\nuB} $.
This loss is suitable when the target distribution $\target$ has accessible density.
If $\target$ is only accessible via samples, \cite{brunswicEGF} introduced a so-called KL-weak-FM loss which may be put in a generalized $L^q$ form:
\begin{equation}\label{equ:weakFM}
     \LossWKLFM^q (\theta) := b \|\deltafinit(s)\|_{L^q(\nutrain)}-  \mathbb E_{s\sim \target}\log \ftermhat(s)
\end{equation}


\section{$L^p$-universality of Generative Flows}\label{sec:universality}
A key implicit assumption while training Generative Flows is their universality. The notion is rarely stated formally, so we define a universality property for generative flows that suits our purpose. Universality of Generative Flows then depends on architecture choices of the flow as well as corresponding universality statement on the parameterizing neural networks.

A family of generative flows is a set $\Theta$ of pairs $(\ForwardPolicyStar,\foutstar)$,
metrized by the $L^p(\nuB)$ norm on the outflow and the $L^p(\nuB)$ operator norm on the
policy (Appendix \ref{app:notation}). The second is a genuine requirement on the
parameterization: boundedness of $\foutstar$ is easy, finiteness of
$\|\ForwardPolicyStar\|_{L^p(\nuB)}$ is not automatic.

\begin{definition}\label{def:universality} Let $\Theta$ be a family of generative flows on a statespace $(\Statespace,\nuB)$ with $\nuB$ finite and let $p\in [1,\infty]$. The family $\Theta$ is weakly $L^p$-universal if its $\ForwardPolicyStar$ component is bounded as an operator from $L^p(\nuB)$ to $L^p(\nuB)$  and for any distributions $\Finit,\Fterm \ll \nuB$ of equal total mass whose densities lie in $L^p(\nuB)$ (equal masses are forced: the defect integrates to $\Finit(\Statespace)-\Fterm(\Statespace)$, so unequal masses floor the infimum away from $0$)

\begin{equation}
    \inf_{(\ForwardPolicyStar, \foutstar) \in \Theta }\left(\|\deltafinit\|_{L^p(\nuB)} + \|\deltafterm\|_{L^p(\nuB)} \right)=0
\end{equation}
for $(\ForwardPolicyStar, \foutstar)$ as a generative flow from $\Finit$ to $\Fterm$.

The family $\Theta$ is strongly universal if the zero infimum is realized (a flow realizing it for $(\Finit,\Fterm)$ has almost-everywhere-zero residuals, so realizes it for that pair at every $r\in[1,p]$).

\end{definition}


Universality is a statement about the \emph{ergodic theory} of the star forward policy:
$\ForwardPolicyStar$ is \emph{$\nuB$-ergodic} if $\nuB\ForwardPolicyStar=\nuB$ and the
only $\mu\ll\nuB$ it fixes are the multiples of $\nuB$. Here, then, the background measure is also the invariant one; the two part company only in Section~\ref{sec:cv_divergence}, where the invariant measure of the backward chain is written $\lambda$. Invariance makes the density
action $P_\star f := \frac{d((f\nuB)\ForwardPolicyStar)}{d\nuB}$ well defined on
$L^1(\nuB)$, and boundedness of $\|\ForwardPolicyStar\|_{L^p(\nuB)}$ says exactly that
it restricts to a bounded operator of $L^p(\nuB)$ --- automatic for $p=1$, a genuine
hypothesis beyond. Writing $\Pi$ for the projection onto the invariant densities, the
\emph{$L^p$-mixing coefficients} and their sum are
\begin{equation}\label{equ:mixing_coefficients}
    \beta^{(p)}_n := \left\|P_\star^n - \Pi\right\|_{L^p(\nuB)\rightarrow L^p(\nuB)},\qquad B_p := \sum_{n\geq 0}\beta^{(p)}_n .
\end{equation}
Summable $L^p$-mixing means $B_p<+\infty$; for $p=2$ these are the coefficients $\beta_n$ used in Section \ref{sec:cv_divergence} to control the speed of training.

The universality Theorem for Ergodic Generative Flows \cite{brunswicEGF} states $L^1$-universality of EGFs. Its proof may be improved to yield $L^p$-universality, in a quantitative form in which $B_p$ is the only property of the policy that enters.
\begin{theorem}\label{theo:universality_L2}
A generative flow family with a fixed forward policy and a free outflow is weakly
$L^p$-universal provided the policy is ergodic with summable $L^p$-mixing, and is strongly
universal when $p=+\infty$.
\end{theorem}
The summable-mixing hypothesis informally asks that iterating the policy averages functions
out fast enough for the errors to sum; ergodicity alone gives only the weaker Ces\`aro
convergence. It holds in particular whenever the density action has a spectral gap, and the
mixing sum $B_p$ of \eqref{equ:mixing_coefficients} is then finite. What the proof delivers is quantitative. Full
statement and proof: Theorem \ref{theo:universality_L2_full}, Appendix
\ref{app:universality}.

The two regimes part on attainment: at $p=+\infty$ the infimum is reached exactly, while for
$p<\infty$ it is only approached.

Universality is an \emph{existence} statement: mixing produces a minimizer. What it costs to reach that minimizer is a separate question, governed by the same ergodicity of the backward policy read as a speed rather than as a limit --- this is Theorem \ref{theo:training_speed}, and it is the through-line of Section \ref{sec:cv_divergence}.

Interestingly, Theorem \ref{theo:universality_L2} includes the case of graphs and is thus a far-reaching generalization of the classical universality statement for Generative Flows. On a finite path-connected state graph, the parameterization of generative flows by a state flow with a frozen, everywhere-positive backward policy \cite{morozov2025revisiting} attains \emph{exact} flow matching, and its balanced flow is unique up to scale; the target enters as the backward-policy row at the sink. The precise statement (Theorem \ref{theo:universality_graphs}), its complete proof and a discussion of the loop-closure mechanism it rests on are given in Appendix A.

\section{Convergence Bounds for Stable Losses}

 \begin{theorem}\label{theo:RL_CV_bound}
For a generative flow trained on a target of known density with the stable $L^q$
flow-matching loss over a training probability measure that covers the background measure,
$\nuB\ll\nutrain$ with $\dln\in L^\infty(\nutrain)$, and run at inference with the corrected
terminal flow $\widehat\target := (\Finit+\Finstar-\Foutstar)^+$, the loss bounds two things
at once: the sampling error in total variation, and the defect between the total mass of the
initial flow and that of the target. That loss can be driven to zero over an $L^p$-universal
family with $q\leq p$ at an initial flow carrying the target's total mass, and over a family
whose initial mass is itself trained exactly when the target's total mass is a limit of the
masses that family realizes, under technical integrability assumptions.
\end{theorem}
The target's total mass is unknown in practice, so the mass clause is met by training rather
than imposed: an initial flow whose total mass is itself a parameter has that mass driven to
$\target(\Statespace)$ by the same loss that drives the residual --- the measure-theoretic
form of a GFlowNet learning the partition function $Z$ --- and the condition for it, that the
family realize masses approaching $\target(\Statespace)$, is necessary as well as sufficient.

The constant in the sampling-error bound is
$C=\frac{2}{\target(\Statespace)}\left\|\dln\right\|^{1/q}_{L^\infty(\nutrain)}\nuB(\Statespace)^{1-1/q}$,
and the mass defect carries its own; three things about $C$ are worth stating plainly. It
needs no normalization of $\target$, which may be an unnormalized reward as is standard. It
is scale-consistent: replacing
$\nuB$ by $c\nuB$ divides the loss by $c$ and multiplies $C$ by $c$, as it must, since
the left-hand side does not see $\nuB$. And it is \emph{not} free of the size of the state
space --- on a finite space with counting measure and uniform $\nutrain$ it reads
$2N/\target(\Statespace)$, so the effective bound carries a factor
$N^{1-1/q}$, which is a H\"older step and is attained when the residual is spread evenly.
At $q=1$ that factor is absent and the bound is dimension-free; Section \ref{sec:experiments}
measures both regimes. Full statement and proof: Theorem \ref{theo:RL_CV_bound_full},
Appendix \ref{app:stable}.

Two objects the practitioner chooses rather than inherits sit inside this bound: the
initial flow and the training measure. Appendix \ref{app:finit} works out the first --- the
certificate-optimal mass overshoots the partition function, undershoot is the dearer error,
and the optimal shape solves a one-sided Chebyshev problem --- and Appendix \ref{app:replay}
the second, where a replay buffer is read as a device shaping $\nutrain$ and the two pipelines
return opposite verdicts: uniform reweighting is optimal for the certificate's constant,
backward replay for the training rate.

\begin{theorem} \label{theo:IL_CV_bound_body}
For a probability initial flow, a normalized target and a forward policy whose star inflow
has a density, trained with the KL-weak-FM loss at any $q\in[1,+\infty]$ and an admissible
weight $b\geq1+\|\dln\|_{L^{q^*}(\nutrain)}$, the loss is bounded below by the self-entropy
$\mathcal H$ of the target --- with equality as soon as the family contains flows dominating
the target up to a factor $1-\varepsilon$ for every $\varepsilon>0$, attained at exact
matching, which strong $L^p$-universality supplies when the initial flow and the target have
densities in $L^p(\nuB)$ --- and the sampling error is
$O\bigl(\sqrt{\LossWKLFM^q-\mathcal H}\bigr)$ as the loss approaches $\mathcal H$.
\end{theorem}
This is the target-by-samples counterpart of Theorem \ref{theo:RL_CV_bound}: the loss never
uses the density of $\target$, only samples from it. Reaching the self-entropy may require a large $0$-flow component --- expected, since
reaching some regions of the state space can take arbitrarily many moves. Full
statement and proof: Theorem \ref{theo:IL_CV_bound}, Appendix \ref{app:stable}.

\section{Convergence Bounds of Divergence losses}\label{sec:cv_divergence}

Training losses for Generative flows based on divergence and generalized divergences were introduced in \cite{brunswic2024theory} to show that this wide family of losses is generically unstable.
Further exploration was then conducted \cite{hu2024beyond,silva2024divergence} in the acyclic case. 
Instability of these losses implies that analogues of Theorem \ref{theo:RL_CV_bound} for these losses \textbf{cannot exist}; the following theorem makes this precise.

\begin{theorem}\label{theo:no_bound_divergence}
    Let $\mathcal L(F)$ be a divergence-based FM-loss on a directed graph $\mathcal G$, with
    generator continuous at $1$, depending on a target $\target$ and on a training distribution $\nutrain$; the supremum below ranges over the pairs $(\mathcal G,\target)$ with $\#\mathcal G = N$.
\begin{align*}
    &\forall\, N\geq 2,\quad \forall\, \varepsilon > 0, \\
    &\hspace{1cm}
    \inf_{\nutrain}\; \sup_{\#\mathcal G = N}\; \sup_{\substack{
      F \\ \mathcal L(F) \leq \varepsilon
    }} \; \TV\bigl(s_\tau \| \target\bigr) = 1
\end{align*}
\end{theorem}

To go beyond this limitation one must restrict the setting. Proposition 4.2 of
\citet{silvageneralization} bounds the total variation by an idealized trajectory-balance
residual on finite acyclic graphs of bounded horizon, with constants depending on the state
space and the training distribution --- bounded independently of the trained model on each
graph, but not uniformly in its depth; Appendix
\ref{sec:silva_comparison} restates it with explicit constants and compares it with ours. Such
restrictions fit autoregressive LLM finetuning
\cite{bartoldson2025trajectory,yu2024flow,wang2025gflowgr,ho2024proof,takase2024gflownet},
but a growing trend goes beyond autoregressive generation
\cite{zhu2025llada,nie2025large,sahoo2024simple,cetin2025large,von2025generalized}, where it
does not apply.

\citet{morozov2025revisiting} show that the detailed-balance loss is stable in a weaker sense
on finite graphs when the flow is parameterized with a fixed backward policy. Theorem
\ref{theo:db_stable_frozen} generalizes this quantitatively, from the summable $L^2$-mixing
hypothesis of our universality Theorem \ref{theo:universality_L2}. Before developing the
framework we record a counter-example to the naive generalization. Consider the graph

\begin{equation}\label{eq:counterexample_graph}
\xymatrix{
s_0 \ar[r]& 1  \ar[drrr]\ar[r]& 2  \ar[drr] \ar `ur^l[l]`^d[l] [l]  \ar[r] & \cdots \ar[r]& n \ar[d]\ar[r]& \cdots \ar[r] & 2^n \ar `ur^l[ll]`^d[ll] [ll]  \ar[r] \ar[dll] & \cdots  \\ &&&& s_f&&
}
\end{equation}
on which the uniform backward policy sends backward trajectories to $+\infty$, so the bound of
\citet{morozov2025revisiting}, resting on their expected length, says nothing. A policy that
damps the doubling edge repairs their condition and not ours:
$\BackwardPolicy(j\rightarrow 2j)=\frac{1}{2(j+1)}$ is positive on every edge and, when
backward trajectories enter the graph at states of mean $4.5$, gives $\bar\sigma=9$ exactly,
whereas \emph{every} constant doubling probability diverges --- and
its own mixing sum is infinite all the same.\footnote{Appendix \ref{app:doubling} proves both
halves for the family $\varepsilon(j)=c/(j+1)^s$, of which this policy is the member
$c=\tfrac12$, $s=1$: the $9$ is its closed-form backward length
(Proposition \ref{prop:doubling_length}), and $\widehat B=+\infty$ wherever the family is
positive recurrent (Theorem \ref{theo:doubling_unbounded}). The divergence is priced there on
truncations to $2^{12}$ states.}

\subsection{From Existence to Speed: the Backward Policy Sets the Clock}\label{sub:speed}

Theorem \ref{theo:universality_L2} and, on graphs, Theorem \ref{theo:universality_graphs}
settle \emph{existence}: an ergodic backward policy with summable mixing makes the
parameterization universal, and on a finite graph the frozen-backward family contains an
exactly balanced flow, unique up to scale. They say nothing about how long training takes to
find it. That is what the rest of this section answers, and the answer has a single currency.

Let $G$ be a finite path-connected marked graph, let the frozen backward policy be positive on the edges of the loop closure $\widehat G$ (Definition \ref{def:loop_closure}), and for a state $x$ let $\sigma(x)$ be the time the backward trajectory issued from $x$ spends in the graph, that is, its hitting time of $s_0$. Put
\begin{equation}\label{eq:morozov_variables}
    \sigma_*:=\max_x\mathbb E(\sigma\mid X_0=x),\qquad \bar\sigma:=\mathbb E_{x\sim\widehat\BackwardPolicy(s_f\rightarrow\cdot)}\mathbb E(\sigma\mid X_0=x),
\end{equation}
$\bar\sigma$ being the expected backward-trajectory length of \citet{morozov2025revisiting} (the maximum in $\sigma_*$ running over every state of the loop closure, $s_f$ included), and let $N(x)$ be the expected number of visits to $x$ made by one backward trajectory drawn from the target row $\widehat\BackwardPolicy(s_f\rightarrow\cdot)$, with $N_{\min}:=\min_xN(x)$.

\begin{theorem}[training speed under a frozen backward policy]\label{theo:training_speed}
On a finite marked graph with a frozen backward policy positive on the edges of its loop
closure, and for a training measure of uniformly positive density, gradient training of the
state-space balance loss generated by $(\log x)^2$ --- the generator of the practical DB and
TB losses --- converges to a balanced flow from
\emph{every} positive initialization, first at a universal \L ojasiewicz rate, then
exponentially at a rate fixed by the time backward trajectories spend in the graph. Both
phases hold for the gradient flow and for discrete gradient descent, the latter with a step
size chosen from the initialization --- no step size fixed in advance works from every
initialization. No mixing, spectral-gap or
aperiodicity hypothesis is used for this state-space loss; the DB instance inherits a local
rate through the edge lift wherever the lifted mixing sum is finite (Theorems
\ref{theo:db_stable_frozen} and \ref{theo:local_convergence}).
\end{theorem}
Here $\lambda$ is the backward chain's invariant measure, which is its visit ratio,
\begin{equation}\label{eq:occupation}
    \lambda(x)=\frac{N(x)}{2+\bar\sigma},\qquad\text{so}\qquad \lambda_{\min}=\frac{N_{\min}}{2+\bar\sigma},
\end{equation}
and $\varrho_\sigma$ is the second-phase rate, per unit mass at the crossover:
\begin{equation}\label{eq:morozov_rate}
    \varrho_\sigma\;=\;\frac{g''(1)\,w_{\min}\,\lambda_{\min}}{\sigma_*^{\,2}}\;=\;\frac{2\,w_{\min}\,N_{\min}}{\sigma_*^{\,2}\,(2+\bar\sigma)} ,
\end{equation}
Its dependence on the graph and the policy runs through $\sigma_*$, $\bar\sigma$ and
$N_{\min}$ alone. Full statement, with the
\L ojasiewicz envelope and the explicit crossover time, and proof: Theorem
\ref{theo:training_speed_full}, Appendix \ref{app:divergence}.

The last sentence is the point of the reformulation. What governs training is not how long
the backward chain takes to forget its starting point but how long it takes to \emph{leave the
graph}, and \eqref{eq:occupation} converts one into the other. Three consequences: aperiodicity
leaves the hypotheses, which matters because a leveled --- autoregressive --- graph is periodic
and has infinite mixing sum; the hypothesis becomes checkable, since $\bar\sigma$ and the
$N(x)$ are expectations over backward trajectories of the model one is already training; and
the rate degenerates on the graph
\eqref{eq:counterexample_graph} at every doubling probability --- $\sigma_*$ is a maximum over
states --- and not only in the window where Morozov's bound fails (\S\ref{sec:doubling_cost}). The price is that
$\widehat B_\sigma$ is a conservative envelope: \eqref{eq:morozov_rate} is a guarantee, not a
prediction, and Section \ref{sec:experiments} measures how conservative. Remark
\ref{rem:visit_ratio} separates the regime each constant governs.

\subsection{First Variation of Balance Losses}

Everything rests on the first variation of the $g$-divergence balance loss
$\mathcal L_{g,\nu}(\mu) := \int g\left(\frac{d(\mu T)}{d\mu}\right)d\nu$, for an
\emph{arbitrary} Markov kernel $T$. Four measures meet in its gradient: the background
measure $\nuB$, which carries the parameterization; the weighting measure $\nu$, which is
the training measure $\nutrain$ in practice --- we keep the generic letter because the
statements below hold for any $\nu$ with the positivity each one names and a bounded density,
the latter automatic on a finite state space and not on a continuous one; an invariant measure
$\lambda$ of $T$, the one of the four that the practitioner neither chooses nor samples; and
$\nuG$, any finite reference measure equivalent to $\lambda$, which fixes the geometry the
gradient is read in rather than the objective.

\begin{theorem}\label{theo:first_variation}
For a continuously differentiable generator, the $g$-divergence balance loss of an arbitrary
Markov kernel is differentiable at every flow whose balance ratio is bounded away from $0$ and
$\infty$, and --- as soon as the kernel admits a non-zero finite invariant measure equivalent
to that flow --- its gradient is the reversal of the kernel, minus the ratio, applied to $g'$
of that ratio.
\end{theorem}
The reversal here is conjugation in $L^2(\lambda)$ against the kernel's invariant measure: an
involution sending $T$ to the kernel $T^\lambda$ with $\lambda\otimes T^\lambda = T\otimes\lambda$,
and exchanging a backward policy with the forward policy of the \emph{balanced} flow (Lemma
\ref{lem:adjoint}) --- the forward/backward duality of generative flows. Explicitly,
$\nabla_{\mu}^{\nuG}\mathcal L_{g,\nu} = \frac{d\nuG}{d\lambda}\cdot(T^\lambda-r)\left[g'(r)\frac{d\nu}{d\mu}\lambda\right]$
with $r=\frac{d(\mu T)}{d\mu}$. Full statement and proof: Theorem
\ref{theo:first_variation_full}, Appendix \ref{app:divergence}; the rates below are read at
the invariant measure of the chain in play.

FM, DB and TB are all specializations of it. The DB and TB ratios live on edges and trajectories rather than
states, but they are the same balance loss for the same frozen backward policy lifted to a
power of the state space, so every convergence statement below transfers to them with explicit
constants (Appendix \ref{app:divergence}).

\subsection{Gradient Descent Diffuses the Flow by the Backward Policy}\label{sub:gd_diffusion}

We now linearize the gradient flow of these losses at the balanced solution and read off the convergence speed of training from the mixing of $\BackwardPolicy$. Throughout, $(\Statespace',\lambda',T)$ denotes a chain together with its invariant measure: either $(\Statespace,\lambda,\BackwardPolicy)$ for the FM loss, or its edge lift $(\Statespace^2,\lambda_2,K_2)$ for the DB loss. Then $P$ is the density action of $T$ on $L^2(\lambda')$, $P^\dagger$ its adjoint --- the density action of the reversal $T^{\lambda'}$, by Lemma \ref{lem:adjoint} --- $A := P - I$, and $M_w$ multiplication by $w=\frac{d\nutrain}{d\lambda'}$.

\begin{theorem}[gradient descent as backward diffusion]\label{theo:gd_diffusion}
Near the balanced flow, and for any generator of class $C^3$ near $1$ with non-degenerate curvature there, the
gradient field of a balance loss linearizes to the backward policy's diffusion applied twice, so that training relaxes the current flow guess by diffusing
it through the frozen backward policy --- and for the FM loss with a reversible policy and
uniform training weight ($w\equiv1$) it is exactly the square
of that policy's heat flow, each of its eigenmodes decaying at $g''(1)(1-\beta)^2$.
\end{theorem}
Full statement and proof: Theorem \ref{theo:gd_diffusion_full}, Appendix
\ref{app:divergence}.

The convergence speed of training is therefore a property of the spectrum of that policy.

\begin{theorem}[convergence speed under frozen backward policy]\label{theo:db_stable_frozen}
For a training measure of uniformly positive density, a backward policy with summable
mixing makes gradient descent on the \emph{linearized} FM and DB losses with frozen backward
policy converge to the balanced ray at an explicit exponential rate, provided the step size
stays below an explicit cap; the linearization conserves the total flow normalization along
the way. At vanishing training weight the rate is vacuous.
\end{theorem}
The rate is $\varrho = g''(1)w_{\min}/\widehat B^{\,2}$ with $\widehat B$ the mixing sum,
which for the FM loss is $B$ itself and for DB is $1+B$. Full statement and proof:
Theorem \ref{theo:db_stable_frozen_full}, Appendix \ref{app:divergence}.

\begin{theorem}[local exponential convergence]\label{theo:local_convergence}
On a finite state space, and from an explicit neighbourhood of the unit-mass balanced flow,
the rate above upgrades to the full nonlinear dynamics: halved for the gradient flow, and
quartered for gradient descent below an explicit step.
\end{theorem}
Full statement and proof: Theorem \ref{theo:local_convergence_full}, Appendix
\ref{app:divergence}.

These are statements about the \emph{dynamics}: the balanced ray attracts, exponentially and
at an explicit rate. They are not statements about the loss landscape in the sense of
\citet{brunswic2024theory}, whose Definition 3 calls a loss \emph{stable} when adding a
$0$-flow never decreases it,

\begin{equation}\label{eq:def_stability}
    \mathcal L(F+F_0)\;\geq\;\mathcal L(F) \qquad \text{for every $0$-flow $F_0$ that is a subflow of } F,
\end{equation}
and whose Theorem 3 declares the FM, DB and TB losses unstable. One might hope that freezing
the backward policy repairs \eqref{eq:def_stability}, since it collapses the balanced set to a
single ray. It does not, and the obstruction is the very mechanism that makes training work.

\begin{proposition}[freezing does not restore stability]\label{prop:frozen_unstable}
On every finite path-connected graph carrying a directed cycle that avoids both the source
and the sink, for every admissible generator ($g$ continuous, vanishing at $1$ and positive
elsewhere) and every training measure of positive density, the FM and DB losses with an
everywhere-positive frozen backward policy remain unstable in the sense of
\eqref{eq:def_stability}:
some flow is strictly improved by \emph{adding} a $0$-flow.
\end{proposition}
Full statement and proof: Proposition \ref{prop:frozen_unstable_full}, Appendix
\ref{app:divergence}.

What freezing buys is not
\eqref{eq:def_stability} but the weaker property invoked by \citet{morozov2025revisiting}: the
loss vanishes only on the balanced ray, is coercive transverse to it, and gradient descent
converges there. The three notions --- stability, weak stability, convergence --- are distinct,
and only the last two survive a frozen backward policy.

On trajectories of $\ell$ transitions the TB Hessian at balance is at most that of DB times
$\ell$ and a mixing factor, so a DB step cap divided by both remains valid for TB; no rate is
proved for TB (Proposition \ref{prop:tb_hessian}). Summable mixing, sufficient for the FM
and DB rates, is
the hypothesis of Theorem~\ref{theo:universality_L2} at $p=2$: \emph{the same ergodicity of the backward
policy provides both the universality of the parameterization and the convergence speed of its
training}. Practice, however, often trains the backward policy rather than freezing it
\citep{morozov2025revisiting}; Appendix \ref{app:pbtrain} translates that objective into the
spectral variables used here, where flow regularization becomes a slow drift along the
balanced manifold toward an optimal-transport endpoint, every constant above pays a computable
price along the way, and an entropic term repairs it. Remark \ref{rem:visit_ratio} resolves the speed state by state.

Whether convergence extends beyond the neighbourhood of Theorem \ref{theo:local_convergence} depends on the generator, and the next theorem separates two classes of it: the strictly unimodal generators, and those whose $g'$ vanishes on bands on either side of $1$ --- bands of zero mobility $g''$, on which the diffusion of Theorem \ref{theo:gd_diffusion} freezes. The two classes do not exhaust the generators: bands on one side of $1$ only fall under neither, and within the first, global convergence is proved for $(\log x)^2$ and $(x-1)^2$ only, a bounded strictly unimodal generator not being covered. For $g=(\log x)^2$, the generator of the practical DB and TB losses, one may fear a milder failure on cyclic graphs: a flow stuck in a \emph{distant equilibrium} between the frozen-policy constraint and the circulation-inflating force of Theorem \ref{theo:no_bound_divergence}. The following dichotomy rules this out.

\begin{theorem}[global dichotomy for frozen-backward training]\label{theo:global_dichotomy}
If $g$ is strictly unimodal --- as $(\log x)^2$ and $(x-1)^2$ are --- then every critical point
of the loss is balanced, on every graph and for every training measure with positive density,
and --- for $(\log x)^2$ and $(x-1)^2$ --- on a finite state space the gradient flow
converges globally from every positive initialization. If instead
$g'$ vanishes on bands on either side of $1$, then every flow whose ratios lie in those bands
is frozen, for every training measure; and some such generators freeze an open set of flows
at a definite distance from balance, already on the two-state chain with maximal mixing, so
that global convergence fails for them.
\end{theorem}
Full statement and proof: Theorem \ref{theo:global_dichotomy_full}, Appendix
\ref{app:divergence}.

Convex or real-analytic generators exclude this freezing altogether (Remark
\ref{rem:freezing}).

\section{Experiments}\label{sec:experiments}

The bounds above are only as informative as their constants, so we measure them. Every
experiment runs on an \emph{enumerable} state space, which is a design choice rather than a
limitation: the sampling error $\TV(s_\tau\|\target)$ is then a linear solve rather than an
estimate, and a figure testing a bound must not be dominated by Monte-Carlo noise of the same
order as the quantity under test. Flows are tabular elements of $\mathcal M^2(\lambda)$
--- the object Theorems \ref{theo:first_variation}--\ref{theo:local_convergence} actually
describe --- so no network is involved anywhere and the gradient flow measured below is the
one the theory predicts. The parameterized case is correspondingly \emph{not} tested here.
Code, configurations and seeds accompany the paper.

\paragraph{The stable bound holds, and at $q=1$ it does not see the state space.}
Starting from the exactly balanced flow produced by the construction in the proof of Theorem
\ref{theo:universality_L2} --- whose residual and sampling error are both zero to machine
precision, itself a check on that construction --- we perturb multiplicatively over four
orders of magnitude and compare the two sides of Theorem \ref{theo:RL_CV_bound}. Across
$21{,}168$ flows spanning seven environment classes (nine environment--policy families)
and state spaces from $N=6$ to
$N=386{,}880$ there is \textbf{no violation} (Figure \ref{fig:bound}, left). Testing a family
of flows rather than a training trajectory is the stronger test, since the theorem quantifies
over every flow and not only those gradient descent visits.

The tightness is the more interesting half. At $q=1$ the median ratio
$C\mathcal L/\TV$ is $6.0$ at $N=6$ and, at fixed near-balance perturbation size, stays
between $6.0$ and $6.6$ all the way to $N=386{,}880$: the bound sits within one
order of magnitude of the truth and \emph{does not degrade} over a $64{,}000$-fold range of
state-space size. (Pooled over all perturbation sizes the per-$N$ median rises to $32$ at
the largest $N$ --- growth carried by the far-from-balance rows, not by the constant.)
For $q>1$ it grows, with fitted exponents $N^{0.20}$ at $q=2$ and $N^{0.60}$
at $q=\infty$ against the worst case $N^{1-1/q}$ of the H\"older step (Figure \ref{fig:bound},
right). So the $N$-dependence discussed after Theorem \ref{theo:RL_CV_bound} is real but is a
consequence of the choice of $q$, and vanishes at $q=1$.

\begin{figure}[tb]
\centering
\includegraphics[width=0.45\textwidth]{figures/fig1_bound.pdf}\hfill
\includegraphics[width=0.45\textwidth]{figures/fig2_scaling.pdf}
\caption{\textbf{Left:} exactly computed $\TV$ against the right-hand side of Theorem
\ref{theo:RL_CV_bound}, one point per flow; everything lies on or below the diagonal.
\textbf{Right:} the tightness ratio against state-space size, with the fitted exponent per
$q$. At $q=1$ the constant is flat over two decades of $N$.}
\label{fig:bound}
\end{figure}

\paragraph{The KL-weak-FM constant is below $1$ wherever the bound is non-vacuous.}
Theorem \ref{theo:IL_CV_bound_body} gives $\TV = O(\sqrt{\mathcal L - \mathcal H})$, and
its full form names the constant --- $\tfrac1{\sqrt2}$ on the square root, beside a capped
term linear in the gap that puts the guaranteed ratio above $\tfrac1{\sqrt2}$ --- so what
we measure is the constant realized: over $3{,}472$ (flow, $q$) instances on toroidal EGF,
Cayley, hypergrid and product-replacement families ($N$ from $24$ to $336$), the floor
$\mathcal L\geq\mathcal H$ holds on \emph{every} instance, and the measured constant
$\TV/\sqrt{\mathcal L-\mathcal H}$ never exceeds $0.612$. Two disclosures keep this
honest. First, the constant is defined on the $1{,}572$ instances where the loss is
finite; on the other $1{,}900$ it is $+\infty$ because the inferred terminal flow
vanishes on part of the target's support --- the forward-KL mechanism, observed rather
than hypothesized (the bound holds there vacuously). Second, the constant is not one number: the
full-support EGF arm gives $\leq0.163$, a support-restricted EGF variant $0.324$, and the
arms at $q\in\{2,4,\infty\}$ and on non-EGF families $\leq0.612$. Theorem
\ref{theo:IL_CV_bound} restricts neither the exponent nor the parameterization, so that
spread measures the constant rather than the boundary of the proof.

\paragraph{The diffusion picture predicts the training rate; the hitting-time constant does not.}
Running the exact gradient flow of $\mathcal L_{g,\nu}$ in $\mathcal M^2(\lambda)$ and
fitting its exponential rate near balance, the prediction
$\varrho_{\mathrm{true}} := g''(1)w_{\min}/B_{\mathrm{true}}^2$ --- Theorem
\ref{theo:db_stable_frozen} with the exact coercivity constant in place of the mixing sum ---
matches the measured rate to within $1$--$10\%$ on all $20$ distinct chains tested
($23$ runs; three leveled cases are carried in both sweeps), across Cayley
graphs of $\mathfrak S_n$, leveled graphs and product-replacement graphs (Figure
\ref{fig:rates}). The linearization of Theorem \ref{theo:gd_diffusion} is therefore not merely
qualitative.

The guarantee $\varrho_\sigma$ of Theorem \ref{theo:training_speed} is a different matter: valid
everywhere, and between $187$ and $10^{9}$ times too small, the gap widening with $N$. On
leveled graphs, where the mixing sum genuinely diverges and $\varrho_\sigma$ is the only rate
available, the exact coercivity constant is nonetheless $B_{\mathrm{true}}\approx1$ --- the
smallest in the study, with the fastest observed convergence. The mixing route fails precisely
where the underlying coercivity is best. Theorems \ref{theo:db_stable_frozen_full} and
\ref{theo:local_convergence_full} are stated accordingly, for any coercivity constant
\eqref{eq:coercivity} --- a bound $\geq1$ on $\|(I-P)^{-1}\|$ over mean-zero functions ---
which summable mixing and hitting times (Proposition \ref{prop:morozov_rate}) each supply.

\begin{figure}[tb]
\centering
\includegraphics[width=0.72\textwidth]{figures/fig4_rates.pdf}
\caption{Guaranteed against achieved local convergence rate on the $11$ chains of the
Cayley/leveled sweep. The exact
coercivity constant (crosses) lands on the measured rate (circles) everywhere; the
hitting-time guarantee of Theorem \ref{theo:training_speed} (squares) is two and a half to
six orders
of magnitude below. The three rightmost are leveled graphs, where the mixing sum is infinite.}
\label{fig:rates}
\end{figure}

\paragraph{The minimax obstruction, made visible.}
Theorem \ref{theo:no_bound_divergence} is a statement about a construction, and the
construction is small enough to evaluate exactly. On the cyclic graph $\mathcal C_8$ of its
proof, adding circulation $k$ to the base flow leaves the sampler --- and hence the sampling
error --- exactly where it was, at $\TV=0.9857$, while the squared divergence losses fall
thirteen orders of magnitude --- $(\log x)^2$ and $(x-1)^2$ reach $2\cdot10^{-13}$ at
$k=10^6$, and $|x-1|$ falls seven. The stable
loss of Theorem \ref{theo:RL_CV_bound}, being a norm of the \emph{difference} of the two sides
rather than of their ratio, is flat in $k$ (Figure \ref{fig:minimax}). The contrast is in the
shape: one family of losses reports the error and the other does not.

\paragraph{A family where the mixing hypothesis provably holds, uniformly.}
Every graph considered so far has constants that grow with $N$, which leaves open whether the
summable-mixing hypothesis is ever satisfied with a constant independent of the state space.
It is. By the theorem of \citet{lubotzky2001product}, property (T) for $\mathrm{Aut}(F_k)$ ---
known for $k\geq4$ \citep{nitsche2020computer,kaluba2019aut,kaluba2021property} --- gives
the product-replacement graphs $\Gamma_k(G)$ a spectral gap
bounded below uniformly over all finite groups $G$. On the ten groups carried at every $k$
($|G|\leq12$, $N\leq246{,}840$; further groups reach $|G|=60$ and $N=386{,}880$ at
individual $k$), the spread of $B$ across groups is $3.92\times$ at
$k=3$, $2.40\times$ at $k=4$ and $1.93\times$ at $k=5$: it tightens as $k$ crosses the
threshold (Figure \ref{fig:minimax}, centre). At $k=2$, where property (T) fails, five of ten groups
give spectral gap exactly zero.

They also make $B$ a controlled variable rather than a property of the environment: for a
\emph{reversible} backward policy the coercivity constant is exactly the reciprocal spectral
gap, since $I-P$ then has eigenvalues $1-\zeta_i$, an identity confirmed against the
loop-closed constant to within $2\%$ and improving with $N$. Reweighting the moves so as to
keep reversibility, or interpolating with a lazy or restart step, then moves $B$ over a range
of $1154$ in closed form (Figure \ref{fig:minimax}, right). This also answers a practical
objection to the coercivity formulation: for a reversible policy the hypothesis is a single
eigenvalue, computable at $N=4\cdot10^5$ in seconds.

\paragraph{What this study does not reach.}
One gap is taken up in the appendix: Appendix \ref{app:sgd} supplies the
stochastic-gradient term absent from the dynamics above, where balance losses turn out to
be noise-free at their optimum while imitation terms are not. Appendix \ref{app:program}
states, for each result of the paper, the experiment that would ground or refute it best.

\begin{figure}[tb]
\centering
\includegraphics[width=0.33\textwidth]{figures/fig3_minimax.pdf}\hfill
\includegraphics[width=0.63\textwidth]{figures/fig7_pra_scale.pdf}
\caption{\textbf{Left:} on the graph of Theorem \ref{theo:no_bound_divergence}, divergence
losses $\mathcal L_{g,\nu}$ --- $\nu$ the training measure --- vanish along $F_k=\widehat F+k\gamma$ while the sampling error stays pinned at
$0.9857$; the stable loss does not move. \textbf{Centre and right:} $B=1/\mathrm{gap}$ of
$\Gamma_k(G)$ against group order, the spread across groups tightening as $k$ crosses the
property (T) threshold, and the range of $B$ each backward-policy family reaches.}
\label{fig:minimax}
\end{figure}

\section{Conclusion and Limitations}\label{sec:conclusion}

One quantity does double duty in this picture: the ergodicity of the backward policy makes the
parameterization expressive enough to contain a solution, and it sets the clock on which
gradient descent finds it. The same hypothesis that yields universality yields the rate.

The strongest limitation we can state is one our own experiments produced. The hitting-time
guarantee is valid everywhere and is between $187$ and $10^{9}$ times too small --- and it is
worst precisely where the underlying coercivity is best, on leveled graphs where the mixing
sum diverges while convergence is fastest. Since that route enters the proofs only through a
coercivity inequality, Theorems \ref{theo:db_stable_frozen_full} and
\ref{theo:local_convergence_full} hypothesize coercivity directly, summable mixing being one
sufficient condition among others; for a reversible policy the hypothesis is a single
eigenvalue, computable at $N=4\cdot10^5$ in seconds. Second, everything measured here
is tabular: the parameterized case is not tested, and Appendix \ref{app:program} specifies,
for each result, the experiment that would ground or refute it, with success and refutation
criteria fixed in advance.

%\section{Empirical Distributions Training}








\begin{theorem} Let $(\Statespace,\lambda)$ be a compact Riemannian manifold endowed with a positive measure equivalent to Lebesgue measure.
Let $\mu$ be a target distribution and let $\widehat{\mu}$ be the empirical distribution given by $n\in \mathbb N$ samples taken from $\mu$. Let  $(\ForwardPolicyStar, \Foutstar)$ be an EGF on $\Statespace$ trained with $\mathcal L_{} $ on a dequantized distribution obtained by replacing each $\delta_x$ by a distribution $\epsilon(x,\sigma)$ of expectancy $x$ and standard deviation $\sigma$.
Then for any $\eta>0$, there exists $C_\eta>0$ such that:
\begin{equation}
    \mathbb E_{\widehat{\mu}}(W_p(s_\tau || \mu) \leq \mathrm{diam}(\Statespace)\left[\frac{\delta}{1+\delta} + \frac{1}{\sqrt{2}}\sqrt{-\mathcal H-\mathcal L_{CE}+\delta}\right] +2^{1/p}\sigma + \frac{C_\eta}{n^{1/d^*_p(\mu)-\eta}}
\end{equation}


\end{theorem}
\begin{proof}
    Since $W_p$ is a distance on $\Statespace$, we have
    \begin{equation}
        W_p(s_\tau|| \mu ) \leq W_p(s_\tau || \widehat{\mu}^\epsilon) + W_p(\widehat{\mu}^\epsilon|| \widehat{\mu} ) + W_p(\widehat{\mu}||\mu).
    \end{equation}
    Since $\Statespace$ is compact, we have $W_p(s_t||\widehat{\mu}^\epsilon)\leq \mathrm{diam}(\Statespace) \mathrm{TV}(s_t||\widehat{\mu}^\epsilon)$.
    Corollary 3.10 of \cite{brunswicEGF} then gives
    \begin{equation}
         \mathrm{TV}(s_t||\widehat{\mu}^\epsilon) \leq \frac{\delta}{1+\delta} +\frac{1}{\sqrt{2}} \sqrt{-\mathcal H - \mathcal L_{CE}+\log(1+\delta)}
    \end{equation}
\end{proof}


\cite{weed2019sharp}   % A6 cut -- see SUBMISSION-ICLR2027.md

%%%% ---------------------------------------------------------------------
%%%% Required / recommended ICLR 2027 statements.
%%%% Neither counts toward the 9-page limit.  Placed at the end of the main
%%%% text, before references, per the ICLR 2027 author guidelines.
%%%% ---------------------------------------------------------------------
\section*{AI use statement}
% REQUIRED by ICLR 2027.  Author to confirm wording before submission.
Large language models were used substantially in the production of this paper, under
the author's direction and review: they drafted the text of
Appendices~\ref{app:finit}--\ref{app:doubling} from the author's research program and
results; they wrote and tested the numerical audit scripts and figure-generation code
that constitute the experimental evidence of those appendices; they performed
adversarial verification passes over the proofs and the empirical claims, several of
whose findings led to corrections incorporated in this version; and they assisted with
copy-editing, \LaTeX{} and bibliography maintenance, and the implementation and
testing of the experiment library. The research questions, the theorems and their
proof strategies, and the final verification of every mathematical claim are the
author's; every numerical claim is additionally backed by a released audit script
whose exit status re-derives it.

\section*{Reproducibility statement}
% RECOMMENDED by ICLR 2027.
All experiments use enumerable state spaces so that the sampling error
$\TV(s_\tau \| \target)$ is computed exactly rather than estimated; this is
necessary for the figures to test the bounds rather than Monte-Carlo noise.
Anonymized code, configurations and seeds are provided as supplementary material,
including the audit scripts that re-derive every number quoted in the appendices.
Complete proofs of all statements appear in the appendix.

\bibliography{note}
\bibliographystyle{iclr2027_conference}

\appendix
\section{\appendixAtitle}

Proofs of statements from the main text are given as named proofs, in order of appearance: universality first, then the two convergence bounds for stable losses, then the divergence section. One statement is proved here in a form narrower than the one advertised in the main text, and is labelled accordingly: Theorem \ref{theo:universality_L2_body} is the $L^2$ strengthening of Theorem \ref{theo:universality_L2} for EGF parameterizations. Theorem \ref{theo:IL_CV_bound}, the full form of Theorem \ref{theo:IL_CV_bound_body}, constrains neither the parameterization nor the exponent, holding for every $q\in[1,+\infty]$ and every forward policy whose star inflow is dominated by $\nuB$; its floor is reached under a hypothesis of approximate domination of the target by the family, exact flow matching being the borderline case. The final subsection is dedicated to the universality theorem on finite graphs (Theorem \ref{theo:universality_graphs}): statement, complete proof, and a discussion of the loop-closure mechanism.

\subsection{Notation}\label{app:notation}

In this work, we consider generative flows in the measurable setting, encompassing scenarios where the state space $\Statespace$ is either a vector space, $\Statespace = \mathbb R^d$, or potentially a Riemannian manifold \cite{lee2018introduction, gemici2016normalizing}. Both are instances of Polish\footnote{A topological measurable space is Polish if its topology is completely metrizable and it is endowed with its Borel $\sigma$-algebra.} spaces that admit a natural finite background measure $\nuB$ (usually the Lebesgue, Haar \cite{halmos2013measureHaar} or Riemannian volume measures).

    A measure $\mu$ is dominated by $\nu$, denoted $\mu\ll \nu$, if $\nu(A)=0 \Rightarrow \mu(A)=0$ for any measurable $A\subset \Statespace$; the Radon-Nikodym derivative of $\mu$ with respect to $\nu$, denoted $\frac{d\mu}{d\nu}\in L^1(\nu)$, is then the unique $\nu$-integrable function $\varphi$ such that $\varphi \nu = \mu$. A
    Markov kernel $\pi:\mathcal X\rightarrow \mathcal Y$ is a stochastic map denoted $\pi(x)$. The image measure of $\mu$ by $\pi$ denoted $\mu \pi$ is defined by $\mu\pi(A) = \int_{x\in \mathcal X} \mathbb P(\pi(x)\in A) d\mu(x)$ for any measurable $A\subset \mathcal Y$. Their tensor product is the kernel  $\mu\otimes \pi(A\rightarrow B):=\int_{x\in A}\mathbb P(\pi(x)\in B)d\mu(x)$.
    Even when $\mu$ is an unnormalized distribution, we write $x\sim \mu$ for ``the law of $x$ is $\frac{1}{\mu(\Statespace)}\mu$".

    Recall that the Lebesgue space $L^p(\Statespace,\mu)$ for a measure $\mu$ on $\Statespace$  and $p\in [1,+\infty]$ is defined as the set of measurable functions $f:\Statespace \rightarrow \mathbb R$ such that $\int_{s\in \Statespace}|f|^pd\mu(s)<+\infty$. The conjugate exponent of $p\in[1,+\infty]$ is written $p^*$ and is defined by $\frac1p+\frac1{p^*}=1$: $p^*=\frac{p}{p-1}$ on $]1,+\infty[$, $p^*=+\infty$ at $p=1$ and $p^*=1$ at $p=+\infty$, so that H\"older's inequality $\int|fg|d\mu\leq\|f\|_{L^p(\mu)}\|g\|_{L^{p^*}(\mu)}$ holds at both endpoints. The same convention applies to the exponent $q$ of the losses below, whose conjugate exponent is written $q^*$.

A family of generative flows on $(\Statespace,\nuB)$ is a set $\Theta$ whose elements are couples $(\ForwardPolicyStar,\foutstar)$ where $\foutstar : \Statespace \rightarrow \mathbb R_+$  is an integrable scalar field on $\Statespace$, $\foutstar\in L^1(\nuB)$ --- equivalently the star outflow $\Foutstar:=\foutstar\nuB$ is a finite measure, $\Foutstar(\Statespace)<+\infty$ --- and $\ForwardPolicyStar:\Statespace\rightarrow \Statespace$ is a Markov kernel.
For a given $p\in [1,+\infty]$, a family $\Theta$ is metrized by the $L^p(\nuB)$ norm for $\foutstar$ and the operator norm $L^p(\nuB)\rightarrow L^p(\nuB)$ for $\ForwardPolicyStar$:
\begin{eqnarray}
    \|\foutstar\|_{L^p(\nuB)} &:=& \left(\int_{s\in \Statespace} |\foutstar(s)|^pd\nuB(s)\right)^{1/p} \\
    \|\ForwardPolicyStar\|_{L^p(\nuB)} &:=& \sup_{0\neq f\in L^p(\nuB)} \frac{\|f\ForwardPolicyStar\|_{L^p(\nuB)}}{\|f\|_{L^p(\nuB)}}.
\end{eqnarray}
\emph{Four measures, and the conventions they carry.} The \emph{background} measure $\nuB$ is the one flow densities and kernels are parameterized against: $\mathcal M^+(\Statespace,\nuB)$ denotes the non-negative measures dominated by $\nuB$, and $\mathcal M^2(\nuB)$ those with square-integrable density, with inner product $\langle\mu\mid\nu\rangle_{\nuB}:=\int\frac{d\mu}{d\nuB}\frac{d\nu}{d\nuB}\,d\nuB$; in the first-variation theorem it fixes the coordinates and nothing else. The letter $\lambda$ is reserved for an \emph{invariant} measure --- of whichever kernel is in play, and a probability on the loop closure --- assumed to exist, non-zero and finite, and unique up to a positive scalar when that kernel is ergodic, the normalization as a probability on the loop closure pinning that scalar --- the visit-ratio identity \eqref{eq:occupation} is not scale-invariant, and neither is the rate $\varrho=g''(1)w_{\min}/\widehat B^{\,2}$ of Theorem \ref{theo:db_stable_frozen_full}, whose $w_{\min}=\min\frac{d\nu}{d\lambda}$ becomes $w_{\min}/c$ under $\lambda\mapsto c\lambda$ at fixed weighting measure while $\widehat B$ itself does not move, nor the constant $C_\infty=(\min\lambda)^{-1/2}$ of Theorem \ref{theo:local_convergence_full}; it is the measure the reversal $T^\lambda$ and the mixing sum $\widehat B$ are defined against, and the only one of the four that the practitioner neither chooses nor samples. The two coincide in \S\ref{app:universality} and \S\ref{app:stable}, where the star policy is assumed $\nuB$-invariant so that $\nuB$ plays both roles; they part company in \S\ref{app:divergence}, where $\lambda$ is the invariant measure of the backward chain while $\nuB$ remains ambient. The \emph{training} measure is $\nutrain$: the law the replay buffer and the sampling strategy make available, and what the generic weighting measure $\nu$ of a balance loss is in practice --- the letter stays generic because the statements below are quantified over every $\nu$, and Theorem \ref{theo:first_variation_full} asks of it only $\frac{d\nu}{d\mu}\in L^\infty(\mu)$, domination by the current flow guess, which is a constraint on $\nutrain$ and not a property a replay buffer supplies. The fourth measure, $\nuG$, fixes the geometry rather than the objective: writing a flow by its density against $\nuG$ and measuring in $L^2(\nuG)$, $\nabla^{\nuG}_\mu\mathcal L_{g,\nu}$ is the representative of the differential of $\mathcal L_{g,\nu}$ in that inner product. The differential does not depend on $\nuG$; the gradient does, and with it the trajectory of gradient descent --- $\nuG$ is a preconditioner, free up to finiteness and $\nuG\sim\lambda$. The choice matching the parameterization is $\nuG=\nuB$, flows being carried as densities against the background measure, and it is available where $\nuB$ is itself an invariant measure equivalent to $\lambda$: in \S\ref{app:universality}, whose ergodicity hypothesis leaves no other invariant measure dominated by $\nuB$, and in \S\ref{app:doubling}, where $\nuB=\lambda$ outright. Invariance alone does not suffice, Lemma \ref{lem:adjoint} giving $\lambda\ll\nuB$ and not the converse. The choice made throughout this appendix is $\nuG=\lambda$, which sets the prefactor $\frac{d\nuG}{d\lambda}$ to $1$ and puts the gradient in the space where the mixing of $\BackwardPolicy$ acts. That freedom belongs to Theorem \ref{theo:first_variation_full} alone: Proposition \ref{prop:no_distant_equilibrium} and every rate in Theorems \ref{theo:gd_diffusion_full}, \ref{theo:db_stable_frozen_full} and \ref{theo:local_convergence_full} are statements about the gradient read at the invariant measure of the chain in play --- $\nabla^\lambda$ for the state-space loss, $\nabla^{\lambda_2}$ for its edge lift. For a Markov kernel or a bounded operator $\mathcal K$ on measures --- the instances below are $T^\lambda-r$ and $K_2^{\lambda_2}-\widehat r$ --- and a measure $m$, $\mathcal K[m]:=m\mathcal K$ denotes the density action of $\mathcal K$, not its function action on $\frac{dm}{d\lambda}$. The operators $P$, $P^\dagger$, $A$ and $A^\dagger$ act on densities rather than on measures, $[\,\cdot\,]$ marking application there as well. $\Pi$ denotes the mean projection onto the invariant functions of the kernel in play, $f\mapsto\frac{\int f\,d\rho}{\rho(\Statespace)}\mathbf 1$ against that kernel's invariant measure $\rho$. Total variation carries the $\tfrac12$-convention, $\TV(\mu\|\nu):=\tfrac12\int|\frac{d\mu}{d\rho}-\frac{d\nu}{d\rho}|d\rho$ for a common dominating $\rho$; hatted mixing coefficients $\widehat\beta_n:=\|P^n-\Pi\|_{L^2(\lambda)}$ refer to the lifted or backward chain under consideration. Two symbols resemble a measure and are not: $\lamreg$ is the flow-regularization weight of \citet{morozov2025revisiting} (Appendix~\ref{app:pbtrain}), and $\Lambda$ in Appendix~\ref{app:sgd} is a local smoothness constant.

Finiteness of $\|\foutstar\|_{L^1(\nuB)}$ is part of what a generative flow is; the higher $L^p$ norms are not, and most parameterizations of $\foutstar$ yield bounded $\foutstar$ on a compact domain so that they too are easily ensured. However, the definitions above do not guarantee that $\|\ForwardPolicyStar\|_{L^p(\nuB)}$ is finite; the parameterization of the star forward policy should ensure the finiteness of the operator norm of interest.

The following theorem  provides theoretical guarantees on the sampler associated with a generative flow.

    \begin{theorem}
        [\cite{bengio2021flow, brunswic2024theory}] \label{theo:sampling_theorem}
        Let $\Finit$, and $\Fterm$ be unnormalized distributions with $\Fterm\neq0$ and let $(\ForwardPolicyStar, \Foutstar)$ be a generative flow. If the flow-matching constraint
        \begin{equation}
            \Finit + \Foutstar\ForwardPolicyStar = \Fterm + \Foutstar
        \end{equation}
        is satisfied then $\Finit(\Statespace)=\Fterm(\Statespace)$ and:
        \begin{equation}
            \mathbb E(\tau)\leq \frac{\Foutstar(\Statespace)}{\Finit(\Statespace)}+1,\quad \quad \underline s_\tau \sim \Fterm. \nonumber
        \end{equation}

    \end{theorem}




A quantitative version controls the total variation of sampling from target distribution.
    The fundamental remark is that a generative flow $(\ForwardPolicyStar,\foutstar,\Finit,\Fterm)$ may be forced to be flow-matching by replacing $\Finit$ and $\Fterm$: denoting $\Finstar := \Foutstar\ForwardPolicyStar$ the star inflow, the flow $(\ForwardPolicyStar,\foutstar,\Finithat,\Ftermhat)$ is flow-matching where
    \begin{eqnarray}
        \Ftermhat &=& \Fterm +\delta\Fterm \\
        \Finithat &=& \Finit +\delta\Finit \\
        \delta\Fterm &=& \left(\Finit+\Finstar - \Foutstar-\Fterm\right)^+ \\
        \delta\Finit &=& \left(\Finit+\Finstar - \Foutstar-\Fterm\right)^-
    \end{eqnarray}
\begin{theorem}[\cite{brunswicEGF}]


\label{theo:negative_control}
  Let $(\ForwardPolicyStar, \foutstar)$ be a generative flow and let $\Finit$ and $\Fterm$ initial and terminal distributions. Assume that $\Finit\neq 0$, and consider the sample $s_\tau$ of the generative flow Markov chain from $\Finit$ to $\Ftermhat$ then:
    \begin{equation}
        \TV\left(s_\tau \left\| \frac{1}{\Ftermhat(\Statespace )} \Ftermhat\right.\right) \leq  \frac{\delta\Finit(\Statespace)}{\Ftermhat(\Statespace)}
    \end{equation}
    with $\TV$  the total variation.
\end{theorem}

\subsection{Universality}\label{app:universality}

\begin{theorem}[full form of Theorem \ref{theo:universality_L2}]
\label{theo:universality_L2_full}
     Let $(\Statespace,\nuB)$ be a measured Polish space with $\nuB$ finite and let $\ForwardPolicyStar$ be a Markov kernel on $\Statespace$ with finite $L^p(\nuB)\rightarrow L^p(\nuB)$ operator norm. If $\ForwardPolicyStar$ is ergodic on $(\Statespace,\nuB)$ with summable $L^p$-mixing coefficients, then the family $\Theta := \{\ForwardPolicyStar\}\times L^p_+(\Statespace,\nuB)$ is weakly $L^p$-universal.

     More precisely, let $\Finit,\Fterm\ll\nuB$ be distributions with densities in $L^p(\nuB)$, put $\theta := \fterm-\finit$ and let $S:=\sum_{n\geq0}\left(P_\star^n-\Pi\right)$, so that $\|S\theta\|_{L^p(\nuB)}\leq B_p\|\theta\|_{L^p(\nuB)}$. Then for every $\eta>0$ the outflow density $\foutstareta := \left(\eta - S\theta\right)^+$ belongs to $\Theta$ and satisfies
     \begin{equation}\label{equ:universality_quantitative}
         \|\deltafinit\|_{L^p(\nuB)}+\|\delta\fterm\|_{L^p(\nuB)} \;\leq\; 2\left(1+\|\ForwardPolicyStar\|_{L^p(\nuB)}\right)\left\|\left(S\theta-\eta\right)^+\right\|_{L^p(\nuB)} \xrightarrow[\eta\rightarrow+\infty]{} 0 .
     \end{equation}
     Furthermore, if $p=+\infty$ then $\Theta$ is strongly universal: taking $\eta := \|S\theta\|_{L^\infty(\nuB)}$ gives $\deltafinit = \delta\fterm = 0$, and the flow-matching constraint \eqref{equ:FM_const} holds exactly.
\end{theorem}

\begin{proof}
Fix distributions $\Finit,\Fterm\ll\nuB$ with densities in $L^p(\nuB)$ and write $\theta := \fterm-\finit\in L^p(\nuB)$. For a candidate outflow $\Foutstar = \foutstar\nuB$ the flow-matching defect is the signed measure
\begin{equation}
    D := \Finit + \Foutstar\ForwardPolicyStar-\Foutstar-\Fterm, \qquad \frac{dD}{d\nuB} = \left(P_\star-\mathrm{Id}\right)\foutstar-\theta,
\end{equation}
and $\delta\Fterm = D^+$, $\delta\Finit = D^-$ by definition. The positive and negative parts of $D$ having disjoint supports,
\begin{equation}\label{equ:defect_equivalence}
    \|D\|_{L^p(\nuB)}\;\leq\;\|\deltafinit\|_{L^p(\nuB)}+\|\delta\fterm\|_{L^p(\nuB)}\;\leq\;2\|D\|_{L^p(\nuB)},
\end{equation}
so weak $L^p$-universality amounts to showing that $\theta$ lies in the $L^p(\nuB)$-closure of the image of the non-negative cone of $L^p(\nuB)$ under $P_\star-\mathrm{Id}$. (The restriction to densities in $L^p(\nuB)$ is not cosmetic: the residuals differ from $\finit-\fterm$ by a term in the range of $\ForwardPolicyStar-\mathrm{Id}$, so for $\finit-\fterm\notin L^p(\nuB)$ the infimum is $+\infty$ for trivial reasons.) We exhibit an explicit approximating family.

\emph{Step 1: summable mixing inverts $\mathrm{Id}-P_\star$ on mean-zero densities.} Since $\sum_n\beta^{(p)}_n = B_p<+\infty$, the series $S:=\sum_{n\geq0}\left(P^n_\star-\Pi\right)$ converges absolutely in operator norm on $L^p(\nuB)$, with $\|S\|\leq B_p$. Using $\Pi P_\star = P_\star\Pi = \Pi$, its partial sums telescope,
\begin{equation}
    \sum_{n=0}^N\left(P^n_\star-\Pi\right)\left(\mathrm{Id}-P_\star\right) = \sum_{n=0}^N\left(P^n_\star-P^{n+1}_\star\right) = \mathrm{Id}-P^{N+1}_\star,
\end{equation}
and $\left\|P^{N+1}_\star-\Pi\right\| = \beta^{(p)}_{N+1}\rightarrow0$, so that
\begin{equation}\label{equ:poisson}
    S\left(\mathrm{Id}-P_\star\right) = \left(\mathrm{Id}-P_\star\right)S = \mathrm{Id}-\Pi .
\end{equation}
For $p=2$, and with $\ForwardPolicyStar$ replaced by the backward policy, this is Lemma \ref{lem:sigma_mixing}. Now $\Finit$ and $\Fterm$ are both probability distributions, so $\int\theta d\nuB = 0$, that is $\Pi\theta=0$, and \eqref{equ:poisson} solves the Poisson equation \emph{exactly}:
\begin{equation}\label{equ:poisson_solved}
    \left(\mathrm{Id}-P_\star\right)S\theta = \theta,\qquad \|S\theta\|_{L^p(\nuB)}\leq B_p\|\theta\|_{L^p(\nuB)} .
\end{equation}

\emph{Step 2: a $0$-flow restores non-negativity.} Equation \eqref{equ:poisson_solved} says that the \emph{signed} measure $-\left(S\theta\right)\nuB$ satisfies the flow-matching constraint \eqref{equ:FM_const} for the pair $(\Finit,\Fterm)$; it is not an admissible outflow, since $S\theta$ need not have a sign. Invariance supplies the missing freedom: $\nuB\ForwardPolicyStar = \nuB$ gives $\left(\mathrm{Id}-P_\star\right)\mathbf 1 = 0$ --- the constants are $0$-flows in the sense of \citet{brunswic2024theory} --- and $\mathbf 1\in L^p(\nuB)$ because $\nuB$ is finite. For $\eta>0$ set
\begin{equation}
    \nu^\eta_\infty := \left(\eta-S\theta\right)\nuB,\quad \foutstareta := \frac{d\left(\nu^\eta_\infty\right)^+}{d\nuB} = \left(\eta-S\theta\right)^+,\quad \psi_\eta := \frac{d\left(\nu^\eta_\infty\right)^-}{d\nuB} = \left(S\theta-\eta\right)^+,
\end{equation}
(the signed intermediate measure $\nu^\eta_\infty$ of this construction is neither a background nor a training measure; the letter is reused only here) so that $\foutstareta\geq0$ lies in $L^p(\nuB)$ and $\foutstareta\nuB-\nu^\eta_\infty = \psi_\eta\nuB$. By \eqref{equ:poisson_solved} and $\left(\mathrm{Id}-P_\star\right)\eta\mathbf 1 = 0$, the density of $\nu^\eta_\infty\ForwardPolicyStar-\nu^\eta_\infty$ is $-\left(\mathrm{Id}-P_\star\right)\left(\eta\mathbf 1-S\theta\right) = \theta$, that is
\begin{equation}
    \nu^\eta_\infty\ForwardPolicyStar-\nu^\eta_\infty = \Fterm-\Finit .
\end{equation}
Choosing $\Foutstar := \foutstareta\nuB = \nu^\eta_\infty+\psi_\eta\nuB$, the exact part cancels in the defect and only the truncation survives:
\begin{equation}\label{equ:defect_is_truncation}
    D = \left(\Finit-\Fterm\right)+\left(\nu^\eta_\infty\ForwardPolicyStar-\nu^\eta_\infty\right)+\left(\left(\psi_\eta\nuB\right)\ForwardPolicyStar-\psi_\eta\nuB\right) = \left(\psi_\eta\nuB\right)\ForwardPolicyStar-\psi_\eta\nuB .
\end{equation}
Hence $\|D\|_{L^p(\nuB)}\leq\left(1+\|\ForwardPolicyStar\|_{L^p(\nuB)}\right)\|\psi_\eta\|_{L^p(\nuB)}$, which together with \eqref{equ:defect_equivalence} is the quantitative bound \eqref{equ:universality_quantitative}.

\emph{Step 3: the truncation vanishes.} Let $p<+\infty$. Then $0\leq\psi_\eta\leq\left(S\theta\right)^+$ with $\left(S\theta\right)^+\in L^p(\nuB)$, and $\psi_\eta\rightarrow0$ pointwise as $\eta\rightarrow+\infty$; dominated convergence gives $\|\psi_\eta\|_{L^p(\nuB)}\rightarrow0$, so $\Theta$ is weakly $L^p$-universal. Let now $p=+\infty$. For $\eta\geq\|S\theta\|_{L^\infty(\nuB)}$ we have $\psi_\eta = 0$, hence $D=0$ by \eqref{equ:defect_is_truncation}: the infimum is attained and the flow-matching constraint \eqref{equ:FM_const} holds exactly. Since $\nuB$ is finite, $\foutstareta\in L^\infty(\nuB)\subset L^r(\nuB)$ for every $r\in[1,+\infty]$, so the infimum is attained in every $L^r(\nuB)$ and $\Theta$ is strongly universal.
\end{proof}

The $L^2$ strengthening of Theorem \ref{theo:universality_L2_full} used below is the following. Only Step 2 above is revisited: for an EGF the star forward policy is a finite mixture of pushforwards by diffeomorphisms, and the $L^2$ norm of $\left(\psi_\eta\nuB\right)\ForwardPolicyStar$ is then controlled by a change of variables rather than by the operator norm, which yields an explicit constant and a control of the two residuals separately.

\begin{theorem}\label{theo:universality_L2_body}
    Under the same hypotheses as in Theorem \ref{theo:universality_L2_full}, if in addition $\Statespace$ is a differential manifold, if $\nuB$ is finite and comes from a differential volume form and if $\ForwardPolicyStar$ is the policy of an EGF then
    \begin{equation}
        \inf_{\foutstar\in L^2(\nuB)} \left[\|\deltafinit\|_{L^2(\nuB)} + \|\delta\fterm\|_{L^2(\nuB)}\right]= 0
    \end{equation}
\end{theorem}

\begin{proof}
    Continuing where we left at the end of the proof of Theorem \ref{theo:universality_L2_full} --- writing $m$ for the EGF's mixture size, $\forwardalpha^i\leq1$ for its mixture weights, $\Phi_i$ for its diffeomorphisms and $\#$ for pushforward --- for any $\eta$ we have
    \begin{eqnarray}
        \delta\Finit &:=&  \left[(\foutstar \nuB- \nu_\infty^\eta)\ForwardPolicyStar  -  (\foutstar \nuB- \nu_\infty^\eta)\right]^-\\
        \delta\Fterm &:=&  \left[(\foutstar \nuB- \nu_\infty^\eta)\ForwardPolicyStar  -  (\foutstar \nuB- \nu_\infty^\eta)\right]^+
    \end{eqnarray}
    By abuse of notation, for any measure $\mu\ll \nuB$ we write $\|\mu\|_{L^2(\nuB)}$ as a shorthand for $\left\|\frac{d\mu}{d\nuB}\right\|_{L^2(\nuB)}$. Now, write $\psi:= \frac{d(\nu_{\infty}^\eta)^-}{d\nuB}$ so that $\foutstar\nuB - \nu_\infty^\eta = \psi \nuB$ by definition of $\foutstar$. Then
    \begin{eqnarray}
         \frac{1}{2}\left(\|\delta\Finit\|_{L^2(\nuB)} +\|\delta\Fterm\|_{L^2(\nuB)}\right) &\leq& \left\|(\psi \nuB)\ForwardPolicyStar\right\|_{L^2(\nuB)} + \| \psi\|_{L^2(\nuB)}  \\
         &\leq&  \sum_{i=1}^m\|\forwardalpha^i\|_{L^\infty}\left\|\Phi_i\#(\psi \nuB)\right\|_{L^2(\nuB)} + \| \psi\|_{L^2(\nuB)} \\
         &\leq & \left(1+\sum_{i=1}^m\left\| |J\Phi_i^{-1}|^{1/2}\right\|_{L^\infty}\right)\|\psi\|_{L^2(\nuB)} \\
    \end{eqnarray}
    Where we used the manifold $L^p$ bound for change of variable one may find in \citet{stern2010p} Theorem 1. The constant in the right-hand side only depends on $\ForwardPolicyStar$, also we have $\|\psi\|_{L^2(\nuB)} \xrightarrow{\eta\rightarrow +\infty} 0$, we may thus conclude that
    \begin{equation}
        \|\delta\Finit\|_{L^2(\nuB)}+\|\delta\Fterm\|_{L^2(\nuB)}  \xrightarrow{\eta\rightarrow +\infty} 0.
    \end{equation}
    The result follows.

\end{proof}

\subsection{Convergence Bounds for Stable Losses}\label{app:stable}

\begin{theorem}[full form of Theorem \ref{theo:RL_CV_bound}]
\label{theo:RL_CV_bound_full} Let $(\Statespace,\nuB)$ be a Polish space endowed with a non-negative finite background measure, let $\Finit\ll\nuB$ be a non-zero initial flow and let $\target\ll \nuB$ be a non-zero, possibly unnormalized target with known density with respect to $\nuB$. Let $q\in[1,+\infty]$ and let $\nutrain$ be a probability measure with $\nuB \ll \nutrain$ and $\dln\in L^\infty(\nutrain)$.

Write
\begin{equation}\label{equ:stable_constants}
    C_0:=\nuB(\Statespace)^{1-1/q}\left\|\dln\right\|^{1/q}_{L^\infty(\nutrain)},
    \qquad
    C:=\frac{2C_0}{\target(\Statespace)} .
\end{equation}
Let $(\ForwardPolicyStar,\foutstar)$ be a generative flow trained on the target $\Fterm=\target$ with loss $\LFMstableq$, but using $\Fterm=\widehat\target := E^+$ during inference, where $E:=\Finit+\Finstar-\Foutstar$. Write $t:=\target(\Statespace)$. Then:
\begin{enumerate}
    \item the loss controls the sampling error and the initial-mass defect together,
    \begin{equation}\label{equ:stable_bound}
      \TV\left(s_\tau \left\| \frac{\target}{t} \right.\right)\;\leq\;C\LFMstableq ,
    \end{equation}
    and
\begin{equation}\label{equ:mass_floor}
    \left|\Finit(\Statespace)-\target(\Statespace)\right|\;\leq\;C_0\LFMstableq ;
\end{equation}
    \item the infimum of the loss is $0$ at an initial flow carrying the target's total mass. If $\Theta$ is a $L^p$-universal family of generative flows with $q\leq p$, if $\finit$ and $\frac{d\target}{d\nuB}$ lie in $L^p(\nuB)$, if the masses match, $\target(\Statespace)=\Finit(\Statespace)$, and if $\nutrain\ll\nuB$ with $\frac{d\nutrain}{d\nuB}\in L^{q'}(\nuB)$, $q'=\frac{p}{p-q}$ with $q'=1$ when $p=+\infty$ --- only the domination being required when $q=p=+\infty$ --- then
\begin{equation}
    \inf_{(\ForwardPolicyStar,\foutstar)\in \Theta} \LFMstableq = 0 ;
\end{equation}
    \item the mass is the only obstruction over a trainable initial family. Call \emph{initial family} a set $\mathcal I$ of non-zero initial flows $\Finit\ll\nuB$ and write, closures being taken in $\mathbb R$,
\begin{equation}\label{equ:initial_masses}
\begin{aligned}
    Z(\mathcal I)&:=\left\{\Finit(\Statespace)\ :\ \Finit\in\mathcal I\right\}, \\
    Z^p(\mathcal I)&:=\left\{\Finit(\Statespace)\ :\ \Finit\in\mathcal I,\ \finit\in L^p(\nuB)\right\} .
\end{aligned}
\end{equation}
Under the hypotheses of \emph{(2)} on $\Theta$, on $\frac{d\target}{d\nuB}$ and on $\nutrain$,
\begin{equation}\label{equ:joint_infimum}
    t\in\overline{Z^p(\mathcal I)}\ \Longrightarrow\ \inf_{\mathcal I\times\Theta}\LFMstableq = 0\ \Longrightarrow\ t\in\overline{Z(\mathcal I)} ,
\end{equation}
the two conditions coinciding when every member of $\mathcal I$ has its density in $L^p(\nuB)$.
\end{enumerate}
\end{theorem}

The total mass of the target is unknown in practice, so the equality $\target(\Statespace)=\Finit(\Statespace)$ of \emph{(2)} cannot be imposed. Item \emph{(3)} replaces it by a condition on the initial family alone, and that condition is necessary as well as sufficient: by \eqref{equ:mass_floor}, $\Finit(\Statespace)$ converges to $\target(\Statespace)$ along any sequence driving the loss to $0$. The mass defect is then not a floor but a quantity the loss removes, which is the measure-theoretic form of the requirement that a GFlowNet learn the partition function $Z=\target(\Statespace)$. Item \emph{(2)} does not follow from \emph{(3)}: the masses $\left\{\target(\Statespace)+2^{-j}\right\}_{j\geq1}$ satisfy the left-hand condition of \eqref{equ:joint_infimum} with no member of mass $\target(\Statespace)$, and masses bounded away from $\target(\Statespace)$ leave the infimum positive.

\begin{proof}
With $t$ and $E$ as in the statement, and with $z:=\Finit(\Statespace)$, $k:=\frac{d\target}{d\nuB}$ and $e:=\frac{dE}{d\nuB}$,
\begin{equation}
    e \;=\; \finit+\finstar-\foutstar,
\end{equation}
so that $\widehat\target = E^+$ has density $e^+$ and, the flow being trained against $\Fterm=\target$ and $\nutrain$ being a probability measure, the loss \eqref{equ:FMloss} reads $\LFMstableq = \|e-k\|_{L^q(\nutrain)}$.

\emph{Step 0: at inference the terminal flow is its own correction.} Running the sampler with $\Fterm := \widehat\target$ gives
\begin{equation}\label{equ:inference_defect}
    \delta\Fterm  = 0,\qquad \delta\Finit  = E^-,\qquad \Ftermhat = \widehat\target+\delta\Fterm = \widehat\target .
\end{equation}
 Moreover $\Finstar(\Statespace) = \Foutstar\ForwardPolicyStar(\Statespace) = \Foutstar(\Statespace)$, so $E(\Statespace) = z$ and, $E^+$ and $E^-$ having disjoint supports,
\begin{equation}\label{equ:mass_identity}
    \widehat\target(\Statespace) = E^+(\Statespace) = \Finit(\Statespace)+E^-(\Statespace)\;\geq\;  \Finit(\Statespace).
\end{equation}

\emph{Step 1: both error sources are dominated by the $L^1(\nuB)$ flow-matching error.} Write $\Delta:=\|e-k\|_{L^1(\nuB)}$. As $k\geq0$, on $\{e<0\}$ we have $|e-k| = k-e\geq -e$, hence $e^-\leq|e-k|$ pointwise; and $x\mapsto x^+$ being $1$-Lipschitz with $k^+=k$, $|e^+-k|\leq|e-k|$ pointwise. Integrating against $\nuB$,
\begin{equation}\label{equ:two_dominations}
    \delta\Finit(\Statespace) = E^-(\Statespace)\;\leq\;\Delta,\qquad \left\|\widehat\target-\target\right\|_{L^1(\nuB)}\;\leq\;\Delta .
\end{equation}

\emph{Step 2: two total variations.} By the triangular inequality for total variation,
\begin{equation}
    \TV\left(s_\tau\left\|\frac{\target}{\target(\Statespace)}\right.\right)\;\leq\;\underbrace{\TV\left(s_\tau\left\|\frac{\widehat\target}{\widehat\target(\Statespace)}\right.\right)}_{(i)}+\underbrace{\TV\left(\left.\frac{\widehat\target}{\widehat\target(\Statespace)}\right\|\frac{\target}{\target(\Statespace)}\right)}_{(ii)} .
\end{equation}
Term $(i)$ is the sampling error of the corrected flow: $\Finit\neq0$, so Theorem \ref{theo:negative_control} with \eqref{equ:inference_defect}, \eqref{equ:mass_identity} and \eqref{equ:two_dominations} gives
\begin{equation}\label{equ:term_i}
    (i)\;\leq\;\frac{\delta\Finit(\Statespace)}{\Ftermhat(\Statespace)} = \frac{\delta\Finit(\Statespace)}{\Finit(\Statespace)+\delta\Finit(\Statespace)}\;\leq\;\min\left(1,\frac{\Delta}{\Finit(\Statespace)}\right).
\end{equation}
Term $(ii)$ only compares the normalizations of two non-negative finite measures. For $\alpha,\beta$ non-negative and finite, $b:=\beta(\Statespace)>0$ and any $a>0$, writing $\frac{\beta}{b}-\frac{\alpha}{a} = \frac{\beta-\alpha}{a}+\beta\left(\frac1b-\frac1a\right)$ and using $\|\beta\|_{L^1}=b$,
\begin{equation}\label{equ:normalization_lemma}
    \TV\left(\left.\frac{\alpha}{a}\right\|\frac{\beta}{b}\right) = \frac12\left\|\frac{\alpha}{a}-\frac{\beta}{b}\right\|_{L^1}\;\leq\;\frac{\|\alpha-\beta\|_{L^1}+|a-b|}{2a}\;\leq\;\frac{\|\alpha-\beta\|_{L^1}}{a},
\end{equation}
the second inequality under the further hypothesis $a=\alpha(\Statespace)$, by $|a-b|\leq\|\alpha-\beta\|_{L^1}$; so with $\alpha=\target$, $a=t$, $\beta=\widehat\target$ and \eqref{equ:two_dominations}, $(ii)\leq\frac{\Delta}{t}$. For term $(i)$, put $u:=\delta\Finit(\Statespace)=E^-(\Statespace)$, so that $u\leq\Delta$ by \eqref{equ:two_dominations}; as $z>0$, the map $x\mapsto\frac{x}{z+x}$ is non-decreasing on $[0,+\infty[$, so the middle expression of \eqref{equ:term_i} gives $(i)\leq\frac{u}{z+u}\leq\frac{\Delta}{z+\Delta}\leq\frac{\Delta}{z}$, whence
\begin{equation}\label{equ:tv_chain}
    \TV\left(s_\tau\left\|\frac{\target}{t}\right.\right)\;\leq\;\frac{\Delta}{z+\Delta}+\frac{\Delta}{t}\;\leq\;\left(\frac1z+\frac1t\right)\Delta .
\end{equation}
By \eqref{equ:mass_identity}, $\int(e-k)\,d\nuB = E(\Statespace)-t = z-t$, so that
\begin{equation}\label{equ:mass_defect}
    |z-t|\;=\;\left|\int(e-k)\,d\nuB\right|\;\leq\;\Delta ,
\end{equation}
hence $z+\Delta\geq z+|z-t|\geq\max(z,t)\geq t$ and, $\Delta$ being non-negative,
\begin{equation*}
    \TV\left(s_\tau\left\|\frac{\target}{t}\right.\right)\;\leq\;\frac{2\Delta}{t} .
\end{equation*}

\emph{Step 3: from $L^1(\nuB)$ to the training loss.} H\"older's inequality with the conjugate exponents $q$ and $q^*$ against $\nuB$ --- an equality at $q=1$, where $\nuB(\Statespace)^{1-1/q}=1$ --- and an $L^\infty(\nutrain)$ bound absorbing the change of measure,
\begin{eqnarray}
    \Delta &=& \int_{s\in\Statespace}|e-k|(s)d\nuB(s)\;\leq\;\nuB(\Statespace)^{1-1/q}\|e-k\|_{L^q(\nuB)},\\
    \|e-k\|^q_{L^q(\nuB)} &=& \int_{s\in\Statespace}|e-k|^q(s)\dln(s)d\nutrain(s)\;\leq\;\left\|\dln\right\|_{L^\infty(\nutrain)}\left(\LFMstableq\right)^q,
\end{eqnarray}
give $\Delta\leq\nuB(\Statespace)^{1-1/q}\left\|\dln\right\|^{1/q}_{L^\infty(\nutrain)}\LFMstableq = C_0\LFMstableq$. For $q=+\infty$ the same two lines read $\Delta\leq\nuB(\Statespace)\|e-k\|_{L^\infty(\nuB)}$ and $\|e-k\|_{L^\infty(\nuB)}\leq\|e-k\|_{L^\infty(\nutrain)}$, a $\nutrain$-null set being $\nuB$-null, which is again $\Delta\leq C_0\LFMstableq$. That bound with \eqref{equ:mass_defect} is \eqref{equ:mass_floor}.

\emph{Step 3b: the initial mass in the constant.} The map $x\mapsto\frac{x}{z+x}$ being non-decreasing on $[0,+\infty[$ as in Step 2, $\Delta\leq C_0\LFMstableq$ applied to the first member of \eqref{equ:tv_chain} gives
\begin{equation}
    \TV\left(s_\tau\left\|\frac{\target}{t}\right.\right)\;\leq\;\frac{\Delta}{z+\Delta}+\frac{\Delta}{t}\;\leq\;\frac{C_0\LFMstableq}{z+C_0\LFMstableq}+\frac{C_0\LFMstableq}{t} .
\end{equation}
By $z+\Delta\geq t$ of Step 2 and $\Delta\leq C_0\LFMstableq$, $z+C_0\LFMstableq\geq t$, so the right-hand side is at most $\frac{2C_0}{t}\LFMstableq$, which is $C\LFMstableq$ by \eqref{equ:stable_constants}. This is \eqref{equ:stable_bound}.

\emph{Step 4: the infimum of the loss at a matched initial mass.} During training the terminal flow is $\target$, so the defect is $E-\target$, of density $e-k$; the pair $(\Finit,\target)$ has equal total mass by hypothesis, so \eqref{equ:defect_equivalence} and the $L^p$-universality of $\Theta$ in the sense of Definition \ref{def:universality} give $\inf_\Theta\|e-k\|_{L^p(\nuB)} = 0$. For $q\leq p$ with $q<+\infty$, H\"older with the conjugate exponents $\frac pq$ and $q'$ transfers this to the training measure,
\begin{equation}
    \left(\LFMstableq\right)^q = \int_{s\in\Statespace}|e-k|^q(s)\frac{d\nutrain}{d\nuB}(s)d\nuB(s)\;\leq\;\|e-k\|^q_{L^p(\nuB)}\left\|\frac{d\nutrain}{d\nuB}\right\|_{L^{q'}(\nuB)},
\end{equation}
so that $\inf_\Theta\LFMstableq = 0$. For $q=p=+\infty$, $\nutrain\ll\nuB$ gives $\|e-k\|_{L^\infty(\nutrain)}\leq\|e-k\|_{L^\infty(\nuB)}$, a $\nuB$-null set being $\nutrain$-null, and the same conclusion.

\emph{Step 5: the infimum of the loss over an initial family.} Assume $t\in\overline{Z^p(\mathcal I)}$ and fix $\Finit\in\mathcal I$ with $\finit\in L^p(\nuB)$, of mass $z>0$. The distributions $\Finit$ and $\frac zt\target$ are dominated by $\nuB$, have equal total mass $z$ and have their densities in $L^p(\nuB)$, so Definition \ref{def:universality} applies to that pair; the defect of the flow for it is $E-\frac zt\target$, of density $e-\frac ztk$, so \eqref{equ:defect_equivalence} gives $\inf_\Theta\left\|e-\frac ztk\right\|_{L^p(\nuB)} = 0$. As $k\in L^p(\nuB)$, the triangle inequality gives
\begin{equation}
    \inf_\Theta\|e-k\|_{L^p(\nuB)}\;\leq\;\inf_\Theta\left\|e-\frac ztk\right\|_{L^p(\nuB)}+\frac{|z-t|}{t}\left\|k\right\|_{L^p(\nuB)}\;=\;\frac{|z-t|}{t}\left\|k\right\|_{L^p(\nuB)} ,
\end{equation}
and the infimum of $|z-t|$ over $z\in Z^p(\mathcal I)$ is $0$ by hypothesis, so $\inf_{\mathcal I\times\Theta}\|e-k\|_{L^p(\nuB)} = 0$. The H\"older transfer of Step 4 holds at each member of $\mathcal I\times\Theta$ with the same constant, whence $\inf_{\mathcal I\times\Theta}\LFMstableq = 0$. Conversely, let $\inf_{\mathcal I\times\Theta}\LFMstableq = 0$ and let $\left(\Finit^{(n)},\theta_n\right)$ be a minimizing sequence; \eqref{equ:mass_floor} gives $\left|\Finit^{(n)}(\Statespace)-t\right|\leq C_0\LFMstableq\left(\Finit^{(n)},\theta_n\right)\rightarrow0$, and $\Finit^{(n)}(\Statespace)\in Z(\mathcal I)$, so $t\in\overline{Z(\mathcal I)}$. This is \eqref{equ:joint_infimum}, and $Z(\mathcal I)=Z^p(\mathcal I)$ when every member of $\mathcal I$ has its density in $L^p(\nuB)$.
\end{proof}


The floor of the loss \eqref{equ:weakFM} in Theorem \ref{theo:IL_CV_bound} is the \emph{self-entropy} of
the target, taken against the background measure: for a probability $\target\ll\nuB$ of
density $k$,
\begin{equation}
    \mathcal H \;:=\; -\int_{s\in\Statespace}k(s)\log k(s)\,d\nuB(s) .
\end{equation}
This is not a Shannon entropy. Replacing $\nuB$ by $c\nuB$ shifts $\mathcal H$ by $\log c$,
and $\mathcal H<0$ whenever $k>1$ holds $\target$-almost everywhere: the floor is of either
sign.

\begin{theorem} \label{theo:IL_CV_bound} Let $(\Statespace,\nuB)$ be a measured Polish space with $\nuB$ finite and let $\Finit,\target$ be probability distributions absolutely continuous with respect to $\nuB$, with densities $\finit,k$ and with finite self-entropy $\mathcal H$. Let $q\in[1,+\infty]$ and assume $\nuB\sim\nutrain$ with $\dln\in L^{q^*}(\nutrain)$ and $b\geq1+\left\|\dln\right\|_{L^{q^*}(\nutrain)}$. For a generative flow $(\ForwardPolicyStar,\foutstar)$ whose star inflow satisfies $\Finstar\ll\nuB$, define $\LossWKLFM^q(\ForwardPolicyStar,\foutstar)$ as in Equation \eqref{equ:weakFM} and use $\Fterm=\widehat\target:=  \left(\Finit+\Finstar-\Foutstar\right)^+$ during inference. Then:
 \begin{itemize}
     \item for every Markov kernel $\ForwardPolicyStar$ on $\Statespace$ and every $\foutstar\in L^1_+(\Statespace,\nuB)$ with $\Finstar\ll\nuB$, $\LossWKLFM^q(\ForwardPolicyStar,\foutstar)\geq\mathcal H$ and
\begin{multline}\label{equ:il_tv_explicit}
   \TV\left(s_\tau \left\| \target \right.\right) \;\leq\; \frac{1}{\sqrt 2}\sqrt{\LossWKLFM^q-\mathcal H} \\ +\;\min\left(1,\left\|\dln\right\|_{L^{q^*}(\nutrain)}\left(\LossWKLFM^q-\mathcal H\right)\right) ;
\end{multline}
     \item let $\Theta$ be a family of generative flows on $(\Statespace,\nuB)$ whose outflows lie in $L^1_+(\Statespace,\nuB)$ and whose star inflows are dominated by $\nuB$. If $(\ForwardPolicyStar,\foutstar)\in\Theta$ and $\varepsilon\in[0,1)$ satisfy
\begin{equation}\label{equ:il_domination}
    \Finit+\Foutstar\ForwardPolicyStar-\Foutstar\;\geq\;(1-\varepsilon)\target\qquad\text{as measures,}
\end{equation}
then $\LossWKLFM^q(\ForwardPolicyStar,\foutstar)\leq\mathcal H+\log\frac{1}{1-\varepsilon}$. Consequently, if for every $\varepsilon>0$ some member of $\Theta$ satisfies \eqref{equ:il_domination}, then
\begin{equation}
    \inf_{(\ForwardPolicyStar,\foutstar)\in\Theta}\LossWKLFM^q(\ForwardPolicyStar,\foutstar) = \mathcal H ,
\end{equation}
and if some member satisfies \eqref{equ:il_domination} with $\varepsilon=0$ --- equivalently, if $\Theta$ contains a flow satisfying the flow-matching constraint \eqref{equ:FM_const} exactly for the pair $(\Finit,\target)$, in particular if $\finit,k\in L^p(\nuB)$ for some $p\in[1,+\infty]$ and $\Theta$ is strongly $L^p$-universal in the sense of Definition \ref{def:universality} --- the infimum is attained at that flow.
 \end{itemize}
\end{theorem}

\begin{proof}
    Let $\foutstar\in L^1_+(\Statespace,\nuB)$, put $E:=\Finit+\Finstar-\Foutstar$, so that $\widehat\target = E^+$, and write $\delta := \int_{s\in\Statespace }\deltafinit(s)d\nuB(s) = \delta\Finit(\Statespace)$. The density $\deltafinit$ is defined $\nuB$-a.e. and $\nutrain\ll\nuB$, so $\left\|\deltafinit\right\|_{L^{q}(\nutrain)}$ is well defined.

    Running the sampler with $\Fterm := \widehat\target$ gives $\delta\Fterm = 0$, $\delta\Finit = E^-$ and $\Ftermhat = \widehat\target$, as in \eqref{equ:inference_defect}. The measure $\Foutstar$ is finite, so $\Finstar(\Statespace) = \Foutstar(\Statespace)$ and $E(\Statespace) = \Finit(\Statespace) = 1$; the parts $E^+$ and $E^-$ having disjoint supports,
    \begin{eqnarray}
        \widehat\target(\Statespace)\;=\;E^+(\Statespace) &=& \Finit(\Statespace)+E^-(\Statespace) \nonumber \\
        &=&1+\delta \nonumber \\ &=& \int_{s\in \Statespace}\ftermhat(s)d\nuB(s) \label{equ:il_mass_identity}
    \end{eqnarray}
    and in particular $\widehat\target\neq0$ and $\widehat\target(\Statespace)\geq1$.

    From the triangular inequality for total variations, Theorem \ref{theo:negative_control} applied to the generative flow $(\ForwardPolicyStar,\foutstar)$ with $\Finit\neq0$, Pinsker's inequality, $\widehat \target = \Ftermhat$ and $\frac{\delta}{\widehat\target(\Statespace)} = \frac{\delta}{1+\delta}\leq\min(1,\delta)$ by \eqref{equ:il_mass_identity}:
    \begin{eqnarray}
        \TV\left(s_\tau \left\| \target \right.\right)  &\leq&  \TV\left(s_\tau \left\| \frac{\widehat\target}{\widehat\target(\Statespace)} \right.\right)  +  \TV\left(\target\left\| \frac{\widehat\target}{\widehat\target(\Statespace)} \right.\right)   \nonumber\\
        &\leq & \frac{\delta}{\widehat \target(\Statespace)} +  \frac{1}{\sqrt{2}}\sqrt{\KL\left(\target\left\|  \frac{\widehat\target}{\widehat\target(\Statespace)}\right. \right)} \nonumber\\
        &\leq & \min(1,\delta) +  \frac{1}{\sqrt{2}}\sqrt{\KL\left(\target\left\|  \frac{\widehat\target}{\widehat\target(\Statespace)}\right. \right)}. \label{equ:il_tv_split}
    \end{eqnarray}
    Furthermore $\nuB\ll\nutrain$, so $\dln$ is defined, and by H\"older's inequality with the conjugate exponents $q^*$ and $q$:
    \begin{eqnarray}
        \delta &=& \int_{s\in \Statespace }\deltafinit(s)\dln(s)d\nutrain(s)\nonumber\\
        &\leq& \left\|\dln\right\|_{L^{q^*}(\nutrain)}\left\|\deltafinit\right\|_{L^{q}(\nutrain)} \nonumber\\
        &\leq& (b-1)\left\|\deltafinit\right\|_{L^{q}(\nutrain)}. \label{equ:il_holder}
    \end{eqnarray}
    Decomposing the cross-entropy of the target against $\ftermhat$ as the forward Kullback--Leibler divergence plus self-entropy plus normalization ($-\mathbb E_{s\sim\target}\log\ftermhat = \KL\left(\target\left\|\frac{\widehat\target}{\widehat\target(\Statespace)}\right.\right) + \mathcal H - \log\widehat\target(\Statespace)$) and using $\widehat\target(\Statespace) = 1+\delta$ from \eqref{equ:il_mass_identity},
        \begin{eqnarray}
        \LossWKLFM^q  &=&\KL\left(\target\left\| \frac{\widehat\target}{\widehat\target(\Statespace)}\right. \right)  + b\left\|\deltafinit\right\|_{L^{q}(\nutrain)} \nonumber\\ && \quad - \log(1+\delta)+ \mathcal H\label{eq:wFM_strong_bound}
    \end{eqnarray}
 but since, by $\log(1+\delta)\leq\delta$ and \eqref{equ:il_holder},
 \begin{eqnarray}
 b\left\|\deltafinit\right\|_{L^{q}(\nutrain)} - \log(1+\delta) &\geq& b\left\|\deltafinit\right\|_{L^{q}(\nutrain)} -\delta \nonumber\\ &\geq& \left\|\deltafinit\right\|_{L^{q}(\nutrain)} \nonumber\\&\geq& 0   \label{eq:lp_delta}
 \end{eqnarray}

   we deduce that
    \begin{equation}
\LossWKLFM^q -\mathcal H \geq \KL\left(\target\left\| \frac{\widehat\target}{\widehat\target(\Statespace)}\right. \right) \geq 0
    \end{equation}
    so that
    $ \LossWKLFM^q \geq   \mathcal H $ and, by \eqref{equ:il_tv_split},

    \begin{eqnarray}
         \TV\left(s_\tau \left\| \target \right.\right) &\leq& \min(1,\delta) +  \frac{1}{\sqrt{2}}\sqrt{ \LossWKLFM^q  - \mathcal H}.
    \end{eqnarray}
   Furthermore, from equations \eqref{eq:wFM_strong_bound} and \eqref{eq:lp_delta}

   \begin{equation}
       \LossWKLFM^q  -\mathcal H \geq \left\|\deltafinit\right\|_{L^{q}(\nutrain)}
   \end{equation}
   So finally, by \eqref{equ:il_holder} and the monotonicity of $x\mapsto\min(1,x)$:
   \begin{multline}
         \TV\left(s_\tau \left\| \target \right.\right) \;\leq\; \frac{1}{\sqrt{2}}\sqrt{ \LossWKLFM^q  - \mathcal H} \\ +\;\min\left(1,\left\|\dln\right\|_{L^{q^*}(\nutrain)}\left(\LossWKLFM^q  - \mathcal H\right)\right),
    \end{multline}
    which is \eqref{equ:il_tv_explicit} and the first item, the coefficient $\left\|\dln\right\|_{L^{q^*}(\nutrain)}$ depending on neither $\ForwardPolicyStar$ nor $\foutstar$.

It remains to prove the second. Every member of $\Theta$ has its outflow in
$L^1_+(\Statespace,\nuB)$ and its star inflow dominated by $\nuB$, so the first item applies
to each of them and $\inf_\Theta\LossWKLFM^q\geq\mathcal H$. Let
$(\ForwardPolicyStar,\foutstar)\in\Theta$ and $\varepsilon\in[0,1)$ satisfy
\eqref{equ:il_domination} and write $e := \finit+\finstar-\foutstar$ for the density of
$E = \Finit+\Finstar-\Foutstar$: then $e\geq(1-\varepsilon)k\geq0$ $\nuB$-almost everywhere,
so $\deltafinit = e^- = 0$, $\ftermhat = e^+ = e$, and the loss of that flow is
\begin{eqnarray}
    \LossWKLFM^q &=& -\int_{s\in\Statespace} k(s)\log e(s)\,d\nuB(s) \nonumber\\
    &\leq& -\int_{s\in\Statespace} k(s)\log\left((1-\varepsilon)k(s)\right)d\nuB(s) \nonumber\\
    &=& \mathcal H+\log\frac{1}{1-\varepsilon} ,
\end{eqnarray}
the inequality by monotonicity of the logarithm on $\{k>0\}$ and both integrands vanishing
on $\{k=0\}$, the last equality by $\int k\,d\nuB = 1$. The term
$b\left\|\deltafinit\right\|_{L^{q}(\nutrain)}$ vanishes there for every $q$, every
admissible $b$ and every $\nutrain$. If for every $\varepsilon>0$ some member of $\Theta$
satisfies \eqref{equ:il_domination}, then
$\mathcal H\leq\inf_\Theta\LossWKLFM^q\leq\mathcal H+\log\frac{1}{1-\varepsilon}$ for every
$\varepsilon>0$, so $\inf_\Theta\LossWKLFM^q = \mathcal H$. At $\varepsilon = 0$ the measure
$E-\target$ is non-negative and
$E(\Statespace) = \Finit(\Statespace) = 1 = \target(\Statespace)$, so $E = \target$, which is
the flow-matching constraint \eqref{equ:FM_const} for the pair $(\Finit,\target)$;
conversely that constraint gives $E = \target$, hence \eqref{equ:il_domination} with
$\varepsilon = 0$. The display above gives $\LossWKLFM^q\leq\mathcal H$ there, hence
$\LossWKLFM^q = \mathcal H$ by the first item, and the infimum is attained at that flow.
Such a member exists when $\finit,k\in L^p(\nuB)$ and $\Theta$ is strongly
$L^p$-universal: $\Finit$ and $\target$ are distributions dominated by $\nuB$, of equal
total mass and with densities in $L^p(\nuB)$, so Definition \ref{def:universality} applies
to the pair $(\Finit,\target)$ and the element of $\Theta$ realizing the infimum has
$\deltafinit = \delta\fterm = 0$, which by \eqref{equ:defect_equivalence} is
\eqref{equ:FM_const} for that pair.
\end{proof}

\subsection{Convergence Bounds of Divergence Losses}\label{app:divergence}

\emph{The doubling graph of \eqref{eq:counterexample_graph}.} With doubling probability
$c/(j+1)$, the thresholds sit at $c=1$ for finite expected backward length and $c=1/\ln 2$ for
recurrence (Proposition \ref{prop:doubling_phase}), so there is a window in which the condition
of \citet{morozov2025revisiting} fails although the chain remains recurrent.

For a strictly unimodal generator (Theorem \ref{theo:global_dichotomy_full}\emph{(1)}), at fixed $\nu$ the gradient flow preserves $\|\frac{d\mu}{d\lambda}\|_{L^2(\lambda)}$ while the total mass increases strictly off balance: \emph{the inflation is the algorithm} --- the circulation-inflating force is the mass component of the gradient, extinguished exactly at flow-matching, so no stalemate with the constraint can form. Convex or real-analytic generators exclude the freezing of Theorem \ref{theo:global_dichotomy_full}\emph{(2)} altogether (Remark \ref{rem:freezing}). Remark \ref{rem:cycle_no_stalemate} works out a five-vertex graph with a cycle and an over-inflated initialization: its expected backward length is finite, which by Proposition 3.12 of \citet{morozov2025revisiting} bounds the circulation of the balanced flow; the mass identity prevents the stalemate; and the transient cost of a distant initialization is, numerically, quadratic in the overload.

Under the same frozen backward policy, the Hessian at balance of the trajectory-balance loss on trajectories of $\ell$ transitions is positive semi-definite and at most $\ell(1+2\widehat B')$ times that of the detailed-balance loss, $\widehat B'$ the edge-lifted mixing sum without its $n=0$ term (Proposition \ref{prop:tb_hessian}): a step cap derived from the detailed-balance Hessian remains sufficient for trajectory balance once divided by that factor, and no trajectory-balance rate is claimed. The detailed-balance rates rest on a coercivity inequality that summable mixing, the hypothesis of Theorem \ref{theo:universality_L2} at $p=2$, supplies, so \emph{the same ergodicity of the backward policy provides both universality and the training rate}, quantifying the observation of \citet{morozov2025revisiting}, whose variables Theorem \ref{theo:training_speed} restores in place of the mixing constant. Remark \ref{rem:visit_ratio} then resolves the speed by state: heuristically, far from balance an out-of-equilibrium state is corrected at a rate proportional to its visit ratio, the mixing constant taking over near balance.

\begin{lemma}\label{lem:cycle_counterexample}
Let $N\geq2$ and let $\mathcal C_N$ be the graph with states $x_1,\dots,x_N$, cycle edges $x_i\rightarrow x_{i+1 \bmod N}$, an initial edge $s_0\rightarrow x_1$ and a terminal edge $x_i\rightarrow s_f$ for every $i$. Let $\widehat F$ be the unit-mass flow routed $s_0\rightarrow x_1\rightarrow x_2\rightarrow s_f$, let $\gamma$ be the unit circulation on the cycle, equal to $1$ on each edge $x_i\rightarrow x_{i+1\bmod N}$ and to $0$ elsewhere, and set $F_k:=\widehat F+k\gamma$ for every integer $k\geq0$. Then:
\begin{enumerate}
    \item every $F_k$ satisfies the flow-matching constraint \eqref{equ:FM_const} with terminal flow $\delta_{x_2}$, and its sampler satisfies $s_\tau(F_k)\sim\delta_{x_2}$; hence $\TV\bigl(s_\tau(F_k)\,\|\,\target\bigr)=1-\target(x_2)$ for every $k$ and every probability measure $\target$ on $\{x_1,\dots,x_N\}$;
    \item let $\target$ be a probability measure and $\nutrain$ a finite measure on $\{x_1,\dots,x_N\}$, and let $\mathcal L(F):=\int g(\rho_F)\,d\nutrain$ be a divergence-based FM loss in the sense of \citet[\S4.1]{brunswic2024theory}: $\rho_F:=(\finit+f_\leftarrow)/(\target+f_\rightarrow)$ is the ratio of the inflow side to the outflow side of \eqref{equ:FM_const}, $\finit$ being the initial flow and $f_\leftarrow$, $f_\rightarrow$ the flows of $F$ into and out of each state along the cycle edges, and $g:\mathbb R_+\rightarrow\mathbb R_+$ vanishes exactly at $1$. If $g$ is continuous at $1$, then $\mathcal L(F_k)$ is defined for $k\geq1$ and tends to $0$ as $k\rightarrow\infty$.
\end{enumerate}
\end{lemma}

\begin{proof}
\emph{(1)} The flow $\widehat F$ carries mass $1$ on the edges $s_0\rightarrow x_1$, $x_1\rightarrow x_2$ and $x_2\rightarrow s_f$ and nothing elsewhere: at $x_1$ its initial flow $1$ equals its outflow $1$ to $x_2$, at $x_2$ its inflow $1$ from $x_1$ equals its terminal flow $1$, and every other state carries nothing, so $\widehat F$ satisfies \eqref{equ:FM_const} with terminal flow $\delta_{x_2}$. The circulation $\gamma$ brings mass $1$ into each $x_i$, takes mass $1$ out of it, and carries no initial or terminal flow, so it is a $0$-flow in the sense of \citet{brunswic2024theory}, and $F_k=\widehat F+k\gamma$ satisfies \eqref{equ:FM_const} with the same initial and terminal flows. By Theorem \ref{theo:sampling_theorem}, $s_\tau(F_k)\sim\delta_{x_2}$, whence, $\target$ being a probability,
\begin{equation*}
    \TV\bigl(s_\tau(F_k)\,\|\,\target\bigr)=\tfrac12\Bigl(1-\target(x_2)+\sum_{i\neq2}\target(x_i)\Bigr)=1-\target(x_2).
\end{equation*}

\emph{(2)} Write $\widehat f_\leftarrow$, $\widehat f_\rightarrow$ for the flows of $\widehat F$ into and out of each state along the cycle edges, and $a_i:=(\finit+\widehat f_\leftarrow)(x_i)\geq0$, $b_i:=(\target+\widehat f_\rightarrow)(x_i)\geq0$, which do not depend on $k$. The circulation adds $k$ to both the inflow and the outflow of every $x_i$, so for $k\geq1$
\begin{equation*}
    \rho_{F_k}(x_i)=\frac{a_i+k}{b_i+k},\qquad |\rho_{F_k}(x_i)-1|=\frac{|a_i-b_i|}{b_i+k}\leq\frac{\max_j|a_j-b_j|}{k},
\end{equation*}
and $\max_i|\rho_{F_k}(x_i)-1|\rightarrow0$ as $k\rightarrow\infty$. As $g$ is continuous at $1$ with $g(1)=0$, $\max_ig(\rho_{F_k}(x_i))\rightarrow0$, and $0\leq\mathcal L(F_k)\leq\nutrain(\{x_1,\dots,x_N\})\max_ig(\rho_{F_k}(x_i))\rightarrow0$, $\nutrain$ being finite.
\end{proof}

\begin{proof}[Proof of Theorem \ref{theo:no_bound_divergence}]
Fix $N\geq2$, $\varepsilon>0$ and a training distribution $\nutrain$; the supremum over graphs of size $N$ ranges over the data $(\mathcal G,\target)$ of the loss. Take $\mathcal G:=\mathcal C_N$, which has $\#\mathcal G=N$, $\#\mathcal G$ counting the states other than $s_0$ and $s_f$, and the flows $F_k$ of Lemma \ref{lem:cycle_counterexample}, and, for $\delta\in(0,1)$, let $\target(x_1)=1-\delta$ and $\target(x_i)=\delta/(N-1)$ otherwise. By that lemma, $\TV\bigl(s_\tau(F_k)\,\|\,\target\bigr)=1-\target(x_2)=1-\frac{\delta}{N-1}$ for every $k$, and $\mathcal L(F_k)\leq\varepsilon$ for $k$ large.

Therefore $\sup_{\#\mathcal G=N}\sup_{\mathcal L\leq\varepsilon}\TV \geq 1-\frac{\delta}{N-1}$ for every $\delta\in(0,1)$, so it equals $1$; the construction being independent of $\nutrain$, the infimum over $\nutrain$ is also $1$. Theorem 3 of \citet{brunswic2024theory} establishes the instability of the DB and TB losses through the same scale invariance along $0$-flows; the total-variation statement is proved here for FM only.
\end{proof}

\begin{lemma}[$L^2(\lambda)$-conjugation and the $\lambda$-reversal]\label{lem:adjoint}
Let $\Statespace$ be Polish, let $\nuB \in \mathcal M^+(\Statespace)$, let $T$ be a Markov kernel on $\Statespace$, and let $\lambda\in \mathcal M^+(\Statespace,\nuB)$ be non-zero, finite and $T$-invariant. Then:
\begin{enumerate}
    \item the measure $T\otimes\lambda$ on $\Statespace^2$ disintegrates over its first marginal --- which is $\lambda T=\lambda$ --- into a $\lambda$-almost everywhere unique Markov kernel $T^\lambda$, the \emph{$\lambda$-reversal} of $T$, characterized by $\lambda\otimes T^\lambda= T\otimes \lambda$;
    \item $\lambda$ is $T^\lambda$-invariant and $T^\lambda\otimes\lambda = \lambda\otimes T$, so that $T$ is in turn the $\lambda$-reversal of $T^\lambda$: reversal is an involution, $\lambda$-almost everywhere, on the Markov kernels leaving $\lambda$ invariant;
    \item the density actions $\mu\mapsto \mu T$ and $\mu\mapsto \mu T^\lambda$ are mutually adjoint contractions of $\mathcal M^2(\lambda)$, each of operator norm $1$: $\langle \mu T\mid \nu\rangle_\lambda = \langle \mu \mid \nu T^\lambda\rangle_\lambda$ for all $\mu,\nu\in\mathcal M^2(\lambda)$; under $\mu\mapsto\frac{d\mu}{d\lambda}$ the density action of each of $T$ and $T^\lambda$ on $\mathcal M^2(\lambda)$ is the function action of the other on $L^2(\lambda)$;
    \item when $T=\BackwardPolicy$ is a backward policy --- a Markov kernel on $\Statespace$ leaving $\lambda$ invariant, which on a finite marked graph is the loop closure $\widehat\BackwardPolicy$ of Definition \ref{def:loop_closure} and not $\BackwardPolicy$ itself --- its reversal, written $\ForwardPolicy^\lambda:=\BackwardPolicy^\lambda$, is the forward policy that the \emph{balanced} flow of inflow $\lambda$ pairs with $\BackwardPolicy$, and $(\ForwardPolicy^\lambda)^\lambda=\BackwardPolicy$: $L^2(\lambda)$-conjugation carries the backward policy to that forward policy and the forward policy back to the backward one, which is the forward/backward duality of generative flows (Proposition 1 in Appendix B of \citet{brunswic2024theory}). On a finite marked graph with $\lambda>0$ everywhere, $\ForwardPolicy^\lambda(s\rightarrow s')=\lambda(s')\widehat\BackwardPolicy(s'\rightarrow s)/\lambda(s)$ is supported on the edges of the loop closure $\widehat G$, wrap edge included, so $\ForwardPolicy^\lambda(s_f\rightarrow s_0)=1$ and $s_f$ is not terminal for it.
\end{enumerate}
\end{lemma}

\begin{proof} Throughout, $(\lambda\otimes A)(dx\,dy)=\lambda(dx)A(x\rightarrow dy)$ and $(A\otimes\lambda)(dx\,dy)=A(y\rightarrow dx)\lambda(dy)$ for a Markov kernel $A$, so that $A\otimes\lambda=\mathrm{fl}_*(\lambda\otimes A)$, where $\mathrm{fl}(x,y):=(y,x)$ is the coordinate flip of $\Statespace^2$ and $\mathrm{fl}_*$ the image measure it induces.

\emph{(1)} The first marginal of $T\otimes\lambda$ is $X\mapsto \int_{y\in\Statespace}T(y\rightarrow X)\lambda(dy)=\lambda T(X)=\lambda(X)$ by $T$-invariance, and its second marginal is $\lambda$. Since $\Statespace$ is Polish and $\lambda$ is finite, $T\otimes\lambda$ is a finite measure on a standard Borel space, so it disintegrates over its first marginal into a Markov kernel, unique up to a $\lambda$-negligible set \cite{halmos2013measureHaar}; that disintegration is the identity $\lambda\otimes T^\lambda = T\otimes\lambda$.

\emph{(2)} $\lambda T^\lambda$ is the second marginal of $\lambda\otimes T^\lambda=T\otimes\lambda$, which is $\lambda$. Applying the flip identity to $T^\lambda$ and then to $T$,
\begin{equation}\label{eq:reversal_flip}
    T^\lambda\otimes\lambda \;=\; \mathrm{fl}_*(\lambda\otimes T^\lambda) \;=\; \mathrm{fl}_*(T\otimes\lambda) \;=\; \lambda\otimes T,
\end{equation}
so $\lambda\otimes T$ is a disintegration of $T^\lambda\otimes\lambda$ over its first marginal $\lambda$; the disintegration of a finite measure on $\Statespace^2$ over its first marginal being unique up to a $\lambda$-negligible set, $T$ is the $\lambda$-reversal of $T^\lambda$.

\emph{(3)} Let $\mu\in\mathcal M^2(\lambda)$ and let $X\subset\Statespace$ be measurable with $\lambda(X)=0$. Then $\int T(y\rightarrow X)\lambda(dy)=\lambda T(X)=0$, so $T(y\rightarrow X)=0$ for $\lambda$-almost every $y$, and $\mu T(X)=\int T(y\rightarrow X)\mu(dy)=0$ since $\mu\ll\lambda$: hence $\mu T\ll\lambda$, and likewise $\mu T^\lambda\ll\lambda$ using $\lambda T^\lambda=\lambda$ from \emph{(2)}. Write $u:=\frac{d\mu}{d\lambda}$ and $v:=\frac{d\nu}{d\lambda}$ for $\mu,\nu\in\mathcal M^2(\lambda)$, and $Af$ for the function action $x\mapsto\int f(y)A(x\rightarrow dy)$ of a Markov kernel $A$. By Jensen's inequality applied to the probability measures $T^\lambda(x\rightarrow\cdot)$, and then by $\lambda T^\lambda=\lambda$,
\begin{eqnarray}
    \int (T^\lambda u)^2\,d\lambda &\leq& \int T^\lambda(u^2)\,d\lambda\nonumber\\
    &=& \int u^2\,d(\lambda T^\lambda) = \|u\|_{L^2(\lambda)}^2, \label{eq:reversal_contraction}
\end{eqnarray}
and the same bound holds for $T$ by $\lambda T=\lambda$. Consequently
\begin{eqnarray}
    \iint |v(x)|\,|u(y)|\,(\lambda\otimes T^\lambda)(dx\,dy) &=& \int |v|\cdot T^\lambda|u|\,d\lambda\nonumber\\
    &\leq& \|v\|_{L^2(\lambda)}\|u\|_{L^2(\lambda)}
\end{eqnarray}
by Cauchy--Schwarz and \eqref{eq:reversal_contraction}: Fubini's theorem therefore applies to the signed integrand $(x,y)\mapsto v(x)u(y)$ in
\begin{eqnarray}
    \langle \mu T\mid \nu\rangle_\lambda &=& \int v\,d(\mu T)\nonumber\\
    &=& \int\left[\int v(x)\,T(y\rightarrow dx)\right]u(y)\,\lambda(dy)\nonumber\\
    &=& \iint v(x)u(y)\,(T\otimes\lambda)(dx\,dy)\nonumber\\
    &=& \iint v(x)u(y)\,(\lambda\otimes T^\lambda)(dx\,dy)\nonumber\\
    &=& \int v(x)\left[\int u(y)\,T^\lambda(x\rightarrow dy)\right]\lambda(dx)\nonumber\\
    &=& \int v\cdot T^\lambda u\,d\lambda\nonumber\\
    &=& \int u\,d(\nu T^\lambda) \;=\; \langle \mu\mid \nu T^\lambda\rangle_\lambda,
\end{eqnarray}
the last two equalities by the definition of the image measure $\nu T^\lambda$ and by $\nu T^\lambda\ll\lambda$. The map $\mu\mapsto\frac{d\mu}{d\lambda}$ is an isometry of $\mathcal M^2(\lambda)$ into $L^2(\lambda)$ carrying $\langle\cdot\mid\cdot\rangle_\lambda$ to the $L^2(\lambda)$ inner product, which is non-degenerate, so $\mu\mapsto\mu T^\lambda$ is the only map satisfying this identity: it is the adjoint of $\mu\mapsto\mu T$ \cite{brezis2011functional}. Taking $\nu$ arbitrary in $\langle\mu T\mid\nu\rangle_\lambda=\int v\cdot T^\lambda u\,d\lambda$ gives $\frac{d(\mu T)}{d\lambda}=T^\lambda u$ in $L^2(\lambda)$; exchanging the roles of $T$ and $T^\lambda$, legitimate by \eqref{eq:reversal_flip}, gives $\frac{d(\mu T^\lambda)}{d\lambda}=Tu$. With \eqref{eq:reversal_contraction} both density actions are therefore contractions of $\mathcal M^2(\lambda)$, and $\lambda\in\mathcal M^2(\lambda)$ --- non-zero since $\lambda\neq0$, of finite norm since $\lambda$ is finite --- satisfies $\lambda T=\lambda T^\lambda=\lambda$, so both operator norms equal $1$.

\emph{(4)} At $T=\BackwardPolicy$ the characterization of \emph{(1)} reads $\lambda\otimes\ForwardPolicy^\lambda=\BackwardPolicy\otimes\lambda$, which is the balance constraint of the pair $(\lambda,\ForwardPolicy^\lambda)$ with $\BackwardPolicy$ frozen, and \emph{(1)} states that $\ForwardPolicy^\lambda$ is the only Markov kernel satisfying it up to a $\lambda$-negligible set; $(\ForwardPolicy^\lambda)^\lambda=\BackwardPolicy$ is \emph{(2)}. On a finite marked graph the loop closure $\widehat\BackwardPolicy$ is a Markov kernel on $\mathcal V$ (Definition \ref{def:loop_closure}), and where $\lambda>0$ the disintegration of \emph{(1)} is $\ForwardPolicy^\lambda(s\rightarrow s')=\lambda(s')\widehat\BackwardPolicy(s'\rightarrow s)/\lambda(s)$: it is non-negative, sums to $1$ over $s'$ by $\lambda\widehat\BackwardPolicy=\lambda$, and satisfies $\lambda(s)\ForwardPolicy^\lambda(s\rightarrow s')=\widehat\BackwardPolicy(s'\rightarrow s)\lambda(s')$, which is the characterizing identity read at $(s,s')$. It is positive only where $\widehat\BackwardPolicy(s'\rightarrow s)>0$, hence only on edges $s\rightarrow s'$ of $\widehat G$; the wrap edge being the only edge out of $s_f$ in $\widehat G$, the probability measure $\ForwardPolicy^\lambda(s_f\rightarrow\cdot)$ is supported on $\{s_0\}$.
\end{proof}

\begin{theorem}\label{theo:first_variation_full}
 Under the hypotheses of Lemma \ref{lem:adjoint}, let $\mu\in \mathcal M^+(\Statespace,\nuB)$ with $\mu\sim\lambda$, $\frac{d\mu}{d\lambda}\in L^2(\lambda)$, $\mu T\ll\mu$, and let $r:=\frac{d(\mu T)}{d\mu}$ be essentially bounded away from $0$ and $\infty$. For any finite measure $\nu$ with $\frac{d\nu}{d\mu}\in L^\infty(\mu)$, define the associated $g$-divergence balance loss $\mathcal L_{g,\nu}(\mu) := \int_{\Statespace} g(r)\, d\nu$. Then, for any $g:\mathbb R_+^*\rightarrow \mathbb R$ continuously differentiable with locally Lipschitz derivative, and any finite $\nuG\in \mathcal M^+(\Statespace,\nuB)$ with $\nuG\sim\lambda$, the derivative of $\mathcal L_{g,\nu}$ --- taken along directions $\delta\ll\mu$ with $\|\frac{d\delta}{d\mu}\|_{L^\infty(\mu)}<1$ and $\frac{d(\delta T)}{d\mu}\in L^\infty(\mu)$ --- is represented by
 \begin{equation}
      \nabla_{\mu}^{\nuG} \mathcal L_{g,\nu} = \frac{d\nuG}{d\lambda}\cdot\left(T^\lambda-r\right)\left[g'\left(r\right)\frac{d\nu}{d\mu}\lambda \right].
     \end{equation}

\end{theorem}

\begin{proof}
 To begin with,
     \begin{eqnarray}
         \mathcal L_{g,\nu} = \int_{x\in \Statespace} g\left(\frac{d(\mu T)}{d\mu}\right) d\nu(x)
     \end{eqnarray}

     Let $r(x,\mu) := \frac{d(\mu T)}{d\mu}(x)$ and, for an admissible direction $\delta$, write $u := \frac{d\delta}{d\mu}$ and $v := \frac{d(\delta T)}{d\mu}$. Then
\[
r(x,\mu+\delta) \;=\; \frac{r+v}{1+u} \;=\; r \;+\; v\;-\; r\,u+\Delta(x,\delta),
\qquad \Delta = \frac{u\,(ru-v)}{1+u},
\]
so for $\|u\|_{L^\infty(\mu)}\leq\frac12$, $|\Delta|\leq2|u|\,(r|u|+|v|)$: the remainder is second order in $(u,v)$, uniformly; $r$ being essentially bounded and $g'$ locally Lipschitz on the essential range of $r$, the corresponding remainder of $\mathcal L_{g,\nu}$ is second order in $(u,v)$, hence negligible against the linear functional as $\delta\rightarrow0$ within the stated class. Since $\mu\sim\lambda$ and $u\in L^\infty(\mu)$, $\frac{d\delta}{d\lambda}=u\frac{d\mu}{d\lambda}$ with $\frac{d\mu}{d\lambda}\in L^2(\lambda)$ by hypothesis, so $\delta\in\mathcal M^2(\lambda)$; and $g'(r)\frac{d\nu}{d\mu}$ is bounded $\lambda$-almost everywhere --- $r$ is essentially valued in a compact subset of $(0,\infty)$, on which $g'$ is continuous, and $\frac{d\nu}{d\mu}\in L^\infty(\mu)=L^\infty(\lambda)$ --- so $g'(r)\frac{d\nu}{d\mu}\lambda\in\mathcal M^2(\lambda)$, $\lambda$ being finite. Lemma \ref{lem:adjoint}\emph{(3)} applies to that pair and, by the symmetry of $\langle\cdot\mid\cdot\rangle_\lambda$, gives $\left\langle g'(r)\frac{d\nu}{d\mu}\lambda \mid \delta T\right\rangle_\lambda = \left\langle \left[g'(r)\frac{d\nu}{d\mu}\lambda\right]T^\lambda \mid \delta\right\rangle_\lambda$. Hence the first variation of $\mathcal L_{g,\nu}$ is
\begin{eqnarray}
\delta \mathcal L_{g,\nu}
&=& \int_{\Statespace} g'(r(x))\,\delta r(x)\, d\nu(x)\nonumber\\
&=& \int_{\Statespace} g'(r)\!\left[\frac{d\,(\delta T)}{d\mu}- r\frac{d\,\delta}{d\mu}\right] d\nu\nonumber\\
&=&\left\langle g'(r)\frac{d\nu}{d\mu}\lambda \mid \delta T\right\rangle_\lambda \nonumber\\
&&\quad \quad - \left\langle rg'(r)\frac{d\nu}{d\mu}\lambda \mid \delta\right\rangle_\lambda\nonumber\\
&=& \left\langle \left[g'(r)\frac{d\nu}{d\mu}\lambda\right]T^\lambda \mid \delta\right\rangle_\lambda \nonumber\\
&&\quad \quad - \left\langle r g'(r)\frac{d\nu}{d\mu}\lambda \mid \delta\right\rangle_\lambda\nonumber\\
&=&\left\langle \left[g'(r)\frac{d\nu}{d\mu}\lambda\right](T^\lambda-r) \mid \delta\right\rangle_\lambda.
\end{eqnarray}

Write $\psi:=g'(r)\frac{d\nu}{d\mu}$ and $\Phi:=\left[\psi\lambda\right](T^\lambda-r)$, so that $\delta\mathcal L_{g,\nu}=\langle\Phi\mid\delta\rangle_\lambda$, and $\varpi:=\frac{d\nuG}{d\lambda}$, which exists since $\nuG\ll\lambda$ and satisfies $\frac{d\lambda}{d\nuG}=\frac1\varpi$ since $\lambda\ll\nuG$. By Lemma \ref{lem:adjoint}\emph{(3)}, $\frac{d\Phi}{d\lambda}=T\psi-r\psi$, which is bounded $\lambda$-almost everywhere, $T$ being Markov, $\psi$ bounded $\lambda$-almost everywhere and $r$ essentially bounded. Let $G:=\varpi\,\frac{d\Phi}{d\lambda}\,\lambda$, so that $\frac{dG}{d\nuG}=\frac{d\Phi}{d\lambda}$ and, $\nuG$ being finite, $G\in\mathcal M^2(\nuG)$. For an admissible $\delta$, $\frac{d\delta}{d\nuG}=\frac1\varpi\frac{d\delta}{d\lambda}$ and
\begin{equation}
    \int\frac{dG}{d\nuG}\frac{d\delta}{d\nuG}\,d\nuG \;=\; \int\frac{d\Phi}{d\lambda}\frac{d\delta}{d\lambda}\,d\lambda \;=\; \delta\mathcal L_{g,\nu},
\end{equation}
both integrals converging absolutely, $\frac{d\Phi}{d\lambda}$ being bounded and $\frac{d\delta}{d\lambda}$ lying in $L^2(\lambda)\subset L^1(\lambda)$, $\lambda$ being finite. Hence $G=\frac{d\nuG}{d\lambda}\Phi$ represents the derivative of $\mathcal L_{g,\nu}$ in $\langle\cdot\mid\cdot\rangle_{\nuG}$.
\end{proof}

\begin{corollary}[FM gradients]\label{cor:gradient_formulas} Under the assumptions of Theorem \ref{theo:first_variation_full} with $T=\BackwardPolicy$ and $\mu = \Fin$:

 \begin{equation}
      \nabla_{\Fin}^{\lambda} \mathcal L_{g,\Fin} = \left(\ForwardPolicy^\lambda-r\right)\left[g'\left(r\right)\lambda \right].
     \end{equation}

     For $g(x)=(x-1)^2$
      \begin{equation}
      \nabla_{\Fin}^{\lambda} \mathcal L_{g,\Fin} = 2\left(\ForwardPolicy^\lambda-r\right)
      \left[\left(r-1\right)\lambda \right] = 2 \left(\ForwardPolicy^\lambda (r\lambda)- \lambda-r^2\lambda+r\lambda\right).
     \end{equation}

     The case $g(x)=|x-1|$ is analogous when $r$ is moreover essentially bounded away from $1$, with $g'(r)$ replaced by $\mathrm{sign}(r-1)$; $|x-1|$ is not differentiable at $1$, and the hypothesis excludes every flow whose ratio comes arbitrarily close to $1$ on a set of positive measure --- in particular the balanced flow $r\equiv1$, at which $\mathcal L_{g,\Fin}$ vanishes.
\end{corollary}

\begin{proof}
Apply Theorem \ref{theo:first_variation_full} to $T=\BackwardPolicy$, $\mu=\Fin$, with $\nuG=\lambda$ and $\nu=\Fin$, so that $\frac{d\nuG}{d\lambda}=\frac{d\nu}{d\Fin}=1$; this gives the first display. For $g(x)=(x-1)^2$, $g'(r)=2(r-1)$ and the second display follows by expanding and using $\lambda\ForwardPolicy^\lambda=\lambda$, which holds since $\lambda\ForwardPolicy^\lambda$ is the second marginal of $\lambda \otimes \ForwardPolicy^\lambda = \BackwardPolicy\otimes\lambda$. For $g(x)=|x-1|$, let $a>0$ with $|r-1|\geq a$ $\Fin$-almost everywhere, and let $\widetilde g(x):=|x-1|$ for $|x-1|\geq\frac a2$ and $\widetilde g(x):=\frac{(x-1)^2}a+\frac a4$ elsewhere, which is $C^1$ with Lipschitz derivative; $r$ being essentially bounded, admissible directions with $\left\|\frac{d\delta}{d\Fin}\right\|_{L^\infty(\Fin)}$ and $\left\|\frac{d(\delta\BackwardPolicy)}{d\Fin}\right\|_{L^\infty(\Fin)}$ small enough keep the ratio of $\Fin+\delta$ within $\frac a2$ of $r$, so $\mathcal L_{g,\Fin}$ and $\mathcal L_{\widetilde g,\Fin}$ agree near $\Fin$ and have the same derivative there, and $\widetilde g'(r)=\mathrm{sign}(r-1)$ $\Fin$-almost everywhere.
\end{proof}

\begin{definition}[edge lift]\label{def:edge_lift}
Let $\BackwardPolicy$ be a backward policy on $(\Statespace,\lambda)$ with $\lambda$ $\BackwardPolicy$-invariant. The \emph{edge lift} of $\BackwardPolicy$ is the Markov kernel $K_2$ on $\Statespace^2$ given by
\begin{equation}
    K_2 : (s,s') \longmapsto (\BackwardPolicy(s),\, s),
\end{equation}
i.e. the current edge is shifted one $\BackwardPolicy$-step backward along the trajectory. The \emph{edge measure} is $\lambda_2(ds\,ds') := \BackwardPolicy(s'\rightarrow ds)\,\lambda(ds')$.
\end{definition}

\begin{lemma}\label{lem:lift_wellposed}
$\lambda_2$ is $K_2$-invariant, $K_2$ is a contraction of $\mathcal M^2(\lambda_2)$, and the $\mathcal M^2(\lambda_2)$-dual of $K_2$ is the edge lift $(s,s')\mapsto (s', \ForwardPolicy^\lambda(s'))$ of the dual forward policy $\ForwardPolicy^\lambda = \BackwardPolicy^\lambda$.
\end{lemma}

\begin{proof}
For $\mu\in\mathcal M^2(\lambda_2)$, writing $m_1\mu$ for its first marginal, a direct computation gives $\mu K_2(ds\,ds') = \BackwardPolicy(s'\rightarrow ds)(m_1\mu)(ds')$, so that
\begin{equation}\label{eq:muK2_density}
    \frac{d(\mu K_2)}{d\lambda_2}(s,s') = \frac{d(m_1\mu)}{d\lambda}(s').
\end{equation}
Taking $\mu=\lambda_2$, whose first marginal is $\lambda\BackwardPolicy=\lambda$, yields invariance. By the defining identity $\lambda\otimes\ForwardPolicy^\lambda = \BackwardPolicy\otimes\lambda$, under $\lambda_2$ the conditional law of $s'$ given $s$ is $\ForwardPolicy^\lambda(s\rightarrow\cdot)$, so that $\frac{d(m_1\mu)}{d\lambda}(s) = \int\frac{d\mu}{d\lambda_2}(s,s')\,\ForwardPolicy^\lambda(s\rightarrow ds')$: the map $\frac{d\mu}{d\lambda_2}\mapsto \frac{d(m_1\mu)}{d\lambda}$ is the conditional expectation given the first coordinate, hence an $L^2$ contraction; both marginals of $\lambda_2$ equal $\lambda$, so \eqref{eq:muK2_density} gives $\|\mu K_2\|_{\mathcal M^2(\lambda_2)}\leq \|\mu\|_{\mathcal M^2(\lambda_2)}$. Finally, the duality identity $\lambda_2\otimes K_2^* = K_2 \otimes \lambda_2$ for $K_2^*(s,s') := (s',\ForwardPolicy^\lambda(s'))$ follows by writing both sides as the law of three consecutive states of the stationary backward chain; by Lemma \ref{lem:adjoint}\emph{(1)}, applied to $K_2$ and its invariant measure $\lambda_2$ on $\Statespace^2$, $K_2^*$ is therefore the $\lambda_2$-reversal of $K_2$, and by Lemma \ref{lem:adjoint}\emph{(3)} its density action is the $\mathcal M^2(\lambda_2)$-adjoint of that of $K_2$.
\end{proof}

\begin{proposition}[DB as a lifted balance loss]\label{prop:db_lift}
Let $\mu\in \mathcal M^2(\lambda_2)$ be an edge flow, let $F:=m_1\mu$ and let $\ForwardPolicy^\mu$ be the disintegration $\mu = F\otimes \ForwardPolicy^\mu$. Then for any $g$ and any training distribution $\widehat\nu$ on $\Statespace^2$,
\begin{equation}
    \int_{\Statespace^2} g\left(\frac{d(\mu K_2)}{d\mu}\right)d\widehat\nu
    = \int g\left(\frac{F(ds')\BackwardPolicy(s'\rightarrow ds)}{F(ds)\ForwardPolicy^\mu(s\rightarrow ds')}\right)d\widehat\nu,
\end{equation}
the $g$-divergence detailed-balance loss with frozen backward policy (the usual orientation of the ratio corresponds to replacing $g$ by $x \mapsto g(1/x)$, which is immaterial for $g$ symmetric such as $(\log x)^2$ and exchanges the two branches of $(x-1)^2$ up to second order at $1$). In particular $\mu K_2=\mu$ if and only if $(F,\ForwardPolicy^\mu)$ satisfies detailed balance with respect to $\BackwardPolicy$.
\end{proposition}

\begin{proof}
By Equation \eqref{eq:muK2_density}, $\frac{d(\mu K_2)}{d\lambda_2}(s,s') = \frac{dF}{d\lambda}(s')$ while $\frac{d\mu}{d\lambda_2}(s,s') = \frac{F(ds)\ForwardPolicy^\mu(s\rightarrow ds')}{\lambda(ds')\BackwardPolicy(s'\rightarrow ds)}$; their ratio is the announced edge ratio, and it equals $1$ $\mu$-almost everywhere exactly when $F\otimes\ForwardPolicy^\mu$ coincides with the backward edge flow.
\end{proof}

\begin{corollary}[DB gradient]\label{cor:db_gradient}
Under the hypotheses of Theorem \ref{theo:first_variation_full} applied on $(\Statespace^2,\lambda_2, K_2)$, the DB loss with frozen backward policy satisfies, with $\widehat r := \frac{d(\mu K_2)}{d\mu}$,
\begin{equation}
    \nabla^{\lambda_2}_{\mu} \mathcal L_{\mathrm{DB},g,\widehat\nu}
    = \left(K_2^{\lambda_2} - \widehat r\right)\left[ g'(\widehat r)\, \frac{d\widehat\nu}{d\mu}\,\lambda_2 \right]
\end{equation}
where $K_2^{\lambda_2}$ is the forward edge kernel of Lemma \ref{lem:lift_wellposed}.
\end{corollary}

\begin{proof}
By Lemma \ref{lem:lift_wellposed}, $(\Statespace^2,\lambda_2,K_2)$ satisfies the hypotheses of Theorem \ref{theo:first_variation_full}; the formula is its specialization $T=K_2$, applied to the DB loss of Proposition \ref{prop:db_lift}.
\end{proof}

\begin{lemma}[mixing transfer]\label{lem:lift_mixing}
In the setting of Definition \ref{def:edge_lift}, let $\beta_n:=\|\BackwardPolicy^n-\Pi\|_{L^2(\lambda)}$ and $\widehat\beta_n:=\|K_2^n-\Pi_2\|_{L^2(\lambda_2)}$ be the $L^2$-mixing coefficients of the density actions of $\BackwardPolicy$ and of its edge lift $K_2$, where $\Pi$ and $\Pi_2$ are the mean projections against $\lambda$ and $\lambda_2$. Then $\widehat\beta_n = \beta_{n-1}$ for every $n\geq 1$.
\end{lemma}

\begin{proof}
By Lemma \ref{lem:lift_wellposed} and Lemma \ref{lem:adjoint}\emph{(3)}, the density actions of $\BackwardPolicy$ on $\mathcal M^2(\lambda)$ and of $K_2$ on $\mathcal M^2(\lambda_2)$ are the $L^2$-adjoints of their function actions; $\Pi$ and $\Pi_2$ being orthogonal projections, $\beta_n$ and $\widehat\beta_n$ are also the norms of $\BackwardPolicy^n-\Pi$ and $K_2^n-\Pi_2$ acting on functions, and we work with the function actions. For $f\in L^2(\lambda_2)$ set $g_f(s') := \int f(z,s')\BackwardPolicy(s'\rightarrow dz)$, and write $\mathrm{pr}_1(s,s'):=s$. By construction of $K_2$, $(K_2f)(s,s')=\int f(z,s)\BackwardPolicy(s\rightarrow dz)=g_f(s)$ and $K_2(\varphi\circ\mathrm{pr}_1)=(\BackwardPolicy\varphi)\circ\mathrm{pr}_1$ for $\varphi\in L^2(\lambda)$, so that $K_2^n f = (\BackwardPolicy^{n-1} g_f)\circ\mathrm{pr}_1$ for $n\geq1$; and $\int f\,d\lambda_2=\int g_f\,d\lambda$ with $\lambda_2(\Statespace^2)=\lambda(\Statespace)$, so that $\Pi_2 f = \Pi g_f$. Since the first marginal of $\lambda_2$ is $\lambda\BackwardPolicy=\lambda$,
\begin{equation}\label{eq:lift_mixing_identity}
    \|K_2^nf - \Pi_2 f\|_{L^2(\lambda_2)} = \|(\BackwardPolicy^{n-1}-\Pi)g_f\|_{L^2(\lambda)}\qquad(n\geq1).
\end{equation}
We prove the two inequalities separately. Jensen's inequality applied to the probability measures $\BackwardPolicy(s'\rightarrow\cdot)$ gives $\int g_f^2\,d\lambda\leq\iint f(z,s')^2\,\BackwardPolicy(s'\rightarrow dz)\,\lambda(ds')=\|f\|^2_{L^2(\lambda_2)}$, so \eqref{eq:lift_mixing_identity} gives $\|K_2^nf - \Pi_2 f\|_{L^2(\lambda_2)}\leq\beta_{n-1}\|f\|_{L^2(\lambda_2)}$, that is $\widehat\beta_n\leq\beta_{n-1}$. Conversely, let $\varphi\in L^2(\lambda)$ and $f(z,s'):=\varphi(s')$. Then $g_f=\varphi$, $\BackwardPolicy(s'\rightarrow\cdot)$ being a probability, and $\|f\|_{L^2(\lambda_2)}=\|\varphi\|_{L^2(\lambda)}$, the second marginal of $\lambda_2$ being $\lambda$; so \eqref{eq:lift_mixing_identity} gives $\|(\BackwardPolicy^{n-1}-\Pi)\varphi\|_{L^2(\lambda)}\leq\widehat\beta_n\|\varphi\|_{L^2(\lambda)}$, and the supremum over $\varphi$ gives $\beta_{n-1}\leq\widehat\beta_n$.
\end{proof}

\begin{lemma}[coercivity transfer]\label{lem:lift_coercivity}
In the setting of Definition \ref{def:edge_lift}, let $P$ and $P_2$ be the density actions $\varphi\mapsto\frac{d((\varphi\lambda)\BackwardPolicy)}{d\lambda}$ of $\BackwardPolicy$ on $L^2(\lambda)$ and $h\mapsto\frac{d((h\lambda_2)K_2)}{d\lambda_2}$ of its edge lift $K_2$ on $L^2(\lambda_2)$, and let $\Pi$ and $\Pi_2$ be the mean projections against $\lambda$ and $\lambda_2$. If $C\geq0$ satisfies $\|\varphi-\Pi\varphi\|_{L^2(\lambda)}\leq C\,\|(I-P)\varphi\|_{L^2(\lambda)}$ for every $\varphi\in L^2(\lambda)$, then
\begin{equation*}
    \|h-\Pi_2h\|_{L^2(\lambda_2)}\;\leq\;(1+C)\,\|(I-P_2)h\|_{L^2(\lambda_2)}\qquad\text{for every }h\in L^2(\lambda_2).
\end{equation*}
\end{lemma}

\begin{proof}
Norms of functions on $\Statespace$ are taken in $L^2(\lambda)$, and norms of functions on $\Statespace^2$ in $L^2(\lambda_2)$. Write $\mathrm{pr}_2(s,s'):=s'$ and, for $h\in L^2(\lambda_2)$, $k_h(s):=\int h(s,z)\,\ForwardPolicy^\lambda(s\rightarrow dz)$, where $\ForwardPolicy^\lambda$ is the dual forward policy of Lemma \ref{lem:lift_wellposed}. By Definition \ref{def:edge_lift} and Lemma \ref{lem:adjoint}\emph{(1)}, $\lambda_2(ds\,dz)=\lambda(ds)\,\ForwardPolicy^\lambda(s\rightarrow dz)$, so $k_h$ is the density against $\lambda$ of the first marginal of $h\lambda_2$, and \eqref{eq:muK2_density} gives $P_2h=k_h\circ\mathrm{pr}_2$. The same identity gives $\int k_h\,d\lambda=\int h\,d\lambda_2$ and $\lambda_2(\Statespace^2)=\lambda(\Statespace)$, so $\Pi_2h=\Pi k_h$; and, with Jensen's inequality applied to the probability measures $\ForwardPolicy^\lambda(s\rightarrow\cdot)$, $\|k_h\|^2\leq\iint h(s,z)^2\,\ForwardPolicy^\lambda(s\rightarrow dz)\,\lambda(ds)=\|h\|^2$. By Lemma \ref{lem:adjoint}\emph{(3)}, $P$ is the function action of $\ForwardPolicy^\lambda$: $(P\varphi)(s)=\int\varphi(z)\,\ForwardPolicy^\lambda(s\rightarrow dz)$.

Let $h\in L^2(\lambda_2)$ and $v:=(I-P_2)h=h-k_h\circ\mathrm{pr}_2$. Then $k_v=k_h-Pk_h=(I-P)k_h$, and
\begin{equation*}
    h-\Pi_2h\;=\;v+\bigl(k_h-\Pi k_h\bigr)\circ\mathrm{pr}_2 .
\end{equation*}
The second marginal of $\lambda_2$ being $\lambda$, $\|(k_h-\Pi k_h)\circ\mathrm{pr}_2\|=\|k_h-\Pi k_h\|\leq C\|(I-P)k_h\|=C\|k_v\|\leq C\|v\|$, and the triangle inequality gives $\|h-\Pi_2h\|\leq(1+C)\|v\|=(1+C)\|(I-P_2)h\|$.
\end{proof}

Under summable mixing of the backward policy, the constant $C=\sum_{n\geq0}\beta_n$ that Lemma \ref{lem:sigma_mixing} provides, $\beta_n$ as in Lemma \ref{lem:lift_mixing}, lifts to $1+C$, which is also the value of the edge-lifted mixing sum in Theorem \ref{theo:db_stable_frozen_full}.

\begin{theorem}[full form of Theorem \ref{theo:gd_diffusion}]\label{theo:gd_diffusion_full}
Under the hypotheses of Lemma \ref{lem:adjoint}, let $T$ be ergodic; write $P$ for the density action $h\mapsto\frac{d((h\lambda)T)}{d\lambda}$ of $T$ on $L^2(\lambda)$, $P^\dagger$ for its adjoint, $A:=P-I$, and $M_w$ for the multiplication by a non-negative $w\in L^\infty(\lambda)$, and let $\nu:=w\lambda$. Let $a\in(0,1)$ and let $g:\mathbb R_+^*\rightarrow\mathbb R$ be $C^3$ on $[1-a,1+a]$ with $g(1)=g'(1)=0<g''(1)$; set
\begin{gather*}
    \Gamma_3:=\sup_{[1-a,1+a]}|g'''|,\qquad C_g:=g''(1)+\frac{a\Gamma_3}2,\\
    C_4:=\|w\|_{L^\infty}\bigl(4g''(1)+3\Gamma_3\bigr),\qquad C_5:=\frac{16}3\,C_g\|w\|_{L^\infty},\qquad K:=2C_4+C_5.
\end{gather*}
Then for every $\varepsilon\in(0,a/4]$ and every $h\in L^\infty(\lambda)$ with $\|h\|_{L^\infty(\lambda)}\leq\varepsilon$, the ratio $r:=\frac{d(\mu T)}{d\mu}$ of $\mu:=(1+h)\lambda$ is $r=\frac{1+Ph}{1+h}$ and satisfies, $\lambda$-almost everywhere,
\begin{equation}\label{eq:gd_ratio_bounds}
    |r-1|\leq\tfrac43\,|Ah|,\qquad |r-1|\leq\tfrac{2a}3,\qquad |g'(r)|\leq C_g\,|r-1|;
\end{equation}
$\mathcal L_{g,\nu}$ has at $\mu$ a gradient $\nabla^{\lambda}\mathcal L_{g,\nu}(\mu)$ in the sense of Theorem \ref{theo:first_variation_full}, with density $P^\dagger\psi-r\psi$ against $\lambda$, where $\psi:=g'(r)\frac{w}{1+h}$; and
\begin{equation}\label{eq:linearized_flow}
\begin{aligned}
    \nabla^{\lambda}\mathcal L_{g,\nu}(\mu) &= g''(1)\, \left(P^\dagger - I\right)\!\left[\,w\,(P-I)h\,\right]\lambda + \mathrm{err}(h),\\
    \|\mathrm{err}(h)\|_{\mathcal M^2(\lambda)}&\leq K\varepsilon\,\|(P-I)h\|_{L^2(\lambda)}\leq 2K\varepsilon\,\|h\|_{L^2(\lambda)}.
\end{aligned}
\end{equation}
The linearization of the gradient flow $\dot\mu_t = -\nabla^{\lambda}\mathcal L_{g,\nu}(\mu_t)$ at $\lambda$ is therefore
\begin{equation}
    \dot h = -\,g''(1)\, A^\dagger M_w A\, h .
\end{equation}
In particular for $w\equiv1$ and $T$ reversible with respect to $\lambda$,
$\dot h = -g''(1)(I-P)^2 h$: training is the \emph{square} of the backward-policy heat flow $\dot h = -(I-P)h$, and each eigenmode $P\varphi = \beta\varphi$ of the backward policy decays at rate $g''(1)(1-\beta)^2$.
\end{theorem}

\begin{proof}
All norms without index are $L^2(\lambda)$ norms. By Lemma \ref{lem:adjoint}\emph{(3)}, $P$ is the function action of the Markov kernel $T^\lambda$ and $P^\dagger$ that of $T$; hence $P\mathbf 1=\mathbf 1$, $\|Pf\|_{L^\infty(\lambda)}\leq\|f\|_{L^\infty(\lambda)}$, $P^\dagger$ preserves $\lambda$-integrals, and by \eqref{eq:reversal_contraction} $\|P\|\leq1$, $\|P^\dagger\|\leq1$, $\|A\|\leq2$ and $\|A^\dagger\|\leq2$. Fix $\varepsilon\in(0,a/4]$ and $h$ with $\|h\|_{L^\infty(\lambda)}\leq\varepsilon$, and set $R(x):=g'(x)-g''(1)(x-1)$ for $x\in[1-a,1+a]$.
\begin{itemize}
    \item \emph{The ratio is $r=\frac{1+Ph}{1+h}$; it satisfies \eqref{eq:gd_ratio_bounds}, $|r-1|\leq\frac{8\varepsilon}3$ and $(r-1)^2\leq4\varepsilon|Ah|$, and $|R(r)|\leq\frac{\Gamma_3}2(r-1)^2$.} Indeed, $\varepsilon\leq\frac a4<\frac14$ gives $\frac34\leq1+h\leq\frac54$, so $\mu\sim\lambda$; by Lemma \ref{lem:adjoint}\emph{(3)}, $\frac{d(\mu T)}{d\lambda}=P(1+h)=1+Ph$, hence $r=\frac{1+Ph}{1+h}$ and $r-1=\frac{Ah}{1+h}$. Therefore $|r-1|\leq\frac43|Ah|$ and, as $|Ah|\leq|Ph|+|h|\leq2\varepsilon$, $|r-1|\leq\frac{2\varepsilon}{1-\varepsilon}\leq\frac{8\varepsilon}3\leq\frac{2a}3$, and $(r-1)^2\leq\frac43|Ah|\cdot\frac{8\varepsilon}3\leq4\varepsilon|Ah|$. As $r$ takes its values in $[1-a,1+a]$, on which $g$ is $C^3$, Taylor's formula at $1$ with $g'(1)=0$ gives $g'(r)=g''(1)(r-1)+R(r)$ with $|R(r)|\leq\frac{\Gamma_3}2(r-1)^2$; hence $|g'(r)|\leq\bigl(g''(1)+\frac{\Gamma_3}2|r-1|\bigr)|r-1|\leq C_g|r-1|$.
    \item \emph{The gradient $\nabla^\lambda\mathcal L_{g,\nu}(\mu)$ is defined and its density against $\lambda$ is $D:=P^\dagger\psi-r\psi$.} Indeed, let $\widetilde g:\mathbb R_+^*\rightarrow\mathbb R$ be equal to $g$ on $[1-a,1+a]$, to $g(1-a)+g'(1-a)(x-1+a)$ at $x<1-a$ and to $g(1+a)+g'(1+a)(x-1-a)$ at $x>1+a$. Its derivative is $g'$ on $[1-a,1+a]$, where $g'$ is $C^2$, and constant on either side, so $\widetilde g$ is continuously differentiable with Lipschitz derivative on $\mathbb R_+^*$, and $\widetilde g'(r)=g'(r)$, $r$ taking its values in $[1-a,1+a]$ by \eqref{eq:gd_ratio_bounds}. Next, $\mu\in\mathcal M^+(\Statespace,\nuB)$ since $\mu\ll\lambda\ll\nuB$; $\mu\sim\lambda$ and $\frac{d\mu}{d\lambda}=1+h\in L^2(\lambda)$, $\lambda$ being finite; $\mu T\ll\lambda\sim\mu$ by Lemma \ref{lem:adjoint}\emph{(3)}; $r$ takes its values in $[1-\frac{2a}3,1+\frac{2a}3]\subset(0,+\infty)$; and $\nu$ is finite with $\frac{d\nu}{d\mu}=\frac w{1+h}\in L^\infty(\mu)$. Theorem \ref{theo:first_variation_full} applied to $\widetilde g$ with $\nuG=\lambda$ gives $\nabla^\lambda\mathcal L_{\widetilde g,\nu}(\mu)=\bigl(T^\lambda-r\bigr)[\psi\lambda]$. For a direction $\delta$ admissible in that theorem, write $u$ and $v$ for the densities of $\delta$ and $\delta T$ against $\mu$; the ratio of $\mu+\delta$ is $\frac{r+v}{1+u}$, and if $\|u\|_{L^\infty(\mu)}+\|v\|_{L^\infty(\mu)}\leq\frac a6$ then, by $|r-1|\leq\frac{2a}3$ and $a<1$, $\bigl|\frac{r+v}{1+u}-1\bigr|\leq\frac{2a/3+a/6}{1-a/6}\leq a$. Hence $\mathcal L_{g,\nu}(\mu+\delta)=\mathcal L_{\widetilde g,\nu}(\mu+\delta)$ for all such $\delta$, the two losses have the same derivative at $\mu$, and $\nabla^\lambda\mathcal L_{g,\nu}(\mu)=\bigl(T^\lambda-r\bigr)[\psi\lambda]$. By Lemma \ref{lem:adjoint}\emph{(3)} the density of $(\psi\lambda)T^\lambda$ against $\lambda$ is the function action of $T$ on $\psi$, which is $P^\dagger\psi$.
    \item \emph{$\psi=g''(1)\,w\,Ah+e$ with $|e|\leq C_4\varepsilon|Ah|$ pointwise.} Indeed, $\psi=g''(1)\,w\,\frac{r-1}{1+h}+\frac{R(r)w}{1+h}$ and $\frac{r-1}{1+h}=\frac{Ah}{(1+h)^2}=Ah+Ah\bigl(\frac1{(1+h)^2}-1\bigr)$. As $|h|\leq\varepsilon\leq\frac14$, $\bigl|\frac1{(1+h)^2}-1\bigr|=\frac{|h|\,|2+h|}{(1+h)^2}\leq\frac{9\varepsilon/4}{9/16}=4\varepsilon$; and $|R(r)|\leq\frac{\Gamma_3}2(r-1)^2$ with $(r-1)^2\leq4\varepsilon|Ah|$. Hence $|e|\leq4g''(1)\|w\|_{L^\infty}\varepsilon|Ah|+\frac43\cdot\frac{\Gamma_3}2\cdot4\varepsilon|Ah|\,\|w\|_{L^\infty}\leq C_4\varepsilon|Ah|$.
    \item \emph{$|(r-1)\psi|\leq C_5\varepsilon|Ah|$ pointwise.} Indeed, by \eqref{eq:gd_ratio_bounds} and $\frac1{1+h}\leq\frac43$, $|(r-1)\psi|\leq\frac43C_g\|w\|_{L^\infty}(r-1)^2$, and $(r-1)^2\leq4\varepsilon|Ah|$.
    \item \emph{The expansion \eqref{eq:linearized_flow} holds.} Indeed, $D=A^\dagger\psi-(r-1)\psi=A^\dagger\bigl[g''(1)\,w\,Ah\bigr]+E$ with $E:=A^\dagger e-(r-1)\psi$, and the bounds $|e|\leq C_4\varepsilon|Ah|$ and $|(r-1)\psi|\leq C_5\varepsilon|Ah|$ give $\|E\|\leq\|A^\dagger\|\,\|e\|+\|(r-1)\psi\|\leq(2C_4+C_5)\,\varepsilon\,\|Ah\|=K\varepsilon\|Ah\|\leq2K\varepsilon\|h\|$. The measure $\mathrm{err}(h):=E\lambda$ satisfies \eqref{eq:linearized_flow}.
    \item \emph{The linearization is $\dot h=-g''(1)A^\dagger M_wAh$, and in the reversible case with $w\equiv1$ each eigenmode $P\varphi=\beta\varphi$ decays at rate $g''(1)(1-\beta)^2$.} Indeed, along $\mu_t=(1+h_t)\lambda$ the flow reads $\dot h_t=-D$, and \eqref{eq:linearized_flow} at every $\varepsilon\in(0,a/4]$ with $\varepsilon\geq\|h_t\|_{L^\infty(\lambda)}$ bounds the remainder by $2K\|h_t\|_{L^\infty(\lambda)}\|h_t\|$, so the part of $-D$ linear in $h$ is $-g''(1)A^\dagger M_wAh$. If $T$ is reversible with respect to $\lambda$, that is $\lambda\otimes T=T\otimes\lambda$, the uniqueness in Lemma \ref{lem:adjoint}\emph{(1)} gives $T^\lambda=T$, hence $P=P^\dagger$ and $A^\dagger A=(I-P)^2$; for $P\varphi=\beta\varphi$, $(I-P)^2\varphi=(1-\beta)^2\varphi$, so $h_t=e^{-g''(1)(1-\beta)^2t}\varphi$ solves $\dot h=-g''(1)(I-P)^2h$.
\end{itemize}
\end{proof}

\begin{lemma}[mixing controls coercivity]\label{lem:sigma_mixing}
Let $T$ be a Markov kernel on a Polish space $\Statespace$ leaving a non-zero finite measure $\lambda$ invariant, $P$ its density action $h\mapsto\frac{d((h\lambda)T)}{d\lambda}$ on $L^2(\lambda)$, $\Pi$ the $\lambda$-mean projection $h\mapsto\lambda(\Statespace)^{-1}\int h\,d\lambda$, and $\widehat\beta_n:=\|P^n-\Pi\|_{L^2(\lambda)}$. Assume $\widehat B := \sum_{n\geq0}\widehat\beta_n < +\infty$ (summable $L^2$-mixing). Then $S := \sum_{n\geq0}(P^n - \Pi)$ converges in operator norm, $S(I-P) = I - \Pi$, and consequently
\begin{equation}
    \forall h\in L^2(\lambda),\quad \|h - \Pi h\|_{L^2(\lambda)} \;\leq\; \widehat B\,\|(I-P)h\|_{L^2(\lambda)}.
\end{equation}
If moreover $\lambda$ is not a multiple of a Dirac mass --- in particular if $\Statespace$ is finite with $\#\Statespace\geq2$ and $\lambda>0$ everywhere --- then $\widehat\beta_0=1$ and $\widehat B\geq1$.
\end{lemma}

\begin{proof}[Proof of Lemma \ref{lem:sigma_mixing}]
$\|P^n - \Pi\| = \widehat\beta_n$ is summable so $S$ converges absolutely. Since $P$ preserves $\lambda$-integrals and fixes constants, $\Pi P = P \Pi = \Pi$, whence $(P^n-\Pi)(I-P) = P^n - P^{n+1}$ and the partial sums telescope to $I - P^{N+1} \rightarrow I-\Pi$. Then $\|h-\Pi h\| = \|S(I-P)h\| \leq \widehat B\|(I-P)h\|$. If $\lambda$ is not a multiple of a Dirac mass, then, $\lambda$ being a finite measure on the Polish space $\Statespace$, it charges a Borel set $X$ with $0<\lambda(X)<\lambda(\Statespace)$, so $\mathbf 1_X-\Pi\mathbf 1_X\neq0$; hence $I-\Pi$ is a non-zero orthogonal projection, $\widehat\beta_0=\|I-\Pi\|=1$, and $\widehat B\geq\widehat\beta_0=1$.
\end{proof}

\begin{theorem}[full form of Theorem \ref{theo:db_stable_frozen}]\label{theo:db_stable_frozen_full}
Let $T$, $\lambda$, $P$, $A$, $w$, $M_w$ and $\nu$ be as in Theorem \ref{theo:gd_diffusion_full}, and let $g$ be twice differentiable at $1$ with $g''(1)>0$. Assume moreover that $w \geq w_{\min} > 0$ and that $\widehat B\geq1$ is a \emph{coercivity constant} of $T$, $\Pi$ denoting the $\lambda$-mean projection $h\mapsto\lambda(\Statespace)^{-1}\int h\,d\lambda$:
\begin{equation}\label{eq:coercivity}
    \|h-\Pi h\|_{L^2(\lambda)}\;\leq\;\widehat B\,\|(I-P)h\|_{L^2(\lambda)}\qquad\text{for every }h\in L^2(\lambda).
\end{equation}
Set
\begin{equation}
    \varrho := \frac{g''(1)\, w_{\min}}{\widehat B^{\,2}},\qquad H:=g''(1)\,A^\dagger M_wA.
\end{equation}
Then along the linearized gradient flow $\dot h_t=-Hh_t$ of $\mathcal L_{g,\nu}$ at the balanced solution: the component $\Pi h_t$ (the total flow normalization) is conserved, and
\begin{equation}
    \|h_t - \Pi h_t\|_{L^2(\lambda)} \leq e^{-\varrho\, t}\, \|h_0 - \Pi h_0\|_{L^2(\lambda)}.
\end{equation}
For the linearized gradient descent $h_{k+1}=h_k-\varepsilon Hh_k$ with step $\varepsilon \leq \left(4 g''(1)\|w\|_{L^\infty}\right)^{-1}$, the contraction factor per step is $1-\varepsilon\varrho$. For the FM loss $T=\BackwardPolicy$, and for the DB loss $T=K_2$ with $\lambda_2$ in place of $\lambda$ (Definition \ref{def:edge_lift}). When $\BackwardPolicy$ has summable $L^2$-mixing, $B:=\sum_{n\geq0}\beta_n<\infty$ with $\beta_n$ as in Lemma \ref{lem:lift_mixing}, and its invariant measure $\lambda$ is not a multiple of a Dirac mass, $\widehat B:=\sum_{n\geq0}\widehat\beta_n$, where $\widehat\beta_n:=\|P^n-\Pi\|_{L^2(\lambda)}$, is a coercivity constant of $T$ with $\widehat B\geq1$, equal to $B$ for the FM loss and to $1+B$ for the DB loss: \emph{a backward policy with summable mixing makes the linearized gradient descent on the FM and DB losses with frozen backward policy converge to the balanced ray, at rate at least $g''(1)w_{\min}/(1+B)^2$} (the nonlinear upgrade, at half the rate for the gradient flow and from an explicit neighbourhood, is Theorem \ref{theo:local_convergence_full}).
\end{theorem}

\begin{proof}
Note that $H\succeq 0$, $M_w$ being non-negative. $\Pi H = 0$ since $P^\dagger$ preserves integrals, so $\Pi h_t$ is conserved; and for $h_\perp := h - \Pi h$, using $Ah = Ah_\perp$ and \eqref{eq:coercivity},
\begin{equation}
    \langle h, Hh\rangle \geq g''(1) w_{\min}\|Ah_\perp\|^2 \geq \varrho\, \|h_\perp\|^2,
\end{equation}
whence $\frac{d}{dt}\|h_\perp\|^2 = -2\langle h,Hh\rangle \leq -2\varrho\|h_\perp\|^2$ and Gr\"onwall concludes. For the discrete case, $\|A\|\leq 2$ gives $\|H\| \leq 4g''(1)\|w\|_{L^\infty}$, so $\varepsilon\|H\|\leq1$ and $\|(I-\varepsilon H)h_\perp\| \leq (1-\varepsilon\varrho)\|h_\perp\|$.

It remains to prove the assertion on summable mixing. Assume $B<\infty$ and that $\lambda$ is not a multiple of a Dirac mass; $\lambda$ being a finite measure on the Polish space $\Statespace$, it charges a Borel set $X$ with $0<\lambda(X)<\lambda(\Statespace)$. For the FM loss, $T=\BackwardPolicy$, so $\widehat\beta_n=\beta_n$ and $\widehat B=B$; and $\widehat\beta_0=\|I-\Pi\|_{L^2(\lambda)}=1$, $I-\Pi$ being an orthogonal projection that is non-zero, since $\mathbf 1_X$ is not $\lambda$-almost everywhere constant. For the DB loss, $T=K_2$ on $(\Statespace^2,\lambda_2)$, and Lemma \ref{lem:lift_mixing} gives $\widehat B=\widehat\beta_0+\sum_{n\geq1}\beta_{n-1}=\widehat\beta_0+B$, where $\widehat\beta_0=\|I-\Pi_2\|_{L^2(\lambda_2)}=1$: $I-\Pi_2$ is an orthogonal projection, and it is non-zero, since $f(s,s'):=\mathbf 1_X(s')$ is not $\lambda_2$-almost everywhere constant, the second marginal of $\lambda_2$ being $\lambda$. In both cases $\widehat B<\infty$ and $\widehat\beta_0=1$, so Lemma \ref{lem:sigma_mixing} applied to $T$ gives \eqref{eq:coercivity}, and $\widehat B\geq\widehat\beta_0=1$; finally $\widehat B\leq1+B$, so $\varrho\geq g''(1)w_{\min}/(1+B)^2$.
\end{proof}

\begin{theorem}[full form of Theorem \ref{theo:local_convergence}]\label{theo:local_convergence_full}
In the setting of Theorem \ref{theo:db_stable_frozen_full}, assume moreover that $g(1)=g'(1)=0$, that $g$ is $C^3$ on $[1-a,1+a]$ for some $a\in(0,1)$, and that the state space $\Statespace$ is finite, with $\lambda$ a probability and $\lambda(x)>0$ for every $x\in\Statespace$; set $\lambda_{\min}:=\min_{x\in\Statespace}\lambda(x)$ and $C_\infty:=\lambda_{\min}^{-1/2}$, so that $C_\infty\geq1$. For the DB loss, $T=K_2$ (Definition \ref{def:edge_lift}) is taken on the set $E$ of pairs $(s,s')\in\Statespace^2$ with $\BackwardPolicy(s'\rightarrow s)>0$, which is finite when $\Statespace$ is, which $K_2$ maps into itself, and on which $\lambda_2(s,s')=\BackwardPolicy(s'\rightarrow s)\lambda(s')>0$ when $\lambda>0$; $E$ plays the role of $\Statespace$ and the probability $\lambda_2$ that of $\lambda$ --- in particular $\lambda_{\min}:=\min_{e\in E}\lambda_2(e)$ and $C_\infty:=\lambda_{\min}^{-1/2}$. Then there are explicit $\varepsilon_0\in\bigl(0,a/(16C_\infty)\bigr]$ and $C\geq1$, depending only on $g''(1)$, $\sup_{[1-a,1+a]}|g'''|$, $a$, $w_{\min}$, $\|w\|_{L^\infty}$, $\widehat B$ and $C_\infty$, such that for every initialization $\mu_0=(1+h_0)\lambda$ with $\|h_0\|_{L^2(\lambda)}\leq\varepsilon_0$, the \emph{nonlinear} gradient flow $\dot\mu_t=-\nabla^{\lambda}\mathcal L_{g,\nu}(\mu_t)$ converges to a balanced flow $c_\infty\lambda$ with $|c_\infty-1-\Pi h_0|\leq C\|h_0-\Pi h_0\|^2$ and
\begin{equation}
    \|h_t-(c_\infty-1)\|_{L^2(\lambda)} \;\leq\; 2\,e^{-\varrho t/2}\,\|h_0-\Pi h_0\|_{L^2(\lambda)}.
\end{equation}
There is moreover an explicit $\gamma_0>0$, depending only on the same quantities, such that for every such $h_0$ and every step $0<\gamma\leq\gamma_0$ the gradient descent $h_{k+1}:=h_k-\gamma D_k$, $D_k$ the density against $\lambda$ of $\nabla^{\lambda}\mathcal L_{g,\nu}\bigl((1+h_k)\lambda\bigr)$, is well defined and satisfies $\|h_{k+1}-\Pi h_{k+1}\|_{L^2(\lambda)}\leq(1-\gamma\varrho/4)\|h_k-\Pi h_k\|_{L^2(\lambda)}$ for every $k\geq0$.
\end{theorem}

\begin{proof}
All norms are $L^2(\lambda)$; set $\Gamma_3:=\sup_{[1-a,1+a]}|g'''|$, $C_g:=g''(1)+\frac{a\Gamma_3}2$, $m_t:=\Pi h_t$ (a scalar, $T$ being ergodic), $h^\perp:=h-m$, and recall $Ah=Ah^\perp$, $\|P\|\leq1$, $\|A\|\leq2$, $\|A^\dagger\|\leq2$. Since $\Statespace$ is finite, $\lambda_{\min}f(x)^2\leq\lambda(x)f(x)^2\leq\|f\|^2$ for every function $f$ and every $x\in\Statespace$, so that $\|f\|_{L^\infty(\lambda)}\leq C_\infty\|f\|$; and $C_\infty\geq1$ because $\lambda_{\min}\leq\lambda(\Statespace)=1$. For the DB loss the same holds on $E$ with $\lambda_2$, which is a probability, its second marginal being $\lambda$.

\emph{Step 1 (expansion of the gradient field).} Assume $\|h\|_{L^\infty}\leq\varepsilon\leq a/4$, and let $K$ be the constant of Theorem \ref{theo:gd_diffusion_full} for $g$ and $w$. By that theorem, the ratio at $(1+h)\lambda$ is $r=\frac{1+Ph}{1+h}$; \eqref{eq:gd_ratio_bounds} gives, pointwise, $|r-1|\leq\frac43|Ah^\perp|$, $|r-1|\leq\frac{2a}3$ and $|g'(r)|\leq C_g|r-1|$; the gradient density at $(1+h)\lambda$ is $D=P^\dagger\psi-r\psi$ with $\psi:=g'(r)\frac w{1+h}$; and \eqref{eq:linearized_flow} reads
\begin{eqnarray}
    D &=& A^\dagger\bigl[g''(1)\,w\,Ah^\perp\bigr]+E,\label{eq:gradient_expansion}\\
    \|E\|&\leq& K\varepsilon\|Ah^\perp\|.\nonumber
\end{eqnarray}

\emph{Step 2 (energy estimate).} Using $\langle h^\perp, A^\dagger[g''(1)wAh^\perp]\rangle=g''(1)\langle Ah^\perp,wAh^\perp\rangle\geq g''(1)w_{\min}\|Ah^\perp\|^2$ and $\|h^\perp\|\leq\widehat B\|Ah^\perp\|$, which is \eqref{eq:coercivity}, for $\varepsilon\leq\varepsilon_1:=\frac{g''(1)w_{\min}}{2K\widehat B}$,
\begin{eqnarray*}
    \frac{d}{dt}\frac12\|h^\perp\|^2=-\langle h^\perp,D\rangle&\leq&-\frac{g''(1)w_{\min}}2\|Ah^\perp\|^2\\
    &\leq&-\frac\varrho2\|h^\perp\|^2,
\end{eqnarray*}
whence $\|h_t^\perp\|\leq e^{-\varrho t/2}\|h_0^\perp\|$ and $\int_0^\infty\|Ah_s^\perp\|^2ds\leq\|h_0^\perp\|^2/(g''(1)w_{\min})$.

\emph{Step 3 (drift of the mass).} Since $P^\dagger$ preserves $\lambda$-integrals, $\dot m=-\int D\,d\lambda=\int(r-1)\psi\,d\lambda$; by the pointwise bounds of Step 1 and $\frac1{1+h}\leq\frac43$, $|(r-1)\psi|\leq\frac43C_g\|w\|_{L^\infty}(r-1)^2\leq\frac{64}{27}C_g\|w\|_{L^\infty}|Ah^\perp|^2$, so $|\dot m|\leq C_6\|Ah^\perp\|^2$ with $C_6:=4C_g\|w\|_{L^\infty}$. By Step 2, $m_t$ converges to some $m_\infty$ with $|m_\infty-m_t|\leq\frac{4C_6}\varrho e^{-\varrho t}\|h_0^\perp\|^2$ and $|m_t-m_0|\leq C_7\|h_0^\perp\|^2$, $C_7:=C_6/(g''(1)w_{\min})$.

\emph{Step 4 (the flow from $h_0$ exists for all times, is unique, and stays in $\{\|h\|_{L^\infty}\leq\frac\varepsilon2\}$).} Set $\varepsilon:=\min\bigl(\tfrac a4,\varepsilon_1\bigr)$ and $\varepsilon_0:=\min\bigl(\tfrac\varepsilon{2C_\infty(2+C_7)},\tfrac\varrho{4C_6},1\bigr)$; since $C_7\geq0$ and $\varepsilon\leq a/4$, $\varepsilon_0\leq\frac\varepsilon{4C_\infty}\leq\frac a{16C_\infty}$. Write $\mathcal O:=\{h\in\mathbb R^{\Statespace}:\|h\|_{L^\infty}<\frac a4\}$, an open set. For $h\in\mathcal O$, Theorem \ref{theo:gd_diffusion_full}, applied at the same $a$ with $\frac a4$ in the role of its $\varepsilon$, gives the gradient density $D(h)=P^\dagger\psi-r\psi$ at $(1+h)\lambda$, with $r=\frac{1+Ph}{1+h}$ taking its values in $[1-\frac{2a}3,1+\frac{2a}3]\subset(1-a,1+a)$, on which $g'$ is $C^2$, and $1+h>\frac34$; hence $D$ is a $C^1$ function of $h\in\mathcal O$, and by the Cauchy--Lipschitz theorem the equation $\dot h=-D(h)$, which is the gradient flow in densities, has from every $h_0\in\mathcal O$ a unique maximal solution $(h_t)_{t\in[0,t_{\max})}$ in $\mathcal O$. If $\|h_0\|\leq\varepsilon_0$, then $\|h_0\|_{L^\infty}\leq C_\infty\varepsilon_0\leq\frac\varepsilon4$, so $h_0\in\mathcal O$, and on any interval $[0,t]\subset[0,t_{\max})$ where $\|h_s\|_{L^\infty}\leq\varepsilon$ the estimates above give $\|h_t\|_{L^\infty}\leq C_\infty(|m_t|+\|h_t^\perp\|)\leq C_\infty\varepsilon_0(2+C_7)\leq\varepsilon/2$: by continuity this holds on all of $[0,t_{\max})$, so the solution remains in $\{h:\|h\|_{L^\infty}\leq\frac\varepsilon2\}$, a compact subset of $\mathcal O$ since $\frac\varepsilon2\leq\frac a8$. A maximal solution of a locally Lipschitz equation that remains in a compact subset of its open domain is defined for all times, so $t_{\max}=+\infty$, and Steps 1--3 hold globally in time. Finally, let $(h'_t)_{t\in[0,T)}$, $T\in(0,+\infty]$, be a solution of the gradient flow from $h_0$, and $t'$ the supremum of the $t\in[0,T)$ with $h'_s\in\mathcal O$ for every $s\in[0,t]$. By uniqueness $h'_t=h_t$ for $t<t'$. If $t'<T$, continuity at $t'$ gives $\|h'_{t'}\|_{L^\infty}\leq\frac\varepsilon2<\frac a4$, so $h'$ remains in $\mathcal O$ on a neighbourhood of $t'$, against the definition of $t'$. Hence $t'=T$, and $h'_t=h_t$ for every $t\in[0,T)$.

\emph{Step 5 (conclusion).} With $c_\infty:=1+m_\infty$ and $C:=\max(1,C_7)$, $|c_\infty-1-\Pi h_0|\leq C_7\|h_0^\perp\|^2\leq C\|h_0^\perp\|^2$ and
\begin{eqnarray*}
    \|h_t-(c_\infty-1)\|&\leq&|m_\infty-m_t|+\|h_t^\perp\|\\
    &\leq& e^{-\varrho t/2}\|h_0^\perp\|\Bigl(1+\frac{4C_6}\varrho\varepsilon_0\Bigr)\\
    &\leq&2e^{-\varrho t/2}\|h_0^\perp\|.
\end{eqnarray*}

\emph{Step 6 (gradient descent).} Set $L:=2g''(1)\|w\|_{L^\infty}+K\varepsilon$ and $\gamma_0:=\frac{g''(1)w_{\min}}{2L^2}$, which depend only on the quantities listed in the statement, let $0<\gamma\leq\gamma_0$, and write $m_k:=\Pi h_k$. Since $L\geq2g''(1)\|w\|_{L^\infty}\geq2g''(1)w_{\min}$, $\gamma\leq(8g''(1)w_{\min})^{-1}$; since $\widehat B\geq1$, $\varrho\leq g''(1)w_{\min}$; hence
\begin{equation}\label{eq:gd_step_cap}
    \gamma L^2\leq\frac{g''(1)w_{\min}}2,\qquad\frac{\gamma\varrho}2\leq\frac1{16}.
\end{equation}
We first prove that if $\|h_k\|_{L^\infty}\leq\varepsilon$, then $D_k$ is defined and
\begin{align}
    \|h^\perp_{k+1}\|^2+\gamma\,\frac{g''(1)w_{\min}}2\,\|Ah^\perp_k\|^2&\leq\|h^\perp_k\|^2,\label{eq:gd_energy}\\
    \|h^\perp_{k+1}\|&\leq\Bigl(1-\frac{\gamma\varrho}4\Bigr)\|h^\perp_k\|,\label{eq:gd_contraction}\\
    |m_{k+1}-m_k|&\leq2C_7\bigl(\|h^\perp_k\|^2-\|h^\perp_{k+1}\|^2\bigr).\label{eq:gd_drift}
\end{align}
By Step 1, $D_k$ is defined, and \eqref{eq:gradient_expansion} with $\|A^\dagger\|\leq2$ gives $\|D_k\|\leq L\|Ah^\perp_k\|$. As $\Pi h^\perp_k=0$, $h^\perp_{k+1}=h^\perp_k-\gamma D_k^\perp$ with $\langle h^\perp_k,D_k^\perp\rangle=\langle h^\perp_k,D_k\rangle$ and $\|D^\perp_k\|\leq\|D_k\|$; the inequality $\langle h^\perp,D\rangle\geq\frac{g''(1)w_{\min}}2\|Ah^\perp\|^2$ of Step 2 and \eqref{eq:gd_step_cap} then give
\begin{eqnarray*}
    \|h^\perp_{k+1}\|^2&\leq&\|h^\perp_k\|^2-\gamma\bigl(g''(1)w_{\min}-\gamma L^2\bigr)\|Ah^\perp_k\|^2\\
    &\leq&\|h^\perp_k\|^2-\gamma\,\frac{g''(1)w_{\min}}2\,\|Ah^\perp_k\|^2,
\end{eqnarray*}
which is \eqref{eq:gd_energy}. By \eqref{eq:coercivity}, $\frac{g''(1)w_{\min}}2\|Ah^\perp_k\|^2\geq\frac\varrho2\|h^\perp_k\|^2$, so \eqref{eq:gd_energy} gives $\|h^\perp_{k+1}\|^2\leq(1-\frac{\gamma\varrho}2)\|h^\perp_k\|^2$; as $\sqrt{1-x}\leq1-\frac x2$ for $x\in[0,1]$ and $x=\frac{\gamma\varrho}2\leq\frac1{16}$ by \eqref{eq:gd_step_cap}, this is \eqref{eq:gd_contraction}. Finally $m_{k+1}-m_k=-\gamma\int D_k\,d\lambda$, and the bound $|\int D\,d\lambda|\leq C_6\|Ah^\perp\|^2$ of Step 3 together with \eqref{eq:gd_energy} gives $|m_{k+1}-m_k|\leq\gamma C_6\|Ah^\perp_k\|^2\leq\frac{2C_6}{g''(1)w_{\min}}\bigl(\|h^\perp_k\|^2-\|h^\perp_{k+1}\|^2\bigr)$, which is \eqref{eq:gd_drift}.

It remains to prove, by induction on $k$, that every $k\geq0$ satisfies
\begin{equation}\label{eq:gd_invariant}
    \|h_k\|_{L^\infty}\leq\varepsilon,\qquad \|h^\perp_k\|\leq\|h^\perp_0\|,\qquad |m_k-m_0|\leq2C_7\bigl(\|h^\perp_0\|^2-\|h^\perp_k\|^2\bigr).
\end{equation}
Note that $|m_0|\leq\|h_0\|$ by the Cauchy--Schwarz inequality, $\lambda$ being a probability, and that $\|h^\perp_0\|\leq\|h_0\|\leq\varepsilon_0\leq1$. At $k=0$, $\|h_0\|_{L^\infty}\leq C_\infty\varepsilon_0\leq\varepsilon$ by Step 4, and the other two conditions hold with equality. If \eqref{eq:gd_invariant} holds at $k$, then \eqref{eq:gd_contraction} and \eqref{eq:gd_drift} hold at $k$: the first gives $\|h^\perp_{k+1}\|\leq\|h^\perp_k\|\leq\|h^\perp_0\|$, and the second, added to the third condition of \eqref{eq:gd_invariant}, gives $|m_{k+1}-m_0|\leq2C_7(\|h^\perp_0\|^2-\|h^\perp_{k+1}\|^2)\leq2C_7\varepsilon_0$. Hence, pointwise,
\begin{equation*}
    |h_{k+1}|\leq|m_{k+1}|+C_\infty\|h^\perp_{k+1}\|\leq C_\infty\bigl(\varepsilon_0+2C_7\varepsilon_0+\varepsilon_0\bigr)\leq2C_\infty\varepsilon_0(2+C_7)\leq\varepsilon,
\end{equation*}
the last inequality by the definition of $\varepsilon_0$ in Step 4. This closes the induction: $h_k$ is well defined and \eqref{eq:gd_contraction} holds at every $k\geq0$.
\end{proof}

\begin{proposition}[full form of Proposition \ref{prop:frozen_unstable}]\label{prop:frozen_unstable_full}
Let $\mathcal G$ be a finite path-connected marked graph carrying a directed cycle $\gamma$ disjoint from $\{s_0,s_f\}$, let $\widehat\BackwardPolicy$ be an everywhere-positive frozen backward policy on its loop closure $\widehat G$ (Definition \ref{def:loop_closure}), let $\lambda$ be the associated balanced flow and let $g$ be continuous with $g(1)=0$ and $g>0$ elsewhere. Write $F^\star(u\rightarrow v):=\lambda(v)\widehat\BackwardPolicy(v\rightarrow u)$ for the balanced edgeflow, $\mathbf 1_\gamma$ for the unit circulation along $\gamma$, $\mathcal L_{g,\nu}(F):=\mathcal L_{g,\nu}(\mu_F)=\int g\bigl(\frac{d(\mu_F\widehat\BackwardPolicy)}{d\mu_F}\bigr)d\nu$ for the flow-matching loss of an edgeflow $F$ with frozen backward policy $\widehat\BackwardPolicy$, where $\mu_F(v):=\sum_uF(u\rightarrow v)$ is the inflow of $F$, and $\mathcal L_{\mathrm{DB},g,\widehat\nu}$ for the detailed-balance loss of Proposition \ref{prop:db_lift} with frozen backward policy $\widehat\BackwardPolicy$. Then for every training measure $\nu$ of positive density, every edge training measure $\widehat\nu$ with positive density against the edge measure $\lambda_2$ of $\widehat\BackwardPolicy$ --- which is positive exactly on the edges of $\widehat G$ --- and every
\begin{equation}
    0<\varepsilon\leq\tfrac12\min_{(u\rightarrow v)\in\gamma}F^\star(u\rightarrow v),
\end{equation}
the edgeflow $F:=F^\star-\varepsilon\mathbf 1_\gamma$ and the $0$-flow $F_0:=\varepsilon\mathbf 1_\gamma$ satisfy $0\leq F_0\leq F$ and
\begin{equation}
    \mathcal L_{g,\nu}(F+F_0)=0<\mathcal L_{g,\nu}(F),\qquad \mathcal L_{\mathrm{DB},g,\widehat\nu}(F+F_0)=0<\mathcal L_{\mathrm{DB},g,\widehat\nu}(F).
\end{equation}
Hence the FM and DB losses with frozen backward policy remain \emph{unstable} in the sense of \eqref{eq:def_stability}, on every such graph and for every such generator and training measures.
\end{proposition}

\begin{proof}
Write $r$ for the ratio $\frac{d(\mu\widehat\BackwardPolicy)}{d\mu}$ of a state flow $\mu$, so that $\mathcal L_{g,\nu}(F)=\int g(r)\,d\nu$ at $\mu=\mu_F$. Since the rows of $\widehat\BackwardPolicy$ sum to $1$, $\mu_{F^\star}=\lambda$, and $\mu_{\mathbf 1_\gamma}=\partial\gamma:=\sum_{v\in\gamma}\delta_v$.

\emph{Admissibility.} $\lambda>0$ and $\widehat\BackwardPolicy>0$ give $F^\star>0$ on every edge, so $\varepsilon$ is well defined and positive, and $2\varepsilon\mathbf 1_\gamma\leq F^\star$ edgewise; hence $F=F^\star-\varepsilon\mathbf 1_\gamma\geq\varepsilon\mathbf 1_\gamma=F_0\geq0$, and $F\geq\frac12F^\star>0$ on every edge of $\widehat G$. The outflow of $F^\star$ at $u$ is $\sum_v\lambda(v)\widehat\BackwardPolicy(v\rightarrow u)=\lambda(u)$, by $\lambda\widehat\BackwardPolicy=\lambda$, which is its inflow $\mu_{F^\star}(u)$; so $F^\star$ is a circulation on $\widehat G$, and so is $\mathbf 1_\gamma$, $\gamma$ being a directed cycle. Hence $F$ is a non-negative circulation on $\widehat G$ whose restriction to the edges of $\mathcal G$ is a flow in the sense of Definition 1 of \citet{brunswic2024theory}; $F_0=\varepsilon\mathbf 1_\gamma$ is a flow carrying no flow into $s_f$, since $\gamma$ avoids $s_f$, i.e. a $0$-flow in the sense of Definition 2 of \citet{brunswic2024theory}; and $F_0\leq F$.

\emph{$\mathcal L_{g,\nu}(F)>0$.} We have $\mu_F=\lambda-\varepsilon\,\partial\gamma$, and $\varepsilon\,\partial\gamma(v)=\sum_u\varepsilon\mathbf 1_\gamma(u\rightarrow v)\leq\frac12\sum_uF^\star(u\rightarrow v)=\frac12\lambda(v)$, so $\mu_F\geq\frac12\lambda>0$ and the ratio $r$ at $\mu_F$ is defined. As $\gamma$ avoids $s_f$, $\partial\gamma(s_f)=0<\lambda(s_f)$ while $\partial\gamma\neq0$, so $\partial\gamma$ is not proportional to $\lambda$ and neither is $\mu_F$. The loop-closed chain being ergodic, its invariant measures are exactly the multiples of $\lambda$, so $\mu_F\widehat\BackwardPolicy\neq\mu_F$: the ratio $r$ differs from $1$ on a set of positive $\nu$-measure, and $g>0$ off $1$ gives $\mathcal L_{g,\nu}(F)>0$.

\emph{$\mathcal L_{g,\nu}(F+F_0)=0$.} $F+F_0=F^\star$, whose state flow $\lambda$ is $\widehat\BackwardPolicy$-invariant, so $r\equiv1$ and $g(1)=0$.

\emph{The DB case.} Let $F'$ be an edgeflow positive on every edge of $\widehat G$, and $\mu^{\mathrm{out}}_{F'}(v):=\sum_wF'(v\rightarrow w)$ its outflow, which is the first marginal of Proposition \ref{prop:db_lift}. By that proposition $\mathcal L_{\mathrm{DB},g,\widehat\nu}(F')$ integrates against $\widehat\nu$ the function $g$ of the edge ratio $\frac{\mu^{\mathrm{out}}_{F'}(v)\widehat\BackwardPolicy(v\rightarrow u)}{F'(u\rightarrow v)}$; as $\lambda_2(u,v)=\widehat\BackwardPolicy(v\rightarrow u)\lambda(v)$ is positive exactly when $u\rightarrow v$ is an edge of $\widehat G$ --- $\widehat\BackwardPolicy$ being positive on the edges of $\widehat G$ and supported on in-neighbours, and $\lambda>0$ --- $\widehat\nu$ is positive on every edge of $\widehat G$ and null off them; with $g(1)=0$ and $g>0$ off $1$, the loss vanishes if every edge satisfies $F'(u\rightarrow v)=\mu^{\mathrm{out}}_{F'}(v)\widehat\BackwardPolicy(v\rightarrow u)$, and is positive otherwise. Both $F^\star$ and $F$ are positive on every edge and are circulations, so their outflows equal their inflows: $\mu^{\mathrm{out}}_{F^\star}=\lambda$ and $\mu^{\mathrm{out}}_F=\mu_F$. At $F+F_0=F^\star$ every edge satisfies the balance by construction, so the DB loss is $0$ there. At $F$ some edge does not: balance on every edge would give, summing over $v$ at each $u$, $\mu_F(u)=\sum_v\mu_F(v)\widehat\BackwardPolicy(v\rightarrow u)$, that is $\mu_F\widehat\BackwardPolicy=\mu_F$, contradicting $\mu_F\widehat\BackwardPolicy\neq\mu_F$. Hence $\mathcal L_{\mathrm{DB},g,\widehat\nu}(F)>0$, and the same pair $(F,F_0)$ witnesses DB instability.

Hence \eqref{eq:def_stability} fails at the pair $(F,F_0)$ for both losses. In state-flow coordinates the same computation reads: $t\mapsto\mathcal L_{g,\nu}(\lambda-\varepsilon\partial\gamma+t\,\partial\gamma)$ drops from a positive value at $t=0$ to $0$ at $t=\varepsilon$, $\partial\gamma$ being the state-flow signature of the circulation.
\end{proof}

\begin{proposition}[TB gradient]\label{prop:tb_gradient}
Let $\widehat{\Statespace}$ be a finite set, $\BackwardPolicy$ a Markov kernel on $\widehat{\Statespace}$ and $\lambda$ a $\BackwardPolicy$-invariant probability with $\lambda(s)>0$ for every $s\in\widehat{\Statespace}$, and let $E:=\{(s,s')\in\widehat{\Statespace}^2:\BackwardPolicy(s'\rightarrow s)>0\}$, which the edge lift $K_2$ of Definition \ref{def:edge_lift} maps into itself and on which the edge measure $\lambda_2$ is positive; write $K_2^{\lambda_2}$ for the $\lambda_2$-reversal of $K_2$ (Lemma \ref{lem:lift_wellposed}). Identify a measure on $E$ with its vector of weights and write $\langle\mu\mid\delta\rangle_{\lambda_2}:=\sum_{e\in E}\mu(e)\delta(e)/\lambda_2(e)$. For an edge flow $\mu\in(0,+\infty)^E$ write $F(s):=\sum_{s':(s,s')\in E}\mu(s,s')$ for its first marginal, $\ForwardPolicy^\mu(s\rightarrow s'):=\mu(s,s')/F(s)$, and $\widehat r(e):=(\mu K_2)(e)/\mu(e)$ for $e\in E$. Let $\ell\geq1$, let $\widehat\nu$ be a finite measure on $\widehat{\Statespace}^{\ell+1}$ charging only trajectories $\tau=(x_0,\dots,x_\ell)$ whose edges $e_k:=(x_{k-1},x_k)$ all lie in $E$, let $\widehat\nu_k$ be the law of $e_k$ under $\widehat\nu$, let $g:\mathbb R_+^*\rightarrow\mathbb R$ be continuously differentiable, and define the TB loss with frozen backward policy
\begin{equation}
    \mathcal L_{\mathrm{TB},g,\widehat\nu}(\mu) := \int g\left(\rho(\tau)\right) d\widehat\nu(\tau),
    \qquad \rho(\tau) := \prod_{k=1}^{\ell} \widehat r(e_k).
\end{equation}
Then $F>0$ on $\widehat{\Statespace}$, $\widehat r$ takes its values in $(0,+\infty)$, and the telescoping identity
$\rho(\tau) = \frac{F(x_\ell)\prod_{k=1}^\ell\BackwardPolicy(x_k\rightarrow x_{k-1})}{F(x_0)\prod_{k=1}^\ell\ForwardPolicy^\mu(x_{k-1}\rightarrow x_k)}$
holds, which identifies $\rho$ with the (inverted) trajectory balance ratio; on the loop closure of a marked graph (Theorem \ref{theo:universality_graphs}), $F(x_0)$ plays the role of $Z$ when $x_0=s_0$, $Z$ being the target mass of Theorem \ref{theo:universality_graphs}\emph{(3)}, and the target enters through the frozen row at the sink, $\widehat\BackwardPolicy(s_f\rightarrow\cdot)=\target/Z$. Moreover $\mathcal L_{\mathrm{TB},g,\widehat\nu}$ is $C^1$ on $(0,+\infty)^E$, and its gradient in $\langle\cdot\mid\cdot\rangle_{\lambda_2}$ is
\begin{equation}\label{eq:tb_gradient}
    \nabla^{\lambda_2}_{\mu} \mathcal L_{\mathrm{TB},g,\widehat\nu}
    = \sum_{k=1}^{\ell} \left(K_2^{\lambda_2} - \widehat r\right)\left[\psi_k\, \lambda_2\right],
\end{equation}
\begin{equation}
    \psi_k := \frac{1}{\widehat r}\,\mathbb E_{\tau\sim\widehat\nu}\bigl( g'(\rho)\rho \mid e_k \bigr)\frac{d\widehat\nu_k}{d\mu}\quad\text{where }\widehat\nu_k>0,\qquad \psi_k:=0\quad\text{where }\widehat\nu_k=0.
\end{equation}
For $\ell = 1$ this is the formula of Corollary \ref{cor:db_gradient}.
\end{proposition}

\begin{proof}
By Definition \ref{def:edge_lift}, $K_2$ moves $(s,s')$ to $(z,s)$ with probability $\BackwardPolicy(s\rightarrow z)$, and $(z,s)\in E$ exactly when that probability is positive, so $K_2$ maps $E$ into itself; $\lambda_2(s,s')=\BackwardPolicy(s'\rightarrow s)\lambda(s')$ is positive on $E$, $\lambda$ being positive. For $e=(s,s')\in E$ and $\mu\in(0,+\infty)^E$, $(\mu K_2)(e)=\sum_{z:(s',z)\in E}\mu(s',z)\,\BackwardPolicy(s'\rightarrow s)=F(s')\BackwardPolicy(s'\rightarrow s)$.
\begin{itemize}
    \item \emph{$F>0$, $\widehat r>0$, and the telescoping identity holds.} Indeed, for $s\in\widehat{\Statespace}$, $\lambda(s)=\sum_z\lambda(z)\BackwardPolicy(z\rightarrow s)>0$ provides a $z$ with $(s,z)\in E$, so $F(s)\geq\mu(s,z)>0$. Hence, for $e=(s,s')\in E$, $\widehat r(e)=\frac{F(s')\BackwardPolicy(s'\rightarrow s)}{F(s)\ForwardPolicy^\mu(s\rightarrow s')}>0$, using $\mu(s,s')=F(s)\ForwardPolicy^\mu(s\rightarrow s')$. In the product over the edges $e_k=(x_{k-1},x_k)$ of $\tau$ the factors $F(x_1),\dots,F(x_{\ell-1})$ cancel, which is the telescoping identity.
    \item \emph{$\mathcal L_{\mathrm{TB},g,\widehat\nu}$ is $C^1$ on $(0,+\infty)^E$.} Indeed, each $F(s)$ is a linear form in $\mu$, positive on $(0,+\infty)^E$ since $F>0$ there, so each $\widehat r(e)$ and each $\rho(\tau)$ is a $C^\infty$ function of $\mu$ with values in $(0,+\infty)$; $g$ being $C^1$ there, each $g(\rho(\tau))$ is $C^1$, and $\mathcal L_{\mathrm{TB},g,\widehat\nu}$ is a finite sum of them.
    \item \emph{The gradient is \eqref{eq:tb_gradient}.} Indeed, let $\delta\in\mathbb R^E$. As $\mu\mapsto\mu K_2$ is linear, the derivative of $\log\widehat r(e)$ along $\delta$ is $\frac{(\delta K_2)(e)}{(\mu K_2)(e)}-\frac{\delta(e)}{\mu(e)}$, and that of $\rho(\tau)$ is $\rho(\tau)$ times the sum of these over the edges of $\tau$. Set $a_k(e):=\sum_{\tau:e_k=e}\widehat\nu(\tau)\,g'(\rho(\tau))\rho(\tau)$, which vanishes where $\widehat\nu_k$ does and equals $\mathbb E_{\tau\sim\widehat\nu}\bigl(g'(\rho)\rho\mid e_k=e\bigr)\widehat\nu_k(e)$ elsewhere, so that $\psi_k=\frac{a_k}{\widehat r\mu}=\frac{a_k}{\mu K_2}$ on $E$. The derivative of $\mathcal L_{\mathrm{TB},g,\widehat\nu}$ along $\delta$ is therefore
\begin{eqnarray}
    \sum_{k=1}^\ell\sum_{e\in E}a_k(e)\left[\frac{(\delta K_2)(e)}{(\mu K_2)(e)}-\frac{\delta(e)}{\mu(e)}\right]
    &=& \sum_{k=1}^\ell \left\langle \psi_k\lambda_2 \,\middle|\, \delta K_2\right\rangle_{\lambda_2} - \left\langle \widehat r\psi_k\lambda_2 \,\middle|\, \delta\right\rangle_{\lambda_2} \nonumber \\
    &=& \sum_{k=1}^\ell \left\langle \left(\psi_k\lambda_2\right) K_2^{\lambda_2} - \widehat r\psi_k\lambda_2 \,\middle|\, \delta\right\rangle_{\lambda_2},\nonumber
\end{eqnarray}
the second equality by Lemma \ref{lem:lift_wellposed} and Lemma \ref{lem:adjoint}\emph{(3)}, applied to $K_2$ and its invariant measure $\lambda_2$; this is \eqref{eq:tb_gradient}.
    \item \emph{For $\ell=1$ the formula is that of Corollary \ref{cor:db_gradient}.} Indeed, at $\ell=1$, $\rho(\tau)=\widehat r(e_1)$, so $\psi_1=g'(\widehat r)\frac{d\widehat\nu_1}{d\mu}$ and \eqref{eq:tb_gradient} is the display of that corollary for the training distribution $\widehat\nu_1$.
\end{itemize}
\end{proof}

\begin{proposition}[TB Hessian]\label{prop:tb_hessian}
In the setting of Proposition \ref{prop:tb_gradient}, let $g(x) = (\log x)^2$ and let $\widehat\nu$ be the window measure
\begin{equation*}
    \widehat\nu(x_0,\dots,x_\ell):=\lambda(x_\ell)\prod_{k=1}^{\ell}\BackwardPolicy(x_k\rightarrow x_{k-1}).
\end{equation*}
Write $P$ for the density action $h\mapsto\frac{d((h\lambda_2)K_2)}{d\lambda_2}$ of $K_2$ on $L^2(\lambda_2)$, $Q:=P^\dagger$, $A:=P-I$, $\Pi_2$ for the $\lambda_2$-mean projection, $\widehat\beta_m:=\|K_2^m-\Pi_2\|_{L^2(\lambda_2)}$ as in Lemma \ref{lem:lift_mixing}, $\widehat B' := \sum_{m\geq1}\widehat\beta_m$, and
\begin{equation*}
    \Theta_\ell := I + \sum_{m=1}^{\ell-1}\left(1-\tfrac m\ell\right)\left(Q^m + (Q^\dagger)^m\right),\qquad
    H_{\mathrm{TB}} := \ell\, g''(1)\, A^\dagger\, \Theta_\ell\, A,\qquad H_{\mathrm{DB}}:=g''(1)\,A^\dagger A .
\end{equation*}
Then every edge marginal $\widehat\nu_k$ is $\lambda_2$, and for every norm $\|\cdot\|$ on $\mathbb R^E$, as $h\rightarrow0$ in $\mathbb R^E$,
\begin{equation}
    \mathcal L_{\mathrm{TB},g,\widehat\nu}\bigl((1+h)\lambda_2\bigr) = \tfrac12\langle h, H_{\mathrm{TB}}h\rangle_{L^2(\lambda_2)} + O\bigl(\|h\|^3\bigr).
\end{equation}
The operator $\Theta_\ell$ is self-adjoint and positive semi-definite on $L^2(\lambda_2)$, and $\Theta_\ell\preceq(1+2\widehat B')\,I$ on the mean-zero functions; consequently
\begin{equation}
    0\;\preceq\; H_{\mathrm{TB}} \;\preceq\; \ell\,(1 + 2\widehat B')\, H_{\mathrm{DB}},
\end{equation}
where $H_{\mathrm{DB}}$ is $H_{\mathrm{TB}}$ at $\ell=1$, at which $\mathcal L_{\mathrm{TB},g,\widehat\nu}$ is the DB loss $\mathcal L_{\mathrm{DB},g,\lambda_2}$ of Proposition \ref{prop:db_lift}. If $\widehat{\Statespace}$ has at least two states, $\widehat B'\geq1$.
\end{proposition}

\begin{proof}
All inner products and norms without index are those of $L^2(\lambda_2)$. The measure $\lambda_2$ is a probability, its second marginal being $\lambda$. By Lemma \ref{lem:lift_wellposed} and Lemma \ref{lem:adjoint}\emph{(3)}, $Q$ is the function action of $K_2$ and $P$ that of $K_2^{\lambda_2}$, both Markov kernels; hence $P\mathbf 1=\mathbf 1$, $|Ph|\leq\|h\|_{L^\infty}$, $\|P\|=\|Q\|\leq1$, and $\int Ph\,d\lambda_2=\langle h,Q\mathbf 1\rangle=\int h\,d\lambda_2$.
\begin{itemize}
    \item \emph{Every $\widehat\nu_k$ is $\lambda_2$, and under $\widehat\nu$ the edges $e_\ell,e_{\ell-1},\dots,e_1$ form a Markov chain with kernel $K_2$ started from $\lambda_2$.} Indeed, $\widehat\nu$ is the law of $(x_0,\dots,x_\ell)$ when $x_\ell\sim\lambda$ and, for $k=\ell,\dots,1$, $x_{k-1}$ is drawn from $\BackwardPolicy(x_k\rightarrow\cdot)$ given $(x_k,\dots,x_\ell)$. By the invariance of $\lambda$ every $x_k$ has law $\lambda$, so $e_k=(x_{k-1},x_k)$ has law $\BackwardPolicy(x_k\rightarrow x_{k-1})\lambda(x_k)=\lambda_2(e_k)$; and given $(e_k,\dots,e_\ell)$ the edge $e_{k-1}=(x_{k-2},x_{k-1})$ has $x_{k-2}$ drawn from $\BackwardPolicy(x_{k-1}\rightarrow\cdot)$, which is the law $K_2(e_k\rightarrow\cdot)$ of Definition \ref{def:edge_lift}.
    \item \emph{For every $u\in L^2(\lambda_2)$, $\mathbb E_{\tau\sim\widehat\nu}\bigl(\sum_{k=1}^\ell u(e_k)\bigr)^2=\ell\langle u,\Theta_\ell u\rangle$; hence $\Theta_\ell\succeq0$.} Indeed, the edges forming a stationary $K_2$-chain, $\mathbb E_{\tau\sim\widehat\nu}\bigl(u(e_k)^2\bigr)=\|u\|^2$ and, for $j<k$, $e_j$ is the state at time $k-j$ of the $K_2$-chain started from $e_k\sim\lambda_2$, so $\mathbb E_{\tau\sim\widehat\nu}\bigl(u(e_j)u(e_k)\bigr)=\langle Q^{k-j}u,u\rangle$. The difference $m=k-j\in\{1,\dots,\ell-1\}$ occurs for $\ell-m$ pairs, and $\langle Q^mu,u\rangle=\langle u,Q^mu\rangle=\frac12\langle u,(Q^m+(Q^\dagger)^m)u\rangle$, so
    \begin{equation}\label{eq:tb_window_variance}
        \mathbb E_{\tau\sim\widehat\nu}\Bigl(\sum_{k=1}^\ell u(e_k)\Bigr)^2=\ell\|u\|^2+\sum_{m=1}^{\ell-1}(\ell-m)\langle u,(Q^m+(Q^\dagger)^m)u\rangle=\ell\langle u,\Theta_\ell u\rangle .
    \end{equation}
    The operator $\Theta_\ell$ is self-adjoint by construction, and the left-hand side is non-negative.
    \item \emph{$\Theta_\ell\preceq(1+2\widehat B')\,I$ on the mean-zero functions.} Indeed, $\Pi_2$ is self-adjoint and $Q^m-\Pi_2=(P^m-\Pi_2)^\dagger$, so $\|Q^m-\Pi_2\|=\widehat\beta_m$. For $\Pi_2u=0$, $|\langle Q^mu,u\rangle|=|\langle(Q^m-\Pi_2)u,u\rangle|\leq\widehat\beta_m\|u\|^2$, and likewise for $(Q^\dagger)^m$, so
    \begin{equation}\label{eq:tb_theta_bound}
        \langle u,\Theta_\ell u\rangle\leq\Bigl(1+2\sum_{m=1}^{\ell-1}\bigl(1-\tfrac m\ell\bigr)\widehat\beta_m\Bigr)\|u\|^2\leq(1+2\widehat B')\|u\|^2\qquad\text{whenever }\Pi_2u=0.
    \end{equation}
    \item \emph{The expansion holds.} Indeed, let $\|h\|_{L^\infty}\leq\frac12$ and $\mu:=(1+h)\lambda_2\in(0,+\infty)^E$. By Lemma \ref{lem:adjoint}\emph{(3)}, $\frac{d(\mu K_2)}{d\lambda_2}=P(1+h)=1+Ph$, so $\widehat r=\frac{1+Ph}{1+h}$ and $u:=\log\widehat r=Ah+R$ with $R:=\bigl[\log(1+Ph)-Ph\bigr]-\bigl[\log(1+h)-h\bigr]$; as $|\log(1+t)-t|\leq\sum_{n\geq2}\frac{|t|^n}n\leq\frac{t^2}{2(1-|t|)}\leq t^2$ for $|t|\leq\frac12$, $|R|\leq2\|h\|_{L^\infty}^2$. For $g=(\log x)^2$, $g(\rho(\tau))=\bigl(\sum_ku(e_k)\bigr)^2$, so by \eqref{eq:tb_window_variance}
\begin{equation*}
    \mathcal L_{\mathrm{TB},g,\widehat\nu}(\mu)=\ell\langle u,\Theta_\ell u\rangle=\ell\langle Ah,\Theta_\ell Ah\rangle+\ell\langle 2Ah+R,\Theta_\ell R\rangle .
\end{equation*}
As $\|\Theta_\ell\|\leq1+2\sum_{m=1}^{\ell-1}(1-\frac m\ell)=\ell$, $\|Ah\|\leq2\|h\|_{L^\infty}$ and $\|R\|\leq2\|h\|^2_{L^\infty}\leq\|h\|_{L^\infty}$, the last term is at most $\ell^2\cdot5\|h\|_{L^\infty}\cdot2\|h\|_{L^\infty}^2=10\ell^2\|h\|^3_{L^\infty}$. Since $g''(1)=2$, $\ell\langle Ah,\Theta_\ell Ah\rangle=\frac12\langle h,H_{\mathrm{TB}}h\rangle$; all norms on $\mathbb R^E$ being equivalent, the remainder is $O(\|h\|^3)$ for each of them.
    \item \emph{$0\preceq H_{\mathrm{TB}}\preceq\ell(1+2\widehat B')H_{\mathrm{DB}}$, and at $\ell=1$ the operator $H_{\mathrm{TB}}$ is $H_{\mathrm{DB}}$ and the loss is $\mathcal L_{\mathrm{DB},g,\lambda_2}$.} Indeed, $\langle h,H_{\mathrm{TB}}h\rangle=\ell g''(1)\langle Ah,\Theta_\ell Ah\rangle\geq0$ by \eqref{eq:tb_window_variance}; $Ah$ is mean-zero, $P$ preserving $\lambda_2$-integrals, so \eqref{eq:tb_theta_bound} gives $\langle h,H_{\mathrm{TB}}h\rangle\leq\ell(1+2\widehat B')g''(1)\|Ah\|^2=\ell(1+2\widehat B')\langle h,H_{\mathrm{DB}}h\rangle$. At $\ell=1$, $\Theta_1=I$, $\widehat\nu(x_0,x_1)=\lambda(x_1)\BackwardPolicy(x_1\rightarrow x_0)=\lambda_2(x_0,x_1)$ and $\rho(\tau)=\widehat r(e_1)$, so $\mathcal L_{\mathrm{TB},g,\widehat\nu}(\mu)=\int g(\widehat r)\,d\lambda_2$, which is the loss of Proposition \ref{prop:db_lift} for the training distribution $\lambda_2$.
    \item \emph{If $\widehat{\Statespace}$ has at least two states, $\widehat B'\geq1$.} Indeed, by Lemma \ref{lem:lift_mixing}, $\widehat\beta_1=\beta_0=\|I-\Pi\|_{L^2(\lambda)}$, and $I-\Pi$ is an orthogonal projection, non-zero since the indicator of one state is not $\lambda$-almost everywhere constant, $\lambda$ being positive at every state; so $\widehat B'\geq\widehat\beta_1=1$.
\end{itemize}
\end{proof}

\begin{remark}\label{rem:tb_vs_db}
Proposition \ref{prop:tb_hessian} compares TB and DB under the same frozen backward policy and the stationary window measure. Heuristically, one trajectory performs $\ell$ detailed-balance updates at once; what is proved is one-sided, $0\preceq H_{\mathrm{TB}}\preceq\ell(1+2\widehat B')H_{\mathrm{DB}}$, so a step cap derived from the norm of $H_{\mathrm{DB}}$ remains sufficient for TB once divided by $\ell(1+2\widehat B')$, and no TB rate is claimed. On an eigenfunction of $Q$ with eigenvalue $e^{i\theta}$ the Rayleigh quotient of $\Theta_\ell$ is the Fej\'er kernel $\sum_{|m|<\ell}(1-\frac{|m|}\ell)e^{im\theta}$: it vanishes at $\theta=\pi$ for $\ell$ even, equals $\frac1\ell$ there for $\ell$ odd, and tends to $\ell$ as $\theta\rightarrow0$. On the counter-example graph \eqref{eq:counterexample_graph} with the damping $\BackwardPolicy(j\rightarrow2j)=c/(j+1)$, the condition of \citet{morozov2025revisiting} and summable mixing fail on different sets: backward trajectories have infinite expected length exactly for $1\leq c<2$, whereas the state-space mixing sum $\widehat B$ is $+\infty$ at every $c$ for which the loop-closed chain is positive recurrent (Theorem \ref{theo:doubling_unbounded}). At $c=\frac14$, where $\bar\sigma=\bar\jmath/(1-c)=6$ for a target row of mean $\bar\jmath=\frac92$ (Proposition \ref{prop:doubling_length}), the rate of Theorem \ref{theo:db_stable_frozen_full} on the truncation at an even $K\geq K_0$ satisfies $\varrho_K\leq g''(1)w_{\min}/(c_8^2K)\rightarrow0$, with $c_8$ and $K_0$ the constants of Theorem \ref{theo:doubling_main}\emph{(5)}: in the flow-matching instance at every coercivity constant \eqref{eq:coercivity} of the truncation, $\widehat B_K=\|S\|\geq c_8\sqrt K$ being the smallest of them, and the mixing sum $\sum_n\widehat\beta_n$ of the truncated state chain, when finite, being one of them (Lemma \ref{lem:sigma_mixing}) and at least $\widehat B_K$ (Lemma \ref{lem:doubling_operator}\emph{(3)}); in the detailed-balance instance at the constant $1+\widehat B_K$ that Lemma \ref{lem:lift_coercivity} supplies. Summable $L^2$-mixing, which supplies the coercivity hypothesis \eqref{eq:coercivity} of these rates (Lemma \ref{lem:sigma_mixing}), is exactly the hypothesis of Theorem \ref{theo:universality_L2} at $p=2$: heuristically, \emph{the same ergodicity property of the backward policy provides both the universality of the parameterization and the speed of its training}, and the expected-backward-length condition of \citet{morozov2025revisiting} is the $L^1$ shadow of $\widehat B<\infty$. Proposition \ref{prop:morozov_rate} replaces $\widehat B$ by the hitting-time constant $\widehat B_\sigma$ in Theorems \ref{theo:db_stable_frozen_full} and \ref{theo:local_convergence_full} and in Proposition \ref{prop:no_distant_equilibrium}\emph{(3)}, in their flow-matching instance, and by $1+\widehat B_\sigma$ in the detailed-balance instance of the two theorems.
\end{remark}

The next two propositions carry the detailed content of Theorem \ref{theo:global_dichotomy}; Remark \ref{rem:freezing} records which generators cannot freeze and Remark \ref{rem:cycle_no_stalemate} works an example.

\begin{proposition}[nonlinear freezing: local stability is sharp]\label{prop:nonlinear_freezing}
Let $0<a<b<1<c<d$, $U^-:=(a,b)$ and $U^+:=(c,d)$. Let $\chi(t):=e^{-1/t}$ for $t>0$ and $\chi(t):=0$ for $t\leq0$, let $\phi_{p,q}(t):=\chi(p-t)+\chi(t-q)$ for $p<q$, and set
\begin{equation*}
    m(t):=\frac{2\,\phi_{a,b}(t)\,\phi_{c,d}(t)}{\phi_{a,b}(1)\,\phi_{c,d}(1)}\quad(t\in\mathbb R),\qquad g(x):=\int_1^x m(t)(t-1)\,dt\quad(x>0).
\end{equation*}
Then $m$ is $C^\infty$, $0\leq m\leq C:=\frac8{\phi_{a,b}(1)\phi_{c,d}(1)}$, $m(1)=2$, and $m$ vanishes exactly on $[a,b]\cup[c,d]$; $g$ is a $C^\infty$ admissible generator --- $g(1)=0$, $g>0$ elsewhere, $g''(1)=2$, $|g(x)|\leq C(1+x^2)$ --- with $g'\equiv0$ on $[a,b]\cup[c,d]$; and:
\begin{enumerate}
    \item for every $\tilde g:\mathbb R_+^*\rightarrow\mathbb R$ continuously differentiable with locally Lipschitz derivative, positive on $U^-\cup U^+$ and with $\tilde g'\equiv0$ there, $g$ included, and every ergodic $(\widehat{\Statespace},\lambda,T)$, every $\mu$ whose ratio $r=\frac{d(\mu T)}{d\mu}$ takes values in $U^-\cup U^+$ $\lambda$-almost everywhere is a critical point of $\mathcal L_{\tilde g,\nu}$ for \emph{every} training measure $\nu$; on a finite state space the set of such $\mu$ is open, and $\mathcal L_{\tilde g,\nu}$ is locally constant on it, and positive when $\nu\neq0$;
    \item on the two-state chain $T(i\rightarrow j)=\frac12$, $\lambda=(\frac12,\frac12)$ --- for which $\widehat\beta_n=0$ for all $n\geq1$, $\widehat B=1$, and the rate $\varrho=2w_{\min}$ of Theorem \ref{theo:db_stable_frozen_full} is maximal --- the choice $(a,b,c,d)=(\frac58,\frac78,\frac{11}8,\frac{13}8)$ yields a nonempty open set of flows, on which every ratio deviates from $1$ by at least $\frac18$, where gradient flow and gradient descent (any step, any $\nu$) are stationary;
    \item let $\varepsilon>0$, let $\tilde g:\mathbb R_+^*\rightarrow\mathbb R$ be continuously differentiable with locally Lipschitz derivative and $|\tilde g'|\leq\varepsilon$ on $U^-\cup U^+$, let $(\widehat{\Statespace},\lambda,T)$ be ergodic with $\lambda$ a probability, let $\nu=w\lambda$ with $w\in L^\infty(\lambda)$ non-negative, and set $C':=2(1+d)\|w\|_{L^\infty(\lambda)}$; then on the set of flows $\mu=(1+h)\lambda$ with $\|h\|_{L^\infty(\lambda)}\leq\frac12$ whose ratio takes values in $U^-\cup U^+$ $\lambda$-almost everywhere, the gradient $\nabla^\lambda\mathcal L_{\tilde g,\nu}$ has norm at most $C'\varepsilon$ in $\mathcal M^2(\lambda)$, so a trajectory of the gradient flow moves a distance $\eta$ in $\mathcal M^2(\lambda)$ in a time at least $\eta/(C'\varepsilon)$ while it remains in that set.
\end{enumerate}
In particular, no bound on the time the gradient flow takes to bring the loss below a given positive level, finite and depending only on that level, on $g''(1)$, on the mixing coefficients and on $w$, holds for every admissible generator.
\end{proposition}

\begin{proof}
\emph{The properties of $m$ and $g$.} The function $\chi$ is $C^\infty$ on $\mathbb R$, and $0\leq\chi\leq1$, since $-1/t<0$ for $t>0$. Hence $\phi_{p,q}$ is $C^\infty$ with values in $[0,2]$; its two terms being non-negative, $\phi_{p,q}(t)=0$ if and only if $\chi(p-t)=\chi(t-q)=0$, that is $p\leq t\leq q$. As $1\notin[a,b]\cup[c,d]$, $\phi_{a,b}(1)\phi_{c,d}(1)>0$, so $m$ is $C^\infty$ with $0\leq m\leq C$ and $m(1)=2$, and $m(t)=0$ if and only if $\phi_{a,b}(t)=0$ or $\phi_{c,d}(t)=0$, that is $t\in[a,b]\cup[c,d]$. The integrand $t\mapsto m(t)(t-1)$ being $C^\infty$, $g$ is $C^\infty$ on $\mathbb R_+^*$ with $g'(x)=m(x)(x-1)$, which vanishes on $[a,b]\cup[c,d]$; $g(1)=0$ and $g''(1)=m'(1)\cdot0+m(1)=2$. For $x>1$, $g(x)>0$: the integrand is continuous and non-negative on $[1,x]$ and positive on $(1,\min(x,c))$, where $m>0$. For $x<1$, $g(x)=\int_x^1m(t)(1-t)\,dt>0$: the integrand is continuous and non-negative on $[x,1]$ and positive on $(\max(x,b),1)$. Finally $0\leq m\leq C$ gives $|g(x)|\leq C\bigl|\int_1^x|t-1|\,dt\bigr|=\frac C2(x-1)^2\leq C(1+x^2)$.

\emph{(1)} For a direction $\delta$ admissible in Theorem \ref{theo:first_variation_full}, the expansion $r(x,\mu+s\delta)=r(x)+s\bigl[\frac{d(\delta T)}{d\mu}-r\frac{d\delta}{d\mu}\bigr](x)+o(s)$ and dominated convergence give
\begin{equation*}
    \frac d{ds}\mathcal L_{\tilde g,\nu}(\mu+s\delta)\Big|_{s=0}=\int \tilde g'(r)\Bigl[\frac{d(\delta T)}{d\mu}-r\frac{d\delta}{d\mu}\Bigr]d\nu=0,
\end{equation*}
since $\tilde g'(r)=0$ $\nu$-almost everywhere ($\nu\ll\mu\sim\lambda$): $\mu$ is critical for every $\nu$. On a finite state space, $\mu\mapsto r(\mu)$ is continuous on $\{\mu>0\}$, so band-valued ratios persist under small perturbations and the set is open; $\tilde g'$ vanishing identically on each band makes $\tilde g$ constant on each of the two intervals, so every $\tilde g(r(x))$, hence $\mathcal L_{\tilde g,\nu}$, is locally constant --- and positive when $\nu\neq0$, since $\tilde g>0$ on the bands, so that $\tilde g(r)>0$ $\nu$-almost everywhere.

\emph{(2)} On the two-state chain, the density action is $Ph=\bigl(\int h\,d\lambda\bigr)\mathbf 1=\Pi h$, so $P^n-\Pi=0$ for every $n\geq1$: $\widehat\beta_n=0$, $\widehat B=\widehat\beta_0=1$ and $\varrho=g''(1)w_{\min}=2w_{\min}$. For $\mu^*=(1,\frac12)$: $\mu^*T=(\frac34,\frac34)$, so $r=(\frac34,\frac32)\in U^-\times U^+$ for $U^-=(\frac58,\frac78)$ and $U^+=(\frac{11}8,\frac{13}8)$. By continuity, a whole neighbourhood of $\mu^*$ has band-valued ratios, and \emph{(1)} applies: the gradient of $\mathcal L_{g,\nu}$ vanishes there for every $\nu$, so the gradient flow is constant and gradient descent is fixed for every step. Every ratio in $(\frac58,\frac78)\cup(\frac{11}8,\frac{13}8)$ deviates from $1$ by at least $\frac18$.

\emph{(3)} Let $\mu=(1+h)\lambda$ be in the set. Then $\frac12\leq1+h\leq\frac32$, so $\mu\sim\lambda$ and $\frac{d\nu}{d\mu}=\frac w{1+h}\in L^\infty(\mu)$, and $r$ takes its values in $(a,d)$; by Theorem \ref{theo:first_variation_full} with $\nuG=\lambda$ and Lemma \ref{lem:adjoint}\emph{(3)}, the density of $\nabla^\lambda\mathcal L_{\tilde g,\nu}(\mu)$ against $\lambda$ is $D=P^\dagger\psi-r\psi$ with $\psi:=\tilde g'(r)\frac w{1+h}$. Pointwise $|\psi|\leq2\|w\|_{L^\infty(\lambda)}\varepsilon$, so $\|\psi\|_{L^2(\lambda)}\leq2\|w\|_{L^\infty(\lambda)}\varepsilon$, $\lambda$ being a probability; as $\|P^\dagger\|\leq1$ and $0<r<d$, $\|D\|_{L^2(\lambda)}\leq(1+d)\|\psi\|_{L^2(\lambda)}\leq C'\varepsilon$. Hence $\|\mu_t-\mu_0\|_{\mathcal M^2(\lambda)}\leq C'\varepsilon\,t$ as long as the trajectory remains in the set, and moving a distance $\eta$ within it takes a time at least $\eta/(C'\varepsilon)$.

\emph{The last assertion.} On the instance of \emph{(2)}, where $g''(1)=2$ and $\widehat\beta_n=0$ for $n\geq1$, let $\mu^*:=(1,\frac12)$, whose ratio $(\frac34,\frac32)$ is band-valued, and let $\nu\neq0$; the gradient flow of $\mathcal L_{g,\nu}$ started at $\mu^*$ is constant, so its loss is $\mathcal L_{g,\nu}(\mu^*)>0$ at every time by \emph{(1)}, and it never falls below any level in $\bigl(0,\mathcal L_{g,\nu}(\mu^*)\bigr)$.
\end{proof}

\begin{remark}\label{rem:freezing}
\emph{(i)} Convexity excludes the mechanism: if $g$ is convex, $g'$ is nondecreasing, and $g'\equiv0$ on $[\alpha,\beta]\subset(0,1)$ would make $g$ nondecreasing on $[\beta,1]$, contradicting $g(\beta)>0=g(1)$ (symmetrically above $1$). Moreover, a convex differentiable $g$ has $g'\leq0$ on $(0,1)$, so $|g'(x_0)|\leq\varepsilon$ at a single $x_0<1$ forces $|g'|\leq\varepsilon$ on $[x_0,1]$, hence $g\leq\varepsilon(1-x_0)$ there. Real-analytic generators exclude exact freezing too, by the identity theorem: $g'\equiv0$ on an interval would make $g$ constant on $(0,+\infty)$. Heuristically, the obstruction lives between $C^\infty$ and analytic, and outside convexity. The generators admitted by \citet{brunswic2024theory} are not required to be convex, and $(\log x)^2$ is not convex, $g''(x)=2(1-\log x)/x^2$ being negative for $x>e$; it saturates one-sidedly only ($g'\rightarrow0$ at $+\infty$ but $|g'|\rightarrow\infty$ at $0^+$), and it is strictly unimodal, so it has no spurious critical point (Proposition \ref{prop:no_distant_equilibrium}\emph{(1)}). \emph{(ii)} For $\delta\in(0,\frac14)$, take $(a,b,c,d)=(1-2\delta,1-\delta,1+\delta,1+2\delta)$ in Proposition \ref{prop:nonlinear_freezing}. On the two-state chain of its item \emph{(2)}, every flow $\mu=(1,1-\eta)$ with $2\delta<\eta<\frac{4\delta}{1+4\delta}$ has ratios $r=\bigl(1-\frac\eta2,\frac{1-\eta/2}{1-\eta}\bigr)\in U^-\times U^+$, so it is frozen by item \emph{(1)}, at distance $O(\delta)$ from the balanced ray; at $\delta=\frac1{10}$, $\eta=\frac14$ gives $r=(0.875,1.1667)$. Every admissible generator $C^3$ on $[1,1+\delta]$ with $g''(1)=2$ and $g'\equiv0$ on $[1+\delta,1+2\delta]$ has $g'(1)=g'(1+\delta)=0$, so $\int_1^{1+\delta}g''=0$ with $g''(1)=2$; as $g''(x)\geq2-(x-1)\sup_{[1,1+\delta]}|g'''|$ on $[1,1+\delta]$, integrating gives $0\geq2\delta-\frac{\delta^2}2\sup_{[1,1+\delta]}|g'''|$, whence $\sup_{[1,1+\delta]}|g'''|\geq4/\delta$; and since the frozen flows above do not converge, the radius $\varepsilon_0$ of Theorem \ref{theo:local_convergence_full} must depend on $g$ beyond $g''(1)$. \emph{(iii)} The freezing configurations are invisible to the linearization at balance. With $\phi_{p,q}$ as in Proposition \ref{prop:nonlinear_freezing} and $0<\theta<\min(1-b,c-1)$, the functions $\zeta:=\phi_{a,b}\phi_{c,d}$ and $\xi:=\phi_{1-\theta,1+\theta}$ are $C^\infty$ and non-negative, with the disjoint zero sets $[a,b]\cup[c,d]$ and $[1-\theta,1+\theta]$, so $m:=2\zeta/(\zeta+\xi)$ is $C^\infty$ with values in $[0,2]$, equals $2$ on $[1-\theta,1+\theta]$ and vanishes exactly on $[a,b]\cup[c,d]$. The generator $g(x):=\int_1^xm(t)(t-1)\,dt$ is then admissible, with $g''(1)=2$ and $g'\equiv0$ on the bands, so item \emph{(1)} of that proposition applies to it, and it equals $(x-1)^2$ on $[1-\theta,1+\theta]$ --- which the generator constructed in the proposition does on no neighbourhood of $1$. \emph{(iv)} For every admissible differentiable strictly unimodal generator, every critical point of $\mathcal L_{g,\nu}$ is balanced when $\nu$ has positive density (Proposition \ref{prop:no_distant_equilibrium}\emph{(1)}); global convergence is proved, on finite state spaces, for $(\log x)^2$ and $(x-1)^2$ only (Proposition \ref{prop:no_distant_equilibrium}\emph{(3)}), a bounded strictly unimodal generator not being covered; infinite state spaces are not addressed.
\end{remark}

\begin{proposition}[strictly unimodal generators admit no distant equilibrium]\label{prop:no_distant_equilibrium}
Let $g$ be admissible and differentiable with $\mathrm{sign}\,g'(x)=\mathrm{sign}(x-1)$ for $x\neq1$ (\emph{strictly unimodal}; e.g. $(\log x)^2$ and $(x-1)^2$), let $(\widehat{\Statespace},\lambda,T)$ be ergodic, and write $u:=\frac{d\mu}{d\lambda}$ and $r:=\frac{d(\mu T)}{d\mu}$.
\begin{enumerate}
    \item \emph{Mass identity.} For any $\nu$ with $\frac{d\nu}{d\mu}>0$ $\mu$-a.e., the total mass of $\nabla^{\lambda}\mathcal L_{g,\nu}(\mu)$ is $\int g'(r)(1-r)\frac{d\nu}{d\mu}d\lambda\leq0$, with equality if and only if $\mu$ is balanced. Every critical point of $\mathcal L_{g,\nu}$ is therefore balanced: combined with Proposition \ref{prop:nonlinear_freezing}, spurious equilibria can arise when $g'$ vanishes on bands on either side of $1$, and strict unimodality excludes them.
    \item \emph{The inflation is the algorithm.} For fixed $\nu=w\lambda$ the loss is scale-invariant, so the gradient flow preserves $\|u_t\|_{L^2(\lambda)}$, while by \emph{(1)} the total mass $\mu_t(\widehat{\Statespace})$ strictly increases off balance: training is a monotone ascent of the mass on a sphere of $\mathcal M^2(\lambda)$, whose unique maximizer (Cauchy--Schwarz) is the balanced flow. The circulation-inflating force is the mass component of the gradient, and it is extinguished exactly at flow-matching: no stalemate with the frozen-policy constraint can occur.
    \item \emph{Global convergence, quantitatively.} Let $\widehat{\Statespace}$ be finite, let $\lambda$ be a probability with $\lambda_{\min}:=\min_{x}\lambda(x)>0$, let $g=(\log x)^2$ or $g=(x-1)^2$, and let $\nu=w\lambda$ with $w\geq w_{\min}>0$. Then $|g'(x)(1-x)|\geq\frac{\delta^2}2$ whenever $\delta\in(0,\frac12]$ and $|x-1|\geq\delta$, whence, for every $\delta\in(0,\frac12]$ and every $\mu\sim\lambda$,
    \begin{equation*}
        \bigl\|\nabla^{\lambda}\mathcal L_{g,\nu}(\mu)\bigr\|_{\mathcal M^2(\lambda)}\;\geq\;\frac{w_{\min}\lambda_{\min}^{1/2}}{\|u\|_{L^2(\lambda)}}\,\frac{\delta^2}2\;\lambda\bigl(|r-1|\geq\delta\bigr);
    \end{equation*}
    the gradient flow $\dot\mu_t=-\nabla^{\lambda}\mathcal L_{g,\nu}(\mu_t)$ from any $\mu_0\sim\lambda$, with $u_0:=\frac{d\mu_0}{d\lambda}$, has a unique solution $(\mu_t)_{t\geq0}$ with $\mu_t\sim\lambda$, which converges to the balanced flow of the sphere $\{\|u\|_{L^2(\lambda)}=\|u_0\|_{L^2(\lambda)}\}$, and whose ratios approach $1$ in explicit time: for every $\delta\in(0,\frac12]$ some $t\leq\mathcal L_{g,\nu}(\mu_0)/c(\delta)^2$, with $c(\delta):=w_{\min}\lambda_{\min}^{3/2}\delta^2/(2\|u_0\|_{L^2(\lambda)})$, has $|r_t-1|<\delta$ at every state, $r_t$ being the ratio of $\mu_t$. If moreover $\widehat B\geq1$ satisfies \eqref{eq:coercivity}, $\varepsilon_0\in(0,\frac12]$, $m_0:=\mu_0(\widehat{\Statespace})$ and $\delta=\delta_0:=\varepsilon_0m_0/(\widehat B\|u_0\|_{L^2(\lambda)})$, then at that time $\mu_t=m_t(1+h)\lambda$ with $m_t:=\mu_t(\widehat{\Statespace})$, $\Pi h=0$ and $\|h\|_{L^2(\lambda)}\leq\varepsilon_0$. Along the flow, $\|u_t\|_{L^2(\lambda)}=\|u_0\|_{L^2(\lambda)}$ and $\mu_t(\widehat{\Statespace})$ is non-decreasing.
\end{enumerate}
\end{proposition}

\begin{proof}
\emph{(1)} The gradient density is $D=P^\dagger\varphi-r\varphi$ with $\varphi:=g'(r)\frac{d\nu}{d\mu}$. Since $P\mathbf1=\mathbf1$ ($\lambda$ is $T$-invariant), $\int P^\dagger\varphi\,d\lambda=\langle\varphi,P\mathbf1\rangle=\int\varphi\,d\lambda$, so
\begin{equation*}
    \int D\,d\lambda=\int\varphi(1-r)\,d\lambda=\int g'(r)(1-r)\tfrac{d\nu}{d\mu}\,d\lambda .
\end{equation*}
Strict unimodality makes $g'(r)(1-r)<0$ wherever $r\neq1$, so the integrand is nonpositive and the integral vanishes if and only if $r=1$ holds $\frac{d\nu}{d\mu}\lambda$-a.e., i.e.\ $\mu$-a.e., i.e.\ $\mu T=\mu$. At a critical point $D=0$, so the mass vanishes and $\mu$ is balanced.

\emph{(2)} For fixed $\nu$, $r(c\mu)=r(\mu)$ gives $\mathcal L_{g,\nu}(c\mu)=\mathcal L_{g,\nu}(\mu)$, whence $0=\frac{d}{dc}\mathcal L(c\mu)|_{c=1}=\langle\nabla\mathcal L,\mu\rangle_{\mathcal M^2(\lambda)}=\int Du\,d\lambda$: the gradient is orthogonal to the radial direction, so $\frac d{dt}\|u_t\|^2_{L^2(\lambda)}=-2\int Du\,d\lambda=0$. Meanwhile $\frac d{dt}\mu_t(\widehat{\Statespace})=-\int D\,d\lambda>0$ off balance by \emph{(1)}. On the sphere $\|u\|=\|u_0\|$, Cauchy--Schwarz gives $\int u\,d\lambda\leq\|u\|_{L^2(\lambda)}$ with equality if and only if $u$ is constant, i.e.\ $\mu$ is balanced: the mass is a strict Lyapunov function whose unique maximizer on the sphere is the balanced flow.

\emph{(3)} For $g=(\log x)^2$, $|g'(x)(1-x)|=2|\log x|\frac{|x-1|}x$. For $x\leq1-\delta$: $|\log x|\geq\delta$ and $\frac{1-x}x\geq\delta$, so the product is $\geq2\delta^2$. For $x\geq1+\delta$: $\log x\geq\frac\delta{1+\delta}$ and $\frac{x-1}x\geq\frac\delta{1+\delta}$, so the product is $\geq\frac{2\delta^2}{(1+\delta)^2}\geq\frac{\delta^2}2$ for $\delta\leq\frac12$. For $g=(x-1)^2$, $|g'(x)(1-x)|=2(x-1)^2\geq2\delta^2\geq\frac{\delta^2}2$ whenever $|x-1|\geq\delta$. For $\mu\sim\lambda$, $u(x)^2\lambda(x)\leq\|u\|^2$ gives $u\leq\|u\|\lambda_{\min}^{-1/2}$ pointwise, so $\frac{d\nu}{d\mu}=\frac wu\geq w_{\min}\lambda_{\min}^{1/2}/\|u\|$, and since $\lambda$ is a probability measure, the Cauchy--Schwarz inequality and \emph{(1)} give
\begin{equation}\label{eq:unimodal_gradient_lower}
    \|D\|_{L^2(\lambda)}\;\geq\;\Bigl|\int D\,d\lambda\Bigr|\;\geq\;\frac{w_{\min}\lambda_{\min}^{1/2}}{\|u\|}\,\frac{\delta^2}2\,\lambda\bigl(|r-1|\geq\delta\bigr).
\end{equation}
Existence: write $\mathcal U:=\{u\in\mathbb R^{\widehat{\Statespace}}:u>0\}$, identify a flow $\mu\sim\lambda$ with its density $u\in\mathcal U$, and set $\mathcal L(u):=\mathcal L_{g,\nu}(u\lambda)$, $\mathcal L_0:=\mathcal L(u_0)$, $m_0:=\Pi u_0=\mu_0(\widehat{\Statespace})>0$ and $p_{\min}:=\min\{T(y\rightarrow z):T(y\rightarrow z)>0\}$. The chain $T$ is irreducible: if the set $C$ of states accessible from some state were a proper subset of $\widehat{\Statespace}$, then $T(x\rightarrow C)=1$ for $x\in C$, so $\lambda(C)=\sum_x\lambda(x)T(x\rightarrow C)$ would give $T(x\rightarrow C)=0$ for every $x\notin C$, $\lambda$ being positive, and $\mathbf 1_C\lambda$ would be a $T$-invariant measure that is not a multiple of $\lambda$ --- at $z\in C$, $(\mathbf 1_C\lambda)T(z)=(\lambda T)(z)=\lambda(z)$, no state outside $C$ leading into $C$, and at $z\notin C$ it vanishes, $C$ being closed under accessibility --- against ergodicity. For $u\in\mathcal U$ and $z\in\widehat{\Statespace}$, $(Pu)(z)=\sum_y\lambda(y)u(y)T(y\rightarrow z)/\lambda(z)>0$, since $\lambda(z)=\sum_y\lambda(y)T(y\rightarrow z)>0$ provides a $y$ with $T(y\rightarrow z)>0$; hence $r(u)=Pu/u$ takes its values in $(0,+\infty)$, and $r(u)$, $\mathcal L(u)$ and the gradient density $D(u)=P^\dagger\varphi-r\varphi$, $\varphi=g'(r)\frac wu$, are $C^\infty$ functions of $u\in\mathcal U$. By the Cauchy--Lipschitz theorem the equation $\dot u=-D(u)$, which is the gradient flow in densities, has a unique maximal solution $(u_t)_{t\in[0,t_{\max})}$ in $\mathcal U$. Along it:
\begin{itemize}
    \item \emph{$\|u_t\|=\|u_0\|$.} Indeed, $\langle D(u),u\rangle=\langle\varphi,Pu\rangle-\langle\varphi,r(u)u\rangle=0$ since $r(u)u=Pu$, so $\frac d{dt}\|u_t\|^2=-2\langle D(u_t),u_t\rangle=0$.
    \item \emph{$\mathcal L(u_t)\leq\mathcal L_0$ and $\Pi u_t\geq m_0$.} Indeed, $D(u)$ being the density of the gradient of $\mathcal L$ (Theorem \ref{theo:first_variation_full} with $\nuG=\lambda$), $\frac d{dt}\mathcal L(u_t)=-\|D(u_t)\|^2\leq0$; and $\frac d{dt}\Pi u_t=-\int D(u_t)\,d\lambda\geq0$ by \emph{(1)}.
    \item \emph{$r(u_t)\leq R_{\max}$}, where $R_{\max}:=e^{M}$ with $M:=\max\bigl(1,\sqrt{\mathcal L_0/(w_{\min}\lambda_{\min})}\bigr)$ for $g=(\log x)^2$, and $R_{\max}:=1+\sqrt{\mathcal L_0/(w_{\min}\lambda_{\min})}$ for $g=(x-1)^2$. Indeed, $\mathcal L(u_t)=\sum_x w(x)\lambda(x)g(r(u_t)(x))$ is a sum of non-negative terms, so $g(r(u_t)(x))\leq\mathcal L_0/(w_{\min}\lambda_{\min})$ at every $x$ since $\mathcal L(u_t)\leq\mathcal L_0$; for $(\log x)^2$ this gives $\log r(u_t)(x)\leq M$, and for $(x-1)^2$ it gives $r(u_t)(x)-1\leq\sqrt{\mathcal L_0/(w_{\min}\lambda_{\min})}$.
    \item \emph{$u_t\geq u_{\min}:=m_0\bigl(\lambda_{\min}p_{\min}/R_{\max}\bigr)^{\#\widehat{\Statespace}-1}$ at every state.} Indeed, if $T(y\rightarrow z)>0$, then $R_{\max}u_t(z)\geq r(u_t)(z)u_t(z)=(Pu_t)(z)\geq\lambda(y)u_t(y)T(y\rightarrow z)/\lambda(z)\geq\lambda_{\min}p_{\min}u_t(y)$, $\lambda(z)$ being at most $1$. Some state $y_0$ has $u_t(y_0)\geq\Pi u_t\geq m_0$, $\lambda$ being a probability; by irreducibility every state is reached from $y_0$ along at most $\#\widehat{\Statespace}-1$ transitions of positive probability; and $\lambda_{\min}p_{\min}/R_{\max}\leq1$.
\end{itemize}
Hence the solution never leaves $\mathcal K:=\{u\in\mathbb R^{\widehat{\Statespace}}:\|u\|=\|u_0\|,\ u\geq u_{\min}\}$, a compact subset of $\mathcal U$; a maximal solution of a locally Lipschitz equation that remains in a compact subset of its open domain is defined for all times, so $t_{\max}=+\infty$. Convergence: by LaSalle's principle, the trajectory staying in $\mathcal K$, its $\omega$-limit is contained in $\{\frac d{dt}\mathrm{mass}=0\}$, which by \emph{(1)} and \emph{(2)} is the single balanced point of the sphere: the flow converges to it. Quantitatively, let $\delta\in(0,\frac12]$. At a time $t$ with $\max|r_t-1|\geq\delta$ one has $\lambda(|r_t-1|\geq\delta)\geq\lambda_{\min}$, so \eqref{eq:unimodal_gradient_lower} with $\|u_t\|=\|u_0\|$ gives $\frac{d}{dt}\mathcal L(u_t)=-\|D(u_t)\|^2\leq-c(\delta)^2$. If $\max|r_t-1|\geq\delta$ held at every $t\leq t^*:=\mathcal L_0/c(\delta)^2$, then $\mathcal L(u_{t^*})\leq0$, hence $\mathcal L(u_{t^*})=0$ and $r_{t^*}\equiv1$, $w\lambda$ being positive at every state and $g$ vanishing only at $1$, a contradiction; so some $t\leq t^*$ has $|r_t-1|<\delta$ at every state. Let now $\widehat B$, $\varepsilon_0$ and $\delta_0$ be as in the statement. Then $m_0=\Pi u_0\leq\|u_0\|$ by the Cauchy--Schwarz inequality, $\lambda$ being a probability, so $\delta_0\leq\varepsilon_0\leq\frac12$; let $t$ have $|r_t-1|<\delta_0$ at every state and write $u:=u_t$. As $(P-I)u=(r_t-1)u$, $\|(I-P)u\|\leq\delta_0\|u\|=\delta_0\|u_0\|$, and \eqref{eq:coercivity} gives $\|u-\Pi u\|\leq\widehat B\delta_0\|u_0\|=\varepsilon_0m_0$. The mass $m_t=\Pi u_t$ satisfies $m_t\geq m_0>0$, so $h:=u/m_t-1$ satisfies $\mu_t=m_t(1+h)\lambda$, $\Pi h=0$ and $\|h\|=\|u-\Pi u\|/m_t\leq\varepsilon_0$. Finally, $\|u_t\|=\|u_0\|$ holds along the solution, and $\frac d{dt}\mu_t(\widehat{\Statespace})=\frac d{dt}\Pi u_t=-\int D(u_t)\,d\lambda\geq0$ by \emph{(1)}.
\end{proof}

\begin{remark}[a five-vertex cycle: no stalemate between constraint and inflation]\label{rem:cycle_no_stalemate}
Take the marked graph with vertices $\{s_0,x_1,x_2,x_3,s_f\}$, edges $s_0\rightarrow x_1$, the cycle $x_1\rightarrow x_2\rightarrow x_3\rightarrow x_1$, and $x_3\rightarrow s_f$, loop-closed as in Definition \ref{def:loop_closure}; freeze the backward policy $\BackwardPolicy(x_1\rightarrow x_3)=p$, $\BackwardPolicy(x_1\rightarrow s_0)=1-p$ with $p\in(0,1)$, all other rows being deterministic --- at $p=1$ the edge $s_0\rightarrow x_1$ carries no backward probability, at $p=0$ the edge $x_3\rightarrow x_1$ carries none, and in either case the mixing sum $\sum_n\widehat\beta_n$ is $+\infty$. The expected backward-trajectory length is $\frac3{1-p}<\infty$: the condition of \citet{morozov2025revisiting} holds; by their Proposition 3.12 the flow induced by the frozen backward policy has total flow through the internal states $F(s_f)\cdot\frac3{1-p}$, $F(s_f)$ its flow into the sink, so the balanced flow carries no exploding circulation. The invariant measure gives equal mass $q$ to the three cycle states and $(1-p)q$ to $s_0$ and $s_f$. Over-inflate the cycle: $u:=\frac{d\mu}{d\lambda}$ equals $M>1$ on $\{x_1,x_2,x_3\}$ and $1$ elsewhere. The ratios are then $r(x_1)=r(x_2)=r(s_f)=1$, $r(x_3)=p+\frac{1-p}M$ and $r(s_0)=M$; as $M\rightarrow\infty$, $r(x_3)\rightarrow p$, and heuristically the frozen-policy constraint appears as a persistent leak ratio $p$ at the junction $x_3$. For $g=(\log x)^2$ the inflation force $|g'(r)r|=2|\log r|$ equals $2\log M$ at the starved source $s_0$, is largest there over the states for every $M>1$, and increases with $M$. The configuration is not a critical point of $\mathcal L_{g,\nu}$ for any $\nu$ of positive density (Proposition \ref{prop:no_distant_equilibrium}\emph{(1)}, $r(x_3)$ being different from $1$). At fixed $\nu$ the gradient field is homogeneous of degree $-1$, so, for a flow $\mu_0\sim\lambda$ with gradient flow $(\mu_t)_{t\geq0}$ and for $\alpha>0$, the gradient flow from $\alpha\mu_0$, which exists by Proposition \ref{prop:no_distant_equilibrium}\emph{(3)}, is $t\mapsto\alpha\,\mu_{t/\alpha^2}$: a trajectory started at norm $\Theta(M)$ runs $\Theta(M^2)$ times slower than the normalized one. Numerically, the transient is of order $M^2$ --- its duration $t_M$, the first time every ratio lies within $\frac1{10}$ of $1$, has $t_M/M^2\approx0.720$ for $M=10,\dots,320$ at $p=\frac12$ and $w\equiv1$ --- after which the exponential phase of Theorem \ref{theo:local_convergence} takes over. As $p\rightarrow1$, the expected backward-trajectory length, the mixing sum $\sum_n\widehat\beta_n$ --- with $P$ the density action of the loop-closed backward chain on $L^2(\lambda)$ and $\Pi$ its $\lambda$-mean projection, the test density $v=\mathbf 1_{x_1}-\mathbf 1_{x_2}$ has $\Pi v=0$ and, for $p<1$, $\|P^3v-v\|_{L^2(\lambda)}^2\leq2(1-p)\|v\|_{L^2(\lambda)}^2$, so, $P$ being an $L^2(\lambda)$-contraction and $P^{3k}v-v=\sum_{j<k}P^{3j}(P^3v-v)$, $\widehat\beta_{3k}\geq\|P^{3k}v\|_{L^2(\lambda)}/\|v\|_{L^2(\lambda)}\geq1-k\sqrt{2(1-p)}\geq\frac12$ for $0\leq k\leq K_p:=\lfloor(8(1-p))^{-1/2}\rfloor$, whence $\sum_{n\geq0}\widehat\beta_n\geq\frac12K_p$ --- and the constant $\widehat B_\sigma$ of Proposition \ref{prop:morozov_rate} tend to $+\infty$, while the rate $\varrho$ of Theorem \ref{theo:db_stable_frozen_full} read at the mixing sum $\widehat B=\sum_n\widehat\beta_n$, the rate $\varrho_\sigma$ of Proposition \ref{prop:morozov_rate} and the constant of Corollary \ref{cor:global_lojasiewicz} tend to $0$; numerically, the smallest constant satisfying \eqref{eq:coercivity} stays below $1.62$.
\end{remark}

\begin{theorem}[full form of Theorem \ref{theo:global_dichotomy}]\label{theo:global_dichotomy_full}
Let $(\widehat{\Statespace},\lambda,T)$ be ergodic and write $r=\frac{d(\mu T)}{d\mu}$.
\begin{enumerate}
    \item If $g$ is \emph{strictly unimodal} --- $\mathrm{sign}\,g'(x)=\mathrm{sign}(x-1)$ for $x\neq1$, as for $(\log x)^2$ and $(x-1)^2$ --- then the total mass of $\nabla^{\lambda}\mathcal L_{g,\nu}$ equals $\int g'(r)(1-r)\frac{d\nu}{d\mu}d\lambda\leq0$, with equality only at balance: \emph{every critical point of $\mathcal L_{g,\nu}$ is balanced}, for every training measure with positive density, on every graph, cycles included. On a finite state space with $\nu=w\lambda$, $w\geq w_{\min}>0$, and for $g=(\log x)^2$ or $g=(x-1)^2$, the gradient flow converges to a balanced flow from \emph{every} initialization $\mu_0\sim\lambda$ (an explicit off-balance gradient lower bound).
    \item If instead $g'$ vanishes on bands on either side of $1$, then every flow with band-valued ratios is a critical point of $\mathcal L_{g,\nu}$ for \emph{every} $\nu$; admissible $C^\infty$ generators of this kind with $g''(1)=2$ exist --- the generator constructed in Proposition \ref{prop:nonlinear_freezing} is one --- and some of them freeze, already on the two-state chain with maximal mixing, an open set of flows whose every ratio deviates from $1$ by at least $\frac18$: global convergence fails for those (Proposition \ref{prop:nonlinear_freezing}\emph{(2)}).
\end{enumerate}
\end{theorem}

\begin{proof}
Item \emph{1} is Proposition \ref{prop:no_distant_equilibrium}, items \emph{(1)} and \emph{(3)}; item \emph{2} is Proposition \ref{prop:nonlinear_freezing}, items \emph{(1)} and \emph{(2)}.
\end{proof}

Heuristically, a bounded strictly unimodal generator can stall toward $r=\infty$; item \emph{1} of Theorem \ref{theo:global_dichotomy_full} does not cover it.

The constant $\widehat B$ of Lemma \ref{lem:sigma_mixing} is an operator quantity; on graphs, \citet{morozov2025revisiting} control training through a probabilistic one, the expected time spent in the graph by backward trajectories. The next proposition rewrites every flow-matching and detailed-balance rate of the main text in those variables, and is the engine of Theorem \ref{theo:training_speed}. The route is the Poisson equation at the regeneration state $s_0$: it supplies the coercivity of Lemma \ref{lem:sigma_mixing} directly from hitting times, and --- unlike the mixing route --- needs no aperiodicity.

\begin{proposition}[the training rate in the variables of Morozov et al.]\label{prop:morozov_rate}
Let $(\widehat G,\widehat\BackwardPolicy)$ be the loop closure (Definition \ref{def:loop_closure}) of a finite path-connected marked graph with a backward policy positive on its edges, $\lambda$ its invariant probability and $P$ the density action of the backward chain. For $x\in\mathcal V$ let $\sigma(x)$ be the hitting time of $s_0$ by the backward chain started at $x$ --- \emph{the time spent in the graph by the backward trajectory of $x$} --- and set
\begin{equation*}
    \sigma_*:=\max_{x}\mathbb E(\sigma\mid X_0=x),\qquad
    \bar\sigma:=\mathbb E_{x\sim\widehat\BackwardPolicy(s_f\rightarrow\cdot)}\,\mathbb E(\sigma\mid X_0=x),
\end{equation*}
$\bar\sigma$ being the expected backward-trajectory length of \citet{morozov2025revisiting}, and let $N(x)$ be the expected number of visits to $x$ of one backward trajectory drawn from the target row $\widehat\BackwardPolicy(s_f\rightarrow\cdot)$, counted from $s_f$ and stopped on reaching $s_0$ --- which it cannot leave, and which is the only state from which the backward chain enters $s_f$, the wrap being the only edge out of $s_f$ in $\widehat G$ --- so that $N(s_0)=N(s_f)=1$. Then:
\begin{enumerate}
    \item $\displaystyle\lambda(x)=\frac{N(x)}{2+\bar\sigma}$ for every $x\in\mathcal V$; in particular $\lambda(s_0)=\lambda(s_f)=\frac1{2+\bar\sigma}$;
    \item for every $h\in L^2(\lambda)$,
    \begin{equation*}
        \|h-\Pi h\|_{L^2(\lambda)}\;\leq\;\widehat B_\sigma\,\|(I-P)h\|_{L^2(\lambda)},
    \end{equation*}
    \begin{equation*}
        \widehat B_\sigma:=\sigma_*\sqrt{\frac{2+\bar\sigma}{\min_{x}N(x)}}\;;
    \end{equation*}
    \item $\widehat B_\sigma\geq1$; hence, for the flow-matching loss on $\widehat G$, $\widehat B:=\widehat B_\sigma$ meets the hypothesis of Theorems \ref{theo:db_stable_frozen_full} and \ref{theo:local_convergence_full} and of Proposition \ref{prop:no_distant_equilibrium}\emph{(3)} --- a constant $\widehat B\geq1$ satisfying \eqref{eq:coercivity} for the backward chain --- and the rate $\varrho$ of those theorems is then
    \begin{equation*}
        \varrho_\sigma\;=\;\frac{g''(1)\,w_{\min}\,\min_x N(x)}{\sigma_*^2\,\bigl(2+\bar\sigma\bigr)}.
    \end{equation*}
    For the detailed-balance loss, $\widehat B:=1+\widehat B_\sigma$ meets the hypothesis of the same two theorems for $T=K_2$, the edge lift of $\widehat\BackwardPolicy$ (Definition \ref{def:edge_lift}), with $\lambda_2$ in place of $\lambda$, and their rate is then $g''(1)w_{\min}/(1+\widehat B_\sigma)^2$; the radius and the step of Theorem \ref{theo:local_convergence_full} still depend on $\min_E\lambda_2$, through $C_\infty$.
\end{enumerate}
The constant $\widehat B_\sigma$ is finite without any aperiodicity assumption: on a leveled graph, where every trajectory has the same length $t_m$ (the autoregressive case), the loop-closed chain is periodic and the mixing sum $\sum_{n\geq0}\|P^n-\Pi\|_{L^2(\lambda)}$ is $+\infty$, while $\widehat B_\sigma\leq (t_m+1)\sqrt{(2+t_m)/\min_xN(x)}$.
\end{proposition}

\begin{proof}
\emph{(1)} The state $s_0$ is a regeneration state: an excursion of the backward chain from $s_0$ visits $s_0$, then $s_f$, then a sample $x$ of the target row, then runs the backward trajectory of $x$ until it hits $s_0$; its length is $\tau^{\mathrm{ret}}=2+\sigma(x)$, so $\mathbb E(\tau^{\mathrm{ret}}\mid X_0=s_0)=2+\bar\sigma$, and its expected number of visits to a state $y$ is $N(y)$. The occupation identity for positive recurrent irreducible chains \cite{levin2017markov} gives $\lambda(y)=\mathbb E(\#\{0\leq n<\tau^{\mathrm{ret}}:X_n=y\}\mid X_0=s_0)/\mathbb E(\tau^{\mathrm{ret}}\mid X_0=s_0)=N(y)/(2+\bar\sigma)$.

\emph{(2)} Let $Q$ be the function action of the backward chain, the $L^2(\lambda)$-adjoint of $P$. Fix $h$ and set $f:=(I-Q)h$; then $\int f\,d\lambda=0$ by invariance. Define $u(s_0):=0$ and, for $x\neq s_0$,
\begin{equation*}
    u(x):=\mathbb E\Bigl(\sum_{n=0}^{\sigma(x)-1}f(X_n)\,\Bigm|\,X_0=x\Bigr).
\end{equation*}
For $x\neq s_0$, conditioning on the first step gives $u(x)=f(x)+\mathbb E(u(X_1)\mid X_0=x)$, i.e.\ $(I-Q)u=f$ at $x$; at $s_0$, the occupation identity of \emph{(1)} applied to $f$ yields $\mathbb E\bigl(\sum_{n=0}^{\tau^{\mathrm{ret}}-1}f(X_n)\bigm| X_0=s_0\bigr)=(2+\bar\sigma)\int f\,d\lambda=0$, whence $\mathbb E(u(X_1)\mid X_0=s_0)=-f(s_0)$ and $(I-Q)u=f$ at $s_0$ as well. Since the chain is irreducible on a finite space, $h-u$ is $Q$-harmonic, hence constant, so $h-\Pi h=u-\Pi u$ and
\begin{eqnarray*}
    \|h-\Pi h\|_{L^2(\lambda)}&\leq&\|u\|_{L^2(\lambda)}\leq\|u\|_{L^\infty}\\
    &\leq&\sigma_*\|f\|_{L^\infty}\leq\frac{\sigma_*}{\sqrt{\lambda_{\min}}}\,\|f\|_{L^2(\lambda)},
\end{eqnarray*}
and $\lambda_{\min}=\min_xN(x)/(2+\bar\sigma)$ by \emph{(1)}. This proves the coercivity bound for $I-Q$; both $I-P$ and $I-Q$ preserve the mean-zero subspace and are adjoint to each other on it, so their restricted inverses have the same norm, and the bound holds for $I-P$.

\emph{(3)} By \emph{(1)}, $\widehat B_\sigma=\sigma_*\lambda_{\min}^{-1/2}$; $\lambda_{\min}\leq1$, $\lambda$ being a probability, and $\sigma_*\geq\mathbb E(\sigma\mid X_0=s_f)\geq1$, $s_f$ being distinct from $s_0$; hence $\widehat B_\sigma\geq1$. By \emph{(2)}, $\widehat B_\sigma$ satisfies \eqref{eq:coercivity} for the backward chain, and $g''(1)w_{\min}/\widehat B_\sigma^{\,2}=\varrho_\sigma$ by the definition of $\widehat B_\sigma$. For the detailed-balance loss, Lemma \ref{lem:lift_coercivity} with $C:=\widehat B_\sigma$, whose hypothesis for $\widehat\BackwardPolicy$ is \emph{(2)}, shows that $1+\widehat B_\sigma$, which is at least $1$, satisfies \eqref{eq:coercivity} for $K_2$ on $L^2(\lambda_2)$, and $g''(1)w_{\min}/(1+\widehat B_\sigma)^2$ is the corresponding rate by definition.

\emph{The last assertion.} On a leveled graph all excursions have length exactly $t_m+2$, so the chain is periodic of period $t_m+2$ and $\|P^{n}-\Pi\|$ does not tend to $0$, whence the mixing sum is $+\infty$; while $\sigma\leq t_m$ deterministically on the graph's states and $\sigma(s_f)\leq t_m+1$, so $\sigma_*\leq t_m+1$ and $\bar\sigma\leq t_m$.
\end{proof}

\begin{corollary}[global \L ojasiewicz inequality for $(\log x)^2$]\label{cor:global_lojasiewicz}
In the setting of Proposition \ref{prop:no_distant_equilibrium}\emph{(3)} --- $g=(\log x)^2$, finite state space, $\nu=w\lambda$ with $w_{\min}\leq w\leq\|w\|_{L^\infty}$ --- set $M:=\max\bigl(1,\sqrt{\mathcal L(\mu_0)/(w_{\min}\lambda_{\min})}\bigr)$ and
\begin{equation*}
\kappa:=\frac{w_{\min}\,\lambda_{\min}^{1/2}}{\|u_0\|_{L^2(\lambda)}\,\|w\|_{L^\infty}\,M}.
\end{equation*}
Then, along the gradient flow,
\begin{equation*}
    -\frac{d\mathcal L}{dt}\;\geq\;\kappa^2\,\mathcal L^2,
    \quad\text{hence}\quad
    \mathcal L(\mu_t)\leq\bigl(\mathcal L(\mu_0)^{-1}+\kappa^2 t\bigr)^{-1}.
\end{equation*}
\end{corollary}
\begin{proof}
For $x\in(0,e^{M}]$ with $M\geq1$ one has $2M|x-1|\geq x|\log x|$: for $x\leq1$ this is $x\log\frac1x\leq1-x$; for $x\in[1,2]$, $x\log x\leq x(x-1)\leq2M(x-1)$; for $x\geq2$, $x\log x\leq Mx\leq2M(x-1)$. Hence
\begin{equation*}
|g'(x)(1-x)|=2|\log x|\tfrac{|x-1|}x\geq\tfrac{g(x)}M \;\text{ on }(0,e^M].
\end{equation*}
Along the flow $g(r_t(x))\leq\mathcal L(\mu_t)/(w_{\min}\lambda_{\min})\leq M^2$, so $r_t\leq e^{M}$ everywhere, and the mass identity of Proposition \ref{prop:no_distant_equilibrium}\emph{(1)} gives
\begin{eqnarray*}
    \|\nabla\mathcal L\|\geq\Bigl|\int D\,d\lambda\Bigr|&\geq&\frac{w_{\min}}{u_{\max}}\,\frac1M\int g(r)\,d\lambda\\
    &\geq&\frac{w_{\min}}{u_{\max}\|w\|_{L^\infty}M}\,\mathcal L,
\end{eqnarray*}
with $u_{\max}\leq\|u_0\|\lambda_{\min}^{-1/2}$ on the invariant sphere; conclude with $-\mathcal L'=\|\nabla\mathcal L\|^2$.
\end{proof}

\begin{theorem}[full form of Theorem \ref{theo:training_speed}]\label{theo:training_speed_full}
In the setting above, let $g=(\log x)^2$ --- the generator of the practical DB and TB losses --- let $\nu=w\lambda$ with $w\geq w_{\min}>0$, and let $p_{\min}$ be the smallest positive transition probability of $\widehat\BackwardPolicy$. Then the invariant measure of the backward chain is the visit ratio,
\begin{equation}
    \lambda(x)=\frac{N(x)}{2+\bar\sigma},\qquad\text{so}\qquad \lambda_{\min}=\frac{N_{\min}}{2+\bar\sigma},
\end{equation}
and, from \emph{every} initialization $\mu_0\sim\lambda$ --- with $u_0:=\frac{d\mu_0}{d\lambda}$, $m_0:=\mu_0(\mathcal V)$, and $\mathcal L(\mu_0)^{-1}:=+\infty$ if $\mu_0$ is balanced --- the gradient flow of $\mathcal L_{g,\nu}$ converges to the balanced flow of its sphere in the two phases \emph{1} and \emph{2}, and gradient descent with a step chosen from the initialization converges to a balanced flow as \emph{3} states:
\begin{enumerate}
    \item \emph{global phase} --- the loss obeys a \L ojasiewicz inequality $-\dot{\mathcal L}\geq\kappa^2\mathcal L^2$, hence decays at the universal rate
    \begin{equation}\label{eq:global_phase}
        \mathcal L(\mu_t)\;\leq\;\bigl(\mathcal L(\mu_0)^{-1}+\kappa^2t\bigr)^{-1},
        \qquad
        \kappa:=\frac{w_{\min}\,\lambda_{\min}^{1/2}}{\|u_0\|_{L^2(\lambda)}\,\|w\|_{L^\infty}\,M},
    \end{equation}
    with $M:=\max\bigl(1,\sqrt{\mathcal L(\mu_0)/(w_{\min}\lambda_{\min})}\bigr)$;
    \item \emph{local phase} --- let $\varepsilon_0$ be the radius of Theorem \ref{theo:local_convergence_full} for $g=(\log x)^2$ at $a=\frac12$, with $\widehat B_\sigma$ in place of $\widehat B$ and $C_\infty=\lambda_{\min}^{-1/2}$; at some time
    \begin{equation*}
        t_1\;\leq\;T_0:=\frac{4\,\mathcal L(\mu_0)\,\|u_0\|_{L^2(\lambda)}^6\,\sigma_*^{\,4}}{w_{\min}^2\,\varepsilon_0^4\,m_0^4\,\lambda_{\min}^{5}},
    \end{equation*}
    the flow rescaled to unit mass enters the neighbourhood of Theorem \ref{theo:local_convergence_full}, and from $t_1$ on the flow converges exponentially at rate $\varrho_\sigma/(2m_1^2)$ --- $m_1\in[m_0,\|u_0\|_{L^2(\lambda)}]$ the flow's mass at $t_1$, so exactly $\varrho_\sigma/2$ once the entry point is rescaled to unit mass --- where
    \begin{equation}
        \varrho_\sigma\;=\;\frac{g''(1)\,w_{\min}\,\lambda_{\min}}{\sigma_*^{\,2}}\;=\;\frac{2\,w_{\min}\,N_{\min}}{\sigma_*^{\,2}\,(2+\bar\sigma)} ;
    \end{equation}
    \item \emph{gradient descent} --- write $\|\cdot\|:=\|\cdot\|_{L^2(\lambda)}$, write $D(u)$ for the density against $\lambda$ of $\nabla^\lambda\mathcal L_{g,\nu}(u\lambda)$ when $u>0$, and set $u_{\min}:=m_0\bigl(\lambda_{\min}\,p_{\min}\,e^{-M}\bigr)^{\#\mathcal V-1}$. There is an explicit $\gamma_*>0$ such that, for every step $0<\gamma\leq\gamma_*$, there is an explicit integer $k_0(\gamma)$ and the gradient descent $u_{k+1}:=u_k-\gamma D(u_k)$ is well defined and satisfies:
    \begin{enumerate}
        \item for every $k\geq0$, $u_k\geq u_{\min}$ pointwise, $\mathcal L(u_{k+1}\lambda)\leq\mathcal L(u_k\lambda)-\frac\gamma2\|D(u_k)\|^2$, $\Pi u_{k+1}\geq\Pi u_k$ and $\|u_{k+1}\|^2=\|u_k\|^2+\gamma^2\|D(u_k)\|^2\leq2\|u_0\|^2$;
        \item for every $k\geq0$, $\mathcal L(u_k\lambda)\leq\bigl(\mathcal L(\mu_0)^{-1}+k\gamma\kappa^2/4\bigr)^{-1}$;
        \item for every $k\geq k_0(\gamma)$, $\|u_k/\Pi u_k-1\|\leq\varepsilon_0$, with $\varepsilon_0$ as in \emph{2}, and
        \begin{equation*}
            \|u_{k+1}-\Pi u_{k+1}\|\;\leq\;\Bigl(1-\frac{\gamma\,\varrho_\sigma}{4\,(\Pi u_{k_0(\gamma)})^2}\Bigr)\|u_k-\Pi u_k\|,
        \end{equation*}
        where $m_0\leq\Pi u_{k_0(\gamma)}\leq\sqrt2\,\|u_0\|$;
        \item $\Pi u_k$ increases to some $m_\infty\leq\sqrt2\|u_0\|$, $m_\infty\lambda$ is balanced, and for every $k\geq k_0(\gamma)$, $m_\infty-\Pi u_k\leq\|u_k-\Pi u_k\|^2/\Pi u_{k_0(\gamma)}$ and $\|u_k-m_\infty\|\leq\frac98\|u_k-\Pi u_k\|$.
    \end{enumerate}
    No step bound independent of the initialization exists: for every $\gamma>0$ and every $\mu_0\sim\lambda$ that is not balanced, $su_0-\gamma D(su_0)$ takes a negative value for every small enough $s>0$.
\end{enumerate}
Every constant's dependence on the graph and the backward policy runs through $\sigma_*$, $\bar\sigma$ and $N_{\min}$ alone, except that $u_{\min}$ and $\gamma_*$ depend in addition on $p_{\min}$ and on $\#\mathcal V$; the initialization enters through $\|u_0\|$, $m_0$ and $\mathcal L(\mu_0)$, and the training measure through $w_{\min}$ and $\|w\|_{L^\infty}$. No mixing, spectral-gap or aperiodicity hypothesis is used.
\end{theorem}

\begin{proof}
The three ingredients are in place and the proof is their assembly. The backward chain of a finite path-connected marked graph with a backward policy positive on the edges of its loop closure is irreducible with a positive invariant probability $\lambda$, and the balanced flow exists and is unique up to scale, by Theorem \ref{theo:universality_graphs}\emph{(1)--(2)}; \eqref{eq:occupation} is Proposition \ref{prop:morozov_rate}\emph{(1)}, which also turns $\lambda_{\min}$ into $N_{\min}/(2+\bar\sigma)$ wherever it occurs below.

\emph{Global phase.} The generator $g=(\log x)^2$ is strictly unimodal, so Proposition \ref{prop:no_distant_equilibrium}\emph{(3)} applies: the gradient flow from any $\mu_0\sim\lambda$ exists for all times and converges to the balanced flow of its sphere. Inequality \eqref{eq:global_phase}, with the stated $M$ and $\kappa$, is Corollary \ref{cor:global_lojasiewicz}.

\emph{Local phase.} By Proposition \ref{prop:morozov_rate}\emph{(3)}, $\widehat B_\sigma\geq1$ satisfies \eqref{eq:coercivity} for the backward chain, the rate of Theorem \ref{theo:local_convergence_full} with $\widehat B_\sigma$ in place of $\widehat B$ is $\varrho_\sigma$, which is \eqref{eq:morozov_rate} since $g''(1)=2$ for $g=(\log x)^2$, and $\widehat B_\sigma=\sigma_*\lambda_{\min}^{-1/2}$ by \emph{(1)} of that proposition; and $\varepsilon_0\leq\lambda_{\min}^{1/2}/32\leq\frac12$ by Theorem \ref{theo:local_convergence_full} at $a=\frac12$ and $C_\infty=\lambda_{\min}^{-1/2}$. Proposition \ref{prop:no_distant_equilibrium}\emph{(3)} with $\widehat B=\widehat B_\sigma$ and $\delta=\delta_0:=\varepsilon_0m_0/(\widehat B_\sigma\|u_0\|)$ provides a time $t_1\leq\mathcal L(\mu_0)/c(\delta_0)^2$ at which $\mu_{t_1}=m_1(1+h)\lambda$, with $m_1:=\mu_{t_1}(\mathcal V)$, $\Pi h=0$ and $\|h\|\leq\varepsilon_0$; and
\begin{equation*}
    \frac{\mathcal L(\mu_0)}{c(\delta_0)^2}\;=\;\frac{4\,\mathcal L(\mu_0)\,\|u_0\|^2}{w_{\min}^2\,\lambda_{\min}^3\,\delta_0^4}\;=\;T_0,
\end{equation*}
the last equality by $\widehat B_\sigma=\sigma_*\lambda_{\min}^{-1/2}$. The same proposition gives $m_1\geq m_0$, the mass being non-decreasing, and $m_1=\Pi u_{t_1}\leq\|u_{t_1}\|=\|u_0\|$ by the Cauchy--Schwarz inequality, $\lambda$ being a probability. By its uniqueness clause the trajectory after $t_1$ is the gradient flow started at $\mu_{t_1}$. For $s>0$ the ratio of $s\mu$ is that of $\mu$, so by Theorem \ref{theo:first_variation_full} with $\nuG=\lambda$, $\nabla^\lambda\mathcal L_{g,\nu}(s\mu)=\nabla^\lambda\mathcal L_{g,\nu}(\mu)/s$, $\frac{d\nu}{d(s\mu)}$ being $\frac1s\frac{d\nu}{d\mu}$; hence $\theta\mapsto\mu_{t_1+m_1^2\theta}/m_1$ is the gradient flow started at $(1+h)\lambda$. Theorem \ref{theo:local_convergence_full} gives the exponential convergence of that flow at rate $\varrho_\sigma/2$ in $\theta$, that is at rate $\varrho_\sigma/(2m_1^2)$ in $t=t_1+m_1^2\theta$.

\emph{Gradient descent.} For $u>0$ write $r(u):=Pu/u$, where $(Pu)(z)=\sum_y\lambda(y)u(y)\widehat\BackwardPolicy(y\rightarrow z)/\lambda(z)$, write $\mathcal L(u):=\mathcal L_{g,\nu}(u\lambda)$ and $\psi(u):=g'(r(u))\,w/u$, and let $P^\dagger$ be the $L^2(\lambda)$-adjoint of $P$, the function action $(P^\dagger f)(y)=\sum_z\widehat\BackwardPolicy(y\rightarrow z)f(z)$ (Lemma \ref{lem:adjoint}\emph{(3)}); by Theorem \ref{theo:first_variation_full} with $\nuG=\lambda$, $D(u)=P^\dagger\psi(u)-r(u)\psi(u)$. The generator $g=(\log x)^2$ is $C^\infty$ on $(0,\infty)$ with $g(1)=g'(1)=0$ and $g''(1)=2$; let $\varepsilon_0$ and $\gamma_0$ be the constants of Theorem \ref{theo:local_convergence_full} for this $g$ at $a=\frac12$, with $\widehat B_\sigma$ in place of $\widehat B$ (Proposition \ref{prop:morozov_rate}\emph{(3)}) and $C_\infty=\lambda_{\min}^{-1/2}$, so that $\varepsilon_0\leq\lambda_{\min}^{1/2}/32\leq\frac1{32}$ and the rate of that theorem is $\varrho_\sigma$. Set $\theta:=\lambda_{\min}p_{\min}e^{-M}$, $M':=M+\ln3$, $b_1:=2M'e^{M'}$, $b_2:=2(1+M')e^{2M'}$,
\begin{equation*}
    b_3:=\frac{4\|w\|_{L^\infty}}{u_{\min}^2}\bigl(1+e^{M'}\bigr)\bigl(b_2(1+e^{M'})+2b_1\bigr),\qquad
    \gamma_*:=\min\Bigl(\frac1{b_3},\ \frac{\|u_0\|^2}2\,\mathcal L(\mu_0)^{-1},\ \gamma_0\,m_0^2\Bigr),
\end{equation*}
\begin{equation*}
    \delta_1:=\frac{\varepsilon_0\,m_0}{\sqrt2\,\widehat B_\sigma\|u_0\|},\qquad
    k_0(\gamma):=\Bigl\lceil\frac{16}{\gamma\,\kappa^2\,\lambda_{\min}\,w_{\min}\,\delta_1^2}\Bigr\rceil
    =\Bigl\lceil\frac{32\,\|u_0\|^4\,\|w\|_{L^\infty}^2\,M^2\,\sigma_*^{\,2}}{\gamma\,w_{\min}^3\,\lambda_{\min}^3\,\varepsilon_0^2\,m_0^2}\Bigr\rceil,
\end{equation*}
the last equality by $\widehat B_\sigma=\sigma_*\lambda_{\min}^{-1/2}$, which is Proposition \ref{prop:morozov_rate}\emph{(1)}. Note that $\theta\leq1$, hence $u_{\min}\leq m_0$; that $m_0=\Pi u_0\leq\|u_0\|$; that $\widehat B_\sigma\geq1$, since $\lambda_{\min}\leq1$ and $\sigma_*\geq\mathbb E(\sigma\mid X_0=s_f)\geq1$, $s_f$ being distinct from $s_0$; and that $b_3\geq32\|w\|_{L^\infty}/u_{\min}^2$, since $b_2\geq2$ and $1+e^{M'}\geq2$. If $\mu_0$ is balanced, then $r(u_0)\equiv1$ and $g'(1)=0$ give $D(u_0)=0$, so $u_k=u_0=m_0\geq u_{\min}$ for every $k$, and \emph{(a)}--\emph{(d)} hold with $\Pi u_{k_0(\gamma)}=m_\infty=m_0$. Assume from now on that $\mu_0$ is not balanced, and fix $0<\gamma\leq\gamma_*$.
\begin{itemize}
    \item \emph{Every $u>0$ satisfies $\langle D(u),u\rangle=0$, $\Pi D(u)\leq0$ and $D(su)=D(u)/s$ for every $s>0$, and $u^+:=u-\gamma D(u)$ satisfies $\|u^+\|^2=\|u\|^2+\gamma^2\|D(u)\|^2$ and $\Pi u^+\geq\Pi u$.} Indeed, $r(u)u=Pu$ gives $\langle D(u),u\rangle=\langle\psi(u),Pu\rangle-\langle\psi(u),r(u)u\rangle=0$; $\Pi D(u)\leq0$ is Proposition \ref{prop:no_distant_equilibrium}\emph{(1)}; $r(su)=r(u)$ and $\psi(su)=\psi(u)/s$ give $D(su)=D(u)/s$; and the last two claims follow by expanding the square and by linearity of $\Pi$.
    \item \emph{Every $u>0$ with $\mathcal L(u)\leq\mathcal L(\mu_0)$ and $\Pi u\geq m_0$ satisfies $e^{-M}\leq r(u)\leq e^M$ and $u\geq u_{\min}$.} Indeed, $\lambda(y)w(y)\bigl(\log r(u)(y)\bigr)^2\leq\mathcal L(u)$ at every state $y$ gives $(\log r(u))^2\leq\mathcal L(\mu_0)/(w_{\min}\lambda_{\min})\leq M^2$. If $\widehat\BackwardPolicy(y\rightarrow z)>0$, then $\widehat\BackwardPolicy(y\rightarrow z)\geq p_{\min}$ and $\lambda(z)\leq1$ give $r(u)(z)\,u(z)=(Pu)(z)\geq\lambda(y)u(y)\widehat\BackwardPolicy(y\rightarrow z)\geq\lambda_{\min}p_{\min}u(y)$, hence $u(z)\geq\theta u(y)$. Let $y^*$ maximize $u$, so that $u(y^*)\geq\Pi u\geq m_0$. The backward chain being irreducible (Theorem \ref{theo:universality_graphs}\emph{(1)}), every state $z$ is reached from $y^*$ along at most $\#\mathcal V-1$ transitions of positive probability, and $\theta\leq1$ gives $u(z)\geq\theta^{\#\mathcal V-1}u(y^*)\geq u_{\min}$.
    \item \emph{Every $u>0$ with $\mathcal L(u)\leq\mathcal L(\mu_0)$ and $\Pi u\geq m_0$ satisfies $u^+\geq u/2$ and $\mathcal L(u^+)\leq\mathcal L(u)-\frac\gamma2\|D(u)\|^2$.} Indeed, such a $u$ satisfies $u\geq u_{\min}$ and $e^{-M}\leq r(u)\leq e^M$ by the positivity floor, and $|g'(x)|=2|\log x|/x\leq2Me^M$ for $x\in[e^{-M},e^M]$; so $|P^\dagger\psi(u)|\leq\max|\psi(u)|\leq2M\|w\|_{L^\infty}e^M/u_{\min}$ and $|r(u)\psi(u)|=2|\log r(u)|\,w/u\leq2M\|w\|_{L^\infty}/u$, whence $\gamma|D(u)|/u\leq2\gamma M\|w\|_{L^\infty}(e^M+1)/u_{\min}^2\leq\frac12$, because $\gamma\leq1/b_3$ and $b_3u_{\min}^2/(4\|w\|_{L^\infty})\geq2b_1(1+e^{M'})\geq M(e^M+1)$. For $s\in[0,\gamma]$, $u_s:=u-sD(u)$ therefore satisfies $u/2\leq u_s\leq3u/2$; $P$ being positive, $r_s:=r(u_s)$ lies in $[r(u)/3,3r(u)]\subset[e^{-M'},e^{M'}]$, where $|g'|\leq b_1$ and $|g''(x)|=2|1-\log x|/x^2\leq b_2$; and $u_s\geq u_{\min}/2$. Hence $f(s):=\mathcal L(u_s)$ is twice differentiable on $[0,\gamma]$. With $v:=-D(u)$ and $\chi_s:=(Pv-r_sv)/u_s$, one has $\frac{d}{ds}r_s=\chi_s$ and $\frac{d^2}{ds^2}r_s=-2\chi_sv/u_s$, so that
    \begin{equation*}
        f'(s)=\int w\,g'(r_s)\,\chi_s\,d\lambda,\qquad f''(s)=\int w\Bigl(g''(r_s)\,\chi_s^2-2g'(r_s)\,\frac{\chi_s\,v}{u_s}\Bigr)d\lambda,
    \end{equation*}
    and $f'(0)=\langle\psi(u),Pv-r(u)v\rangle=\langle D(u),v\rangle=-\|D(u)\|^2$. Now $|\chi_s|\leq2(P|v|+e^{M'}|v|)/u_{\min}$ pointwise and $\|P|v|\|\leq\|v\|$ by Lemma \ref{lem:adjoint}\emph{(3)}, so the triangle inequality gives $\|P|v|+e^{M'}|v|\|\leq(1+e^{M'})\|v\|$ and the Cauchy--Schwarz inequality gives $\int(P|v|+e^{M'}|v|)|v|\,d\lambda\leq(1+e^{M'})\|v\|^2$; therefore $|f''|\leq b_3\|D(u)\|^2$ on $[0,\gamma]$. Taylor's formula and $\gamma b_3\leq1$ give $\mathcal L(u^+)=f(\gamma)\leq\mathcal L(u)-\gamma\|D(u)\|^2+\frac{\gamma^2b_3}2\|D(u)\|^2\leq\mathcal L(u)-\frac\gamma2\|D(u)\|^2$, and $u^+=u_\gamma\geq u/2$.
    \item \emph{Every $u_k$ is well defined and lies in $\mathcal U_\gamma:=\bigl\{u>0:\ \mathcal L(u)\leq\mathcal L(\mu_0),\ \Pi u\geq m_0,\ \|u\|^2+2\gamma\mathcal L(u)\leq\|u_0\|^2+2\gamma\mathcal L(\mu_0)\bigr\}$, and \emph{(a)} holds.} Indeed, $u_0\in\mathcal U_\gamma$. If $u_k\in\mathcal U_\gamma$, the positivity floor and the descent inequality at $u=u_k$ give $u_k\geq u_{\min}$, $u_{k+1}\geq u_k/2>0$ and $\mathcal L(u_{k+1})\leq\mathcal L(u_k)-\frac\gamma2\|D(u_k)\|^2$, and the identities for $D$ give $\Pi u_{k+1}\geq\Pi u_k$ and
    \begin{equation*}
        \|u_{k+1}\|^2+2\gamma\mathcal L(u_{k+1})\leq\|u_k\|^2+\gamma^2\|D(u_k)\|^2+2\gamma\mathcal L(u_k)-\gamma^2\|D(u_k)\|^2,
    \end{equation*}
    so $u_{k+1}\in\mathcal U_\gamma$. Finally $\gamma\leq\frac{\|u_0\|^2}2\mathcal L(\mu_0)^{-1}$ gives $\|u_k\|^2\leq\|u_0\|^2+2\gamma\mathcal L(\mu_0)\leq2\|u_0\|^2$.
    \item \emph{Assertion \emph{(b)} holds.} Indeed, for $x\in(0,e^M]$, $x|\log x|\leq2M|x-1|$: for $x\leq1$ this is $x\log\frac1x\leq1-x$; for $1\leq x\leq2$, $x\log x\leq x(x-1)\leq2M(x-1)$; for $x\geq2$, $x\log x\leq Mx\leq2M(x-1)$. Hence $|g'(x)(1-x)|=2|\log x|\,|x-1|/x\geq g(x)/M$ on $(0,e^M]$, which contains the values of $r(u_k)$ by the positivity floor, $u_k$ lying in $\mathcal U_\gamma$. By Proposition \ref{prop:no_distant_equilibrium}\emph{(1)}, the Cauchy--Schwarz inequality, $\max u_k\leq\lambda_{\min}^{-1/2}\|u_k\|\leq\sqrt2\,\lambda_{\min}^{-1/2}\|u_0\|$ and $\mathcal L(u_k)\leq\|w\|_{L^\infty}\int g(r(u_k))\,d\lambda$,
    \begin{eqnarray*}
        \|D(u_k)\|\;\geq\;|\Pi D(u_k)|&=&\int\bigl|g'(r(u_k))\bigl(1-r(u_k)\bigr)\bigr|\,\frac w{u_k}\,d\lambda\\
        &\geq&\frac{w_{\min}}{M\max u_k}\int g(r(u_k))\,d\lambda\;\geq\;\frac\kappa{\sqrt2}\,\mathcal L(u_k).
    \end{eqnarray*}
    With \emph{(a)}, $\mathcal L(u_{k+1})\leq\mathcal L(u_k)\bigl(1-\frac{\gamma\kappa^2}4\mathcal L(u_k)\bigr)$. If $\mathcal L(u_{k+1})>0$, then $\frac{\gamma\kappa^2}4\mathcal L(u_k)<1$ and $\mathcal L(u_{k+1})^{-1}\geq\mathcal L(u_k)^{-1}\bigl(1-\frac{\gamma\kappa^2}4\mathcal L(u_k)\bigr)^{-1}\geq\mathcal L(u_k)^{-1}+\frac{\gamma\kappa^2}4$; if $\mathcal L(u_j)=0$, then $\mathcal L(u_k)=0$ for every $k\geq j$ by \emph{(a)}. Induction on $k$ gives \emph{(b)}.
    \item \emph{For every $k\geq k_0(\gamma)$, $|r(u_k)-1|\leq\delta_1\leq\frac1{32}$ pointwise and $\|u_k-\Pi u_k\|\leq\varepsilon_0m_0\leq\varepsilon_0\Pi u_k$; in particular the first claim of \emph{(c)} holds.} Indeed, $\delta_1\leq\varepsilon_0\leq\frac1{32}$ since $\widehat B_\sigma\geq1$ and $m_0\leq\|u_0\|$. By \emph{(b)} and the choice of $k_0(\gamma)$, $\mathcal L(u_k)\leq4/(k\gamma\kappa^2)\leq\lambda_{\min}w_{\min}\delta_1^2/4$, so $|\log r(u_k)|\leq\delta_1/2$ at every state, and $|r(u_k)-1|\leq e^{\delta_1/2}-1\leq\delta_1$ since $e^x-1\leq2x$ on $[0,1]$. As $u_k>0$, $(P-I)u_k=(r(u_k)-1)u_k$, whence $\|(P-I)u_k\|\leq\delta_1\|u_k\|\leq\sqrt2\,\delta_1\|u_0\|$ by \emph{(a)}; Proposition \ref{prop:morozov_rate}\emph{(2)} gives $\|u_k-\Pi u_k\|\leq\sqrt2\,\widehat B_\sigma\delta_1\|u_0\|=\varepsilon_0m_0$, and $m_0\leq\Pi u_k$ by \emph{(a)}.
    \item \emph{Assertion \emph{(c)} holds.} Indeed, put $m:=\Pi u_{k_0(\gamma)}$; by \emph{(a)} and the Cauchy--Schwarz inequality, $m_0\leq m\leq\|u_{k_0(\gamma)}\|\leq\sqrt2\|u_0\|$. Let $h_j:=u_{k_0(\gamma)+j}/m-1$ for $j\geq0$. As $D(u)=D(u/m)/m$ for $u>0$, $h_{j+1}=h_j-\frac\gamma{m^2}D(1+h_j)$, where $\frac\gamma{m^2}\leq\frac{\gamma_*}{m_0^2}\leq\gamma_0$; and $\|h_0\|=\|u_{k_0(\gamma)}-\Pi u_{k_0(\gamma)}\|/m\leq\varepsilon_0$ by the entry bound. The gradient-descent clause of Theorem \ref{theo:local_convergence_full}, whose hypotheses hold with $\widehat B_\sigma$ in place of $\widehat B$ by Proposition \ref{prop:morozov_rate}\emph{(3)}, gives $\|h_{j+1}-\Pi h_{j+1}\|\leq\bigl(1-\frac{\gamma\varrho_\sigma}{4m^2}\bigr)\|h_j-\Pi h_j\|$ for every $j\geq0$; multiplying by $m$ gives the contraction of \emph{(c)}.
    \item \emph{Assertion \emph{(d)} holds.} Indeed, put $m:=\Pi u_{k_0(\gamma)}$, so that $m_0\leq m\leq\sqrt2\|u_0\|$ by \emph{(c)}, write $\xi_k:=u_k-\Pi u_k$, and let $k\geq k_0(\gamma)$. By \emph{(c)}, $\|\xi_{k+1}\|\leq\|\xi_k\|$; moreover $\frac{\gamma\varrho_\sigma}{4m^2}\leq\frac1{64}$, since $\gamma\leq1/b_3\leq u_{\min}^2/(32\|w\|_{L^\infty})\leq m^2/(32\|w\|_{L^\infty})$ and $\varrho_\sigma\leq2w_{\min}\leq2\|w\|_{L^\infty}$; with $m^2\leq2\|u_0\|^2$, this gives $\|\xi_k\|\leq\bigl(1-\frac{\gamma\varrho_\sigma}{8\|u_0\|^2}\bigr)^{k-k_0(\gamma)}\|\xi_{k_0(\gamma)}\|\rightarrow0$. By the entry bound $\|\xi_{k_0(\gamma)}\|\leq\varepsilon_0m_0$, $\|\xi_k\|\leq\varepsilon_0m$, so $|\xi_k|\leq\lambda_{\min}^{-1/2}\varepsilon_0m\leq m/32$ pointwise and, $\Pi u_k\geq m$ by \emph{(a)}, $u_k\geq\frac{31}{32}m$. By \emph{(a)} the mass $\Pi u_k$ is nondecreasing, and it is bounded by $\|u_k\|\leq\sqrt2\|u_0\|$, so it converges to some $m_\infty\leq\sqrt2\|u_0\|$; with $\|\xi_k\|\rightarrow0$, $u_k\rightarrow m_\infty$ in $L^2(\lambda)$, and $m_\infty\lambda$ is balanced, $\lambda$ being invariant. Next, $r(u_k)\geq\frac{31}{32}$ by the entry bound $|r(u_k)-1|\leq\frac1{32}$, so $|\log r(u_k)|\leq\frac{32}{31}|r(u_k)-1|$; and $(P-I)\Pi u_k=0$ gives $r(u_k)-1=(P-I)\xi_k/u_k$. Hence $|\log r(u_k)|\leq(\frac{32}{31})^2|(P-I)\xi_k|/m$ and, as $\|P-I\|\leq2$ by Lemma \ref{lem:adjoint}\emph{(3)} and $(\frac{32}{31})^4\leq\frac65$,
    \begin{equation*}
        \mathcal L(u_k)\;\leq\;\|w\|_{L^\infty}\Bigl(\frac{32}{31}\Bigr)^4\frac{\|(P-I)\xi_k\|^2}{m^2}\;\leq\;\frac{24}5\,\|w\|_{L^\infty}\,\frac{\|\xi_k\|^2}{m^2}.
    \end{equation*}
    Summing $\|u_{i+1}\|^2=\|u_i\|^2+\gamma^2\|D(u_i)\|^2$ over $k\leq i<j$ and letting $j\rightarrow\infty$, with $\|u_j\|^2=(\Pi u_j)^2+\|\xi_j\|^2\rightarrow m_\infty^2$, gives $m_\infty^2=(\Pi u_k)^2+\|\xi_k\|^2+\gamma^2\sum_{i\geq k}\|D(u_i)\|^2$; summing the descent inequality of \emph{(a)} gives $\gamma^2\sum_{i\geq k}\|D(u_i)\|^2\leq2\gamma\mathcal L(u_k)$. With $\gamma\leq m^2/(32\|w\|_{L^\infty})$, therefore,
    \begin{equation*}
        m_\infty^2-(\Pi u_k)^2\;\leq\;2\gamma\mathcal L(u_k)+\|\xi_k\|^2\;\leq\;\Bigl(\frac3{10}+1\Bigr)\|\xi_k\|^2,
    \end{equation*}
    and $m_\infty+\Pi u_k\geq2m$ gives $m_\infty-\Pi u_k\leq\frac{13}{20}\|\xi_k\|^2/m\leq\|\xi_k\|^2/m$. Finally $\|u_k-m_\infty\|\leq\|\xi_k\|+m_\infty-\Pi u_k\leq(1+\frac{13}{20}\varepsilon_0)\|\xi_k\|\leq\frac98\|\xi_k\|$, since $\|\xi_k\|\leq\varepsilon_0m$ and $\varepsilon_0\leq\frac1{32}$.
\end{itemize}
It remains to prove the last claim of \emph{3}. Let $\gamma'>0$ and let $\mu_0\sim\lambda$ be not balanced. By Proposition \ref{prop:no_distant_equilibrium}\emph{(1)}, $\Pi D(u_0)<0$, so $D(u_0)\neq0$; as $\langle D(u_0),u_0\rangle=0$ and $u_0>0$, there is a state $y$ with $D(u_0)(y)>0$. As $D(su_0)=D(u_0)/s$, $su_0(y)-\gamma'D(su_0)(y)=su_0(y)-\gamma'D(u_0)(y)/s<0$ for every $s>0$ with $s^2<\gamma'D(u_0)(y)/u_0(y)$.

No step invoked summable mixing, a spectral gap or aperiodicity: Proposition \ref{prop:morozov_rate} derives the coercivity from hitting times at the regeneration state $s_0$, which a finite irreducible chain possesses whether or not it is aperiodic.
\end{proof}

\begin{remark}[the speed is carried by visit ratios, then by mixing]\label{rem:visit_ratio}
Convergence itself uses no mixing constant (Theorem \ref{theo:training_speed_full}). Heuristically, the finiteness of $\widehat B_\sigma$ on a periodic loop closure (Proposition \ref{prop:morozov_rate}) is a convenience rather than a gain --- periodicity is removable by lazification or by an aperiodic perturbation of the wrap, each of which changes the loss, after which some finite mixing constant yields convergence --- and the constants matter for the \emph{speed}, the two phases of the dynamics reading them in two different places. \emph{Far from balance}, the mass identity gives a lower bound state by state: in the setting of Theorem \ref{theo:training_speed_full}, with $u_{\max}:=\max_x\frac{d\mu}{d\lambda}(x)$, for a region $U$ on which $|r-1|\geq\delta$ with $\delta\in(0,\frac12]$, the contribution of $U$ to the gradient mass is
\begin{equation*}
    \int_U|g'(r)(1-r)|\tfrac{d\nu}{d\mu}\,d\lambda\;\geq\;\frac{w_{\min}}{u_{\max}}\cdot\frac{\delta^2}2\cdot\frac{\sum_{x\in U}N(x)}{2+\bar\sigma}\,:
\end{equation*}
\emph{the contribution of an out-of-equilibrium state to this lower bound is proportional to its visit ratio $N(x)/(2+\bar\sigma)$}. Rarely visited unbalanced states are the slow modes of the global phase in the sense of this lower bound, which suggests, heuristically, equilibration on times of order $(2+\bar\sigma)/N(x)$ --- a lower bound on the gradient mass bounds that time from neither side; Corollary \ref{cor:global_lojasiewicz} aggregates the mass identity into the global $\mathcal L^{-1}$-decay in time. \emph{Near balance} the mixing constant takes over: the exponential rate is $\varrho=g''(1)w_{\min}/\widehat B^2$ for the linearized gradient flow (Theorem \ref{theo:db_stable_frozen_full}) and $\varrho/2$ for the nonlinear one (Theorem \ref{theo:local_convergence_full}). Heuristically, the visitation measure no longer discriminates there, all ratios being close to $1$, while the rate is spectral, sensitive to the \emph{shape} of the residual perturbation, which $\widehat B$ measures and visit counts do not. The constant $\widehat B_\sigma$ of Proposition \ref{prop:morozov_rate}, built from times spent in the graph and visit counts, bounds the coercivity constant --- $\widehat B_\sigma\geq\|S\|$ wherever the operator $S$ of Lemma \ref{lem:sigma_mixing} is defined --- and not the mixing sum $\widehat B$, which is $+\infty$ on a leveled graph where $\widehat B_\sigma$ is finite.
\end{remark}


\subsection{Universality on Finite Graphs: the Frozen-Backward Family}\label{app:universality_graphs}

This subsection specializes the universality theory to the discrete setting in which generative flows were introduced \cite{bengio2021flow}, and makes precise the statement announced in the main text: on a finite path-connected state graph, the parameterization of generative flows by a state flow with a frozen backward policy \cite{morozov2025revisiting} attains \emph{exact} flow matching, and the balanced flow is unique up to scale.

Throughout, a \emph{marked graph} $G=(\mathcal V,\mathcal E)$ is a finite directed graph with two distinguished vertices: a source $s_0$ with no incoming edge and a sink $s_f$ with no outgoing edge. Its internal states are $\Statespace:=\mathcal V\setminus\{s_0,s_f\}$, and an internal state $x$ is \emph{terminating} if $x\rightarrow s_f\in\mathcal E$, the edge $x\rightarrow s_f$ being then a \emph{terminating edge}. A \emph{backward policy} on $G$ is a collection of probability distributions $\BackwardPolicy(s\rightarrow\cdot)$, one for each $s\in\mathcal V\setminus\{s_0\}$, supported on the in-neighbours of $s$: $\BackwardPolicy(s\rightarrow s')>0$ only if $s'\rightarrow s\in \mathcal E$.

\begin{definition}[path-connectedness]\label{def:path_connected}
A marked graph $G$ is \emph{path-connected} if for any $s\in \mathcal V$ there exists a path $$s_0\rightarrow s_1 \rightarrow \cdots \rightarrow s\rightarrow \cdots \rightarrow s_f$$ from $s_0$ to $s_f$ via $s$.
\end{definition}

\begin{definition}[loop closure]\label{def:loop_closure}
The \emph{loop closure} of a marked graph $G$ is $\widehat G := (\mathcal V,\ \mathcal E\cup\{s_f\rightarrow s_0\})$. A backward policy $\BackwardPolicy$ on $G$ extends uniquely to a backward policy $\widehat\BackwardPolicy$ on $\widehat G$ by $\widehat\BackwardPolicy(s_0\rightarrow s_f):=1$, the wrap edge $s_f\rightarrow s_0$ being the only incoming edge of $s_0$ in $\widehat G$. The \emph{backward chain} of $(G,\BackwardPolicy)$ is the Markov chain on $\mathcal V$ with kernel $\widehat\BackwardPolicy$.
\end{definition}

\begin{equation}\label{eq:loop_closure}
        \xymatrix{
        s_0 \ar[r]&\Statespace \ar@(ur,ul)\ar[r] & s_f &&
        s_0 \ar[r]&\Statespace \ar@(ur,ul)\ar[r]& s_f \ar `dr_l[ll]`_u[ll]_{\mathrm{wrap}} [ll]
        }
\end{equation}

\begin{theorem}[universality on finite graphs with frozen backward policy]\label{theo:universality_graphs}
Let $G=(\mathcal V,\mathcal E)$ be a finite path-connected marked graph endowed with the counting measure $\mu$, let $\BackwardPolicy$ be a backward policy on $G$ such that $\BackwardPolicy(s\rightarrow s')>0$ for every edge $s'\rightarrow s$ of $G$, and let $\widehat\BackwardPolicy$ be its loop closure. Fix a total mass $Z>0$. Then:
\begin{enumerate}
    \item the backward chain $\widehat\BackwardPolicy$ is irreducible on $\mathcal V$ and admits a unique invariant probability measure $\lambda$; moreover $\lambda>0$ everywhere and $\lambda(s_0)=\lambda(s_f)$;
    \item the frozen-backward family on the loop closure,
    \begin{equation*}
        \Theta_{\BackwardPolicy}:=\left\{(\ForwardPolicy,\fout)\ :\ (\fout\mu)\otimes \ForwardPolicy = \widehat\BackwardPolicy\otimes(\fout\mu)\right\},
    \end{equation*}
    is, apart from the trivial element $\fout\equiv0$, exactly the ray $\fout\mu = c\,\lambda$, $c>0$, with $\ForwardPolicy=\widehat\BackwardPolicy^{\lambda}$ the $\lambda$-reversal of $\widehat\BackwardPolicy$ (Lemma \ref{lem:adjoint}), given on edges by $\ForwardPolicy(s\rightarrow s') = \lambda(s')\widehat\BackwardPolicy(s'\rightarrow s)/\lambda(s)$;
    \item for $c=Z/\lambda(s_0)$, cutting the wrap edge turns this element into a generative flow on $G$ satisfying the flow-matching constraint \eqref{equ:FM_const} \emph{exactly}, with initial flow of total mass $Z$ carried by the edges out of $s_0$ and terminal flow
    \begin{equation*}
        \Fterm = \target := Z\, \BackwardPolicy(s_f\rightarrow\cdot),
    \end{equation*}
    positive on every terminating state; in particular its sampler satisfies $s_\tau\sim \target/Z$ by Theorem \ref{theo:sampling_theorem}.
\end{enumerate}
Consequently, for every target $\target$ of mass $Z$ carried by the terminating states of $G$ and positive at each of them, freezing the sink row $\BackwardPolicy(s_f\rightarrow\cdot):=\target/Z$ of a backward policy $\BackwardPolicy$ positive on the edges of $G$ produces a frozen-backward family whose unique balanced flow samples $\target/Z$ exactly: the universality infimum vanishes and is attained, simultaneously for every $p\in[1,+\infty]$ --- the family is strongly universal over terminal targets.
\end{theorem}

\begin{proof}
\emph{(1)} Every vertex of $\widehat G$ has an incoming edge: internal vertices and $s_f$ lie on source-to-sink paths, and the wrap edge feeds $s_0$; so $\widehat\BackwardPolicy$ is a genuine Markov kernel on $\mathcal V$. Note that a path of $G$ uses the edge $s_0\rightarrow s_f$ only if it is that edge alone, $s_0$ having no incoming edge and $s_f$ no outgoing one. For irreducibility, let $s,s'\in\mathcal V$. Choose a path $p: s_0\rightarrow\cdots\rightarrow s$ in $G$ as follows: for $s\neq s_f$, the initial segment ending at $s$ of a path from $s_0$ to $s_f$ via $s$ (path-connectedness), which does not use the edge $s_0\rightarrow s_f$ since it does not end at $s_f$; for $s=s_f$ and $\Statespace\neq\emptyset$, a path from $s_0$ to $s_f$ via some $y\in\Statespace$ (path-connectedness), which does not use the edge $s_0\rightarrow s_f$ since it visits $y$; for $s=s_f$ and $\Statespace=\emptyset$, the edge $s_0\rightarrow s_f$. If $s'\neq s_0$, choose a path $p': s'\rightarrow\cdots\rightarrow s_f$ in $G$ (path-connectedness), which does not use the edge $s_0\rightarrow s_f$ since it does not start at $s_0$. The backward chain started at $s$ realizes the reversal of $p$ (from $s$ down to $s_0$) with positive probability, each reversed edge carrying $\BackwardPolicy>0$ by hypothesis, which settles $s'=s_0$; if $s'\neq s_0$, it then moves $s_0\rightarrow s_f$ with probability $1$, and realizes the reversal of $p'$ (from $s_f$ down to $s'$) with positive probability. Hence $s'$ is accessible from $s$: the chain is irreducible, and this argument consumes the positivity of $\BackwardPolicy$ on the edge $s_0\rightarrow s_f$ only if $\Statespace=\emptyset$. A finite irreducible Markov chain admits a unique invariant probability measure, which is everywhere positive \cite{levin2017markov}; note that aperiodicity is not needed. Finally, the only backward transition into $s_f$ is from $s_0$ (the wrap edge is the only edge out of $s_f$ in $\widehat G$) and it has probability $1$, so stationarity at $s_f$ reads $\lambda(s_f)=\lambda(s_0)\widehat\BackwardPolicy(s_0\rightarrow s_f)=\lambda(s_0)$.

\emph{(2)} The constraint is an equality of measures on $\mathcal V^2$: in coordinates $(u,v)$ standing for a transition $u\rightarrow v$, it reads $\fout(u)\ForwardPolicy(u\rightarrow v) = \widehat\BackwardPolicy(v\rightarrow u)\fout(v)$. Summing over $v$ gives, at every $u\in \mathcal V$,
\begin{equation*}
\fout(u) = \sum_{v\in\mathcal V}\fout(v)\,\widehat\BackwardPolicy(v\rightarrow u),
\end{equation*}
i.e. $\fout\mu$ is invariant under the backward chain, whence $\fout\mu = c\,\lambda$ for some $c>0$ by \emph{(1)}. The constraint then becomes $(c\lambda)\otimes\ForwardPolicy = \widehat\BackwardPolicy\otimes(c\lambda)$, whose unique solution is the $\lambda$-reversal of Lemma \ref{lem:adjoint}; on a finite space the disintegration is the elementary formula displayed in the statement, well-defined since $\lambda>0$. Conversely this pair satisfies the constraint by definition of the reversal.

\emph{(3)} Define the edge flow $e(u\rightarrow v):= c\,\lambda(v)\,\widehat\BackwardPolicy(v\rightarrow u)$ for every edge $u\rightarrow v$ of $\widehat G$: it is the edge measure of the flow of \emph{(2)}. At every vertex,
\begin{eqnarray}
\sum_{v} e(u\rightarrow v) &=& c\sum_v \lambda(v)\widehat\BackwardPolicy(v\rightarrow u) = c\,\lambda(u),\nonumber\\
\sum_{u} e(u\rightarrow v) &=& c\,\lambda(v)\sum_u\widehat\BackwardPolicy(v\rightarrow u) = c\,\lambda(v),\nonumber
\end{eqnarray}
the first by stationarity of $\lambda$, the second because $\widehat\BackwardPolicy(v\rightarrow\cdot)$ is a probability: $e$ is a circulation, with throughput $c\lambda(u)$ at every vertex $u$. The wrap edge carries $e(s_f\rightarrow s_0) = c\,\lambda(s_0)\,\widehat\BackwardPolicy(s_0\rightarrow s_f) = c\,\lambda(s_0) = Z$. Cut it, declaring its flow to be an initial flow $Z$ injected at $s_0$ and a terminal flow $Z$ collected at $s_f$; conservation at every vertex is untouched. In the source--sink reading of $G$, the initial flow is $\Finit(v) = e(s_0\rightarrow v) = c\lambda(v)\BackwardPolicy(v\rightarrow s_0)$, of total mass $c\lambda(s_0)=Z$, and the terminal flow is, at every $x\in\Statespace$,
\begin{eqnarray*}
\Fterm(x) \;=\; e(x\rightarrow s_f) &=& c\,\lambda(s_f)\,\BackwardPolicy(s_f\rightarrow x)\\
&=& Z\,\BackwardPolicy(s_f\rightarrow x) \;=\; \target(x),
\end{eqnarray*}
using $\lambda(s_f)=\lambda(s_0)$; it is positive on every terminating state since $\BackwardPolicy$ is positive on every terminating edge $x\rightarrow s_f$. Explicitly, on the internal state space $\Statespace$ the generative flow has star outflow $\foutstar(u)= c\lambda(u) - e(u\rightarrow s_f)$ and star forward policy $\ForwardPolicyStar(u\rightarrow v) := e(u\rightarrow v)/\foutstar(u)$ (at states where $\foutstar$ vanishes the trajectory terminates almost surely and the policy is immaterial), so that $\foutstar(u)\ForwardPolicyStar(u\rightarrow v)=e(u\rightarrow v)$ and, at every $v\in\Statespace$,
\begin{eqnarray*}
\lefteqn{\Finit(v) + \bigl((\foutstar\mu)\ForwardPolicyStar\bigr)(v)}\\
&=& \Finit(v) + \sum_{u \in \Statespace} e(u\rightarrow v)\\
&=& c\,\lambda(v) \;=\; \Fterm(v) + \foutstar(v),
\end{eqnarray*}
which is the flow-matching constraint \eqref{equ:FM_const}. By Theorem \ref{theo:sampling_theorem}, the sampler terminates almost surely with $s_\tau \sim \Fterm/Z = \target/Z$.

For the final statement, let $\target$ be a target of mass $Z$ carried by the terminating states and positive at each of them; a terminating state then exists, so $\Statespace\neq\emptyset$. The row $\BackwardPolicy(s_f\rightarrow\cdot):=\target/Z$ is a probability supported on the in-neighbours of $s_f$, so the frozen policy is a backward policy on $G$; it is positive on every edge of $G$ except on the edge $s_0\rightarrow s_f$, where it vanishes, if that edge is in $\mathcal E$. The proofs of \emph{(1)}--\emph{(3)} consume positivity on no other edge: that of \emph{(1)} consumes it on $s_0\rightarrow s_f$ only if $\Statespace=\emptyset$, that of \emph{(2)} consumes only \emph{(1)}, and that of \emph{(3)} consumes only \emph{(1)}, \emph{(2)}, the value $\widehat\BackwardPolicy(s_0\rightarrow s_f)=1$ and the positivity of the sink row on the terminating edges. Hence \emph{(1)}--\emph{(3)} hold for the frozen policy, with terminal flow $Z\cdot\target/Z=\target$: the balanced flow matches the target exactly, every reasonable loss vanishes on it, and all residuals $\delta\Finit$, $\delta\Fterm$ are zero --- in every $L^p$ norm at once, the state space being finite.
\end{proof}

\begin{remark}[the loop closure is necessary]\label{rem:loop_closure_necessary}
On $G$ itself, i.e. without the wrap edge, the frozen-backward family --- with $\ForwardPolicy$ a probability row at every state of $\mathcal V$, $s_f$ included, each row a distribution on $\mathcal V$ and not required to be supported on $\mathcal E$ --- is trivial. Indeed the marginal argument of \emph{(2)}, which sums that row at each state, forces $\fout(s_0)=0$ ($\BackwardPolicy$ being undefined at $s_0$) and $\fout\mu$ to be invariant under the backward kernel restricted to $\mathcal V\setminus\{s_0\}$, which is substochastic: from every vertex, the backward chain leaks to $s_0$ within $\#\mathcal V$ steps with probability bounded below, by reversing a source-to-vertex path as in the proof of \emph{(1)}. Iterating the invariance identity gives $\fout = 0$. The loop closure is therefore not a convenience of the proof: it is the mechanism that makes the frozen-backward parameterization nonempty, and it is where the target enters --- the reward is the backward-policy row at the sink, $\target = Z\BackwardPolicy(s_f\rightarrow\cdot)$, consistently with the wrap-transition reading of the trajectory-balance loss in Proposition \ref{prop:tb_gradient}.
\end{remark}

\begin{remark}[relation to Theorem \ref{theo:universality_L2} and to training]\label{rem:graphs_vs_L2}
Theorem \ref{theo:universality_graphs} is more precise than Theorem \ref{theo:universality_L2} applied to the loop-closed chain --- exactness, and uniqueness of the balanced flow --- and less demanding: it needs only the irreducibility of a finite chain, whereas the summable-mixing hypothesis of Theorem \ref{theo:universality_L2} fails when the loop closure is periodic (the loop closure of a single directed path is a directed cycle). When $\widehat\BackwardPolicy$ is moreover aperiodic, the finite irreducible chain is geometrically ergodic, its $L^2$-mixing coefficients are summable, and Theorem \ref{theo:db_stable_frozen_full} gives the rate $g''(1)w_{\min}/\widehat B^2$ of the linearized gradient descent toward the balanced ray of item \emph{(2)}. Proposition \ref{prop:morozov_rate} removes the aperiodicity hypothesis, replacing $\widehat B$ by the hitting-time constant $\widehat B_\sigma$. Heuristically, the uniqueness up to scale in \emph{(2)} is the static counterpart of the conservation of $\Pi h_t$ in Theorem \ref{theo:db_stable_frozen_full}: the balanced locus is one-dimensional, and the linearized training moves transversally to it, whereas along the nonlinear gradient flow of a strictly unimodal generator the total mass increases (Proposition \ref{prop:no_distant_equilibrium}\emph{(2)}).
\end{remark}

\begin{remark}[targets vanishing at a terminating state]\label{rem:partial_support}
If the target $\target$ vanishes on some terminating states, let $G'$ be the marked subgraph of $G$ induced on $s_0$, $s_f$ and the states from which a directed path of $G$ reaches $\mathrm{supp}\,\target$, $\mathrm{supp}\,\target$ included, with the terminating edges $y\rightarrow s_f$, $\target(y)=0$, deleted. It is again path-connected, and a backward policy on $G$ positive on its edges restricts to one on $G'$ with only its sink row renormalized. Theorem \ref{theo:universality_graphs} applied to $G'$ gives a flow which, extended by zero, satisfies the flow-matching constraint \eqref{equ:FM_const} on $G$ exactly, with terminal flow $\target$; it belongs to the frozen-backward family of $G'$, not to that of $G$. Alternatively, if $\target$ has mass $Z$ and $\mathrm{unif}$ is the uniform probability on the terminating states, the targets $(1-\varepsilon)\target+\varepsilon Z\,\mathrm{unif}$, positive at every terminating state, realize $\target$ as a limit of exactly matched flows: the universality infimum over the union of the frozen-backward families of $G$, one for each backward policy positive on its edges, still vanishes, and it is not attained, every element of that union being either the trivial element or a flow whose terminal flow $Z'\BackwardPolicy(s_f\rightarrow\cdot)$, $Z'>0$, is positive on every terminating state.
\end{remark}
   % proofs.tex issues \appendix itself
\section{\appendixBtitle}\label{sec:silva_comparison}

Three works of Silva and coauthors are related to the bounds of the present paper.
\citet{silvageneralization} prove a bound on the total variation of the terminal distribution by a trajectory-balance residual (their Proposition 4.2), which is, to our knowledge, the only loss-to-distance bound in this line of work; we analyze it in detail below.
\citet{silva2025gflownets} bound the total variation induced by a \emph{flow perturbation} of an exactly balanced model (their Theorem 1) --- a sensitivity analysis at the balanced manifold rather than a bound by the value of a training loss.
\citet{silva2024divergence} prove an equality of expected on-policy gradients between the TB loss and the reverse Kullback--Leibler divergence (their Proposition 1), with no value bound.
None of the three addresses the behaviour of a minimizing \emph{sequence} of models --- the question addressed by Theorems \ref{theo:gd_diffusion} and \ref{theo:db_stable_frozen}, which describe where gradient descent goes, and how fast, near the balanced flow, and by Proposition \ref{prop:frozen_unstable}, which delimits what that convergence does \emph{not} give.

\subsection{Setting}

All three works operate on a finite state space carrying a pointed directed acyclic state graph in which every trajectory terminates within a fixed horizon $t_m$ (the ``finitely absorbing'' axiom of \citet{silvageneralization}, Definition A.1, from which acyclicity is derived), with total variation taken between distributions on the set $\mathcal X$ of terminal states; in \citet{silvageneralization} the backward policy $p_B$ is moreover \emph{fixed to the uniform policy} throughout.
In the notation of the present paper: $\Statespace$ finite, $\BackwardPolicy$ frozen, and the trajectory expectations of their losses are taken over \emph{backward} trajectories $\tau\sim p_B(\cdot\mid x)$ --- the same frozen-backward-policy structure as in Theorem \ref{theo:first_variation} and its specializations.
When the loop-closed backward policy is moreover aperiodic it has summable $L^2$-mixing, and for a training measure of positive density Theorem \ref{theo:db_stable_frozen} applies to the FM and DB losses on such a graph and supplies what the static bounds below do not: a convergence rate for the linearized training dynamics.

\subsection{The bound of \citet{silvageneralization} with explicit constants}

Proposition 4.2 of \citet{silvageneralization} states, for $q_{E,T}$ the uniform distribution on $\mathcal X$ and $p_{E,T}$ a training distribution on $\mathcal X$,
\begin{eqnarray}
&&\TV(p_T,\pi) \nonumber \\ && \quad \lesssim\;
\Bigl[\bigl(1+\chi^2(q_{E,T}\|p_{E,T})\bigr) \, \mathcal E(p_F)\Bigr]^{1/2},
\label{eq:silva_bound}
\end{eqnarray}
where $\pi = R/Z$ is the \emph{normalized} target, $p_T$ the terminal marginal of the model, and
\begin{equation}
\mathcal E(p_F) := \mathbb E_{x\sim p_{E,T}}\,\mathbb E_{\tau\sim p_B(\cdot|x)}\left[\left(\log \frac{p_F(\tau)}{\pi(x)\,p_B(\tau\mid x)}\right)^{\!2}\right].
\end{equation}
The following restatement makes the constant absorbed in $\lesssim$ explicit; the proof follows their scheme (H\"older, Jensen over backward trajectories, a log-Lipschitz step).

\begin{proposition}[explicit-constant form of Prop.\ 4.2 of \citet{silvageneralization}]\label{prop:silva_explicit}
Let $\mathcal X$ be finite, let $p_{E,T}$ and $\pi$ have full support on $\mathcal X$, assume $p_F(\tau)>0$ for every trajectory $\tau$ with $p_B(\tau\mid x)>0$, and assume $p_B(\cdot\mid x)$ dominates the forward trajectory distribution into each $x$. Set
\begin{equation}
    M := \max_{\tau,x:\,p_B(\tau\mid x)>0} \frac{p_F(\tau)}{p_B(\tau\mid x)}, \qquad M' := \max\left(M,\ \|\pi\|_{\infty}\right).
\end{equation}
Then
\begin{eqnarray}
&&\TV(p_T,\pi) \nonumber\\ &&\;\; \leq\; \frac{|\mathcal X|}{2}\, M'\,
\Bigl[\bigl(1+\chi^2(q_{E,T}\|p_{E,T})\bigr)\,\mathcal E(p_F)\Bigr]^{1/2}.
\end{eqnarray}
\end{proposition}
\begin{proof}
Write $\phi(x) := |\pi(x)-p_T(x)|$, so $\TV(p_T,\pi) = \frac{|\mathcal X|}{2}\,\mathbb E_{x\sim q_{E,T}}[\phi(x)]$. Since $p_{E,T}$ has full support, the Cauchy--Schwarz inequality with the change of measure $q_{E,T} = \frac{q_{E,T}}{p_{E,T}}\,p_{E,T}$ gives
\begin{equation}
    \mathbb E_{q_{E,T}}[\phi] \leq \left(\mathbb E_{p_{E,T}}[\phi^2]\right)^{1/2}\bigl(1+\chi^2(q_{E,T}\|p_{E,T})\bigr)^{1/2}.
\end{equation}
By the domination assumption, $p_T(x) = \mathbb E_{\tau\sim p_B(\cdot|x)}\bigl[p_F(\tau)/p_B(\tau\mid x)\bigr]$, so Jensen's inequality yields
\begin{equation}
    \phi(x)^2 \leq \mathbb E_{\tau\sim p_B(\cdot|x)}\left[\left(\frac{p_F(\tau)}{p_B(\tau\mid x)} - \pi(x)\right)^{\!2}\right].
\end{equation}
Finally, for $a := p_F(\tau)/p_B(\tau\mid x) \in (0, M]$ and $b := \pi(x)\in(0,\|\pi\|_\infty]$, the mean value theorem applied to $\exp$ on $(-\infty, \log M']$ gives $|a-b| \leq M'\,|\log a - \log b|$, whence $(a-b)^2 \leq M'^2\bigl(\log\frac{p_F(\tau)}{\pi(x)p_B(\tau|x)}\bigr)^2$. Combining the three displays proves the claim.
\end{proof}

\begin{remark}[hypotheses used by the bound]\label{rem:silva_hypotheses}
The proof above isolates the hypotheses under which Equation \eqref{eq:silva_bound} operates; they are used, but not all stated, in \citet{silvageneralization}.
\begin{enumerate}
    \item \emph{Full support of $p_{E,T}$.} The change-of-measure step divides by $p_{E,T}$ under a uniformly weighted sum; off full support, $\chi^2(q_{E,T}\|p_{E,T}) = +\infty$ and the right side of Proposition \ref{prop:silva_explicit} is $+\infty$, or $\infty\cdot0$ when $\mathcal E(p_F)=0$. Heuristically, $\chi^2$ plays here, in $L^2$ form, the role of the factor $\|\dln\|_{L^\infty}$ in Theorem \ref{theo:RL_CV_bound}.
    \item \emph{Full support of the trained $p_F$.} If $p_F(\tau)=0$ for a trajectory $\tau$ into a terminal state $x$ with $p_{E,T}(x)>0$ and $p_B(\tau\mid x)>0$, then $\mathcal E(p_F)=+\infty$: the bound holds but carries no information. \citet{silvageneralization} remove this with the $\alpha$-mixture realizability device of their Lemma B.1, which they invoke only for their Theorem 5.2.
    \item \emph{Domination by the backward policy.} The identity $p_T(x)=\mathbb E_{p_B(\cdot|x)}[p_F/p_B]$ requires $p_B(\cdot\mid x)$ to dominate forward trajectories into $x$ --- automatic for the uniform $p_B$ that \citet{silvageneralization} fix throughout, an assumption under the extension to learnable $p_B$ they mention.
    \item \emph{The constant depends on the trained model, boundedly on each fixed graph.} $M=\max_{\tau,x} p_F(\tau)/p_B(\tau\mid x)$ is a maximum over the \emph{learned} forward policy; it is quantified in Remark \ref{rem:silva_model_constant} below.
    \item \emph{Constant of the log-Lipschitz step.} The step $(a-b)^2\lesssim (\log a - \log b)^2$ is justified in Lemma D.1 of \citet{silvageneralization} for arguments in $(0,1]$ after rescaling by $M$; since $\pi(x) \leq M$ is not implied by the definition of $M$ (on a star, $s_o$ with $K\geq3$ terminal children, the uniform forward policy has $M=1/K$, while a target with $\pi(x)=\tfrac12$ at one child has $\|\pi\|_\infty=\tfrac12>M$), a sufficient rescaling constant is $M'=\max(M,\|\pi\|_\infty)$, as in Proposition \ref{prop:silva_explicit}.
    \item \emph{Structural axioms.} Finiteness of $\mathcal X$, the uniform horizon $t_m$, and acyclicity (derived from the finitely-absorbing axiom) are standing hypotheses of their framework; the prefactor $\frac{|\mathcal X|}{2}$ is absorbed into $\lesssim$ in their statement.
\end{enumerate}
\end{remark}

\begin{remark}[the constant is not uniform in the depth]\label{rem:silva_model_constant}
Heuristically, the reading of Proposition \ref{prop:silva_explicit} hinges on its constant: the proposition is a statement about a \emph{fixed} model, while a convergence guarantee for training must hold uniformly along a minimizing sequence --- and $M'$ is a maximum over the very policy being trained. On a fixed graph that dependence is bounded: since $p_F(\tau)\leq1$ and $\|\pi\|_\infty\leq1$, every forward policy has
\begin{equation}
    M' \;\leq\; \max\Bigl(1,\ \max_{\tau,x:\,p_B(\tau\mid x)>0}\frac{1}{p_B(\tau\mid x)}\Bigr).
\end{equation}
On a tree, $p_B(\tau\mid x)=1$ and $M\leq1$; on a state graph with several backward paths into a terminal state, $M$ can exceed $1$ but stays at most that bound. Consider the ladder graph of depth $n$: $n+1$ states in a row, two parallel edges at each of the $n$ stages, and the last state $x$ as the only terminal state, so that $\pi(x)=p_{E,T}(x)=1$, $x$ has $2^{n}$ incoming trajectories, the uniform backward policy gives $p_B(\tau\mid x)=2^{-n}$ for each, and $M'\leq 2^{n}$. At the balanced solution, $p_F(\tau)=\pi(x)p_B(\tau\mid x)$ and $M=\|\pi\|_\infty=1$. But already at the level of trajectory distributions (before imposing Markov realizability of $p_F$), a distribution placing mass $c$ on a single trajectory $\tau_0$ has $M \geq c\,2^{n}$, while the corresponding term of the residual $\mathcal E(p_F)$ is
\begin{equation}
    2^{-n}\left(\log\bigl(c\,2^{n}\bigr)\right)^{2} \xrightarrow[n\rightarrow\infty]{} 0:
\end{equation}
the residual weights the trajectory on which the constant is large by $p_B(\tau_0\mid x)=2^{-n}$. With $c:=1-e^{-\sqrt{\varepsilon/2}}$ and the remaining mass $1-c$ spread uniformly over the other trajectories, $\mathcal E(p_F)\rightarrow\varepsilon/2$ as $n\rightarrow\infty$; consequently, over trajectory distributions,
\begin{equation}
    \sup\{M'(p_F) : \mathcal E(p_F)\leq\varepsilon\} \;\geq\; \bigl(1-e^{-\sqrt{\varepsilon/2}}\bigr)\,2^{n}
\end{equation}
for every $\varepsilon>0$ and every $n$ large enough: on $\{\mathcal E\leq\varepsilon\}$ the constant reaches the order of its fixed-graph bound, and no constant is uniform in the depth. This computation does not quantify the growth of the same supremum over Markov forward policies. On a directed cycle entered at one state and terminal at every state, Lemma \ref{lem:cycle_counterexample} rules out any model-free bound tending to $0$ with the loss for the divergence-based FM losses of \citet{brunswic2024theory} with a generator continuous at $1$ (Remark \ref{rem:path_space}), a class that does not contain the idealized residual $\mathcal E(p_F)$.
\end{remark}

Remark \ref{rem:silva_model_constant} bounds $M'$ on each fixed graph independently of the model, by a quantity that can grow with the depth. Whether a model-free bound can be independent of the state space is settled by the following proposition, on a sequence of DAGs: a per-graph model-free bound always exists (item \emph{(i)}), but its constant necessarily degenerates with the state space (item \emph{(ii)}), even when every degeneracy channel of Proposition \ref{prop:silva_explicit} other than $|\mathcal X|$ is switched off.

\begin{proposition}[the state-space dependence is unavoidable]\label{prop:silva_no_uniform}
Let $w_{\min} := \min\{\,p_{E,T}(x)\,p_B(\tau\mid x) : p_B(\tau\mid x)>0\,\}$.
\begin{enumerate}
    \item On every finite pointed DAG, for every full-support forward policy,
    \begin{equation}
        \TV(p_T,\pi) \;\leq\; \tfrac12\left(e^{\sqrt{\mathcal E(p_F)/w_{\min}}}-1\right).
    \end{equation}
    \item For every $\varepsilon>0$ there exist a finite pointed DAG with uniform $p_B$, a target $\pi$, the \emph{uniform} training distribution $p_{E,T}$ (so $\chi^2(q_{E,T}\|p_{E,T})=0$) and a full-support forward policy $p_F$ with $M'\leq 1$ such that
    \begin{equation}
        \mathcal E(p_F)\leq\varepsilon \qquad\text{and}\qquad \TV(p_T,\pi)=\tfrac18 .
    \end{equation}
\end{enumerate}
In particular, no bound $\TV\leq F\bigl(\mathcal E,\chi^2,M'\bigr)$ with $F$ independent of the state space, non-decreasing in its third argument and $F(\mathcal E,0,1)\rightarrow0$ as $\mathcal E\rightarrow0$ can hold: the factor $1/w_{\min}$ in \emph{(i)} --- of order $|\mathcal X|\,e^{c\,t_m}$ in general --- is not an artifact of the proof.
\end{proposition}
\begin{proof}
\emph{(i)} Every trajectory term of $\mathcal E$ is bounded below by $w_{\min}\,\ell(\tau)^2$ with $\ell(\tau) := \log\frac{p_F(\tau)}{\pi(x)p_B(\tau|x)}$, so $|\ell(\tau)|\leq s := \sqrt{\mathcal E/w_{\min}}$ for every $\tau$ with $p_B(\tau\mid x)>0$. Since $p_T(x) = \pi(x)\,\mathbb E_{\tau\sim p_B(\cdot|x)}[e^{\ell(\tau)}]$, every ratio $p_T(x)/\pi(x)$ lies in $[e^{-s},e^{s}]$ and $\TV = \frac12\sum_x \pi(x)\bigl|\tfrac{p_T(x)}{\pi(x)}-1\bigr| \leq \frac12(e^{s}-1)$.

\emph{(ii)} Take the star DAG: $s_o$ with children $\mathcal X = \{x^*, x'\} \cup \{b_1,\dots,b_K\}$, each terminal linked to $s_f$; each $x\in\mathcal X$ has the single parent $s_o$, so $p_B(\tau\mid x)=1$ and the unique trajectory into $x$ is $(s_o,x)$. Set $\pi(x^*)=\pi(x')=\frac14$ and $\pi(b_i)=\frac1{2K}$, and take the forward policy $p_F(s_o\rightarrow x^*)=\frac38$, $p_F(s_o\rightarrow x')=\frac18$, $p_F(s_o\rightarrow b_i)=\frac1{2K}$. Then $\TV(p_T,\pi) = \frac12(\frac18+\frac18)=\frac18$ for every $K$, while the log-ratios vanish on the $b_i$ and
\begin{equation}
    \mathcal E(p_F) = \frac{\left(\log\frac32\right)^2 + \left(\log 2\right)^2}{K+2} \;\leq\; \frac{1}{K+2},
\end{equation}
which is $\leq\varepsilon$ for $K$ large. Here $M = \max_x p_F(s_o\rightarrow x) = \frac38$ and $\|\pi\|_\infty=\frac14$ for $K\geq2$, the filler weight $\frac1{2K}=\frac12$ being the largest of both at $K=1$, so $M'\leq 1$ for every $K$, and $p_{E,T}=q_{E,T}$ gives $\chi^2=0$.
\end{proof}

\begin{remark}[the two degeneracy channels]\label{rem:silva_two_channels}
The two degeneracy channels are distinct. In the star family of \emph{(ii)} the failure is pure $|\mathcal X|$-averaging: a $\Theta(1)$ misallocation on two terminals is diluted by the uniform training distribution over $K+2$ terminals, with $M'\leq1$ and $\chi^2=0$. The corridor of Remark \ref{rem:silva_model_constant} is the complementary channel, through $M'$, on deep graphs. Conversely, when $\mathcal X$ is small and $p_{E,T}$ covers it, item \emph{(i)} is a genuine model-free bound: the depleted side of any misrouting is then visible in the residual. It is the averaging over a large terminal space --- the regime GFlowNets are designed for --- that hides it.
\end{remark}

\begin{remark}[the hidden size is the path space]\label{rem:path_space}
Heuristically, the graph-dependence of all these bounds is measured by the size of the \emph{path space}. In Proposition \ref{prop:silva_no_uniform}\emph{(i)}, since for a backward policy of full support $p_B(\cdot\mid x)$ is a probability charging each of the $P(x)$ trajectories into $x$, one has $\min_\tau p_B(\tau\mid x)\leq 1/P(x)$ and therefore
\begin{equation}
    \frac{1}{w_{\min}} \;\geq\; \max_{x\in\mathcal X}\ \frac{P(x)}{p_{E,T}(x)} ,
\end{equation}
and the constant $M'$ of Proposition \ref{prop:silva_explicit} satisfies $\sup_{p_F}M'\geq\max_{x\in\mathcal X}P(x)$, the supremum running over full-support forward policies. When a cycle lies on the trajectories into a state $x$, these are infinitely many, and, for a backward policy of full support, both $\inf_\tau p_B(\tau\mid x)=0$ and $\sup_{p_F}M'=+\infty$ follow from that. The infimum vanishes because $p_B(\cdot\mid x)$ is a probability charging each of them. The supremum is infinite because $p_B(\cdot\mid x)$ is summable: along the trajectories $\tau_n$ that wind $n$ times around the cycle, $\liminf_{n\rightarrow\infty} n\,p_B(\tau_n\mid x)=0$, while for each $n$ some full-support forward policy, leaving the cycle with probability $1/(n+2)$ at the state where $\tau_n$ leaves it, has $p_F(\tau_n)\geq c/(n+2)$, with $c>0$ depending on the graph, the cycle and $x$ but not on $n$ --- $M'$ being read with the formula of Proposition \ref{prop:silva_explicit}. The infimum alone would not give the supremum, since every forward policy has $p_F(\tau_n)\leq 1/(n+1)$. Let $\mathcal C_N$ be the directed cycle on states $x_1,\dots,x_N$ of Lemma \ref{lem:cycle_counterexample}, entered at $x_1$ and terminal at every state: on $\mathcal C_N$ with a fixed training distribution, $|\mathcal X|=N$ and $\chi^2(q_{E,T}\|p_{E,T})$ do not depend on the flow, while on the counter-example graph \eqref{eq:counterexample_graph}, whose terminal set is infinite, $|\mathcal X|=\infty$ and $\chi^2(q_{E,T}\|p_{E,T})$ is undefined, there being no uniform distribution on an infinite set. For the divergence-based FM losses of \citet{brunswic2024theory} with a generator continuous at $1$, acyclicity cannot be dropped: on $\mathcal C_N$, for a target $\target\neq\delta_{x_2}$, Lemma \ref{lem:cycle_counterexample} supplies flows whose loss tends to $0$ for every training distribution while their terminal law stays $\delta_{x_2}$, at total variation $1-\target(x_2)>0$ from $\target$, so no model-free bound tending to $0$ with the loss holds. The idealized residual $\mathcal E(p_F)$ lies outside that class, and whether the acyclicity hypothesis of \citet{silvageneralization} can be dropped from their bound the lemma does not decide; what this remark establishes for that bound is the graph-dependence of its constants through the path space.
\end{remark}

\begin{remark}[relation to the bounds of this paper]\label{rem:silva_vs_ours}
Two further points of comparison. \emph{(i)} The residual $\mathcal E(p_F)$ of \citet{silvageneralization} involves the normalized target $\pi = R/Z$ with the true partition function: it is an idealized form of the TB objective, distinct from the trained TB loss, in which the normalization $Z_\theta = F(s_o)$ is learned; one expects relating the two to require a control of $|\log Z_\theta - \log Z|$, which \citet{silvageneralization} do not provide. \emph{(ii)} Theorem \ref{theo:no_bound_divergence} asserts that, for the divergence-based losses of \citet{brunswic2024theory} with a generator continuous at $1$, no bound $\TV\leq F(\mathcal L)$ with a model-free $F$ tending to $0$ at $0$ holds on the graphs of a given size, whatever the training distribution. On the directed cycle $\mathcal C_N$ of Lemma \ref{lem:cycle_counterexample}, the constants $M'$ and $1/w_{\min}$ lie outside the acyclic setting of Propositions \ref{prop:silva_explicit} and \ref{prop:silva_no_uniform}; read with the same formulas, $\sup_{p_F}M'=+\infty$ (Remark \ref{rem:path_space}), while $|\mathcal X|$ and $\chi^2$ stay fixed with the graph and the training distribution; on the counter-example graph \eqref{eq:counterexample_graph}, cyclic as well and with no finite horizon, the terminal set is infinite and $\chi^2$ is undefined too. Their bound thus quantifies the finite acyclic bounded-horizon regime; Lemma \ref{lem:cycle_counterexample} shows that on $\mathcal C_N$ no model-free bound tending to $0$ with the loss survives for the divergence-based FM losses --- whether one survives for the idealized residual $\mathcal E(p_F)$, which lies outside that class, it does not decide --- and Theorems \ref{theo:gd_diffusion} and \ref{theo:db_stable_frozen} supply the dynamic statement missing from their bound, namely the convergence rate of the linearized training under a mixing backward policy.
\end{remark}

\subsection{The two companion papers}

\citet{silva2024divergence} prove (their Proposition 1) that on a finitely absorbing acyclic space, the expected \emph{on-policy} TB gradient equals twice the gradient of the reverse KL divergence, $\nabla_\theta\,\mathbb E_{\tau\sim P_F}[\mathcal L_{TB}] = 2\nabla_\theta\,\KL(P_F\|P_B)$, by the score-function identity. This is a gradient identity, not a value bound, and is therefore compatible with Theorem \ref{theo:no_bound_divergence}. It is the on-policy, trajectory-space counterpart of Theorem \ref{theo:first_variation}, which computes the gradient in operator form $(T^\lambda - r)[\,\cdot\,]$ for an arbitrary training measure and frozen backward policy; the two agree at the balanced flow.

\citet{silva2025gflownets} prove (their Theorem 1) that perturbing the flow of a single node of an \emph{exactly balanced} model by $\delta$ and propagating the excess forward moves the terminal distribution by at most $\frac{\delta}{Z+\delta}(1-\min_{x\in D}\pi(x))$, where $D$ is the set of terminal descendants of the perturbed node; the bound is linear in $\delta$ and tight on trees. This prices a perturbation of the balanced solution in total variation; it is complementary to Theorem \ref{theo:gd_diffusion}, which quantifies the \emph{decay} of such perturbations under gradient descent. Their Proposition 1 --- forward-looking losses with unrestricted terminal flows vanish on models with $p_T \propto R\tilde F$, so that a zero loss does not imply correctness --- is a further instance, at zero loss, of the uniqueness and stability failures that Theorem \ref{theo:no_bound_divergence} formalizes at small loss.

\subsection{Summary}

\begin{center}
\small
\begin{tabular}{@{}p{0.21\columnwidth}p{0.33\columnwidth}p{0.34\columnwidth}@{}}
\hline
 & bounds what & constants; dynamics \\
\hline
\citet{silvageneralization} Prop.\ 4.2 & TV by an idealized TB residual (true $Z$), backward-sampled & $\frac{|\mathcal X|}{2}M'\sqrt{1+\chi^2}$, depth-dependent via $M'$; static \\[2pt]
\citet{silva2025gflownets} Thm.\ 1 & TV under a one-node flow perturbation of a balanced model & linear in $\delta/(Z+\delta)$; static, at the balanced manifold \\[2pt]
This paper & TV by the trained stable / KL-wFM losses (Thms.\ \ref{theo:RL_CV_bound}, \ref{theo:IL_CV_bound_body}); no bound for divergence losses in general (Thm.\ \ref{theo:no_bound_divergence}) & model-free constants; dynamic: linearized training rate from the mixing of $\BackwardPolicy$ (Thm.\ \ref{theo:db_stable_frozen}) \\
\hline
\end{tabular}
\end{center}

% Appendix C -- the initial flow: mass, shape, and the constant of the stable bound.
\section{\appendixCtitle}\label{app:finit}

The constant of Theorem~\ref{theo:RL_CV_bound} ties the guarantee to two objects the
practitioner chooses rather than inherits: the initial flow $\Finit$ and the training
measure $\nutrain$. This appendix works out the first; Appendix~\ref{app:replay} the
second. Throughout, $\Finit\ll\nuB$ has mass $z:=\Finit(\Statespace)$ and density
$\finit$, the target $\target\ll\nuB$ has density $k$ and mass
$t:=\target(\Statespace)$, the trained residual is $E:=\Finit+\Finstar-\Foutstar$ with
density $e$ and $L^1(\nuB)$ error $\Delta:=\|e-k\|_{L^1(\nuB)}$, and the stable loss is
$\LFMstableq=\|e-k\|_{L^q(\nutrain)}$. In this notation Theorem~\ref{theo:RL_CV_bound}
reads, for the inference-time terminal flow $\widehat\target=E^+$,
\begin{equation}\label{equ:paper_bound}
    \TV\left(s_\tau\left\|\frac{\target}{t}\right.\right)\;\leq\;C\,\LFMstableq,\qquad
    C=\frac{2}{t}\left\|\dln\right\|^{1/q}_{L^\infty(\nutrain)}\nuB(\Statespace)^{1-1/q},
\end{equation}
and Theorem~\ref{theo:RL_CV_bound_full} gives the model-free floor
$\LFMstableq\geq|z-t|\big/\bigl(\nuB(\Statespace)^{1-1/q}\|\dln\|^{1/q}_{L^\infty(\nutrain)}\bigr)$.
Write
\begin{equation}\label{equ:dual_norm}
    \Dq \;:=\; \left\|\dln\right\|_{L^{q^*}(\nutrain)} .
\end{equation}

\subsection{The dual mass floor}

\begin{proposition}[dual mass floor]\label{prop:dual_floor}
Assume $\nuB\ll\nutrain$ and $\dln\in L^{q^*}(\nutrain)$. Then for every generative flow
and every initial flow of mass $z$,
\begin{equation}\label{equ:dual_floor}
    \LFMstableq\;\geq\;\frac{|z-t|}{\Dq},
\end{equation}
and
\begin{equation}\label{equ:norm_comparison}
    \Dq\;\leq\;\nuB(\Statespace)^{1-1/q}\left\|\dln\right\|^{1/q}_{L^\infty(\nutrain)},
\end{equation}
so \eqref{equ:dual_floor} sharpens the floor \eqref{equ:mass_floor}. The two floors
agree at $q=1$; for $q>1$ the inequality \eqref{equ:norm_comparison} is strict unless
$\dln$ takes, $\nutrain$-essentially, only the values $0$ and
$\|\dln\|_{L^\infty(\nutrain)}$ --- in particular they agree when $\nutrain$ is
proportional to $\nuB$.
\end{proposition}

\begin{proof}
As in the proof of Theorem~\ref{theo:RL_CV_bound_full}, $\int(e-k)\,d\nuB=E(\Statespace)-t=z-t$. Writing the
integral against $\nutrain$ and applying H\"older with exponents $(q,q^*)$,
\begin{equation*}
    |z-t| \;=\; \left|\int(e-k)\,\dln\,d\nutrain\right|
    \;\leq\;\|e-k\|_{L^q(\nutrain)}\left\|\dln\right\|_{L^{q^*}(\nutrain)}
    \;=\;\LFMstableq\cdot\Dq ,
\end{equation*}
which is \eqref{equ:dual_floor}. For \eqref{equ:norm_comparison}, put $g:=\dln\geq0$ and
note $\int g\,d\nutrain=\nuB(\Statespace)$; then
\begin{equation*}
    \Dq^{q^*}=\int g^{q^*}d\nutrain=\int g^{q^*-1}\cdot g\,d\nutrain
    \;\leq\;\|g\|_{L^\infty(\nutrain)}^{q^*-1}\,\nuB(\Statespace),
\end{equation*}
and $\frac{q^*-1}{q^*}=\frac1q$, $\frac{1}{q^*}=1-\frac1q$ give
\eqref{equ:norm_comparison}. Equality in the middle step forces $g\in\{0,\|g\|_{L^\infty(\nutrain)}\}$
$\nutrain$-a.e.; at $q=1$, $q^*=\infty$ and both sides are $\|g\|_{L^\infty(\nutrain)}$.
\end{proof}

The floor \eqref{equ:mass_floor} is the composition of \eqref{equ:dual_floor} with
\eqref{equ:norm_comparison}; keeping the dual norm $\Dq$ undoes the lossy step. The same
quantity $1+\Dq$ already appears as the admissible weight $b$ of
Theorem~\ref{theo:IL_CV_bound}, so no new hypothesis is introduced.

On finite state spaces the floor is not merely a bound --- it is the exact minimum, and it
does not see the shape of $\Finit$ at all.

\begin{proposition}[exactness; the minimal loss sees only the mass]\label{prop:floor_exact}
Let $\Statespace$ be finite, let $\nuB$ be everywhere positive and invariant under a
$\nuB$-ergodic star policy $\ForwardPolicyStar$, and let $\nutrain$ be everywhere
positive. Fix any initial flow $\Finit$ of mass $z$. Then, minimizing over non-negative
outflows,
\begin{equation}\label{equ:floor_attained}
    \min_{\foutstar\geq0}\;\LFMstableq\left(\ForwardPolicyStar,\foutstar\right)
    \;=\;\frac{|z-t|}{\Dq}\,.
\end{equation}
For $q\in(1,\infty)$ a minimizer is explicit: with $g=\dln$, the extremal residual profile
\begin{equation}\label{equ:extremal_profile}
    e-k \;=\; \varphi_q\;:=\;\frac{z-t}{\int g^{q^*}d\nutrain}\;g^{q^*-1},
\end{equation}
is realized by $\foutstar=\eta\mathbf 1-S\rho$ with $\rho:=\varphi_q-(\finit-k)$,
$S:=(\mathrm{Id}-P_\star+\Pi)^{-1}-\Pi$ (which agrees with
$\sum_{n\geq0}(P_\star^n-\Pi)$ of \eqref{equ:poisson_solved} whenever the mixing is
summable, but exists for every finite $\nuB$-ergodic chain, periodic included), and
any $\eta\geq\max_x S\rho(x)$. At $q=\infty$ the extremal profile is the constant
$\varphi_\infty=(z-t)/\nuB(\Statespace)$; at $q=1$ it is a point mass at any state
maximizing $g$, and the minimum is still \eqref{equ:floor_attained}.
\end{proposition}

\begin{proof}
\emph{Lower bound:} Proposition~\ref{prop:dual_floor}. \emph{Realizability:} on a finite
state space, $\nuB$-ergodicity makes $1$ a simple eigenvalue of $P_\star$, so every
eigenvalue of $\mathrm{Id}-P_\star+\Pi$ is either $1$ (on constants) or $1-\beta\neq0$
(on the complement), the resolvent $S=(\mathrm{Id}-P_\star+\Pi)^{-1}-\Pi$ is defined ---
no aperiodicity is needed, unlike for the series form --- and
$(\mathrm{Id}-P_\star)S=\mathrm{Id}-\Pi$ by direct algebra ($\Pi P_\star=P_\star\Pi=\Pi$).
The residual
of a flow is $e-k=(\finit-k)+(P_\star-\mathrm{Id})\foutstar$. Given any target profile
$\varphi$ with $\int\varphi\,d\nuB=z-t$, the function $\rho=\varphi-(\finit-k)$ is
mean-zero ($\int\rho\,d\nuB=(z-t)-(z-t)=0$, so $\Pi\rho=0$), and
$\foutstar=\eta\mathbf 1-S\rho$ satisfies
$(P_\star-\mathrm{Id})\foutstar=-(P_\star-\mathrm{Id})S\rho=\rho$, hence $e-k=\varphi$;
choosing $\eta\geq\max_xS\rho(x)$ makes $\foutstar\geq0$. Thus \emph{every} profile of the
correct mean is reachable, and the minimization of $\LFMstableq=\|\varphi\|_{L^q(\nutrain)}$
over $\{\int\varphi\,d\nuB=z-t\}$ is the norm of the linear functional
$\varphi\mapsto\int\varphi g\,d\nutrain$ on $L^q(\nutrain)$: its value is
$|z-t|/\|g\|_{L^{q^*}(\nutrain)}$, attained by the H\"older extremizer, which is
\eqref{equ:extremal_profile} for $q\in(1,\infty)$ (one checks
$\int\varphi_q\,d\nuB=z-t$ directly), the stated constant for $q=\infty$, and the point
mass at an argmax of $g$ for $q=1$ (finite space, so the sup is attained).
\end{proof}

\begin{corollary}[separation of roles]\label{cor:separation}
On a finite state space with a $\nuB$-ergodic star policy and unrestricted non-negative
outflow, the infimum of the stable loss is a function of $\Finit(\Statespace)$ alone. The
shape of $\Finit$ is invisible to the loss; its entire footprint is on the \emph{outflow
mass} a minimizer requires, hence on the length of sampling trajectories
(\S\ref{sec:shape}) --- and, for restricted families $\Theta$, on whether the minimizer is
representable at all.
\end{corollary}

\begin{proof}
The value \eqref{equ:floor_attained} depends on the pair $(\Finit,\target)$ only through
the masses $z$ and $t$, which is the first claim, and with it the invisibility of the
shape. The outflow-mass footprint is Proposition~\ref{prop:shape_cost}(1) below ---
proved there at $z=t$ for exactly matching flows, which is the regime the footprint
claim refers to --- whose
minimal mass $\nuB(\Statespace)\max_xS(k-\finit)$ does depend on the shape; the
representability footprint is immediate, since restricting $\foutstar$ to a family
$\Theta$ can only remove the explicit minimizer of
Proposition~\ref{prop:floor_exact}.
\end{proof}

\subsection{The mass of the initial flow: the certificate overshoots $Z$}\label{sec:mass}

The proof of Theorem~\ref{theo:RL_CV_bound_full} passes from the first member of
\eqref{equ:tv_chain}, $\frac{\Delta}{z+\Delta}+\frac{\Delta}{t}$, to the envelope $\frac2t$
of \eqref{equ:stable_bound}; keeping $z$ visible in that same member, in the currency of the
residual rather than of the loss, gives a $z$-dependent constant instead, and that constant
has a preferred $z$.

\begin{proposition}[$z$-carried bound]\label{prop:z_bound}
Under the hypotheses of Theorem~\ref{theo:RL_CV_bound_full}, with the mass
$z:=\Finit(\Statespace)>0$ kept visible and with $\Delta:=\|e-k\|_{L^1(\nuB)}$,
\begin{equation}\label{equ:z_bound}
    \TV\left(s_\tau\left\|\frac{\target}{t}\right.\right)\;\leq\;
    \left(\frac1z+\frac1t\right)\Delta\;\leq\;
    \left(\frac1z+\frac1t\right)
    \nuB(\Statespace)^{1-1/q}\left\|\dln\right\|^{1/q}_{L^\infty(\nutrain)}\LFMstableq .
\end{equation}
At $z=1$ (normalized initial flow, normalized target) this is $\frac1z+\frac1t=1+\frac1t$;
at $z=t$ (the GFlowNet convention $\Finit=Z\mu_0$) it is $\frac2t$, the constant $C$ of
\eqref{equ:paper_bound}, which \eqref{equ:z_bound} exceeds for $z<t$ and undercuts for
$z>t$.
\end{proposition}

\begin{proof}
Steps 0--3 of the proof of Theorem~\ref{theo:RL_CV_bound_full} verbatim: term $(i)$ of
Step~2 is $\leq\Delta/z$ by \eqref{equ:term_i} (which is stated for general $z$), term
$(ii)\leq\Delta/t$ by \eqref{equ:normalization_lemma}, and Step~3 bounds $\Delta$ by
$\nuB(\Statespace)^{1-1/q}\|\dln\|^{1/q}_{L^\infty(\nutrain)}\LFMstableq$.
\end{proof}

Where the initial mass is trained it is an output of the descent rather than a setting, and
Theorem~\ref{theo:RL_CV_bound_full}\emph{(3)} says when the descent can remove its defect;
what is left to the practitioner is which side of $t$ to err on --- in the initialization of
$Z$, or in a one-sided regularizer --- and that is a question about the certificate at a mass
the descent has not yet reached. By \eqref{equ:mass_defect} in the proof of
Theorem~\ref{theo:RL_CV_bound_full} the achievable $L^1(\nuB)$ error is itself floored:
$\Delta\geq|\int(e-k)d\nuB|=|z-t|$, for every flow. Suppose training is able to reach
any $\Delta$ down to a residual budget $\delta\geq0$ but no further --- $\delta$ collects
whatever the family, the optimizer and the finite run leave on the table --- so that the
best certified error at mass $z$ is
\begin{equation}\label{equ:envelope}
    B_\delta(z)\;:=\;\left(\frac1z+\frac1t\right)\max\left(\delta,\;|z-t|\right).
\end{equation}

\begin{proposition}[the certificate-optimal mass overshoots, and undershoot is dearer]
\label{prop:overshoot}
For every $\delta>0$ the map $z\mapsto B_\delta(z)$ is strictly decreasing on
$(0,\,t+\delta]$ and strictly increasing on $[t+\delta,\,\infty)$; its unique minimizer is
\begin{equation}\label{equ:zstar}
    z^*\;=\;t+\delta ,
\end{equation}
the mass that \emph{overshoots} the partition function by exactly the residual budget.
Moreover the penalty is asymmetric: for every $0<s<t$,
\begin{equation}\label{equ:asymmetry}
    B_\delta(t+s)\;<\;B_\delta(t-s):
\end{equation}
undershooting $Z$ by $s$ always certifies strictly worse than overshooting by the same
$s$. In loss currency the same holds with the dual floor: with achievable loss
$\max(\varepsilon,|z-t|/\Dq)$ the kink sits at $z^*=t+\varepsilon\Dq$.
\end{proposition}

\begin{proof}
On $(0,t+\delta]$ where the max is not yet driven by the floor on all of the interval:
for $z\leq t$, $\max(\delta,|z-t|)=\max(\delta,t-z)$ is non-increasing and $\frac1z+\frac1t$
strictly decreasing, so $B_\delta$ strictly decreases; for $t\leq z\leq t+\delta$ the max
equals $\delta$ and the first factor still strictly decreases. On $[t+\delta,\infty)$,
$B_\delta(z)=\left(\frac1z+\frac1t\right)(z-t)=\frac zt-\frac tz$, with derivative
$\frac1t+\frac t{z^2}>0$. The two monotonicities meet at $z^*=t+\delta$. For
\eqref{equ:asymmetry}: if $s\leq\delta$ both maxima equal $\delta$ and
$\frac1{t+s}<\frac1{t-s}$; if $s>\delta$ both maxima equal $s$ and the same comparison of
the first factors applies.
\end{proof}

Two remarks keep this honest. First, the gain over the default $z=t$ is second order:
$B_\delta(t)-B_\delta(z^*)=\delta\left(\frac1t-\frac1{t+\delta}\right)
=\frac{\delta^2}{t(t+\delta)}$, so \eqref{equ:zstar} is a statement about the
\emph{mechanism} --- the certificate is indifferent to overshoot up to the budget, because
extra initial mass enlarges the denominator of term $(i)$ of
\eqref{equ:term_i} --- rather than a large constant saving. Second, and asymmetrically,
the \emph{undershoot} penalty is first order and is not an artifact of the certificate:
the next proposition computes exactly what each regime does to the sampled law.

\begin{proposition}[at the floor: overshoot blurs, undershoot truncates]
\label{prop:blur_vs_drop}
In the setting of Proposition~\ref{prop:floor_exact} with $\nutrain=\nuB/\nuB(\Statespace)$
and $q\in(1,\infty)$, the minimal-loss residual profile is the constant
$\varphi=(z-t)/\nuB(\Statespace)$, and the inference-time terminal flow is
\begin{equation*}
    \widehat\target \;=\;\Bigl(k+\tfrac{z-t}{\nuB(\Statespace)}\Bigr)^{+}\nuB .
\end{equation*}
\begin{itemize}
    \item If $z\geq t$: $\widehat\target=\target+(z-t)\frac{\nuB}{\nuB(\Statespace)}$,
    the flow-matching constraint holds \emph{exactly} for $\widehat\target$, so
    $s_\tau\sim\widehat\target$ by Theorem~\ref{theo:sampling_theorem}, $\deltafinit=0$, and
    \begin{equation}\label{equ:blur}
        \TV\left(s_\tau\left\|\frac\target t\right.\right)
        \;=\;\frac{z-t}{z}\;\TV\left(\frac{\nuB}{\nuB(\Statespace)}
        \left\|\frac\target t\right.\right):
    \end{equation}
    the overshoot mass leaves as a uniform blur, in proportion $\frac{z-t}{z}$.
    \item If $z<t$: $\widehat\target=\bigl(k-\tfrac{t-z}{\nuB(\Statespace)}\bigr)^+\nuB$
    --- a soft-threshold of the target at the water level
    $\ell:=\frac{t-z}{\nuB(\Statespace)}$. Whenever the level exceeds the bottom of
    the target density --- $\ell>\operatorname*{ess\,inf}_{\nuB} k$, i.e.\ as soon as
    the undershoot $t-z$ exceeds $\nuB(\Statespace)\operatorname*{ess\,inf}_{\nuB} k$,
    which for a non-constant target happens throughout an interval of $z$ ---
    every state below the level receives \emph{zero} terminal mass: the missing mass is
    paid first by the lowest-density part of the target (mode drop), and additionally
    $\deltafinit=E^->0$, so the sampler incurs the term-$(i)$ error of
    \eqref{equ:term_i} on top of the truncation. (For a uniform target the threshold is
    never crossed and undershoot only rescales --- the asymmetry of
    Proposition~\ref{prop:overshoot} is then purely the certificate's.)
\end{itemize}
\end{proposition}

\begin{proof}
For uniform $\nutrain$, $g=\dln$ is constant, so the extremizer
\eqref{equ:extremal_profile} is constant, and for $q\in(1,\infty)$ it is the unique
minimal-loss profile (strict convexity of the $L^q$ ball). If $z\geq t$ then
$e=k+\frac{z-t}{\nuB(\Statespace)}\geq0$, so $E^-=0$, $E=E^+=\widehat\target$ and the
corrected flow is flow-matching with terminal flow $\widehat\target$ of mass $z$;
Theorem~\ref{theo:sampling_theorem} gives $s_\tau\sim\widehat\target$ exactly. Then
$\frac{\widehat\target}{z}=\frac tz\cdot\frac\target t+\frac{z-t}{z}\cdot
\frac{\nuB}{\nuB(\Statespace)}$ and
$\TV(\alpha P+(1-\alpha)Q\,\|\,P)=(1-\alpha)\TV(Q\|P)$ yields \eqref{equ:blur}. If $z<t$
the formula for $\widehat\target=E^+$ is the definition of the positive part, and
$E^-=\bigl(\ell-k\bigr)^+\nuB\neq0$ exactly when $\nuB(\{k<\ell\})>0$, i.e.\ when
the level exceeds $\operatorname*{ess\,inf}_{\nuB} k$ --- the stated threshold. The
threshold is not removable: for the uniform target $k\equiv t/\nuB(\Statespace)$ and
any $0<t-z<t$, $e=k-\ell=z/\nuB(\Statespace)>0$ everywhere, so $E^-=0$ and nothing
truncates.
\end{proof}

Read together, Propositions~\ref{prop:overshoot} and \ref{prop:blur_vs_drop} answer the
mass question: \emph{at the certificate level the optimal initial mass overshoots $Z$ by
exactly the residual budget, and the preference for overshoot is not a bound artifact ---
undershoot truncates the target from the bottom while overshoot only dilutes it.} Exactly
at the floor ($\delta=0$) the optimum is $z^*=t$, i.e.\ the GFlowNet convention
$\Finit=Z\mu_0$ with $Z$ learned exactly; the overshoot rule is the first-order correction
when the training residual is not zero.

\subsection{The shape of the initial flow: a one-sided Chebyshev problem}\label{sec:shape}

By Corollary~\ref{cor:separation} the loss cannot see the shape of $\Finit$. What the shape
buys --- or costs --- is \emph{inference}: how much outflow, hence how many sampling steps,
the matching flow must carry. On finite spaces this cost has a closed form.

\begin{proposition}[exact inference cost of a shape mismatch]\label{prop:shape_cost}
In the setting of Proposition~\ref{prop:floor_exact}, let $z=t$ and
$\theta:=k-\finit$ (so $\int\theta\,d\nuB=0$). Consider the flows achieving
\emph{exact} flow matching for the pair $(\Finit,\target)$ with the given star policy.
\begin{enumerate}
    \item The set of exact non-negative outflows is exactly the ray
    $\{\foutstar=\eta\mathbf 1-S\theta:\ \eta\geq\max_xS\theta(x)\}$, and the minimal
    outflow mass is
    \begin{equation}\label{equ:min_outflow}
        \Foutmin(\Statespace)\;=\;\nuB(\Statespace)\,
        \max_{x\in\Statespace}\,S\bigl(k-\finit\bigr)(x)\;\geq\;0,
    \end{equation}
    with equality to $0$ iff $\finit=k$.
    \item For any exact flow, the visit measure of the sampler is
    $\mu=\Fterm+\Foutstar$, so the expected number of visits per trajectory is
    \emph{exactly} $1+\Foutstar(\Statespace)/t$ --- the inequality of
    Theorem~\ref{theo:sampling_theorem} is an identity here; credit where due: this
    equality is Proposition~3.12 of \citet{morozov2025revisiting} (expected trajectory
    length $=$ normalized total flow), independently re-derived here through the cut
    construction --- and its minimum over exact flows is
    $1+\nuB(\Statespace)\max_xS\theta(x)/t$.
    \item Over an admissible class $\mathcal F\subseteq\{\text{densities of mass }t\}$,
    the optimal shape solves the \emph{one-sided Chebyshev problem}
    \begin{equation}\label{equ:chebyshev}
        \min_{\finit\in\mathcal F}\ \max_{x\in\Statespace}\,S\bigl(k-\finit\bigr)(x),
    \end{equation}
    a convex program whenever $\mathcal F$ is convex ($S$ is linear and the max of linear
    functionals convex).
    \item ($L^2$ relaxation and the slow modes.) If $\ForwardPolicyStar$ is reversible
    with spectrum $1=\beta_0>\beta_1\geq\dots$, then $S$ acts on the $\beta$-eigenmode as
    multiplication by $\frac1{1-\beta}$, so
    $\|S\theta\|^2_{L^2(\nuB)}=\sum_{i\geq1}\frac{\langle\theta,v_i\rangle^2_{L^2(\nuB)}}{(1-\beta_i)^2}$:
    mismatch along the \emph{slow} modes of the star policy is what costs outflow and
    trajectory length; fast-mode mismatch is nearly free. Over an affine class the
    relaxed optimum is the $\mathcal F$-projection of $k$ in the seminorm $\|S\cdot\|_{L^2(\nuB)}$,
    which is a genuine seminorm equivalent to $L^2$ on mean-zero functions
    ($\tfrac12\|\varphi\|\leq\|S\varphi\|\leq(1-\beta_1)^{-1}\|\varphi\|$ with $\beta_1$
    the largest sub-unit eigenvalue --- finite by ergodicity alone, since a reversible
    spectrum lies in $[-1,1)$ off the constants).
\end{enumerate}
\end{proposition}

\begin{proof}
(1) Exactness means $(P_\star-\mathrm{Id})\foutstar=\theta$. A particular solution is
$-S\theta$ by the identity $(\mathrm{Id}-P_\star)S=\mathrm{Id}-\Pi$ established in the
proof of Proposition~\ref{prop:floor_exact}; the homogeneous solutions are the fixed
functions of $P_\star$, which by $\nuB$-ergodicity are the constants (a fixed function
$f$ of the density action makes $f\nuB$ a fixed measure $\ll\nuB$, hence a multiple
of $\nuB$). Non-negativity of $\eta\mathbf 1-S\theta$ is $\eta\geq\max S\theta$, and
$\int S\theta\,d\nuB=0$ ($\Pi S=0$) gives the mass formula. The max is $\geq0$ because
$S\theta$ is mean-zero, with equality iff $S\theta=0$ iff $\theta=0$.
(2) The through-flow at each state of a matching flow is
$\Finit+\Finstar=\Fterm+\Foutstar$, and the visit measure solves
$\mu=\Finit+(\mu\odot(1-p_{\mathrm{stop}}))\ForwardPolicyStar$ with
$p_{\mathrm{stop}}=\frac{d\Fterm}{d(\Fterm+\Foutstar)}$; the candidate
$\mu=\Fterm+\Foutstar$ satisfies it --- the continuing measure is then exactly
$\Foutstar$, and $\Finit+\Foutstar\ForwardPolicyStar=\Finit+\Finstar=\mu$ --- and the
absorbing-chain solve is unique. Hence $\mu(\Statespace)=t+\Foutstar(\Statespace)$.
(3) Immediate. (4) For a mean-zero eigenmode, $(\mathrm{Id}-P_\star+\Pi)v_i=(1-\beta_i)v_i$
and $\Pi v_i=0$, so $Sv_i=\frac{v_i}{1-\beta_i}$ for
$i\geq1$; the lower seminorm bound is $\|\varphi\|=\|(\mathrm{Id}-P_\star)S\varphi\|\leq
2\|S\varphi\|$.
\end{proof}

The degenerate optimum of \eqref{equ:chebyshev} --- $\finit=k$, zero outflow, trajectories
of length zero: sample the target directly --- is precisely why the class $\mathcal F$
carries all the content: $\Finit$ must remain \emph{samplable}, and
\eqref{equ:chebyshev} prices the distance between what is samplable and the target
\emph{in the geometry of the star policy}, slow modes first. A warm-started initial
distribution that matches the target's mass across the slow modes (e.g.\ across the
clusters of a two-cluster state space) removes almost the entire inference cost, while
matching the fast detail buys almost nothing --- the numeric instance of
\S\ref{sec:num_finit} makes the two costs differ by a factor $30$ at equal mismatch
amplitude.

\begin{remark}[training does not see a frozen $\Finit$]\label{rem:training}
Two channels one might expect are absent, and the absence is worth recording.
(1)~In the frozen-backward parameterization of Section~\ref{sec:cv_divergence} the initial
flow is an \emph{output} of the construction, not an input: the balanced flow of
Theorem~\ref{theo:universality_graphs} carries
$\Finit(v)=c\,\lambda(v)\,\BackwardPolicy(v\rightarrow s_0)$, so within that family there
is no $\Finit$ to optimize --- the backward policy already chose it.
(2)~For the free-outflow stable pipeline at $q=2$ and a \emph{fixed} initial flow, the loss
$\foutstar\mapsto\|(\finit-k)+(P_\star-\mathrm{Id})\foutstar\|^2_{L^2(\nutrain)}$ is a
linear least-squares problem whose Hessian $2(P_\star-\mathrm{Id})^\dagger
M_{d\nutrain/d\nuB}(P_\star-\mathrm{Id})$ does not involve $\finit$: the convergence
\emph{rate} of gradient descent is $\Finit$-independent, and $\Finit$ enters only through
the distance of the initialization to the solution set (a warm start moves the constant
$\log\|f_0-f^*\|$, not the rate). Likewise Theorem~\ref{theo:IL_CV_bound} fixes $\Finit$ to
a probability, so no mass freedom exists there. The levers of a frozen initial flow are
therefore exactly the two of \S\ref{sec:mass} and \S\ref{sec:shape}: mass (loss floor and
certificate) and shape (inference cost); training dynamics are indifferent.

The second claim is about a frozen $\Finit$ and does not survive a trained mass. With
$\Finit=z\mu_0$, $\mu_0$ a fixed probability and $z$ a parameter --- the setting of
Theorem~\ref{theo:RL_CV_bound_full}\emph{(3)} --- the residual
$zm_0-k+(P_\star-\mathrm{Id})\foutstar$, with $m_0:=\frac{d\mu_0}{d\nuB}$, is jointly linear
in $(z,\foutstar)$, and the $q=2$ Hessian becomes the Gram matrix of
$\{m_0\}\cup\operatorname{range}(P_\star-\mathrm{Id})$ in $L^2(\nutrain)$, whose conditioning
turns on the angle between $m_0$ and that range. The rate is $\Finit$-dependent there,
through the shape $\mu_0$ and not through the mass; what \S\ref{sec:mass} and
\S\ref{sec:shape} price is the certificate at a given mass, not the descent that reaches
it.
\end{remark}

\subsection{Numerical evidence}\label{sec:num_finit}

All computations follow the doctrine of Section~\ref{sec:experiments}: enumerable state
spaces, every quantity computed rather than estimated (the sampled law of $s_\tau$ is a
linear solve). The script \texttt{audit\_finit.py}, released with the code, re-derives
every number below and exits nonzero on any violation.

\paragraph{The dual floor is attained, and dominates the $L^\infty$ floor.}
Over $2{,}000$ random instances (state size in $[6,50)$, random ergodic star policy with
its scaled invariant measure, random $\nutrain$, $q\in\{1,\frac32,2,3,\infty\}$,
$z,t\in[0.3,3]$): the explicit extremal flow of Proposition~\ref{prop:floor_exact} attains
$\LFMstableq=|z-t|/\Dq$ with worst relative error $4.4\cdot10^{-14}$; the dual floor
dominates the floor \eqref{equ:mass_floor} on every instance to machine
precision (worst signed gap $-9\cdot10^{-16}$); $2{,}000$ random
competitor flows never went below it.

\paragraph{The mass envelope has its kink at $t+\delta$, and the asymmetry is real.}
On a $40$-state instance swept over $121$ values of $z$
(Figure~\ref{fig:app_mass}, left): for $z>t$ the corrected flow is exactly flow-matching
and its \emph{exact chain law} equals $\widehat\target/z$ to $4.6\cdot10^{-16}$; the
mixture identity \eqref{equ:blur} holds to $4.9\cdot10^{-17}$. For $z<t$ the terminal
flow is the water-level truncation --- up to $11$ of $40$ states receive zero terminal
mass at the lowest $z$ tested --- and the term-$(i)$ bound of \eqref{equ:term_i} holds on
the exact chain at every $z$. The envelope $B_\delta$ has its argmin at $t+\delta$ to
grid resolution ($2\cdot10^{-4}$), and $B_\delta(t+s)<B_\delta(t-s)$ at all $500$ tested
$s$.

\paragraph{Slow-mode mismatch costs $30\times$ the outflow of fast-mode mismatch.}
On a two-cluster reversible instance (bridge conductance $0.02$;
Figure~\ref{fig:app_mass}, right): the eigen-relation $Sv=v/(1-\beta)$ and the identities
$\Foutmin(\Statespace)=\nuB(\Statespace)\max S\theta$ and
$\mu(\Statespace)/z=1+\Foutstar(\Statespace)/\Finit(\Statespace)$ all hold to
$10^{-9}$ or better. At \emph{equal mismatch amplitude}: mismatch along the slow mode
($\beta=+0.9601$) costs outflow mass $1.1593$ and $2.159$ visits per trajectory; along a
fast mode ($\beta=-0.1857$) it costs $0.0390$ and $1.039$ --- a factor $29.7$ in outflow
for the same size of initialization error, placed differently. Both the Chebyshev and the
$L^2$ objectives were verified convex along $400$ random segments of a mixture class.

\begin{figure*}[tb]
\centering
\includegraphics[width=0.45\linewidth]{figures/figA1_mass.pdf}\hfill
\includegraphics[width=0.45\linewidth]{figures/figA2_shape.pdf}
\caption{\textbf{Left:} the certified-error envelope $B_\delta(z)$ and the exactly
computed $\TV$ of the minimal-loss flow against the initial mass $z$, with the kink at
$z^*=t+\delta$ and the truncation fraction rising only on the undershoot side
(Propositions~\ref{prop:overshoot} and~\ref{prop:blur_vs_drop}).
\textbf{Right:} minimal outflow mass of an exact flow against the mode of the mismatch,
one point per eigenmode of the two-cluster star policy, against the predicted
$1/(1-\beta)$ weighting (Proposition~\ref{prop:shape_cost}(4)); the slow mode costs
$29.7\times$ the fast one at equal amplitude.}
\label{fig:app_mass}
\end{figure*}

% Appendix D -- the replay buffer and the training measure.
\section{\appendixDtitle}\label{app:replay}

The training measure $\nutrain$ enters the theory twice, against two \emph{different}
reference measures. In the stable bound \eqref{equ:paper_bound} it enters against the
background $\nuB$, through $A_q(\nutrain):=\|\dln\|^{1/q}_{L^\infty(\nutrain)}$.
In the training-speed theory of Section~\ref{sec:cv_divergence} it enters against the
backward-chain occupation probability $\lambda$, through $w:=\frac{d\nutrain}{d\lambda}$:
the exponential rate is $\varrho_\sigma=g''(1)w_{\min}\lambda_{\min}/\sigma_*^2$
(Theorem~\ref{theo:training_speed_full}), the \L ojasiewicz constant is
$\kappa=w_{\min}\lambda_{\min}^{1/2}/(\|u_0\|\,\|w\|_{L^\infty}M)$
(Corollary~\ref{cor:global_lojasiewicz}), and the gradient-descent step cap is
$\varepsilon\leq(4g''(1)\|w\|_{L^\infty})^{-1}$
(Theorem~\ref{theo:db_stable_frozen_full}).
A replay buffer is precisely a device that shapes the empirical $\nutrain$; the question
is what it should shape it \emph{into}, and the answer splits by which pipeline carries
the guarantee.

\subsection{The buffer as a $\nutrain$-shaping operator, and coverage}\label{sec:buffer_def}

\begin{definition}\label{def:buffer}
A \emph{replay schedule} is a family of probability measures
$\nutrain = \alpha\,\mu_{\mathrm{on}} + (1-\alpha)\,\rho_{\mathrm{buf}}$, $\alpha\in[0,1]$,
where $\mu_{\mathrm{on}}$ is the on-policy visitation measure of the current sampler and
$\rho_{\mathrm{buf}}$ is the replay measure of a buffer --- e.g.\ uniform over stored
states, importance-reweighted (\emph{inverse-frequency}) replay, or
\emph{backward replay}: draw a stored terminal state and replay the states of a frozen
backward trajectory from it, with multiplicity.
\end{definition}

Coverage is the zeroth-order requirement, and for an empirical buffer on a finite state
space it is a coupon-collector statement: the certificate is infinite until every state
has been trained on at least once.

\begin{lemma}[coverage, and its clock]\label{lem:coverage}
If $\nutrain$ fails to dominate $\nuB$ --- some set of positive $\nuB$-measure is
never trained on --- then $A_q(\nutrain)=+\infty$ and the bound \eqref{equ:paper_bound}
is vacuous; likewise $w_{\min}=0$ and every rate of Section~\ref{sec:cv_divergence}
degenerates. Under backward replay in the setting of Proposition~\ref{prop:morozov_rate},
with its $N$ and $\bar\sigma$, the mean recurrence time of a state $x$ in the concatenated
excursion process is exactly $(2+\bar\sigma)/N(x)$, and the reciprocal of the smallest
visit ratio, $\max_x(2+\bar\sigma)/N(x)$, is a \emph{lower} bound on the expected number
of replayed steps until every state has been replayed, excursions entering the buffer
whole.
\end{lemma}

\begin{proof}
If $\nuB(A)>0=\nutrain(A)$ then $\nuB\not\ll\nutrain$: no density
$\dln$ exists and the essential supremum defining $A_q$ is
$+\infty$ under the usual convention, while any flow whose residual sits inside $A$ is
invisible to $\LFMstableq$, so no finite constant can hold in \eqref{equ:paper_bound}.
On the training side, $\lambda$ is everywhere positive on the loop closure, so
$w=\frac{d\nutrain}{d\lambda}$ has essential infimum $0$ on $A$, and every constant in
the display above vanishes with $w_{\min}$. For the clock: concatenating the i.i.d.\
excursions reproduces the loop-closed chain regenerating at $s_0$ --- the structure
behind the occupation identity, Proposition~\ref{prop:morozov_rate}(1) --- an
irreducible positive-recurrent chain with invariant probability $\lambda$; Kac's
recurrence formula gives mean return time $\mathbb E(\tau_x^+\mid X_0=x)=1/\lambda(x)$, and the
occupation identity turns it into $(2+\bar\sigma)/N(x)$. For coverage, an excursion
replays $s_0$, $s_f$, then the states visited by the backward trajectory of a draw from
$\widehat\BackwardPolicy(s_f\rightarrow\cdot)$ before it hits $s_0$, so by definition of $N$
its expected number of visits to $x$ is $N(x)$. Let $L_i$ be the number of states the
$i$-th excursion replays, $K_x$ the index of the first excursion visiting $x$, and
$p_x:=\mathbb P(\text{an excursion visits }x)$. Then
$\mathbb E(L_i)=\sum_xN(x)=2+\bar\sigma=1/\lambda(s_0)<\infty$, the first equality by
linearity and the others by Proposition~\ref{prop:morozov_rate}(1), $\lambda$ being a
probability; and $0<p_x\leq N(x)$, the upper bound by Markov's inequality and the lower bound since
$N(x)=(2+\bar\sigma)\lambda(x)>0$. The excursions being i.i.d., $K_x$ is a stopping time
for them, geometric with $\mathbb E(K_x)=1/p_x$, and Wald's identity gives
\begin{equation*}
    \mathbb E\Bigl(\sum_{i\leq K_x}L_i\Bigr)=\mathbb E(K_x)\,\mathbb E(L_1)
    =\frac{2+\bar\sigma}{p_x}\geq\frac{2+\bar\sigma}{N(x)}.
\end{equation*}
Excursions entering the buffer whole, the number of replayed steps until every state has
been replayed is $\sum_{i\leq\max_yK_y}L_i$, which is at least $\sum_{i\leq K_x}L_i$ for
every $x$, the $L_i$ being non-negative; taking expectations, then the maximum over $x$,
gives the lower bound.
\end{proof}

Heuristically, the reciprocal visit ratio $(2+\bar\sigma)/N(x)$ is also the time scale of
the slow modes of the global phase of training, which Remark~\ref{rem:visit_ratio} locates
on the rarely visited states.

\subsection{Backward replay realizes the rate-optimal measure}\label{sec:backward_replay}

\begin{proposition}[$\lambda$ is the guarantee-optimal training measure, and backward
replay samples it]\label{prop:backward_replay}
Let $(\widehat G,\widehat\BackwardPolicy)$ be the loop closure of a finite path-connected
marked graph with backward policy positive on its edges, $\lambda$ its invariant
probability. Over probability measures $\nutrain=w\lambda$:
\begin{enumerate}
    \item $w_{\min}\leq1\leq\|w\|_{L^\infty}$, each with equality iff $\nutrain=\lambda$.
    \item Consequently each of the three training-speed constants --- the
    exponential rate $\varrho_\sigma$ (equivalently $\varrho=g''(1)w_{\min}/\widehat B^2$),
    the \L ojasiewicz constant $\kappa$, and the step cap
    $(4g''(1)\|w\|_{L^\infty})^{-1}$ --- is maximized at $\nutrain=\lambda$, uniquely
    among $\lambda$-a.c.\ probabilities.
    \item Backward replay from the target row samples $\lambda$ \emph{exactly}: the
    long-run (renewal) state-frequency of excursions (one visit each to $s_0$ and $s_f$,
    then the backward trajectory of a draw from
    $\widehat\BackwardPolicy(s_f\rightarrow\cdot)$ until it hits $s_0$) is
    $N(x)/(2+\bar\sigma)=\lambda(x)$, by Proposition~\ref{prop:morozov_rate}(1). If the
    buffer's marginal over terminal states is
    $\mu_{\mathrm{buf}}\neq\widehat\BackwardPolicy(s_f\rightarrow\cdot)$,
    the realized measure is the normalized expected-visit vector of the chain killed at
    $s_0$ started from $\mu_{\mathrm{buf}}$ --- computable by one fundamental-matrix
    solve, so the bias of a stale or reward-filtered buffer is a quantity, not a worry.
\end{enumerate}
\end{proposition}

\begin{proof}
(1) $\int w\,d\lambda=\nutrain(\Statespace)=1$, so $\operatorname*{ess\,inf}w\leq1\leq
\operatorname*{ess\,sup}w$, with equality on either side iff $w=1$ $\lambda$-a.e.
(2) $\varrho_\sigma$ and $\varrho$ are increasing in $w_{\min}$; the step cap is
decreasing in $\|w\|_{L^\infty}$; and $\kappa$ is increasing in $w_{\min}$ (both directly
and through $M=\max(1,\sqrt{\mathcal L(\mu_0)/(w_{\min}\lambda_{\min})})$) and decreasing
in $\|w\|_{L^\infty}$. Apply (1).
(3) The first claim is the occupation identity read as a sampling algorithm; the second
is linearity of expected visits in the starting law.
\end{proof}

An honesty clause: $\varrho_\sigma$ is \emph{a guarantee, not a
prediction} (the measured gap is $187$--$10^9\times$, Section~\ref{sec:experiments}), so
Proposition~\ref{prop:backward_replay} optimizes the guarantee. The exact linearized rate
is the smallest positive eigenvalue of $H_\nu=g''(1)A^\dagger M_wA$, and
\S\ref{sec:num_replay} measures where \emph{it} peaks on the $\alpha$-mixtures beside
where the guarantee says; the gap is reported, not hidden.

\subsection{Constant against rate: two pipelines, two verdicts}\label{sec:tension}

\begin{proposition}[stable pipeline: uniform reweighting is bound-optimal, with no
tension]\label{prop:stable_uniform}
Let $q\in[1,\infty]$ and let $\nutrain$ range over probability measures with
$\nuB\ll\nutrain$.
\begin{enumerate}
    \item $\left\|\dln\right\|_{L^\infty(\nutrain)}\geq\nuB(\Statespace)$,
    with equality iff $\nutrain=\nuB/\nuB(\Statespace)$; hence
    $A_q(\nutrain)\geq\nuB(\Statespace)^{1/q}$ with the same equality case, and the
    constant of \eqref{equ:paper_bound} is minimized by uniform-in-$\nuB$ training.
    \item At $q=2$ with free non-negative outflow and $P_\star$ summably $L^2$-mixing
    (so that $B_2<\infty$, Lemma~\ref{lem:sigma_mixing}), gradient descent on
    $\foutstar\mapsto\|(\finit-k)+(P_\star-\mathrm{Id})\foutstar\|^2_{L^2(\nutrain)}$
    contracts the transverse component at rate at least
    $2\inf_\Statespace\frac{d\nutrain}{d\nuB}\big/B_2^2$, and
    $\inf\frac{d\nutrain}{d\nuB}\leq1/\nuB(\Statespace)$ with equality iff
    $\nutrain=\nuB/\nuB(\Statespace)$.
\end{enumerate}
Both the certificate constant and the rate guarantee of the stable pipeline are therefore
optimized by the \emph{same} training measure, the uniform one: for the stable loss a
replay buffer should simply flatten the empirical measure against $\nuB$
(inverse-frequency replay over the visited support), and there is nothing to trade.
\end{proposition}

\begin{proof}
(1) $\int\dln{}d\nutrain=\nuB(\Statespace)$, so the essential sup
is $\geq$ the mean, with equality iff the density is constant.
(2) The loss is quadratic with Hessian $2(P_\star-\mathrm{Id})^\dagger
M_{d\nutrain/d\nuB}(P_\star-\mathrm{Id})$;
$\langle f,Hf\rangle_{L^2(\nuB)}\geq2\inf(\tfrac{d\nutrain}{d\nuB})\|(P_\star-\mathrm{Id})f\|^2_{L^2(\nuB)}
\geq\frac{2\inf(d\nutrain/d\nuB)}{B_2^2}\|f-\Pi f\|^2$ by the coercivity of
Lemma~\ref{lem:sigma_mixing} applied to $P_\star$; and the mean of
$\frac{d\nutrain}{d\nuB}$ against $\nuB/\nuB(\Statespace)$ is
$1/\nuB(\Statespace)$.
\end{proof}

For the \emph{practical} pipeline --- the DB/TB balance losses with frozen backward
policy, $g=(\log x)^2$ --- the two guarantees genuinely pull apart: the rate wants
$\nutrain=\lambda$ (Proposition~\ref{prop:backward_replay}), the certificate wants
$\nutrain\propto\nuB$ (Proposition~\ref{prop:stable_uniform}(1)), and $\lambda$ is
uniform in $\nuB$ only when every state has the same visit ratio. The bridge between
the two loss families is the following lemma, which is where the sink-normalization
subtlety lives.

\begin{lemma}[transfer: balance loss to stable loss, exactly]\label{lem:transfer}
On the loop closure above, let $\mu$ be a positive state flow, $r:=\frac{d(\mu\widehat\BackwardPolicy)}{d\mu}$,
and let the cut generative flow of $\mu$ be defined edge-wise by
$F(u\rightarrow v)=\mu(v)\widehat\BackwardPolicy(v\rightarrow u)$, with
$\Finit:=F(s_0\rightarrow\cdot)$ and terminal flow $F(\cdot\rightarrow s_f)$
(Theorem~\ref{theo:universality_graphs}). Normalize by the sink:
$\widetilde\mu:=\frac{t}{\mu(s_f)}\mu$, $\widetilde u:=\frac{d\widetilde\mu}{d\lambda}$,
where $t:=\target(\Statespace)$ is the target mass.
Then the flow-matching residual of the cut flow of $\widetilde\mu$ against the target
$\target$ is, at every interior state,
\begin{equation}\label{equ:transfer_identity}
    e-k \;=\; -\,(r-1)\,\widetilde u\,\frac{d\lambda}{d\nuB}
\end{equation}
($r$ is scale-invariant, so it may be read on $\mu$). Consequently, for any training
measure $\nutrain$ and any $a\in(0,1)$ with $r$ valued in $[1-a,1+a]$,
\begin{equation}\label{equ:transfer_bound}
    \LFMstableq\;\leq\;\|\widetilde u\|_{L^\infty}
    \left\|\frac{d\lambda}{d\nuB}\right\|_{L^\infty}\|r-1\|_{L^q(\nutrain)},
    \qquad
    \|r-1\|_{L^2(\nutrain)}\leq\sqrt{\frac{\mathcal L_{g,\nutrain}(\mu)}{c_g(a)}},
\end{equation}
with $c_g(a):=\inf_{|x-1|\leq a}g(x)/(x-1)^2>0$ (for $g=(\log x)^2$ and $a=\frac12$,
$c_g=(2\log\frac32)^2\approx0.657$).
\end{lemma}

\begin{proof}
At an interior state $v$, every in-edge of $v$ in $\widehat G$ lies in $G$, so the total
inflow of the cut flow is $\sum_uF(u\rightarrow v)=\mu(v)\sum_u\widehat\BackwardPolicy(v\rightarrow u)=\mu(v)$,
while the outflow through interior edges plus the sink edge is
$\sum_{v'}\mu(v')\widehat\BackwardPolicy(v'\rightarrow v)=(\mu\widehat\BackwardPolicy)(v)=r(v)\mu(v)$.
Hence, splitting off the sink edge $F(v\rightarrow s_f)=\mu(s_f)\widehat\BackwardPolicy(s_f\rightarrow v)
=\mu(s_f)\frac{\target(v)}{t}$ (the frozen sink row is the normalized target,
Proposition~\ref{prop:tb_gradient}),
\begin{equation*}
    E(v)=\mu(v)\bigl(1-r(v)\bigr)+\frac{\mu(s_f)}{t}\,\target(v),
    \qquad
    E-\target=-(r-1)\mu+\Bigl(\frac{\mu(s_f)}{t}-1\Bigr)\target .
\end{equation*}
The second term is killed exactly by the sink normalization $\mu\mapsto\widetilde\mu$,
which leaves $r$ unchanged; \eqref{equ:transfer_identity} is this identity in
$\nuB$-densities. The first inequality of \eqref{equ:transfer_bound} is H\"older with
$L^\infty$ factors; the second is $g(x)\geq c_g(a)(x-1)^2$ on the band, integrated
against $\nutrain$.
\end{proof}

The mass term $\bigl(\frac{\mu(s_f)}{t}-1\bigr)\target$ that the normalization removes is
the initial-flow question of Appendix~\ref{app:finit} arriving from the other side:
$\mu(s_f)$, the terminal mass of the cut flow, \emph{is} its implied partition function
--- at a balanced $\mu$ it equals $\mu(s_0)=F(s_0)$, which plays the role of $Z$ in
Proposition~\ref{prop:tb_gradient}, and off balance $F(s_0)=r(s_0)\,r(s_f)\,\mu(s_f)$ ---
and an unnormalized read-out pays exactly the mass floor of \S\ref{sec:mass}. The balance loss
is scale-invariant, so it can never learn this scale --- normalizing by the sink at
read-out, as \eqref{equ:transfer_identity} prescribes, is the frozen-backward analogue of
learning $Z$, and it is free.

\begin{proposition}[the certificate-optimal mixture: uniform mixed into backward replay
at weight $\sim1/T$]\label{prop:alpha_star}
On the mixture family $\nu_\alpha:=(1-\alpha)\lambda+\alpha\,\nuB/\nuB(\Statespace)$,
$\alpha\in(0,1]$:
\begin{equation*}
    \left\|\frac{d\nuB}{d\nu_\alpha}\right\|_{L^\infty}\leq\frac{\nuB(\Statespace)}{\alpha},
    \qquad w_{\min}(\nu_\alpha)\geq1-\alpha ,
\end{equation*}
so, chaining Theorem~\ref{theo:training_speed_full} (local phase at rate
$\varrho_\sigma(\nu_\alpha)/2\geq(1-\alpha)\widehat\varrho/2$,
$\widehat\varrho:=\varrho_\sigma(\lambda)$), Lemma~\ref{lem:transfer} at $q=2$ and the
$z$-carried bound of Proposition~\ref{prop:z_bound}, the certified error after training
time $T$ has the envelope
\begin{equation}\label{equ:JT}
    J_T(\alpha)\;=\;K\left(\frac{\nuB(\Statespace)}{\alpha}\right)^{1/q}
    e^{-c\,(1-\alpha)\,T},
\end{equation}
with $c$ the realized exponential rate of $\sqrt{\mathcal L_{g,\nu}}$ per unit
$w_{\min}$ and $K$ collecting the $T$- and $\alpha$-independent factors. $J_T$ is
log-convex in $\alpha$ with unique minimizer
\begin{equation}\label{equ:alpha_star}
    \alpha^*(T)\;=\;\min\left(1,\;\frac{1}{q\,c\,T}\right):
\end{equation}
the certificate-optimal buffer replays backward trajectories, tempered by
uniform-in-$\nuB$ reweighting whose weight decays like $1/T$ --- early training buys
coverage for the constant, late training buys rate.
\end{proposition}

\begin{proof}
$\nu_\alpha\geq\alpha\nuB/\nuB(\Statespace)$ gives the density bound;
$\nu_\alpha\geq(1-\alpha)\lambda$ gives $w_{\min}\geq1-\alpha$. For \eqref{equ:JT}, the
local phase contracts $\mathcal L_{g,\nu_\alpha}$ at rate $\varrho_\sigma(\nu_\alpha)/2$
and enters the TV bound under a square root through \eqref{equ:transfer_bound};
absorbing the halvings into $c$ leaves the stated $\alpha$-dependence. Then
$\log J_T(\alpha)=\mathrm{const}+\frac1q\log\frac1\alpha+c T\alpha$ is strictly convex
with critical point $\frac{1}{qcT}$, clipped at $1$.
\end{proof}

Three honest caveats. \emph{(i)} \eqref{equ:JT} chains \emph{guarantees}; the true error
of a trained flow can sit far below it, and the numerics measure the exact-rate optimum
beside the certified one. \emph{(ii)} The rate factor $(1-\alpha)$ uses
$w_{\min}(\nu_\alpha)\geq1-\alpha$, which discards the uniform component's own
contribution to $w$; on graphs where $\nuB/\nuB(\Statespace)$ and $\lambda$ are
close the true rate barely degrades with $\alpha$ and the optimal $\alpha$ is larger.
\emph{(iii)} The crossover $T_0$ of the global phase also depends on $\nu$ through
$w_{\min}$ and $\|w\|_{L^\infty}$, in the same direction --- so the $\alpha$-schedule
\eqref{equ:alpha_star} is conservative about, not contradicted by, the early phase.

\subsection{Positioning}\label{sec:positioning}

\citet{bartoldson2025trajectory} decouple exploration from learning for LLM post-training
--- searcher processes fill a replay buffer, a trainer consumes it off-policy with TB ---
and report large wall-clock speedups; the graphs are autoregressive, hence acyclic and
leveled, and the buffer policy (what to store, what to replay) is an engineering
surface tuned empirically. The statements above are the theory-side complement on
possibly-cyclic graphs: at the level of the guarantees of this paper, the question
``which measure should replay realize?'' has an exact answer --- the backward occupation
measure $\lambda$ for the training rate (and backward replay from terminal draws samples
it with no estimation step, by the occupation identity), tempered toward
uniform-in-$\nuB$ with weight $\sim1/T$ when a total-variation certificate is wanted
at read-out. The inverse-frequency heuristic of prioritized replay appears here with a
proof and a scope: it is exactly right for the \emph{stable} pipeline
(Proposition~\ref{prop:stable_uniform}), and exactly the $\alpha$-component for the
practical one. \citet{morozov2025revisiting}'s on-policy forward sampling, read in this
frame, trains on $\mu_{\mathrm{on}}$ --- and at balance the sampler's visit measure is
the total through-flow $r\mu=c\lambda$, i.e.\ exactly $\lambda$ on the interior states
--- so the frame also explains why on-policy training of a frozen-backward flow behaves
well near convergence and poorly before coverage.

\emph{Against the minimum-flow/OT reading of the same quantity.} The outflow mass that
Appendix~\ref{app:finit} prices is the object \citet{maksimov2026ot} characterize
from the other end: Proposition~3.12 of \citet{morozov2025revisiting} makes
expected trajectory length \emph{equal} to normalized total flow, and
\citet{maksimov2026ot} show that minimizing total flow \emph{over the policy as well},
with $\Finit$ fixed, is a Kantorovich optimal-transport problem with graph shortest-path
costs --- the free-policy endpoint of the question. Proposition~\ref{prop:shape_cost} is
the \emph{fixed-policy} endpoint: with the star policy prescribed and ergodic, the
minimal flow is not a geodesic quantity but a potential-theoretic one,
$\nuB(\Statespace)\max_xS(k-\finit)$, weighted $\frac1{1-\beta}$ on the policy's slow
modes. The ratio of the two --- fixed-policy minimal flow over OT minimal flow --- is
the price of routing transport through a prescribed mixing kernel rather than along
graph geodesics, and it is finite exactly when the policy mixes; making that ratio's
spectral control precise is an open question this comparison surfaces.
Appendix~\ref{app:pbtrain} works out what happens when the policy itself is trained
toward the OT endpoint.

\subsection{Numerical evidence}\label{sec:num_replay}

The script \texttt{audit\_replay.py}, released with the code, computes everything below
exactly (the one simulated quantity is compared to its exact prediction).

\paragraph{$\lambda$ is never beaten on any guarantee.}
Over $10$ random loop-closed marked graphs $\times$ $1{,}000$ random probability
$\nutrain$ each: the occupation identity holds to $4.2\cdot10^{-16}$, and none of the
$10{,}000$ draws beats $\lambda$ on $w_{\min}$, on the \L ojasiewicz factor, or on the
step cap --- nor the uniform measure on the certificate factor
(Proposition~\ref{prop:stable_uniform}(1)).

\paragraph{The guarantee peaks at $\lambda$; the exact rate peaks at an interior
mixture.}
The smallest positive eigenvalue of $H_\nu=g''(1)A^\dagger M_wA$ dominates
$\varrho_\sigma(\nu)$ on every graph and every mixture $\nu_\alpha$ tested (worst margin
positive). On the instance reported in Figure~\ref{fig:app_replay}: the \emph{guarantee}
peaks at $\alpha=0$ ($\nu=\lambda$), as proved --- and the \emph{exact} linearized rate
peaks at $\alpha=0.45$, with conservatism $1588\times$ at $\alpha=0$. The mechanism is
visible in $w=\frac{d\nu_\alpha}{d\lambda}$: uniform tempering makes $w$ large exactly on
rarely-visited states, lifting the slow modes that live there, which a $w_{\min}$-only
bound cannot see. Two scope notes keep the reading honest: the interior optimum is
\emph{shallow} on the shown instance --- the exact rate at $\alpha=0.45$ exceeds the
$\alpha=0$ value by $3.7\%$, so the decisive fact is that tempering costs nothing in
exact rate while buying the certificate, not that it buys much rate --- and on other
graphs of the family the exact-rate optimum can sit at the uniform endpoint
($\alpha=1$) rather than in the interior. Read on the exact rate, the finding is still
\emph{stronger} than Proposition~\ref{prop:alpha_star}, not weaker: tempering backward
replay toward uniform helps the certificate without hurting the exact rate.

\paragraph{Transfer and certificate hold on the exact chain.}
Over $40$ random unbalanced sink-normalized cut flows: the identity
\eqref{equ:transfer_identity} holds to $1.3\cdot10^{-15}$;
$\TV(s_\tau\|\target)\leq C\,\LFMstable$ holds on the exact chain of every instance; the
linearized flow decays at least at $\varrho_\sigma(\nu_\alpha)$ for every tested
$(\alpha,T)$; and the grid argmin of $J_T$ matches \eqref{equ:alpha_star} to
$2.5\cdot10^{-6}$ across four $(q,c,T)$ triples.

\paragraph{Backward replay does sample $\lambda$, and buffer bias is a computable
quantity.}
$20{,}000$ simulated excursions give $\TV(\widehat\nu,\lambda)=0.0013$ (and
$\widehat\nu$ is $\lambda$, not uniform); a reward-squared-biased buffer marginal lands
within $0.0006$ of its fundamental-matrix prediction
(Proposition~\ref{prop:backward_replay}(3)); inverse-frequency reweighting improves
$\sup\frac{d\nuB}{d\nu}$ from $29.6$ to $12.0$ at full coverage.

\begin{figure*}[tb]
\centering
\includegraphics[width=0.9\linewidth]{figures/figA3_replay.pdf}
\caption{\textbf{Left:} certified rate $\varrho_\sigma(\nu_\alpha)$ against the exact
smallest positive Hessian eigenvalue over the mixture family
$\nu_\alpha=(1-\alpha)\lambda+\alpha\,\mathrm{unif}$ (log scale): the guarantee peaks at
$\alpha=0$ as Proposition~\ref{prop:backward_replay} proves, the exact rate peaks at an
interior $\alpha=0.45$, and the vertical gap is the $1588\times$ conservatism the
guarantee admits. \textbf{Right:} the grid argmin of the envelope $J_T$ against the
closed form $\alpha^*(T)=\min(1,1/(qcT))$ of Proposition~\ref{prop:alpha_star}.}
\label{fig:app_replay}
\end{figure*}

% Appendix E -- trained backward policies: flow regularization in spectral variables.
\section{\appendixEtitle}\label{app:pbtrain}

Everything in Section~\ref{sec:cv_divergence} freezes the backward policy; the practice
recommended by \citet{morozov2025revisiting} trains it. Their objective is the log-scale
detailed-balance loss (their Eq.~4) with a learnable $\mathcal P_B$, plus state-flow
regularization (their Eq.~12),
\begin{equation}\label{equ:hse_objective}
    \mathcal L_{\lamreg}(\mathcal F, \mathcal P_B)
    \;=\; \sum_{s\to s'} \widehat\nu(s\to s')
    \Bigl(\log\tfrac{\mathcal F(s)\mathcal P_F(s'|s)}{\mathcal F(s')\mathcal P_B(s|s')}\Bigr)^2
    \;+\; \lamreg\sum_{s}\mathcal F(s)
\end{equation}
(their regularizer is written per transition; summed over a batch it is the total state
flow, which is the form every statement below uses), with reward matching enforced by
the substitution
$\mathcal F(s_f)\mathcal P_B(x|s_f)=\mathcal R(x)$. That substitution does two things at
once in the loop-closure frame of this paper: it pins the total scale ($\mathcal F(s_f) =
\sum_x\mathcal R(x) = Z$ on the balanced set, so the scale ray of the frozen-backward
family --- Theorem~\ref{theo:universality_graphs}(2), the ray whose uniqueness-up-to-scale
Proposition~\ref{prop:frozen_unstable} exploits --- is removed by a boundary condition,
not by a constraint one must add), and it pins the \emph{sink row} of the backward chain to the
normalized target --- exactly the convention
$\widehat{\mathcal P}_B(s_f\to\cdot)=\target/t$ under which
Proposition~\ref{prop:morozov_rate} defines $N$, $\bar\sigma$, $\sigma_*$. Their
Proposition~3.7 then parameterizes the balanced set by the \emph{interior rows of
$\mathcal P_B$ alone}. This appendix translates \eqref{equ:hse_objective} into the
spectral variables of Section~\ref{sec:cv_divergence}: what the $\mathcal P_B$-training
does (a slow drift along the balanced manifold, descending expected trajectory length),
where it goes (the optimal-transport endpoint of \citet{maksimov2026ot}), and what it
costs (every constant of the training-speed theory, at a computable exchange rate).

\subsection{The joint landscape, and the dictionary}\label{sec:joint}

\begin{lemma}[balanced points are critical in all variables]\label{lem:balanced_min}
Let $\mathcal L\geq0$ be $C^1$ in $(\mathcal F,\mathcal P_B)$ and vanish exactly on the
balanced set. Then every balanced point is a global minimum, so
$\nabla_{(\mathcal F,\mathcal P_B)}\mathcal L=0$ there --- including the
$\mathcal P_B$-gradient. If moreover $\mathcal L=\sum\widehat\nu\, g(r)$ with
$g(1)=g'(1)=0$, $g''(1)>0$, the joint Hessian at a balanced point is Gauss--Newton,
$H_{\mathrm{full}} = g''(1)\, J^\dagger M_{\widehat\nu} J$ with $J$ the joint Jacobian of
the residual $r$, and $\ker H_{\mathrm{full}}$ is the tangent space of the balanced
manifold.
\end{lemma}

\begin{proof}
The first claim is the definition of a global minimum of a nonnegative function. For the
second, $\nabla^2[g(r)] = g''(r)\nabla r\otimes\nabla r + g'(r)\nabla^2 r$; at $r\equiv1$
the second term dies with $g'(1)=0$.
\end{proof}

So at balance the only surviving force in \eqref{equ:hse_objective} is the regularizer:
$\mathcal P_B$-training near the balanced manifold is not ``the loss adjusting the
backward policy'' --- the loss is indifferent there --- it is the regularizer
transporting the flow \emph{along} the manifold. What the regularizer descends is
identified exactly by the following dictionary entry.

\begin{lemma}[their Eq.~10 is the occupation identity, summed]\label{lem:eq10_occupation}
On the loop closure with sink row $\target/t$, let $\mathcal F$ be the balanced flow of an
interior-positive $\mathcal P_B$ with $\mathcal F(s_f)=Z$. Then $\mathcal F(s) = Z\,N(s)$
for every state, hence
\begin{equation}\label{equ:eq10}
    \sum_{s\notin\{s_0,s_f\}}\frac{\mathcal F(s)}{\mathcal F(s_f)} \;=\; \bar\sigma,
    \qquad
    \sum_{s}\frac{\mathcal F(s)}{\mathcal F(s_f)} \;=\; 2+\bar\sigma,
    \qquad
    \mathcal F(\Statespace) \;=\; (2+\bar\sigma)\,Z .
\end{equation}
The first display is Eq.~(10) of \citet{morozov2025revisiting} (expected trajectory
length $=$ normalized total interior flow) re-derived in the variables of
Proposition~\ref{prop:morozov_rate}; in particular the state-flow regularizer of
\eqref{equ:hse_objective} equals, on the balanced manifold,
$\lamreg{}Z\,(2+\bar\sigma(\mathcal P_B))$ --- up to an additive constant,
$\lamreg{}Z\,\bar\sigma(\mathcal P_B)$, so \textbf{flow regularization is
gradient descent on the expected backward-trajectory length}.
\end{lemma}

\begin{proof}
The balanced flow of a frozen backward policy is the invariant-measure ray
(Theorem~\ref{theo:universality_graphs}, part (2)): $\mathcal F = c\,\lambda$. The
occupation identity $\lambda = N/(2+\bar\sigma)$
(Proposition~\ref{prop:morozov_rate}(1)) with $\lambda(s_f)=1/(2+\bar\sigma)$ gives
$c = Z(2+\bar\sigma)$ from $\mathcal F(s_f)=Z$, hence $\mathcal F = Z N$. Sum, using
$N(s_0)=N(s_f)=1$.
\end{proof}

The paper's training-speed theorem is stated as a two-phase convergence result, not as
a gradient-domination inequality; the drift analysis below needs the latter, so we
record it as a lemma. It is the coercivity of Proposition~\ref{prop:morozov_rate} read
on the Hessian.

\begin{lemma}[local Polyak--{\L}ojasiewicz inequality]\label{lem:local_pl}
Fix an interior backward policy $p$ with occupation measure $\lambda(p)$ and training
density $w\geq w_{\min}>0$, and consider the frozen-$p$ loss at its quadratic
approximation around the balanced ray, $\mathcal L=\tfrac12\langle h,Hh\rangle$ with
$H=g''(1)A^\dagger M_wA$ on $L^2(\lambda(p))$ (Theorem~\ref{theo:gd_diffusion}). Then,
on the orthogonal complement of $\ker H$,
\begin{equation*}
    \|\nabla\mathcal L\|^2_{L^2(\lambda(p))}\;\geq\;2\,\varrho_\sigma(p)\,\mathcal L .
\end{equation*}
The constant lives in the $L^2(\lambda(p))$ metric, which moves with $p$: along the
drift studied below it degenerates exactly as $\lambda_{\min}(p)\to0$
(Proposition~\ref{prop:constants_path}).
\end{lemma}

\begin{proof}
$\nabla\mathcal L=Hh$ and, $H$ being positive semi-definite,
$\|Hh\|^2\geq\zeta^+_{\min}(H)\,\langle h,Hh\rangle=2\zeta^+_{\min}(H)\,\mathcal L$
on $(\ker H)^\perp$. The coercivity of Proposition~\ref{prop:morozov_rate}(2) bounds the
gap: $\langle h,Hh\rangle=g''(1)\langle Ah,wAh\rangle\geq
g''(1)w_{\min}\|Ah\|^2\geq g''(1)w_{\min}\|h-\Pi h\|^2/\widehat B_\sigma^{\,2}$, so
$\zeta^+_{\min}(H)\geq g''(1)w_{\min}/\widehat B_\sigma^{\,2}=\varrho_\sigma(p)$.
\end{proof}

\begin{proposition}[quasi-static tracking, and the shadow drift]\label{prop:two_timescale}
Consider the gradient flow of $\mathcal L + \lamreg R$, $R := \sum_s
\mathcal F(s)$, on a region where (i) the frozen-$\mathcal P_B$ loss satisfies the local
Polyak--{\L}ojasiewicz inequality of Lemma~\ref{lem:local_pl} with constant
$2\varrho(p)\geq2\varrho_{\min}>0$ at every backward policy $p$ visited, and (ii)
$C_R:=\sup\|\nabla R\|<\infty$. Then
\begin{equation}\label{equ:tracking}
    \mathcal L_t \;\leq\; e^{-\varrho_{\min} t}\,\mathcal L_0
    \;+\; \frac{\lamreg^2\,C_R^2}{2\,\varrho_{\min}} ,
\end{equation}
i.e.\ the trajectory tracks the balanced manifold at distance $O(\lamreg^2)$
in loss. Its slow component follows, to $O(\lamreg^2)$, the shadow ODE
$\dot p = -\lamreg\,P_{\mathrm{tan}}\nabla R = -\lamreg Z\,
P_{\mathrm{tan}}\nabla\bar\sigma(p)$: projected gradient descent on the expected
trajectory length.
\end{proposition}

\begin{proof}
Enlarging the variable space only adds nonnegative terms to the gradient norm, so
pointwise in $p_t$, $\|\nabla_{(\mathcal F,p)}\mathcal L\|^2 \geq
\|\nabla_{\mathcal F}\mathcal L\|^2 \geq 2\varrho(p_t)\mathcal L$ --- the frozen
inequality applied at the instantaneous kernel; no transverse-Hessian decomposition is
needed and the cross-terms of the joint Hessian never enter. Then
$\dot{\mathcal L} = -\langle\nabla\mathcal L, \nabla\mathcal L +
\lamreg\nabla R\rangle \leq -\tfrac12\|\nabla\mathcal L\|^2 +
\tfrac12\lamreg^2\|\nabla R\|^2 \leq -\varrho_{\min}\mathcal L +
\tfrac12\lamreg^2C_R^2$ by Young; Gr\"onwall gives \eqref{equ:tracking}.
The shadow statement is standard singular perturbation for gradient flows on a compact
interior region (the fast variable is slaved at rate $\varrho_{\min}$, the slow one moves
at $O(\lamreg)$); it is used at the level of gradient flow only ---
discrete GD adds the usual $O(\varepsilon)$ terms and is not re-derived here. The
identification $P_{\mathrm{tan}}\nabla R = Z P_{\mathrm{tan}}\nabla\bar\sigma$ is
Lemma~\ref{lem:eq10_occupation} restricted to the manifold.
\end{proof}

\subsection{Where the drift goes: the OT endpoint}\label{sec:drift_target}

\begin{lemma}[normalization reconciliation]\label{lem:normalization}
With reward-pinned terminal flows and total scale $Z$ fixed, minimizing total interior
flow over balanced flows at a \emph{fixed} initial row $\mathcal F(s_0\to u)=L(u)$ is the
linear program of \citet[Eq.~10]{maksimov2026ot}, hence by their Theorem~3.2 a Kantorovich
optimal-transport problem between $L$ and $\mathcal R$ with graph shortest-path costs.
When the initial row is learned (the setting of \eqref{equ:hse_objective}), the relevant
relaxation is $\inf_{L:\,\sum L=Z}\mathrm{OT}(L,\mathcal R)$ --- transport with a free
source marginal, still an LP, in which each reward unit is routed from its nearest
admissible entry.
\end{lemma}

\begin{proof}
A balanced flow with reward-pinned terminal edges and initial row $L$ satisfies exactly
the constraint set of \citet[Eq.~10]{maksimov2026ot}: conservation at interior states,
$\mathcal F(x\to s_f)=\mathcal R(x)$, $\mathcal F(s_0\to u)=L(u)$, non-negative edge
flows; conversely, a feasible point with everywhere-positive throughput is such a flow,
and a general feasible point lies in the \emph{closure} of the family (zero-throughput
states have no defined backward row) --- so the LP's minimum is the infimum over the
closure, the object Proposition~\ref{prop:drift_target}(ii) speaks about. The objective
--- total interior flow --- is the same on both sides, so the optimization problems
coincide and their Theorem~3.2 applies (stated through their reduced problem, Eq.~11,
under their Assumption~3.1). Freeing $L$ at fixed total $\sum L=Z$ replaces the value
by its infimum over source marginals, an LP in $(L,\Pi)$ jointly.
\end{proof}

\begin{proposition}[the drift target]\label{prop:drift_target}
(i) Along the shadow ODE of Proposition~\ref{prop:two_timescale}, $\bar\sigma(p_t)$ is
non-increasing. (ii) The global minimizers of total flow over the closure of the balanced
family are the minimum-flow solutions, supported on the tight (shortest-path) subgraph of
the optimal potentials \citep[Thms~3.2--3.3]{maksimov2026ot}. (iii) Whether the drift
\emph{converges} to a global minimizer is open --- $\bar\sigma(p)$ is a
stochastic-shortest-path policy-optimization landscape, and no claim is made about its
critical points beyond (i).
\end{proposition}

\begin{proof}
(i) is the definition of (projected) gradient descent on $Z\bar\sigma$. (ii) is
Lemma~\ref{lem:normalization} together with the cited theorems; \S\ref{sec:num_pb}
verifies the identification on enumerated instances with an independent transport LP and
a zero-reduced-cost dual certificate.
\end{proof}

\subsection{The spectral price}\label{sec:price}

The endpoint of the drift is a boundary point of the policy polytope: the min-flow
support is generically a strict subgraph, and a backward policy supported on it is
(partly) deterministic. Every constant of Section~\ref{sec:cv_divergence} is a functional
of $\mathcal P_B$; here is what each does along the way. Write $p_t$ for an interior path
with $p_t\to p_1$, $t\to1$.

\begin{proposition}[the constants along the path to minimum flow]\label{prop:constants_path}
Suppose $p_1$'s flow has a support omitting at least one \emph{state}, while every state
still reaches $s_0$ under $p_1$ (no trap). Then, as $t\to1$:
\begin{enumerate}
    \item $N_{\min}(p_t)\to0$ and $\lambda_{\min}(p_t)\to0$: the occupation of every
    off-support state vanishes;
    \item consequently $\widehat B_\sigma = \sigma_*/\lambda_{\min}^{1/2}\to\infty$ and
    $\varrho_\sigma = g''(1)w_{\min}\lambda_{\min}/\sigma_*^2\to0$: \emph{the certified
    rate dies with the off-support occupation};
    \item for any fixed full-support $\nutrain$, $\|w\|_{L^\infty} =
    \|d\nutrain/d\lambda\|_{L^\infty}\to\infty$, so the step cap
    $\varepsilon\leq(4g''(1)\|w\|_{L^\infty})^{-1}$ of
    Theorem~\ref{theo:db_stable_frozen} collapses as well;
    \item $\sigma_*$ and $\bar\sigma$ stay \emph{bounded}, converging to the hitting
    times of the limit policy. The divergence is an occupation phenomenon, not a
    hitting-time one.
\end{enumerate}
\end{proposition}

\begin{proof}
(1) $N$ is the expected-visit vector of the chain killed at $s_0$, started from the sink
row: the solve $(I-Q_p)^\top N = \widehat{\mathcal P}_B(s_f\to\cdot)$ on
$\Statespace\setminus\{s_0\}$. Under the no-trap hypothesis $I-Q_{p_1}$ is invertible
(every state reaches $s_0$), so $N$ is continuous in $p$ on a neighborhood of $p_1$.
Under $p_1$ an excursion starts on the support (the sink row is reward-pinned) and moves
only along the support, so $N(p_1)(x)=0$ for every off-support state; by hypothesis at
least one such state exists, and continuity gives $N_{\min}(p_t)\to0$. Then
$\lambda_{\min}=N_{\min}/(2+\bar\sigma)\to0$ using (4). (2) is substitution into the
definitions of $\widehat B_\sigma$ and $\varrho_\sigma$. (3) At an off-support state
$\nutrain\geq\nu_{\min}>0$ while $\lambda\to0$, so
$\|w\|_{L^\infty}\geq\nu_{\min}/\lambda_{\min}\to\infty$. (4) The hitting times solve
$(I-Q^{s_0}_p)h=\mathbf 1$, again invertible at $p_1$ by the no-trap hypothesis; by
continuity $\sigma_*(p_t)$ and $\bar\sigma(p_t)$ converge to the finite hitting times of
$p_1$.
\end{proof}

\begin{remark}[replay repairs the cap, not the rate]\label{rem:replay_rescue}
The adaptive buffer $\nutrain = \lambda(p_t)$ of Appendix~\ref{app:replay} keeps
$w\equiv1$, so the step cap survives the whole path. The rate does not:
$\varrho_\sigma = g''(1)\lambda_{\min}/\sigma_*^2\to0$ regardless. The speed loss under
flow regularization is intrinsic --- it \emph{is} the occupation mass leaving the
off-geodesic states --- and no training measure can buy it back within the guarantees of
Section~\ref{sec:cv_divergence}. (Do not read this as ``the exact rate dies too'':
$\varrho_\sigma$ is a guarantee, not a prediction. The exact transverse gap is the
subject of the next proposition, and \emph{it} can die for a different reason.)
\end{remark}

\begin{proposition}[gap collapse at a two-component limit]\label{prop:gap_collapse}
Suppose instead $p_t\to p_1$ where the limit chain has $m\geq2$ recurrent classes, and
the invariant measures converge, $\lambda(p_t)\to\lambda_1$ with $\lambda_1>0$
everywhere. Then the exact transverse Hessian gap vanishes:
$\zeta^{\perp}_{\min}\bigl(H(p_t)\bigr)\to0$ for
$H = g''(1)A^\dagger M_wA$, any fixed $w>0$.
\end{proposition}

\begin{proof}
With $m\geq2$ recurrent classes each charged by $\lambda_1$, the per-class invariant
densities $h_j = d\lambda^{(j)}/d\lambda_1$ span an $m$-dimensional kernel of $A_1 = P_1-I$
(density action), and some mean-zero combination $h\neq0$ has
$\langle h, H_1h\rangle=0$ with $\|h\|_{L^2(\lambda_1)}>0$. Since $\lambda_1$ is
everywhere positive, the quadratic forms $H(p_t)$ converge to $H_1$ on a
non-degenerating space, and the smallest transverse eigenvalue converges to $0$.
\end{proof}

\begin{remark}[the positivity hypothesis is necessary, measured]\label{rem:one_leak}
The hypothesis $\lambda_1>0$ everywhere requires the two leak rates --- into and out of
the would-be second class --- to close at comparable speeds, so that the limit measure is
a genuine mixture over both classes. A one-leak degeneration shows the hypothesis is
necessary: closing only the escape route drains the invariant measure onto the trap, the
limit operator \emph{on its own support} has no extra kernel, and the measured gap
stalls ($0.28\to0.014$ and no further on the one-leak path) --- while the two-leak path
takes it to $2.1\cdot10^{-9}$ (Figure~\ref{fig:app_path}, right; both paths are computed
by the released figure script). The distinction is exactly the
difference between ``the chain traps'' (hitting times diverge, $\varrho_\sigma$ dies
through $\sigma_*$) and ``the chain disconnects'' (the exact gap itself dies).
\end{remark}

\begin{remark}[minimum flow is adjacent to collapse]\label{rem:lp_degeneracy}
At a minimum-flow solution the LP leaves the backward rows of \emph{zero-flow states}
completely unconstrained: the min-flow face of the polytope contains kernels whose
unconstrained rows form traps. Flow regularization therefore steers toward a boundary
region where that trapping geometry is not a far-away pathology but a neighboring face.
The released architecture sweep reached it through a collapsed planner; a
flow-regularized trained $\mathcal P_B$ can reach it through the prices above ---
$\lamreg$ is the dial that chooses how close.
\end{remark}

\subsection{The entropic repair}\label{sec:entropic}

\begin{lemma}[trajectory-likelihood regularization]\label{lem:traj_entropy}
For interior-positive $\mathcal P_B$,
\begin{equation*}
    \sum_{s'} N(s')\,H\bigl(\mathcal P_B(\cdot\,|\,s')\bigr)
    \;=\; \mathbb E\bigl[-\log \mathcal P_B(\tau)\bigr],
\end{equation*}
the expected negative log-likelihood of the backward excursion --- so $N$-weighted
conditional entropy is \emph{path} entropy, exactly. (The $\lambda$-weighted variant
differs by the factor $2+\bar\sigma$; and neither is the edge-flow entropy
$\sum_e f_e\log f_e$ of entropic optimal transport, which differs further by the state
marginal term --- the relation to Sinkhorn is precise but not an identity, see
Remark~\ref{rem:entropic_ot}.)
\end{lemma}

\begin{proof}
Both sides are $\sum_{s'}\sum_u (\text{expected visits to } s')\,
\mathcal P_B(u|s')\allowbreak\,(-\log\mathcal P_B(u|s'))$, by Wald / the occupation
decomposition of the excursion.
\end{proof}

\begin{proposition}[soft-SSP keeps the constants alive]\label{prop:soft_ssp}
Regularize the slow objective with the relative entropy of the policy to the uniform
backward proposal, at temperature $\mu$:
\begin{equation*}
    \min_p\ \bar\sigma(p)\;+\;\mu\,\mathbb E_p\Bigl[\textstyle\sum_t
    \log\bigl(d_{x_t}\,p(u_t|x_t)\bigr)\Bigr]
\end{equation*}
(a pure path-entropy bonus --- the object of Lemma~\ref{lem:traj_entropy} --- differs
from this penalty by the trajectory's accumulated log-degrees; the KL-to-uniform form
is the one whose fixed point is stated below). Its optimizer is the soft-Bellman policy
$\pi_\mu(u|x)\propto e^{-V_\mu(u)/\mu}$ over parents, where $V_\mu$
solves $V_\mu(x) = 1-\mu\log\bigl(\tfrac1{d_x}\sum_{u\in\mathrm{pa}(x)}
e^{-V_\mu(u)/\mu}\bigr)$, $V_\mu(s_0)=0$. Then:
\begin{enumerate}
    \item \emph{(fixed point)} on any finite graph where $s_0$ is reachable from every
    state, value iteration from $0$ converges monotonically to a fixed point $V_\mu$
    with $V_\mu\leq(1+\mu\log d_{\max})\,V_{\mathrm{hard}}$;
    \item \emph{(policy floor)} $\pi_\mu(u|x)\geq e^{-\Delta_V/\mu}/d_{\max}$ with
    $\Delta_V$ the largest parent-value spread;
    \item \emph{(the certified floor)} wherever the minimal-occupation state is one soft
    choice off the hard path,
    \begin{equation*}
        \varrho_\sigma(\mu)\;\geq\;
        \frac{g''(1)\,e^{-\Delta_V/\mu}}{d_{\max}\,(2+\bar\sigma(\mu))\,\sigma_*(\mu)^2}
        \;>\;0 ;
    \end{equation*}
    in general $N(x)\geq r_{\min}\,(e^{-\Delta_V/\mu}/d_{\max})^{D}$ for $D$ the
    backward diameter and $r_{\min}$ the smallest positive sink-row mass, giving the
    same shape with a cruder exponent;
    \item \emph{(the price of the repair)} $\bar\sigma(\mu)-\bar\sigma_{\min} \leq
    \mu\,\log d_{\max}\,\bar\sigma(\mu)$.
\end{enumerate}
So the dial is explicit: the excess trajectory length is \emph{linear} in $\mu$ while
$\log\varrho_\sigma(\mu)$ is $\Theta(-1/\mu)$ --- the repair converts the deterministic
limit's exact zero (a void certificate) into a positive floor, at a cost in $\bar\sigma$
one can read off before training.
\end{proposition}

\begin{proof}
(1) The soft-Bellman operator $\mathsf T$ is monotone, and $\mathsf T0=\mathbf 1\geq0$,
so the iterates $V_n$ increase. The bound $V\leq c\,V_{\mathrm{hard}}$ with
$c:=1+\mu\log d_{\max}$ is invariant: using
$-\mu\log\bigl(\tfrac1d\sum e^{-V(u)/\mu}\bigr)\leq\min_uV(u)+\mu\log d$,
\begin{equation*}
    (\mathsf TV)(x)\;\leq\;1+c\min_{\mathrm{pa}(x)}V_{\mathrm{hard}}+\mu\log d_{\max}
    \;\leq\;c\,\bigl(1+\min_{\mathrm{pa}(x)}V_{\mathrm{hard}}\bigr)\;=\;c\,V_{\mathrm{hard}}(x),
\end{equation*}
since $c-1-\mu\log d_{\max}=0$ and $(c-1)\min V_{\mathrm{hard}}\geq0$; $V_{\mathrm{hard}}$
is finite by reachability. Monotone and bounded, $V_n$ converges pointwise; continuity of
$\mathsf T$ makes the limit a fixed point. (2) is immediate from the softmax form.
(4) $V_\mu$ is the optimal value of steps plus a \emph{non-negative} penalty (a
relative entropy per step), so $\mathbb E_{\pi_\mu}[\mathrm{steps}]\leq V_\mu$; and
$V_\mu\leq(1+\mu\log d_{\max})V_{\mathrm{hard}}$ by (1). Taking the sink-row
expectation, $\bar\sigma(\mu)\leq(1+\mu\log d_{\max})\bar\sigma_{\min}$, and
$\bar\sigma(\mu)-\bar\sigma_{\min}\leq\mu\log d_{\max}\,\bar\sigma_{\min}\leq
\mu\log d_{\max}\,\bar\sigma(\mu)$. (3) chains (2) through
$N_{\min}\geq\pi_\mu$-floor at the pivotal choice and
$\lambda_{\min}=N_{\min}/(2+\bar\sigma)$ into Proposition~\ref{prop:morozov_rate}(3);
the general form iterates the floor along a backward path of length at most $D$ from a
sink-row state of mass $\geq r_{\min}$.
\end{proof}

\begin{remark}[relation to entropic OT]\label{rem:entropic_ot}
Path-entropy regularization of a minimum-cost flow is the Schr\"odinger-bridge /
maximum-entropy-RL side of the entropic-OT family, not Sinkhorn itself: Sinkhorn
penalizes the coupling's entropy (edge flows), soft-SSP penalizes the \emph{policy}'s
path measure, and on graphs the two differ by the state-marginal entropy. What matters
here is not the taxonomy but the direction of the exchange: \emph{the certificate
constants of Section~\ref{sec:cv_divergence} are a quantitative reason to prefer the
entropic version of the flow regularizer} --- they are finite on the soft solution and
zero on the hard one, with the exchange rate of Proposition~\ref{prop:soft_ssp}.
\end{remark}

\emph{Summary of the bridge.} In the paper's variables, \eqref{equ:hse_objective}
reads: the loss holds the flow on the balanced manifold at $O(\lamreg^2)$
(Proposition~\ref{prop:two_timescale}); the regularizer walks the manifold by gradient
descent on $\bar\sigma$ (Lemma~\ref{lem:eq10_occupation}); the walk heads toward the
minimum-flow/OT boundary (Proposition~\ref{prop:drift_target}), where the certified
rate, the step cap and (at a disconnecting limit) the exact gap die at computable
speeds (Propositions~\ref{prop:constants_path} and~\ref{prop:gap_collapse}), next to
the min-flow face whose unconstrained backward rows form traps
(Remark~\ref{rem:lp_degeneracy}); and a path-entropy term restores every constant at a
price linear in the temperature (Proposition~\ref{prop:soft_ssp}).

\subsection{Numerical evidence}\label{sec:num_pb}

The script \texttt{audit\_pb\_spectral.py}, released with the code, computes everything
below exactly and exits nonzero on any violation.

\paragraph{The dictionary holds to machine precision.}
Over $12$ random loop-closed marked graphs: the occupation identity to
$4.3\cdot10^{-16}$; $\lambda(s_0)=\lambda(s_f)=1/(2+\bar\sigma)$ to $2.7\cdot10^{-16}$;
the interior sum $\sum\mathcal F/\mathcal F(s_f)=\bar\sigma$ and the full sum
$2+\bar\sigma$ of Lemma~\ref{lem:eq10_occupation} to $1.3\cdot10^{-13}$.

\paragraph{Minimum flow equals optimal transport, with a dual certificate.}
On a two-route instance and a four-terminal diamond: the minimum-interior-flow LP value
equals an independent transport-LP value (cost $=$ interior states on a shortest path)
to $10^{-10}$; complementary slackness against the solver's own dual potentials holds
with zero reduced cost on the support --- the tight-subgraph statement of
\citet[Thm~3.3]{maksimov2026ot} with the \emph{correct} potentials. (A fixed initial row
can force flow down a long route; tightness is relative to the optimal Kantorovich
potentials, not to naive minimum distances --- a naive min-distance version of this
check fails on exactly that instance.)

\paragraph{The spectral price is exact on the two-route instance.}
Along $\mathcal P_B(t)=(1-t)P_{\mathrm{unif}}+tP_{\mathrm{det}}$
(Figure~\ref{fig:app_path}, left): $N_{\min}=(1-t)/2$ \emph{exactly} ($10^{-12}$ at
every $t$); $\widehat B_\sigma$ grows as $(1-t)^{-1/2}$, from $12.65$ to $848.57$ at
$t=0.9999$; $\varrho_\sigma$ dies linearly, $1.25\cdot10^{-2}\to2.78\cdot10^{-6}$;
$\sigma_*$ \emph{decreases}, $4.000\to3.000$; $\|w\|_{L^\infty}$ grows to $1.1\cdot10^4$ for
uniform $\nutrain$ while $w\equiv1$ holds to $10^{-14}$ under adaptive replay
(Remark~\ref{rem:replay_rescue}). On the trap path with both leaks closing together
(Figure~\ref{fig:app_path}, right): the exact transverse gap falls
$2.84\cdot10^{-1}\to2.11\cdot10^{-9}$ and the limit kernel has exactly $2$ weak
invariant components, validating Proposition~\ref{prop:gap_collapse}; the certificate
ordering $\varrho_\sigma\leq\zeta^\perp_{\min}(H)$ holds at every point of both paths.

\paragraph{The soft-SSP floor is exactly tight.}
Fixed-point residual $<10^{-14}$ at every $\mu\in\{0.5,0.2,0.1,0.05,0.02\}$; the policy
floor and the certified $\varrho_\sigma$ floor of Proposition~\ref{prop:soft_ssp}(3)
hold at every $\mu$ --- and the floor is \emph{exactly tight} on this instance
($\varrho_\sigma$ equals it to three digits at every $\mu$, e.g.\ $9.670\cdot10^{-4}$ at
$\mu=0.5$ and $2.07\cdot10^{-45}$ at $\mu=0.02$); $\bar\sigma(\mu)$ decreases to the
hard value $2$ with excess inside the $\mu\log d_{\max}\,\bar\sigma$ bound at every
$\mu$ (Figure~\ref{fig:app_softssp}).

\begin{figure*}[tb]
\centering
\includegraphics[width=0.9\linewidth]{figures/figA5_path.pdf}
\caption{\textbf{Left:} the constants along the interpolation toward the minimum-flow
support (log scales): $\varrho_\sigma$ dies linearly with $N_{\min}=(1-t)/2$,
$\widehat B_\sigma$ diverges as $(1-t)^{-1/2}$, $\sigma_*$ stays bounded
(Proposition~\ref{prop:constants_path}). \textbf{Right:} the exact transverse gap on the
two-leak trap path against the one-leak stall (Proposition~\ref{prop:gap_collapse} and
Remark~\ref{rem:one_leak}).}
\label{fig:app_path}
\end{figure*}

\begin{figure*}[tb]
\centering
\includegraphics[width=0.9\linewidth]{figures/figA6_softssp.pdf}
\caption{The entropic repair's dial (Proposition~\ref{prop:soft_ssp}).
\textbf{Left:} $\varrho_\sigma(\mu)$ against its certified floor
$g''(1)e^{-\Delta_V/\mu}/(d_{\max}(2+\bar\sigma)\sigma_*^2)$ --- exactly tight on this
instance at every temperature. \textbf{Right:} $\bar\sigma(\mu)$'s excess over the hard
value against the linear bound $\mu\log d_{\max}\,\bar\sigma(\mu)$.}
\label{fig:app_softssp}
\end{figure*}

% Appendix F -- stochastic gradients: noise, floors, and the buffer.
\section{\appendixFtitle}\label{app:sgd}

The training theory of Section~\ref{sec:cv_divergence} is exact gradient flow and
full-batch descent; nothing in it analyzes sampled gradients. Yet the one empirical
scaling left open by the trefoil architecture sweep
(\texttt{exp15\_trefoil\_arch\_sweep.py}, released with the code) is a \emph{noise}
scaling: the sub-square-root batch exponent. This appendix sets up the missing term and
gets three things from it: a structural dichotomy (balance losses are noise-free at
their optimum; imitation terms are not), a floor formula with the slow-mode weighting
in exactly one of its two guises, and an honest account of a batch-scaling model that
was registered in advance and then refuted by a longer control run --- with the
protocol that would test it properly (Appendix~\ref{app:program}).

\subsection{The deformation}\label{sec:deformation}

SGD with training measure $\nutrain$, batch $b$ and step $\varepsilon$ on a loss
$\mathcal L$ is $\theta_{k+1} = \theta_k - \varepsilon(\nabla\mathcal L(\theta_k)+\xi_k)$
with $\mathrm{Cov}(\xi_k) = \Sigma_{\nutrain}(\theta_k)/b$, where
$\Sigma_{\nutrain}(\theta) = \mathrm{Cov}_{x\sim\nutrain}(G_x(\theta))$ is the
per-sample gradient covariance. Its diffusion presentation \citep{li2019sme} is
\begin{equation}\label{equ:sme}
    d\theta \;=\; -\nabla\mathcal L(\theta)\,dt
    \;+\; \sqrt{\tfrac{\varepsilon}{b}}\,\Sigma_{\nutrain}(\theta)^{1/2}\,dW .
\end{equation}
The training measure now enters the limiting dynamics \emph{twice}: through the drift
preconditioner $M_w$ of Theorem~\ref{theo:gd_diffusion} ($w = d\nutrain/d\lambda$, the
object all of Appendix~\ref{app:replay} optimizes) and through the diffusion tensor
$\Sigma_{\nutrain}$ --- the deformation the buffer applies to the gradient flow the
theory analyzes. The noise scale is the ratio $\varepsilon/b$ (halving the step buys a
doubled batch), with the caveat that the step cap of
Theorem~\ref{theo:db_stable_frozen} bounds $\varepsilon$ independently of $b$.

\subsection{The dichotomy: balance noise is multiplicative}\label{sec:dichotomy}

\begin{proposition}[balance losses are at interpolation]\label{prop:interpolation}
For any balance loss $\mathcal L_{g,\nu}(\mu)=\int g(r)\,d\nu$ with $g\geq0$,
$g(1)=g'(1)=0$, sampled per state (or per edge, through the lift) --- per-sample loss
$g(r_x)$, $x\sim\nu$:
\begin{enumerate}
    \item at balance every per-sample loss is at its own \emph{global} minimum, so
    $G_x(\theta^*)=0$ for every $x$ --- the exact interpolation regime --- and
    $\Sigma_\nu(\theta^*)=0$ identically, \emph{for every training measure};
    \item near balance the noise is multiplicative with the strong-growth constant
    explicit: $\operatorname{tr}\Sigma_\nu(\theta)\;\leq\;
    2g''(1)\,\max_x\|\nabla r_x\|^2\;\mathcal L(\theta)\,(1+o(1))$;
    \item consequently SGD on a pure balance loss has \emph{no noise floor}: under the
    local PL inequality of Lemma~\ref{lem:local_pl}, the expected loss contracts
    linearly to zero at a rate degraded by at most the strong-growth factor
    \citep[the mechanism of][]{vaswani2019fast} --- batch size rescales the rate, never
    the limit;
    \item in log coordinates ($\mu = e^\theta$ against counting measure) the balance
    Hessian at balance is exactly $g''(1)(P-I)^\top\!\operatorname{diag}(\nu)(P-I)$ with
    $P$ the density action --- the Gauss--Newton face of
    Theorem~\ref{theo:gd_diffusion}'s $H = g''(1)A^\dagger M_wA$.
\end{enumerate}
\end{proposition}

\begin{proof}
(1): $g(r_x)\geq0$ vanishes at $r_x=1$. (2): $\|G_x\|^2 = g'(r_x)^2\|\nabla r_x\|^2$
with $g'(r)^2 = g''(1)^2(r-1)^2(1+o(1))$ and $\mathcal L =
\tfrac{g''(1)}2\int(r-1)^2d\nu\,(1+o(1))$; take $\mathbb E_\nu$ and compare.
(3): one expected-descent line, with $\Lambda$ a local smoothness constant (an upper
bound on $\|\nabla^2\mathcal L\|$ on the neighborhood): $\mathbb E[\mathcal L^+]\leq\mathcal L
-\varepsilon\|\nabla\mathcal L\|^2 + \tfrac{\varepsilon^2\Lambda}2
(\|\nabla\mathcal L\|^2 + \operatorname{tr}\Sigma/b)$, then (2) and PL. (4): expand
$r = 1 + (P-I)h + O(h^2)$ as in Theorem~\ref{theo:gd_diffusion} and apply
Lemma~\ref{lem:balanced_min}'s Gauss--Newton form.
\end{proof}

The training loss of the trefoil architecture sweep is \emph{not} pure
balance: its defect penalty is an exact enumerated sum, and its only sampled term is the
imitation cross-entropy $-\mathbb E_{x\sim\target}\log\fterm(x)$, whose per-sample
gradient at the optimum is nonzero (only its mean vanishes) --- $\Sigma(\theta^*)$ is
the Fisher information of the categorical target, an \emph{additive} noise source. The
dichotomy is the structural content of the batch scaling: the part of the loss this
paper's theory speaks about contributes no floor at all; any floor comes entirely from
the imitation term, and it obeys an Ornstein--Uhlenbeck law.

\subsection{The floor, in two guises}\label{sec:floor}

\begin{proposition}[additive floor]\label{prop:additive_floor}
Linearize \eqref{equ:sme} at an optimum $\theta^*$ with Hessian $H$ and constant noise
$\Sigma = \Sigma(\theta^*)$: $dh = -Hh\,dt + \sqrt{\varepsilon/b}\,\Sigma^{1/2}dW$,
restricted to the complement of $\ker H$, with $\Sigma_\perp$ denoting the restriction
of $\Sigma$ to $(\ker H)^\perp$. The stationary covariance $S$ solves
$HS+SH = \tfrac{\varepsilon}{b}\Sigma_\perp$, and:
\begin{enumerate}
    \item \emph{(loss floor --- no spectral weighting)}\;
    $\mathbb E[\mathcal L-\mathcal L^*] \;\to\; \tfrac12\operatorname{tr}(HS)
    \;=\; \frac{\varepsilon}{4b}\operatorname{tr}\Sigma_\perp$: the stationary excess
    loss does not see the spectrum of $H$ at all;
    \item \emph{(parameter floor --- the slow modes)}\;
    $\operatorname{tr} S = \frac{\varepsilon}{2b}\sum_i
    \frac{\widetilde\Sigma_{ii}}{\zeta_i}$ in $H$'s eigenbasis: parameter (and hence
    TV) error weights each noise mode by the \emph{inverse} eigenvalue --- the
    $\frac1{1-\beta}$ slow-mode weighting of Proposition~\ref{prop:shape_cost}(4),
    reappearing on the noise side. Locally
    $\mathbb E[\TV] \lesssim C_{\mathrm{lip}}\,(\operatorname{tr}S)^{1/2}$ by
    Lipschitz-ness of the exact chain law, not through any $\mathcal L^{1/q}$ route;
    \item \emph{(gauge)} at the mixed balance$+$imitation optimum, $\ker H$ is exactly
    the scale direction, along which the terminal law is invariant --- so the
    restriction in (1)--(2) discards nothing the certificate sees;
    \item \emph{(discrete ground truth)} the discrete stationary equation is
    $S_d = (I-\varepsilon H)S_d(I-\varepsilon H) + \tfrac{\varepsilon^2}b\Sigma_\perp$;
    $S_d = S + O(\varepsilon)$.
\end{enumerate}
\end{proposition}

\begin{proof}
(1): trace the Lyapunov equation: $2\operatorname{tr}(HS) =
\tfrac{\varepsilon}b\operatorname{tr}\Sigma_\perp$; the excess loss is
$\tfrac12\mathbb E[h^\top Hh] = \tfrac12\operatorname{tr}(HS)$. (2): solve the Lyapunov
equation in the eigenbasis, $S_{ij} = \tfrac{\varepsilon}b
\widetilde\Sigma_{ij}/(\zeta_i+\zeta_j)$, and take the trace. (3): both loss terms
are invariant under $\theta\mapsto\theta+c\mathbf 1$, and the exact chain law is
unchanged by a global rescaling of the flow; \S\ref{sec:num_sgd} verifies the kernel is
one-dimensional at the mixed optimum. (4) is direct algebra on the two stationary
equations.
\end{proof}

\subsection{The buffer under noise}\label{sec:buffer_noise}

\begin{proposition}[the replay answer is noise-robust]\label{prop:buffer_robust}
For pure balance losses, the balanced set is $\nutrain$-independent and
$\Sigma_{\nutrain}(\theta^*)=0$ for \emph{every} $\nutrain$: the noise term creates no
new trade-off, and $\nutrain=\lambda$ retains the guarantee-optimality of
Proposition~\ref{prop:backward_replay} unchanged under sampling.
\end{proposition}

\begin{proof}
Immediate from Proposition~\ref{prop:interpolation}(1): the balanced set and the
vanishing of every per-sample gradient there do not depend on the training measure, so
sampling adds no $\nutrain$-dependent term at the optimum, and
Proposition~\ref{prop:backward_replay} applies unchanged.
\end{proof}

For the \emph{mixed} loss the picture genuinely deforms, and the exact curve of
\S\ref{sec:num_sgd} is worth recording: along $\nu_t = (1-t)\lambda + t\,\mathrm{unif}$
on the audit instance, the exact transverse rate \emph{rises}
($5.8\cdot10^{-3}\to4.7\cdot10^{-2}$) and the parameter floor \emph{falls}
($3.1\cdot10^{-7}\to1.6\cdot10^{-7}$) --- both axes prefer tempering toward uniform once
an imitation term is present [computed exactly on one instance; no general theorem
claimed]. This is the third appearance of the same moral: the pure-balance
\emph{guarantees} prefer $\lambda$ (Proposition~\ref{prop:backward_replay}), the exact
rate prefers an interior mixture (\S\ref{sec:num_replay}), the stable pipeline prefers
uniform (Proposition~\ref{prop:stable_uniform}) --- and noise, when it exists at the
optimum, votes with the tempered end. The importance-sampling intuition
``$\nu^*\propto$ root second moment of the per-sample gradient'' is deliberately
\emph{not} used: it belongs to the reweighted-estimator regime (fixed objective),
whereas the buffer here changes the objective itself.

\subsection{The batch model: registered, and refuted by a control}\label{sec:batch_model}

By Propositions~\ref{prop:interpolation} and~\ref{prop:additive_floor}, the sweep's
final TV should decompose as noise floor plus $b$-independent bias,
\begin{equation}\label{equ:batch_model}
    \TV(b) \;\approx\; \sqrt{a\,/\,b + c},
\end{equation}
with $a$ the imitation-noise coefficient
($\propto\varepsilon\cdot$Fisher-along-slow-modes) and $c$ collecting the
finite-$\varepsilon$ floor, model-class and finite-time bias --- inseparable at fixed
learning rate. The sweep's two batch points ($b=512$: TV $0.0345$; $b=4096$: $0.0159$)
determine $(a,c)$ \emph{exactly} --- zero degrees of freedom, so their consistency with
the observed $2.17\times$ improvement against $\sqrt8=2.83$ is a reading, not a
confirmation. The fit gives $a=0.5485$, $c=1.189\cdot10^{-4}$, hence predicted final TVs
$0.0197$ at $b=2048$, $0.0136$ at $8192$, $0.0116$ at $32768$, and an asymptote
$\sqrt c=0.0109$; these were registered before any further run was examined. (The fit
uses the recorded three-significant-figure TVs; full-precision values give $a=0.5471$,
$\sqrt c=0.0107$ --- immaterial, and the registered numbers are kept as registered.)

The refutation arrived through a control rather than the sweep: the long-run control
(\texttt{exp16\_best\_recipe.py}) trains the \emph{same} batch $4096$ for $20{,}000$
steps and reaches TV $0.00574$, still falling --- $1.9\times$ \emph{below} the fitted
asymptote at the same $b$, so no noise-scaling reading survives. The cause is visible in
the fit's inputs: both batch points came from wall-clock-truncated runs still descending
at their last checkpoint (with an $8\%$ step-count mismatch between them), so ``final
TV'' in \eqref{equ:batch_model} was never a stationary quantity, and $c$ absorbed where
training happened to stop. The structural theory of
\S\ref{sec:dichotomy}--\S\ref{sec:floor} is untouched --- it is verified at $0.03\%$
where stationarity actually holds (\S\ref{sec:num_sgd}) --- but a valid test of
\eqref{equ:batch_model} needs arms run to a fixed large step count, not a wall-clock
budget, with truncated points excluded from the fit; experiments (F1) and (F2) of
Appendix~\ref{app:program} specify it, and the sharper first test is the
half-learning-rate arm at fixed $b$ and fixed steps, which isolates the
$\varepsilon$-dependent part of the floor and is immune to the stopping artifact. The
episode is the mass-floor lesson of \S\ref{sec:mass} arriving on the noise side: a floor
fitted from a trajectory that never reached it measures the trajectory, not the floor.

\subsection{Numerical evidence}\label{sec:num_sgd}

The script \texttt{audit\_sgd.py}, released with the code, is fully deterministic by
default (covariances and Lyapunov solves are finite sums and linear solves); an optional
flag adds the seeded stochastic check reported last.

\paragraph{Interpolation, measured.}
Over $5$ random loop-closed graphs: per-sample gradients at balance vanish to
$7.1\cdot10^{-15}$, state by state, by automatic differentiation; the Gauss--Newton
identity of Proposition~\ref{prop:interpolation}(4) holds against the autograd Hessian
to $9.4\cdot10^{-16}$; $\Sigma_\nu$ is PSD; and
$\operatorname{tr}\Sigma/\mathcal L$ is constant over a $4\times$ range of distances to
balance (drift $<10^{-4}$) --- the noise is multiplicative, measured.

\paragraph{The floor algebra is exact.}
At a mixed balance$+$imitation optimum (Newton to $10^{-14}$): $\ker H$ is
one-dimensional and aligned with the scale direction to $10^{-12}$; the exact chain law
is invariant under doubling the flow to $10^{-12}$;
$\operatorname{tr}(HS)=(\varepsilon/2b)\operatorname{tr}\Sigma_\perp$ and the eigenbasis
identity for $\operatorname{tr}S$ both hold at $10^{-11}$; the discrete--continuous
Lyapunov gap halves when $\varepsilon$ halves (ratio $2.000$).

\paragraph{The buffer's two entrances.}
Pure-balance $\Sigma_\nu(\theta^*)$ vanishes at $10^{-13}$ for every $\nu$ in a spiky
random grid; $\varrho_\sigma(\lambda)\geq\varrho_\sigma(\nu)$ on $200$ draws; the
mixed-loss (exact rate, parameter floor) curve of \S\ref{sec:buffer_noise} is printed by
the script, not asserted.

\paragraph{The diffusion presentation is the observed stationary law.}
$64$ tabular SGD chains on the mixed loss ($b=8$, $\varepsilon=0.2/\zeta_{\max}$,
$16{,}000$ steps after burn-in): empirical stationary excess loss $3.130\cdot10^{-5}$
against the predicted $\tfrac12\operatorname{tr}(HS_d)=3.131\cdot10^{-5}$ --- agreement
to $0.03\%$, far inside the statistical tolerance
(Figure~\ref{fig:app_sgdfloor}, right). On these instances \eqref{equ:sme} is not an
analogy; it is the measured stationary law.

\begin{figure*}[tb]
\centering
\includegraphics[width=0.9\linewidth]{figures/figA7_sgdfloor.pdf}
\caption{\textbf{Left:} the batch-model refutation of \S\ref{sec:batch_model}: TV
against steps for the wall-clock-truncated batch-$4096$ sweep arm and for the long-run
control at the same batch, which walks through the fitted asymptote $\sqrt c$ (dashed)
--- the fitted floor was a stopping artifact. \textbf{Right:} where stationarity does
hold, the theory is exact: predicted stationary excess loss
$\tfrac12\operatorname{tr}(HS_d)$ against the empirical mean of $64$ tabular SGD chains
(agreement $0.03\%$; Proposition~\ref{prop:additive_floor}).}
\label{fig:app_sgdfloor}
\end{figure*}

% Appendix G -- the experimental program: what grounds each result, run or specified.
\section{\appendixGtitle}\label{app:program}

Appendices~\ref{app:finit}--\ref{app:sgd} each close with the evidence already in hand:
exact audits (released scripts, every identity computed, exit status as the pass/fail
record) and, where available, trained-system data from the released trefoil
architecture sweep (\texttt{exp15\_trefoil\_arch\_sweep.py}) and its control runs.
This appendix completes the picture in the other direction: for each result, the
experiment that would ground or refute it \emph{best}, whether or not it has been run.
Experiments not yet run are specified here with the same discipline the run ones
followed: enumerable state spaces wherever possible, so every comparison is against an
exactly computed law; \emph{fixed step counts, never wall-clock budgets} (the lesson of
\S\ref{sec:batch_model}: a truncated run measures the trajectory, not the property);
and success \emph{and} refutation criteria stated before the run. Experiments tagged
[exact] need no training and run in CPU minutes on the released audit machinery; those
tagged [trained] train that sweep's network family and carry a rough GPU cost.
Priority order among those not yet run: F2, F1, D1, C1, E1, E2, F3, C2.
Items H1, H2 and H3 have been run; \S\ref{sec:program_doubling} reports their outcomes,
including the two registered criteria they did not meet as written.

\subsection{For Appendix \ref{app:finit} (the initial flow)}

\paragraph{(C1) Trained mass sweep [trained].}
Grounds Propositions~\ref{prop:overshoot} and~\ref{prop:blur_vs_drop} beyond the trivial
point, on a trained system. Protocol: $\Finit=(z/t)\,Z\cdot\mathrm{unif}$ with
$z/t\in\{0.25, 0.5, 0.8, 1, 1.25, 2, 4\}$, fixed $10{,}000$ steps, two seeds, on a
mid-size enumerable environment. Primary quantities: final TV; the truncation signature
(fraction of support states with near-zero terminal mass, and terminal mass in the
bottom decile of target density); the blur signature on the overshoot side. Success: the
asymmetry $\TV(t-s)>\TV(t+s)$ with truncation appearing only under undershoot.
Refutation: a symmetric or reversed penalty --- the overshoot preference would then be a
certificate artifact, and \S\ref{sec:mass} would say so. Cost: $\sim$14 runs,
$\sim$7 GPU-h.

\paragraph{(C2) Warm-start shape [trained, exploratory].}
Grounds Proposition~\ref{prop:shape_cost}(3--4) under training; the open warm-start
question of \S\ref{sec:shape}. Arms: $\Finit\in\{$uniform; a slow-mode-matched coarse
approximation of $\target$ (cluster masses); a fast-detail-matched profile at equal
mismatch amplitude; $\target$ itself as control$\}$, fixed steps. Primary: final outflow
mass, mean trajectory length, TV. Success: the slow-mode warm start cuts outflow and
length at equal accuracy, per the $1/(1-\beta)$ weighting. Refutation: no placement
dependence --- planner adaptation then dominates the initialization, a scope limit worth
knowing. Cost: $\sim$8 runs.

\subsection{For Appendix \ref{app:replay} (the training measure)}

\paragraph{(D1) Replay-schedule sweep on a trained frozen-backward system [trained or
exact].}
Grounds Propositions~\ref{prop:backward_replay} and~\ref{prop:alpha_star}, and the
interior-exact-rate finding of \S\ref{sec:num_replay}, under real training. Protocol: DB
loss, frozen backward policy, enumerable environment ($N\sim10^3$--$10^4$); arms
$\nutrain\in\{$on-policy; backward replay ($\alpha=0$); fixed mixtures
$\alpha\in\{0.2, 0.45, 0.7\}$; uniform ($\alpha=1$); the schedule
$\alpha^*(T)=\min(1,1/(qcT))\}$; fixed $20{,}000$ steps; two seeds; TV computed exactly
per checkpoint. Success: backward replay dominates on-policy in the late phase and
trails it before coverage (Lemma~\ref{lem:coverage}); some interior mixture beats both
endpoints; the schedule matches the best fixed $\alpha$. Refutation: an endpoint
dominating throughout --- the mixture story would then be a guarantee-level artifact.
Cost: $\sim$14 runs (tabular version: CPU hours).

\subsection{For Appendix \ref{app:pbtrain} (trained backward policies)}

\paragraph{(E1) Flow-regularization drift, exact [exact].}
Grounds Proposition~\ref{prop:two_timescale} and Proposition~\ref{prop:drift_target}(i)
under live dynamics --- the only wave with no dynamical evidence yet. Protocol: tabular
gradient descent on $(\mathcal F,\mathcal P_B)$ jointly for
$\mathcal L+\lamreg{}R$ on the two-route instance and one random
$20$--$40$-state loop-closed graph, $\lamreg\in\{10^{-1},3\cdot10^{-2},
10^{-2},3\cdot10^{-3}\}$. Primary: the stationary loss plateau against
$\lamreg$ (predicted log-log slope $2$, from \eqref{equ:tracking});
monotone $\bar\sigma$ descent; the endpoint $\bar\sigma$ against the minimum-flow LP
value. Success: slope $2.0\pm0.2$; $\bar\sigma$ within $5\%$ of the LP value or a
documented interior stall (consistent with Proposition~\ref{prop:drift_target}(iii)).
Refutation: a robustly non-quadratic plateau scaling. Cost: CPU minutes to hours.

\paragraph{(E2) Entropic repair, live [exact].}
Grounds Proposition~\ref{prop:soft_ssp} under training rather than at the fixed point.
Protocol: the (E1) setup plus $\mu\cdot$path-entropy,
$\mu\in\{0.5,0.2,0.1,0.05\}$; measure the trained endpoints' $\varrho_\sigma(\mu)$
against the certified floor, and the $\bar\sigma$ excess against
$\mu\log d_{\max}\,\bar\sigma$. Success: floors hold live, excess linear in $\mu$.
Refutation: an endpoint below its floor --- diagnose non-convergence and report. Cost:
CPU minutes.

\subsection{For Appendix \ref{app:sgd} (stochastic gradients)}

\paragraph{(F1) Fixed-step batch sweep [trained].}
Replaces the refuted fit of \S\ref{sec:batch_model}; grounds
Propositions~\ref{prop:interpolation} and~\ref{prop:additive_floor} jointly. Protocol:
$b\in\{512, 2048, 4096, 8192, 32768\}$, same learning rate, \emph{fixed} $20{,}000$
steps, two seeds, with a stationarity gate (mean $|\Delta\TV|$ over the last $2{,}000$
steps below threshold; extend the run rather than fit through a failing gate); fit
\eqref{equ:batch_model} with all truncated points excluded. Success: monotone $\TV(b)$,
small fit residuals, and the largest-$b$ arm consistent with the fitted $c$. Refutation:
$\TV(32768)$ materially below the small-$b$ fit's floor, or non-monotonicity outside
seed spread. Cost: $\sim$10 runs, $\sim$15--25 GPU-h (the largest batches dominate).

\paragraph{(F2) Half-learning-rate discriminator [trained] --- the first test to run.}
Grounds Proposition~\ref{prop:additive_floor}'s $\varepsilon$-dependence and is immune
to the stopping artifact. Protocol: the control configuration at fixed $b=4096$,
learning rate $\in\{\varepsilon_0,\varepsilon_0/2\}$, fixed $20{,}000$ steps (extend to
$40{,}000$ if the stationarity gate fails), seeds $\{0,1\}$. Primary: stationary final
TV. Success: floors scaling $\approx\sqrt2$ in the learning rate (noise-bound, per the
additive floor $\propto\varepsilon$). Refutation: a ratio $\approx1$ --- the floor is
then bias-bound at these settings and every floor claim gets scoped accordingly; or
persistent non-stationarity, reported as such. Cost: 4--6 runs, $\sim$6 GPU-h.

\paragraph{(F3) Two-guise floor discriminator [exact].}
Separates Proposition~\ref{prop:additive_floor}(1) from (2) --- the unweighted loss
floor against the slow-mode-weighted parameter floor. Protocol: construct two instances
with \emph{equal} $\operatorname{tr}\Sigma_\perp$ but noise aligned with slow versus
fast Hessian modes; run $64$ tabular SGD chains on each. Prediction: \emph{equal} loss
floors (the trace law sees no spectrum), \emph{unequal} parameter/TV floors by the
computed eigenbasis ratio. Success: loss floors within $\sim5\%$, TV ratio within
$\sim15\%$ of the computed value. Refutation: materially different loss floors ---
spectral weighting where the theory says there is none. Cost: CPU minutes.

\subsection{For Appendix \ref{app:doubling} (the doubling graph)}\label{sec:program_doubling}

The three items below were registered in the docstring of
\texttt{exp20\_doubling\_diffusion.py} before the run, and have been run. Each reports what was
found rather than what would be found, and two of the registered criteria are not met as
written; both are recorded here as such.

\paragraph{(H1) The closed forms and the two exponents [exact], run.}
Grounds Propositions~\ref{prop:doubling_drift}, \ref{prop:doubling_length} and
\ref{prop:doubling_exponent} on the family $\varepsilon_{c,s}(j)=c/(j+1)^s$ of
\eqref{eq:doubling_family}. Protocol: expected backward lengths by linear solve, and the decay
exponent of the truncated invariant measure fitted on the dyadic band $[K/16,K/2]$, at eight
values of $c$ from $0.25$ to $2.5$ and truncations to $K=4096$. Success: $\bar\sigma$ matching
the closed form $\bar\jmath/(1-c)$ --- $\bar\jmath$ the mean of the target row --- to a
relative $10^{-6}$ at the largest converged cap; the fitted decay exponent matching the
Cram\'er root $p_*$ of \eqref{eq:doubling_cramer} to $5\%$; and the earlier fitted truncation
exponent of $\bar\sigma$ matching $1-p_*$ to $0.05$ absolute. Refutation: any of the three
outside tolerance --- the closed form is an identity, so a single mismatch kills it.
\emph{Outcome:} confirmed. The closed form holds to a relative $5.8\cdot10^{-10}$ at $c=0.25$
and $9.7\cdot10^{-6}$ at $c=0.5$; at $c=0.75$ and $0.9$ the cap is not converged, and the
residual is identified as truncation by its own decay, $K^{-(p_*-1)}$ to three digits. The
fitted exponents match $p_*$ to three decimals at all eight values of $c$
(\S\ref{sec:doubling_measured}), including the two where $p_*$ is negative and the earlier fit
had clamped it.

\paragraph{(H2) The two bounds on a graph of infinite mixing [exact], run.}
Grounds Theorems~\ref{theo:RL_CV_bound_full} and \ref{theo:IL_CV_bound} where the mixing
hypothesis of the second fails outright. Protocol: $9{,}216$ rows --- eight values of $c$, four
truncations ($K=64$ to $512$), four targets crossed with three initial flows, eight
perturbation scales, three exponents $q$ --- with the sampling error computed exactly at every
row. Success: zero violations of the stable bound; its constant flat in $K$ at fixed
$\nutrain$; the certificate level $\eta$ growing with $K$ on at least the ladder-loaded pairs;
the certificate closing flow matching to $10^{-8}$ and its loss equal to the self-entropy.
Refutation: a violation, which would indict the theorem or the harness and either matters; or
$\eta$ flat in $K$, which would mean the unbounded diffusion operator costs nothing.
\emph{Outcome:} confirmed but for one criterion, which was mis-stated. Zero violations in
$9{,}216$ rows, non-vacuous in $2{,}115$; the certificate closes to $6.0\cdot10^{-8}$ at worst,
a relative $5\cdot10^{-17}$, with its loss equal to the self-entropy; $\eta$ is flat at $11.35$
for the target row at the foot of the ladder and grows like $K^{2.9}$ for the ladder-loaded
targets. The registered ``constant flat in $K$'' criterion cannot be met: with counting
background measure and uniform training measure \eqref{equ:stable_constants} gives $C=3N$ by
construction, measured at $198$, $390$, $774$, $1{,}542$ over that range. The quantity that
carries the meaning is the measured tightness, which is what Section~\ref{sec:experiments}
reports and what is flat here --- $6.33$, $6.00$, $6.59$, $6.47$ at $q=1$.

\paragraph{(H3) The critical direction and the price of unboundedness [exact], run.}
Grounds Theorem~\ref{theo:doubling_unbounded} and Corollary~\ref{cor:doubling_truncation}, and the
reading of the extremal direction as a tail indicator. Protocol: $40$ cells --- eight values of
$c$ crossed with five truncations, $K=64$ to $1024$ --- computing $\widehat B_K$ as a singular
value, the best cut of \eqref{eq:doubling_percut_K}, the extremal direction of $S$, the
certificate level along it, and the exact terminal law of the certificate flow. Success:
$\widehat B_K\sim K^\alpha$ with $\alpha\geq0.9$ at every $c$; the lower bound holding at every
admissible cut; an overlap above $0.8$ between the measured extremal direction and the best
centred tail indicator; the certificate's total variation at machine precision at every cap.
Refutation: $\widehat B_K$ saturating in $K$ for some member of the polynomial family, which
would refute the unboundedness theorem; an overlap near zero, which would refute the mechanism
reading but not the theorem; a total variation growing with $K$, which would contradict
Theorem~\ref{theo:universality_graphs}. \emph{Outcome:} confirmed but for the overlap
threshold. The fitted exponent runs from $0.983$ to $1.139$ over the eight values of $c$; the
lower bound holds in all $40$ cells and its own fitted exponent is $0.473$ to $0.510$, so the
proved growth is $\sqrt K$ against a measured $K$; and the certificate's total variation is
$1.2\cdot10^{-8}$ at worst. The overlap is $0.71$ to $0.84$, above the registered $0.8$ at four
of the eight values of $c$ and below it at the rest (\S\ref{sec:doubling_cost}). The refutation
criterion was an overlap near zero, so the mechanism reading stands; the threshold does not
hold uniformly in $c$, and the appendix says so where it uses the reading.

% Appendix H -- the doubling graph: an unbounded diffusion operator at finite backward length.
\section{\appendixHtitle}\label{app:doubling}

This appendix studies in depth the counter-example graph
\eqref{eq:counterexample_graph} and the limit cases of the universality and convergence
theorems for families of backward policies, see \eqref{eq:doubling_family} below, and prices
what it finds. In particular, some positive recurrent members of that family have an unbounded
diffusion operator while the expected backward-trajectory length of
\citet{morozov2025revisiting} is finite. Theorem~\ref{theo:doubling_main} collects our
theoretical results, \S\ref{sec:doubling_family} proves them, and
\S\ref{sec:doubling_measured} and \S\ref{sec:doubling_cost} measure them on truncations.

\begin{definition}[standing notation and hypotheses of this appendix]\label{def:doubling_setting}
The graph is \eqref{eq:counterexample_graph}, carrying the source $s_0$, the ladder
$1,2,3,\dots$, the sink $s_f$ and the forward edges $s_0\rightarrow1$, $j\rightarrow j+1$,
$2j\rightarrow j$ and $j\rightarrow s_f$; the predecessors of a ladder state $j$ are $j-1$ and
$2j$, so the backward chain at $j$ either \emph{decrements} to $j-1$ or \emph{doubles} to $2j$,
and the loop closure of Definition~\ref{def:loop_closure} adds the wrap $s_f\rightarrow s_0$.
The backward policy is
\begin{equation}\label{eq:doubling_policy}
    \BackwardPolicy(j\rightarrow2j)=\varepsilon(j),\qquad
    \BackwardPolicy(j\rightarrow j-1)=1-\varepsilon(j)\quad(j\geq1),\qquad
    \BackwardPolicy(s_0\rightarrow s_f)=1,
\end{equation}
with $j-1=0$ read as $s_0$. The \emph{target row} $\widehat\BackwardPolicy(s_f\rightarrow\cdot)$
is a probability supported in $\{1,\dots,d\}$ of mean $\bar\jmath$, the integer $d\geq1$ and the
real $\bar\jmath$ being parameters of the setting. The family studied throughout is
\begin{equation}\label{eq:doubling_family}
    \varepsilon_{c,s}(j):=\frac{c}{(j+1)^s},
\end{equation}
and its \emph{standing range} is $s\geq0$, $0<c<2^s$: there $\varepsilon_{c,s}$ takes its values
in $(0,1)$ and the decrement edge out of state $1$ carries positive probability
(Lemma~\ref{lem:doubling_range}), and the loop closure is irreducible
(Lemma~\ref{lem:doubling_irreducible}). Write $\varepsilon_{\max}:=\sup_{j\geq1}\varepsilon(j)$,
which on the standing range is $\varepsilon(1)=c\,2^{-s}<1$
(Lemma~\ref{lem:doubling_range}). The statements below that name
$(c,s)$ are confined to that range unless they widen it explicitly, as
Proposition~\ref{prop:doubling_exponent} does; those that hypothesize $\varepsilon$ directly
are not confined to it.

Write $(X_n)_{n\geq0}$ for the backward chain of the loop closure, $\mathbb P(\,\cdot\mid X_0=x)$
and $\mathbb E(\,\cdot\mid X_0=x)$ for the law of the chain in play started at the state $x$ and
its expectation --- that chain being $(X_n)$ wherever no other is named --- $\sigma$ for the hitting time of $s_0$,
and $\sigma_*$ and $\bar\sigma$ for the two functionals \eqref{eq:morozov_variables},
$\bar\sigma$ being the expected backward-trajectory length of \citet{morozov2025revisiting}.
When the loop-closed chain is positive recurrent, $\lambda$ denotes its invariant probability,
$\lambda_j:=\lambda(\{j\})$ the mass of a ladder state,
$L(m):=\lambda([m,+\infty))=\sum_{j\geq m}\lambda_j$ the tail, and
$f_m:=\mathbf 1_{[m,+\infty)}-L(m)$, for $m>d$, the \emph{centred tail indicator} at the cut $m$.

The background measure is $\nuB=\lambda$ throughout, and the conventions of
\S\ref{app:notation} apply: $T$ is the backward kernel of the loop closure, $P$ its density
action on $L^2(\lambda)$, $\Pi$ the $\lambda$-mean projection
$f\mapsto\bigl(\int f\,d\lambda\bigr)\mathbf 1$ and
$\widehat\beta_n:=\|P^n-\Pi\|_{L^2(\lambda)}$. The forward policy of the balanced flow is the
$\lambda$-reversal of $T$ (Lemma~\ref{lem:adjoint}\emph{(4)}), and its density action $P_\star$
--- the operator written $P^\dagger$ in \S\ref{app:divergence} --- is the \emph{function} action
of $T$ on $L^2(\lambda)$ (Lemma~\ref{lem:adjoint}\emph{(3)}), so that
$(P_\star f)(x)=\mathbb E(f(X_1)\mid X_0=x)$ for $f\in L^2(\lambda)$, and, by
Lemma~\ref{lem:doubling_operator}\emph{(1)}, a contraction of $L^p(\lambda)$ for every
$p\in[1,+\infty]$, on which it acts by the same formula. The orientation is
$\theta:=\fterm-\finit$, and
\begin{equation}\label{eq:doubling_Bhat}
    S:=\left(\mathrm{Id}-P_\star+\Pi\right)^{-1}-\Pi,\qquad
    \widehat B:=\|S\|_{L^2(\lambda)} .
\end{equation}
On the infinite chain the convention is $\widehat B:=+\infty$ when no bounded operator $S$ on
$L^2(\lambda)$ satisfies $S(\mathrm{Id}-P_\star)=\mathrm{Id}-\Pi$.

The \emph{Cram\'er equation} of the family is
\begin{equation}\label{eq:doubling_cramer}
    c\left(2^p-1\right)\;=\;p ,
\end{equation}
and $p_*=p_*(c)$ denotes its non-zero real root, which exists and is unique for every $c>0$
with $c\neq1/\ln2$, is $>1$ exactly when $c<1$ and is $>0$ exactly when $c<1/\ln2$
(Lemma~\ref{lem:doubling_cramer_root}).

For an even integer $K\geq d$, the \emph{truncation at $K$} is the chain obtained from the loop
closure by deleting every doubling edge $j\rightarrow2j$ with $2j>K$, the decrement edge out of
such a state carrying probability $1$; it is a finite irreducible chain on the $K+2$ states
$\{s_0,1,\dots,K,s_f\}$ (Lemma~\ref{lem:doubling_truncation_irreducible}), $\lambda^K$ denotes its
unique invariant probability, $L_K(m):=\lambda^K([m,K])$ its tail and $\widehat B_K$ its
constant \eqref{eq:doubling_Bhat}.
\end{definition}

\begin{lemma}[the standing range]\label{lem:doubling_range}
In the setting of Definition~\ref{def:doubling_setting}, let $s\geq0$ and $0<c<2^s$. Then
$j\mapsto\varepsilon_{c,s}(j)$ is non-increasing on $\{1,2,3,\dots\}$ with values in $(0,1)$,
its supremum is $\varepsilon_{c,s}(1)=c\,2^{-s}<1$, and the decrement edge out of the state $1$
carries the probability $1-c\,2^{-s}>0$.
\end{lemma}

\begin{proof}
For $j\geq1$, $\varepsilon_{c,s}(j)=c\,(j+1)^{-s}$ is positive, $c$ being positive, and
$j\mapsto(j+1)^{-s}$ is non-increasing, $s$ being non-negative: the supremum over $j\geq1$ is
attained at $j=1$ and equals $c\,2^{-s}$, which is $<1$ by $c<2^s$. Hence
$\varepsilon_{c,s}(j)\in(0,1)$ at every $j\geq1$, and by \eqref{eq:doubling_policy} the
decrement edge out of the state $1$ carries $1-\varepsilon_{c,s}(1)=1-c\,2^{-s}>0$.
\end{proof}

\begin{lemma}[positive edges make the loop closure irreducible]\label{lem:doubling_irreducible}
In the setting of Definition~\ref{def:doubling_setting}, let $\varepsilon(j)\in(0,1)$ for every
$j\geq1$. Then the loop-closed backward chain $X$ of \eqref{eq:counterexample_graph} under the
policy \eqref{eq:doubling_policy} is irreducible, and every invariant probability of $X$ is
positive at every state.
\end{lemma}

\begin{proof}
By \eqref{eq:doubling_policy} the doubling edge out of a ladder state $j$ carries the
probability $\varepsilon(j)>0$ and the decrement edge out of it the probability
$1-\varepsilon(j)>0$; the wrap $s_0\rightarrow s_f$ carries probability $1$; and the target
row, a probability supported in $\{1,\dots,d\}$, is positive at some $j_0\in\{1,\dots,d\}$.
\begin{itemize}
    \item \emph{Every state leads to the state $1$.} Indeed, from a ladder state $j$ the $j$
    consecutive decrements reach $s_0$, from $s_0$ the wrap reaches $s_f$, from $s_f$ the target
    row reaches $j_0$, and from $j_0$ the $j_0-1$ consecutive decrements reach $1$.
    \item \emph{The state $1$ leads to every state.} Indeed, every ladder state is reached from
    $1$ by induction: from a ladder state $n\geq1$ the doubling edge $n\rightarrow2n$ followed
    by $n-1$ decrements reaches $n+1$, and $2n\geq n+1$. From the state $1$ the decrement
    reaches $s_0$, and the wrap then reaches $s_f$.
\end{itemize}
Concatenating a path from an arbitrary state to the state $1$ with a path from $1$ to an
arbitrary state, every state leads to every other, so $X$ is irreducible. Let $\lambda$ be an
invariant probability and $y$ a state: $\lambda$ is positive at some state $x$, being a
probability, and irreducibility gives an integer $n\geq0$ with
$\mathbb P(X_n=y\mid X_0=x)>0$, so $\lambda T^n=\lambda$ gives
$\lambda(\{y\})\geq\lambda(\{x\})\,\mathbb P(X_n=y\mid X_0=x)>0$.
\end{proof}

\begin{lemma}[the fixed points of the function action are the constants]\label{lem:doubling_fixed_points}
Let $T$ be an irreducible Markov kernel on a countable state space carrying an invariant
probability $\lambda$ positive at every state, let $P$ be its density action
$Pu:=\frac{d\left((u\lambda)T\right)}{d\lambda}$ and $P_\star$ its function action
$(P_\star v)(x):=\int v\,dT(x,\cdot)$. Then
\begin{equation}\label{eq:doubling_fixed}
    \ker\left(\mathrm{Id}-P_\star\right)\cap L^2(\lambda)
    \;=\;\ker\left(\mathrm{Id}-P\right)\cap L^2(\lambda)
    \;=\;\mathbb R\mathbf 1 .
\end{equation}
\end{lemma}

\begin{proof}
The constants are fixed by $P_\star$, each $T(x,\cdot)$ being a probability, and by $P$,
$\lambda$ being $T$-invariant.
\begin{itemize}
    \item \emph{A fixed point of $P_\star$ in $L^2(\lambda)$ is constant.} Indeed, let
    $f\in L^2(\lambda)$ with $P_\star f=f$. Jensen's inequality for the probability $T(x,\cdot)$
    and the strictly convex $t\mapsto t^2$ gives $(P_\star f)^2\leq P_\star(f^2)$ at every
    state, and $\lambda T=\lambda$ gives
    $\int P_\star(f^2)\,d\lambda=\int f^2\,d\lambda=\int(P_\star f)^2\,d\lambda$; $\lambda$
    being positive at every state, $P_\star(f^2)=(P_\star f)^2$ at every state. Equality in
    Jensen for a strictly convex function makes $f$ constant on $T(x,\cdot)$, of value
    $(P_\star f)(x)=f(x)$: at every state $x$, $f(y)=f(x)$ whenever $T(x,y)>0$, and $f$ is
    constant by irreducibility.
    \item \emph{A fixed point of $P$ in $L^2(\lambda)$ is constant.} Indeed, $P_\star=P^*$ and
    both are contractions of $L^2(\lambda)$ (Lemma~\ref{lem:doubling_operator}\emph{(1)}), and
    a contraction $Q$ of a Hilbert space has
    $\ker(\mathrm{Id}-Q^*)\subseteq\ker(\mathrm{Id}-Q)$ --- if $Q^*g=g$ then
    $\langle Qg,g\rangle=\|g\|^2$ and $\|Qg-g\|^2=\|Qg\|^2-\|g\|^2\leq0$ --- so $Pg=g$ forces
    $P_\star g=g$, and a fixed point of $P_\star$ in $L^2(\lambda)$ is constant.
\end{itemize}
Both kernels are therefore contained in $\mathbb R\mathbf 1$, which is
\eqref{eq:doubling_fixed}.
\end{proof}

\begin{lemma}[the diffusion operator]\label{lem:doubling_operator}
Let $T$ be a Markov kernel on a measurable space carrying an invariant probability $\lambda$,
let $P$ be its density action $Pu:=\frac{d\left((u\lambda)T\right)}{d\lambda}$ and $P_\star$
its function action $(P_\star v)(x):=\int v\,dT(x,\cdot)$, let $\Pi$ be the $\lambda$-mean
projection, let $\widehat\beta_n:=\|P^n-\Pi\|_{L^2(\lambda)}$ and, whenever
$\mathrm{Id}-P_\star+\Pi$ is invertible on $L^2(\lambda)$, put
$S:=\left(\mathrm{Id}-P_\star+\Pi\right)^{-1}-\Pi$ and $\widehat B:=\|S\|_{L^2(\lambda)}$.
Then:
\begin{enumerate}
    \item $P_\star$ is a contraction of $L^p(\lambda)$ for every $p\in[1,+\infty]$ and is the
    $L^2(\lambda)$-adjoint of $P$, so that $\|P_\star^n-\Pi\|_{L^2(\lambda)}=\widehat\beta_n$
    for every $n\geq0$;
    \item if moreover the chain is finite and irreducible then $P_\star$ fixes only the
    constant functions, the operator $\mathrm{Id}-P_\star+\Pi$ is invertible, and
    \begin{equation}\label{eq:doubling_resolvent}
        S(\mathrm{Id}-P_\star)=(\mathrm{Id}-P_\star)S=\mathrm{Id}-\Pi,
        \qquad \Pi S=S\Pi=0 ;
    \end{equation}
    \item if $U:=\sum_{n\geq0}(P_\star^n-\Pi)$ converges in operator norm on $L^2(\lambda)$,
    then $\mathrm{Id}-P_\star+\Pi$ is invertible with inverse $\Pi+U$, the operator $S$ equals
    $U$ and satisfies \eqref{eq:doubling_resolvent}, and
    $\widehat B\leq\sum_{n\geq0}\widehat\beta_n$; in particular, on a finite irreducible chain
    $\widehat B\leq\sum_{n\geq0}\widehat\beta_n$, the inequality being vacuous when its right
    side is $+\infty$.
\end{enumerate}
\end{lemma}

\begin{proof}
Both actions are well defined on $\lambda$-classes: a $\lambda$-null set $Z$ satisfies
$\int T(x,Z)\,\lambda(dx)=(\lambda T)(Z)=\lambda(Z)=0$, so $T(x,Z)=0$ for $\lambda$-almost
every $x$, and $\left((u\lambda)T\right)(Z)=\int T(x,Z)\,u(x)\,\lambda(dx)=0$ for
$u\in L^1(\lambda)$.

\emph{(1)} Let $v\in L^p(\lambda)$ with $p<+\infty$. Jensen's inequality for the probability
$T(x,\cdot)$ and the convex $t\mapsto|t|^p$ gives
$|(P_\star v)(x)|^p\leq\bigl(P_\star\left(|v|^p\right)\bigr)(x)$ for $\lambda$-almost every
$x$, and $\lambda T=\lambda$ gives
$\int P_\star\left(|v|^p\right)d\lambda=\int|v|^p\,d(\lambda T)=\|v\|^p_{L^p(\lambda)}$: hence
$\|P_\star v\|_{L^p(\lambda)}\leq\|v\|_{L^p(\lambda)}$. For $p=+\infty$ the set
$Z:=\{|v|>\|v\|_{L^\infty(\lambda)}\}$ is $\lambda$-null, so $T(x,Z)=0$ and
$|(P_\star v)(x)|\leq\|v\|_{L^\infty(\lambda)}$ for $\lambda$-almost every $x$. In particular
$P_\star$ is bounded on $L^2(\lambda)$, so for $u,v\in L^2(\lambda)$ the function
$(x,y)\mapsto u(x)v(y)$ is $T(x,dy)\lambda(dx)$-integrable ---
$\int|u|\,P_\star|v|\,d\lambda\leq\|u\|_{L^2(\lambda)}\|v\|_{L^2(\lambda)}$ by Cauchy--Schwarz
--- and Fubini gives
\begin{equation}
\begin{aligned}
    \int v\,d\left((u\lambda)T\right)\;&=\;\iint v(y)\,T(x,dy)\,u(x)\,\lambda(dx)\\
    &=\;\langle u,P_\star v\rangle_\lambda .
\end{aligned}
\end{equation}
The same bound makes $v\mapsto\int v\,d\left((u\lambda)T\right)$ a bounded linear functional
on $L^2(\lambda)$ of norm at most $\|u\|_{L^2(\lambda)}$, so it is
$\langle w,\cdot\rangle_\lambda$ for a unique $w\in L^2(\lambda)$ by the Riesz representation
theorem \citep{brezis2011functional}, and taking $v=\mathbf 1_A$ at a measurable $A$ identifies
$w$ with $Pu$: the density action is bounded on $L^2(\lambda)$, and the display reads
$\langle Pu,v\rangle_\lambda=\langle u,P_\star v\rangle_\lambda$, that is $P_\star=P^*$. Since
$\lambda$ is a probability, $\Pi$ is self-adjoint, so $P_\star^n-\Pi=(P^n-\Pi)^*$ and the two
operators have the same norm.

\emph{(2)} An irreducible Markov kernel on a finite state space carries a unique invariant
probability, positive at every state \citep{levin2017markov}, so $\lambda$ is that probability
and $P_\star$ fixes only the constant functions
(Lemma~\ref{lem:doubling_fixed_points}). The space $L^2(\lambda)$ is finite-dimensional and
$\mathrm{Id}-P_\star+\Pi$ is injective on it: $\lambda T=\lambda$ gives $\Pi P_\star=\Pi$ and
$\Pi^2=\Pi$, $\lambda$ being a probability, so $(\mathrm{Id}-P_\star+\Pi)f=0$ forces
$\Pi f=0$ and then $P_\star f=f$, whence $f$ is a constant of zero mean, that is $0$. That
operator is therefore invertible, and $R:=(\mathrm{Id}-P_\star+\Pi)^{-1}$ and $S=R-\Pi$ are
defined. Since $P_\star\mathbf 1=\mathbf 1$ and $\lambda$ is $T$-invariant,
$\Pi P_\star=P_\star\Pi=\Pi$, so $\Pi(\mathrm{Id}-P_\star+\Pi)=(\mathrm{Id}-P_\star+\Pi)\Pi=\Pi$
by $\Pi^2=\Pi$; multiplying the first identity on the right by $R$ and the second on the left by
$R$ gives $\Pi R=R\Pi=\Pi$, whence $\Pi S=\Pi R-\Pi^2=0$ and $S\Pi=R\Pi-\Pi^2=0$. Finally
$\Pi(\mathrm{Id}-P_\star)=(\mathrm{Id}-P_\star)\Pi=0$ and $R(\mathrm{Id}-P_\star+\Pi)=(\mathrm{Id}-P_\star+\Pi)R=\mathrm{Id}$
give $S(\mathrm{Id}-P_\star)=R(\mathrm{Id}-P_\star)=\mathrm{Id}-R\Pi=\mathrm{Id}-\Pi$ and
$(\mathrm{Id}-P_\star)S=(\mathrm{Id}-P_\star)R=\mathrm{Id}-\Pi R=\mathrm{Id}-\Pi$, which is
\eqref{eq:doubling_resolvent}.

\emph{(3)} Suppose the series converges in operator norm, so that
$P_\star^n-\Pi\rightarrow0$ in norm, and put $M:=\mathrm{Id}-P_\star+\Pi$. From
$P_\star\mathbf 1=\mathbf 1$ and $\lambda T=\lambda$ come $P_\star\Pi=\Pi P_\star=\Pi$, and
$\Pi^2=\Pi$, $\lambda$ being a probability; hence $M\Pi=\Pi M=\Pi$ and
$(\mathrm{Id}-P_\star)\Pi=\Pi(\mathrm{Id}-P_\star)=0$. Each term of the series satisfies
$\Pi(P_\star^n-\Pi)=(P_\star^n-\Pi)\Pi=0$, so $\Pi U=U\Pi=0$; and
$(\mathrm{Id}-P_\star)(P_\star^n-\Pi)=(P_\star^n-\Pi)(\mathrm{Id}-P_\star)=P_\star^n-P_\star^{n+1}$,
so both partial sums telescope to $\mathrm{Id}-P_\star^{N+1}\rightarrow\mathrm{Id}-\Pi$ and
$(\mathrm{Id}-P_\star)U=U(\mathrm{Id}-P_\star)=\mathrm{Id}-\Pi$. Therefore
$M(\Pi+U)=\Pi+(\mathrm{Id}-P_\star)U+\Pi U=\mathrm{Id}$ and
$(\Pi+U)M=\Pi+U(\mathrm{Id}-P_\star)+U\Pi=\mathrm{Id}$: the operator $M$ is invertible with
inverse $\Pi+U$, and $S=M^{-1}-\Pi=U$, which with
$(\mathrm{Id}-P_\star)U=U(\mathrm{Id}-P_\star)=\mathrm{Id}-\Pi$ and $\Pi U=U\Pi=0$ is
\eqref{eq:doubling_resolvent}. Then
$\widehat B=\|U\|\leq\sum_n\|P_\star^n-\Pi\|=\sum_n\widehat\beta_n$ by \emph{(1)}. On a finite
irreducible chain $S$ is defined by \emph{(2)}: when $\sum_{n\geq0}\widehat\beta_n=+\infty$ the
bound $\widehat B\leq\sum_{n\geq0}\widehat\beta_n$ holds with an infinite right side, and when
that sum is finite the series $\sum_{n\geq0}(P_\star^n-\Pi)$ converges absolutely in operator
norm by \emph{(1)}, so $S=U$ and the bound is the one just proved.
\end{proof}

\begin{remark}[the two constants written $\widehat B$, and the trajectory reading]\label{rem:doubling_two_constants}
The constant \eqref{eq:doubling_Bhat} is the norm of the diffusion operator, not the mixing sum
$\sum_n\widehat\beta_n$ that Lemma~\ref{lem:sigma_mixing} also writes $\widehat B$, and item
\emph{(3)} of Lemma~\ref{lem:doubling_operator} orders the two on any finite irreducible chain.
They need not coincide: on the two-state chain that flips with probability $3/4$, whose
invariant probability is uniform and whose centred functions are eigenfunctions of its
transition operator at the eigenvalue $-1/2$, $\widehat B=2/3$ while $\widehat\beta_n=2^{-n}$
for every $n\geq0$ and $\sum_n\widehat\beta_n=2$; at the flip probability $1/2$ both equal $1$. Numerically they
also differ on the truncation at which \S\ref{sec:doubling_measured} measures both. On a finite
irreducible chain, wherever the series of that item converges,
$(S\theta)(x)=\sum_{n\geq0}\bigl(\mathbb E(\theta(X_n)\mid X_0=x)-\lambda(\theta)\bigr)$ --- the
operator is read along backward trajectories,
$(P_\star f)(x)=\mathbb E(f(X_1)\mid X_0=x)$ for $f\in L^2(\lambda)$.
\end{remark}

\begin{theorem}[the doubling graph in one statement]\label{theo:doubling_main}
In the setting of Definition~\ref{def:doubling_setting}, with $(c,s)$ in the standing range:
\begin{enumerate}
    \item \emph{(the phase diagram)} the loop-closed backward chain $X$ is irreducible, and its
    recurrence class is
    \begin{center}
    \normalfont
    \begin{tabular}{l@{\quad}l|p{0.38\linewidth}}
        (a) & $s<1$, every $0<c<2^s$   & transient \\
        (b) & $s=1$, $0<c<1$           & positive recurrent \\
        (c) & $s=1$, $1\leq c<1/\ln2$  & null recurrent \\
        (d) & $s=1$, $c=1/\ln2$        & not positive recurrent; null recurrence or
                                         transience not decided \\
        (e) & $s=1$, $1/\ln2<c<2$      & transient \\
        (f) & $s>1$, every $0<c<2^s$   & positive recurrent
    \end{tabular}
    \end{center}
    \item \emph{(the invariant measure at $s=1$)} at $s=1$ and $0<c<1$, with $p_*>1$ the root
    of \eqref{eq:doubling_cramer} and $\lambda$ the invariant probability of $X$, there is a
    constant $C>0$, determined by $c$, $d$ and $\lambda_1,\dots,\lambda_d$, with
    $L(m)=\frac{C}{p_*-1}\,m^{1-p_*}\bigl(1+o(1)\bigr)$;
    \item \emph{(the diffusion operator, on the whole positive recurrent phase)} at every
    $(c,s)$ of the standing range at which $X$ is positive recurrent --- rows \emph{(b)} and
    \emph{(f)} of the table --- with $\lambda$ its invariant probability, and for every
    $p\in[1,+\infty]$,
    $\inf\bigl\{\|(\mathrm{Id}-P_\star)f\|_{L^p(\lambda)}:\Pi f=0,\ \|f\|_{L^p(\lambda)}=1\bigr\}=0$
    and no bounded operator $S$ on $L^p(\lambda)$ satisfies
    $S(\mathrm{Id}-P_\star)=\mathrm{Id}-\Pi$. At $p=2$ this gives $\widehat B=+\infty$ by the
    convention of Definition~\ref{def:doubling_setting}, the series
    $\sum_{n\geq0}(P_\star^n-\Pi)$ does not converge in operator norm, and
    $\sum_{n\geq0}\widehat\beta_n=+\infty$. At $s=1$ and $0<c<1$ the expected
    backward-trajectory length is finite, $\bar\sigma=\bar\jmath/(1-c)$, and there are
    moreover $c_7>0$, depending only on $c$, $d$ and $\lambda_1,\dots,\lambda_{2m_0}$ with
    $m_0$ the integer of Theorem~\ref{theo:doubling_decay}, and an integer $m_2>d$ with
    $\|f_m\|_{L^2(\lambda)}\geq c_7\sqrt m\,\|(\mathrm{Id}-P_\star)f_m\|_{L^2(\lambda)}$ for
    every $m\geq m_2$;
    \item \emph{(an unsolvable flow-matching problem)} at every $(c,s)$ of the standing range
    at which $X$ is positive recurrent, with $\lambda$ its invariant probability, there are
    probability densities $\finit,\fterm\in L^2(\lambda)$ such that no $f\in L^2(\lambda)$
    satisfies the flow-matching equation
    $\left(\mathrm{Id}-P_\star\right)f=\finit-\fterm$ of the loop closure;
    \item \emph{(the truncation)} at $s=1$ and $0<c<1$, the truncation at $K$ is a finite
    irreducible chain on $K+2$ states and $\widehat B_K<+\infty$ for every even integer
    $K\geq d$, and there are $c_8>0$ and an integer $K_0\geq d$, depending only on $c$, $d$
    and $\bar\jmath$, with $\widehat B_K\geq c_8\sqrt K$ for every even $K\geq K_0$.
\end{enumerate}
\end{theorem}

The two conditions come apart on this graph. That of \citet{morozov2025revisiting} is repaired
by any member with $s=1$ and $c<1$, which is what $\bar\sigma=\bar\jmath/(1-c)$ says
(Proposition~\ref{prop:doubling_length}); the summable-mixing hypothesis of
Theorem~\ref{theo:universality_L2} is repaired by no positive recurrent member of
\eqref{eq:doubling_family} --- item \emph{(3)} read at $p=2$ --- and neither is the
coercivity hypothesis \eqref{eq:coercivity} of Theorem~\ref{theo:db_stable_frozen_full}: by
item \emph{(4)} $\mathrm{Id}-P_\star$ does not map $L^2(\lambda)$ onto $\ker\Pi$
(Proposition~\ref{prop:doubling_unsolvable}\emph{(1)}), so its adjoint $\mathrm{Id}-P$
(Lemma~\ref{lem:doubling_operator}\emph{(1)}) is not bounded below on $\ker\Pi$, in the
flow-matching instance and hence in the detailed-balance one --- and at the null
recurrent and transient members there is no invariant probability, hence no $\lambda$, no
$\Pi$, no $\widehat\beta_n$, no $\widehat B$, and no summable mixing or coercivity that could
fail: both hypotheses fail there for want of the objects they quantify over. What item \emph{(4)}
refutes is \emph{exact} $L^2$ flow matching on the infinite chain --- a pair of $L^2$
densities the loop closure cannot balance at all --- and not the weak universality of
Theorem~\ref{theo:universality_L2_full}, whose conclusion is a defect infimum of zero; the
range of $\mathrm{Id}-P_\star$ in $\ker\Pi$ is what
Proposition~\ref{prop:doubling_unsolvable}\emph{(1)} settles, and whether the weak form fails
as well is not decided here. The sufficient condition of
Theorem~\ref{theo:db_stable_frozen_full} supplies no coercivity constant here either: its
summable mixing is, in both instances, $\sum_n\widehat\beta_n<+\infty$, which item \emph{(3)}
refutes at every positive recurrent member. The $c_7\sqrt m$ of \emph{(3)} is not a lower bound on
$\widehat B$, which is $+\infty$ there; it is a
rate along the trial functions $f_m$, and it is the truncation of \emph{(5)} that carries a
lower bound on a finite constant. The boundary of all this is
Remark~\ref{rem:doubling_geometric}: for a geometrically decaying doubling probability the
growth condition \eqref{eq:doubling_star} fails, and no boundedness is claimed for any policy on
this graph.

\subsection{The policy family and its diffusion operator}\label{sec:doubling_family}

This subsection proves the five items of Theorem~\ref{theo:doubling_main}, in six blocks. The
first classifies the family by its drift and gives the expected backward length in closed form
(Propositions~\ref{prop:doubling_drift}, \ref{prop:doubling_length}
and~\ref{prop:doubling_phase}). The second reads stationarity as a balance across a cut of the
ladder and names the exponent at which the invariant measure decays
(Propositions~\ref{prop:doubling_cut} and~\ref{prop:doubling_exponent},
Lemma~\ref{lem:doubling_cramer_root}). The third proves that decay two-sidedly and with no
asymptotic hypothesis (Theorem~\ref{theo:doubling_decay}), and the fourth replaces its two
constants by one, with a rate, and identifies it (Theorem~\ref{theo:doubling_sharp},
Proposition~\ref{prop:doubling_constant}). The fifth proves the diffusion operator unbounded on
every $L^p$ at every positive recurrent member of the family
(Theorem~\ref{theo:doubling_unbounded}, Lemma~\ref{lem:doubling_ramp},
Corollary~\ref{cor:doubling_family}) and exhibits a pair of $L^2$ densities the loop closure
cannot balance (Proposition~\ref{prop:doubling_unsolvable}), and the sixth restores a finite
constant on the
truncation, with a growth in the cap (Corollary~\ref{cor:doubling_truncation}), and assembles
the whole into Theorem~\ref{theo:doubling_main}.

\subsubsection{Drift, excursion length and the phase diagram}

This block proves item \emph{(1)} of Theorem~\ref{theo:doubling_main} and the closed form of the
expected backward-trajectory length at $s=1$. Proposition~\ref{prop:doubling_drift} computes the drift of the
backward chain and finds it independent of the state exactly at $s=1$: there a single constant
separates positive recurrence, null recurrence and transience, while off that exponent
$s>1$ and $s<1$ each carry a single phase. Proposition~\ref{prop:doubling_length} gives the expected backward length
in closed form at $s=1$ and $c<1$, with Lemma~\ref{lem:doubling_excursion} and
Lemma~\ref{lem:doubling_supersolution} --- an excursion bound and a linear supersolution --- as
the two milestones before it. Proposition~\ref{prop:doubling_phase} is the phase diagram.

\begin{proposition}[the drift, and when it is constant]\label{prop:doubling_drift}
In the setting of Definition~\ref{def:doubling_setting}, let $(c,s)$ lie in the standing
range. Then for every $n\geq0$ and every ladder state $j\geq1$,
\begin{equation}\label{eq:doubling_drift}
    \mathbb E\left[X_{n+1}-X_n\mid X_n=j\right]\;=\;c\,(j+1)^{1-s}-1 ,
\end{equation}
and this quantity is independent of $j$ if and only if $s=1$, where it equals $c-1$.
\end{proposition}

\begin{proof}
By \eqref{eq:doubling_policy} the chain moves from $j$ to $2j$ with probability
$\varepsilon(j)$ and to $j-1$ with probability $1-\varepsilon(j)$, so
$\mathbb E\left[X_{n+1}-X_n\mid X_n=j\right]=\varepsilon(j)\,(2j-j)+\bigl(1-\varepsilon(j)\bigr)(-1)=\varepsilon(j)(j+1)-1$;
substituting $\varepsilon(j)=c\,(j+1)^{-s}$ from \eqref{eq:doubling_family} gives
$c\,(j+1)^{1-s}-1$, which is \eqref{eq:doubling_drift}. At $s=1$ this equals $c-1$ for every
$j$; if $s\neq1$ then $j\mapsto c\,(j+1)^{1-s}$ is strictly monotone and $c>0$ on the standing
range, so its values at $j=1$ and at $j=2$ differ.
\end{proof}

\begin{lemma}[the excursion maximum is bounded by the hitting time]\label{lem:doubling_excursion}
In the setting of Definition~\ref{def:doubling_setting}, let $(X_n)_{n\geq0}$ be either the
backward chain of the loop closure or, for an integer $K\geq d$, the backward chain of the
truncation at $K$ --- the loop closure with every doubling edge $j\rightarrow2j$ such that
$2j>K$ deleted, the decrement edge out of such a state carrying probability $1$ --- and let
$\sigma$ be its hitting time of $s_0$. Then $X_{n+1}\geq X_n-1$ at every step issued from a
ladder state, and for every ladder state $j$,
\begin{equation}\label{eq:doubling_excursion}
    \max_{0\leq n\leq\sigma}X_n\;\leq\;\sigma\qquad\mathbb P(\,\cdot\mid X_0=j)\text{-almost surely},
\end{equation}
the left-hand side being read as $\sup_{n\geq0}X_n$ on $\{\sigma=+\infty\}$.
\end{lemma}

\begin{proof}
From a ladder state $j$ the chain moves to $2j\geq j$ or to $j-1$ by
\eqref{eq:doubling_policy}, the doubling edge being absent on the truncation when $2j>K$ and
the decrement then carrying probability $1$; in both chains $X_{n+1}\geq X_n-1$ at every step
issued from a ladder state. Started at a ladder state the chain visits only ladder states
before $\sigma$, the state $s_f$ being entered from $s_0$ alone, so that inequality holds at
every $n<\sigma$.

On $\{\sigma=+\infty\}$ the asserted inequality reads $\sup_{n\geq0}X_n\leq+\infty$. On
$\{\sigma<+\infty\}$ let $n_*\leq\sigma$ attain $\max_{0\leq n\leq\sigma}X_n$. Each of the
$\sigma-n_*$ steps from $n_*$ to $\sigma$ lowers the chain by at most $1$ and
$X_\sigma=s_0$, read as $0$, so $\sigma-n_*\geq\max_{0\leq n\leq\sigma}X_n$ and
$\sigma\geq\max_{0\leq n\leq\sigma}X_n$.
\end{proof}

\begin{lemma}[the linear supersolution]\label{lem:doubling_supersolution}
In the setting of Definition~\ref{def:doubling_setting}, let $s=1$ and $0<c<1$, and put
$h(j):=j/(1-c)$ at every ladder state $j\geq1$ and $h(s_0):=0$. Let $K\geq d$ be an integer
and let the truncation at $K$ be as in Lemma~\ref{lem:doubling_excursion}, with $\bar\sigma_K$
denoting the functional $\bar\sigma$ of \eqref{eq:morozov_variables} formed on it. Then
\begin{equation}\label{eq:doubling_super}
    h(j)\;=\;1+\varepsilon(j)\,h(2j)+\bigl(1-\varepsilon(j)\bigr)\,h(j-1)
    \qquad(j\geq1)
\end{equation}
on the loop closure; on the truncation at $K$, $h(j)\geq1+\mathbb E(h(X_1)\mid X_0=j)$ for every
$1\leq j\leq K$, with equality for $2j\leq K$. Consequently, on both chains,
$\mathbb E(\sigma\mid X_0=j)\leq j/(1-c)<+\infty$ at every ladder state $j$, so that $\sigma<+\infty$
$\mathbb P(\,\cdot\mid X_0=j)$-almost surely, and $\bar\sigma\leq\bar\jmath/(1-c)$ and
$\bar\sigma_K\leq\bar\jmath/(1-c)$.
\end{lemma}

\begin{proof}
\emph{Step 1 (the identity on the loop closure).} Substituting $\varepsilon(j)=c/(j+1)$ from
\eqref{eq:doubling_family} at $s=1$,
\begin{equation}\label{eq:doubling_superstep}
    \varepsilon(j)\,2j+\bigl(1-\varepsilon(j)\bigr)(j-1)
    \;=\;\frac{2cj-c(j-1)}{j+1}+(j-1)\;=\;\frac{cj+c}{j+1}+j-1\;=\;c+j-1 .
\end{equation}
Since $h$ is the linear function $x\mapsto x/(1-c)$ on $\{0,1,2,\dots\}$, with $h(s_0)=0$ read
as $h(0)$, \eqref{eq:doubling_superstep} gives
$1+\varepsilon(j)h(2j)+(1-\varepsilon(j))h(j-1)=1+\frac{c+j-1}{1-c}=\frac j{1-c}=h(j)$, which
is \eqref{eq:doubling_super}.

\emph{Step 2 (the inequality on the truncation).} Let $1\leq j\leq K$. If $2j\leq K$, the
doubling edge out of $j$ is present and the transitions at $j$ are those of the loop closure,
so $1+\mathbb E(h(X_1)\mid X_0=j)=h(j)$ by Step 1. If $2j>K$, the doubling edge has been deleted and
the decrement carries probability $1$, so
$1+\mathbb E(h(X_1)\mid X_0=j)=1+\frac{j-1}{1-c}\leq\frac j{1-c}=h(j)$, the inequality being
$1\leq1/(1-c)$.

\emph{Step 3 (a non-negative supersolution dominates the expected hitting time).} Fix either
chain and put $u_n(x):=\mathbb E(\sigma\wedge n\mid X_0=x)$ at $s_0$ and at every ladder state $x$ of
that chain. Then $u_0=0$ and $u_n(s_0)=0$ for every $n$, and conditioning on the first step
gives $u_{n+1}(j)=1+\mathbb E(u_n(X_1)\mid X_0=j)$ at every ladder state $j$, the chain moving from
$j$ either to a ladder state or to $s_0$. Now $u_0\leq h$, and if $u_n\leq h$ then
$u_{n+1}(j)\leq1+\mathbb E(h(X_1)\mid X_0=j)\leq h(j)$ by Steps 1 and 2 while $u_{n+1}(s_0)=0=h(s_0)$:
by induction $u_n\leq h$ for every $n$. Since $\sigma\wedge n\uparrow\sigma$, monotone
convergence gives $\mathbb E(\sigma\mid X_0=j)=\lim_nu_n(j)\leq h(j)=j/(1-c)<+\infty$ at every ladder
state $j$, and a finite expectation gives $\sigma<+\infty$ $\mathbb P(\,\cdot\mid X_0=j)$-almost surely. The
target row is a probability supported in $\{1,\dots,d\}\subseteq\{1,\dots,K\}$ with mean
$\bar\jmath$ and $h$ is linear there, so averaging $\mathbb E(\sigma\mid X_0=j)\leq h(j)$ over it
gives $\bar\sigma\leq\bar\jmath/(1-c)$ on the loop closure and
$\bar\sigma_K\leq\bar\jmath/(1-c)$ on the truncation at $K$.
\end{proof}

\begin{proposition}[closed-form backward length]\label{prop:doubling_length}
In the setting of Definition~\ref{def:doubling_setting}, let $s=1$ and $0<c<1$. Then
$\mathbb E(\sigma\mid X_0=j)=j/(1-c)$ for every ladder state $j\geq1$, and
$\bar\sigma=\bar\jmath/(1-c)$.
\end{proposition}

\begin{proof}
\emph{Step 1 (the upper bound).} By Lemma~\ref{lem:doubling_supersolution},
$\mathbb E(\sigma\mid X_0=j)\leq j/(1-c)<+\infty$ at every ladder state $j$ and $\sigma<+\infty$
$\mathbb P(\,\cdot\mid X_0=j)$-almost surely.

\emph{Step 2 (the excursion maximum is integrable).} By Lemma~\ref{lem:doubling_excursion},
$\max_{0\leq n\leq\sigma}X_n\leq\sigma$ $\mathbb P(\,\cdot\mid X_0=j)$-almost surely, so
$\mathbb E\bigl(\max_{0\leq n\leq\sigma}X_n\mid X_0=j\bigr)\leq\mathbb E(\sigma\mid X_0=j)<+\infty$ by
Step 1.

\emph{Step 3 (optional stopping).} By Proposition~\ref{prop:doubling_drift} at $s=1$ the
drift is $c-1$ at every ladder state, so $X_{n\wedge\sigma}+(1-c)(n\wedge\sigma)$ is a
martingale and $\mathbb E(X_{n\wedge\sigma}\mid X_0=j)=j-(1-c)\,\mathbb E(n\wedge\sigma\mid X_0=j)$. As
$\sigma<+\infty$ almost surely by Step 1, $X_{n\wedge\sigma}\rightarrow X_\sigma=0$
pointwise, dominated by the integrable $\max_{0\leq n\leq\sigma}X_n$ of Step 2, so the left
side tends to $0$ by dominated convergence, while
$\mathbb E(n\wedge\sigma\mid X_0=j)\rightarrow\mathbb E(\sigma\mid X_0=j)$ by monotone convergence. Hence
$0=j-(1-c)\mathbb E(\sigma\mid X_0=j)$, that is $\mathbb E(\sigma\mid X_0=j)=j/(1-c)$. Averaging over the
target row gives $\bar\sigma=\bar\jmath/(1-c)$.
\end{proof}

The policy exhibited after \eqref{eq:counterexample_graph} in
Section~\ref{sec:cv_divergence}, $\BackwardPolicy(j\rightarrow2j)=\frac1{2(j+1)}$, is the
member $c=\tfrac12$, $s=1$ of the family \eqref{eq:doubling_family}, and
Proposition~\ref{prop:doubling_length} gives its expected backward length as
$\bar\jmath/(1-c)=2\bar\jmath$, which is the $\bar\sigma=9$ stated there exactly when the target
row has mean $\bar\jmath=4.5$ --- the mean every measurement of \S\ref{sec:doubling_measured}
uses.
Theorem~\ref{theo:doubling_unbounded} below gives $\widehat B=+\infty$ for that same policy.

\begin{proposition}[the phase diagram of the family]\label{prop:doubling_phase}
In the setting of Definition~\ref{def:doubling_setting}, let $(c,s)$ lie in the standing
range and let $X$ be the loop-closed backward chain on the infinite graph
\eqref{eq:counterexample_graph}. Then $X$ is irreducible, and:
\begin{enumerate}
    \item if $s>1$, then $X$ is positive recurrent;
    \item if $s<1$, then $X$ is transient;
    \item if $s=1$, then $X$ is positive recurrent for $0<c<1$, with
    $\mathbb E(\sigma\mid X_0=j)=j/(1-c)$ at every ladder state $j\geq1$; for every
    $1\leq c<2$, $\mathbb E(\sigma\mid X_0=j)=+\infty$ at every ladder state $j\geq1$ and $X$
    is not positive recurrent; and $X$ is null recurrent for $1\leq c<1/\ln2$ and transient
    for $1/\ln2<c<2$.
\end{enumerate}
\end{proposition}

\begin{proof}
\emph{Step 1 (irreducibility).} On the standing range $\varepsilon_{c,s}$ takes its values in
$(0,1)$ (Lemma~\ref{lem:doubling_range}), so $X$ is irreducible
(Lemma~\ref{lem:doubling_irreducible}).

\emph{Step 2 (the Lyapunov expansion).} For $t>0$ define $W_t$ on the whole loop closure by
$W_t(j):=(j+1)^{-t}$ at a ladder state $j$ and $W_t(s_0)=W_t(s_f):=1$. There is a constant
$\Xi_t<+\infty$ depending on $t$ alone such that, at every ladder state $j\geq1$,
\begin{equation}\label{eq:doubling_Wdrift}
\begin{aligned}
    \mathbb E\left(W_t(X_1)\mid X_0=j\right)-W_t(j)
    &\;=\;(j+1)^{-t-1}\Bigl[\,t-c\left(1-2^{-t}\right)(j+1)^{1-s}\Bigr]\;+\;\Upsilon_t(j),\\
    |\Upsilon_t(j)|&\;\leq\;\Xi_t\Bigl((j+1)^{-t-2}+\varepsilon(j)(j+1)^{-t-1}\Bigr).
\end{aligned}
\end{equation}
By \eqref{eq:doubling_policy},
$\mathbb E(W_t(X_1)\mid X_0=j)-W_t(j)=\varepsilon(j)\bigl((2j+1)^{-t}-(j+1)^{-t}\bigr)+\bigl(1-\varepsilon(j)\bigr)\bigl(j^{-t}-(j+1)^{-t}\bigr)$.
From $2j+1=2(j+1)\bigl(1-\tfrac1{2(j+1)}\bigr)$ and $j=(j+1)\bigl(1-\tfrac1{j+1}\bigr)$,
\begin{equation*}
\begin{aligned}
    (2j+1)^{-t}&\;=\;2^{-t}(j+1)^{-t}\Bigl(1+\tfrac t{2(j+1)}+O(j^{-2})\Bigr),\\
    j^{-t}&\;=\;(j+1)^{-t}\Bigl(1+\tfrac t{j+1}+O(j^{-2})\Bigr),
\end{aligned}
\end{equation*}
with remainders bounded in terms of $t$ alone. The first product therefore contributes
$\varepsilon(j)(j+1)^{-t}\bigl(2^{-t}-1\bigr)=-c\left(1-2^{-t}\right)(j+1)^{-t-1}(j+1)^{1-s}$
by $\varepsilon(j)=c(j+1)^{-s}$, plus a term $O\bigl(\varepsilon(j)(j+1)^{-t-1}\bigr)$; the
second contributes $t\,(j+1)^{-t-1}$, plus terms $O\bigl((j+1)^{-t-2}\bigr)$ and
$O\bigl(\varepsilon(j)(j+1)^{-t-1}\bigr)$. Collecting the two leading terms into the bracket
and the rest into $\Upsilon_t$ gives \eqref{eq:doubling_Wdrift}.

\emph{Step 3 (positive recurrence for $s>1$).} Define $V_1$ on the whole loop closure by
$V_1(j):=j$ at a ladder state $j$ and $V_1(s_0)=V_1(s_f):=0$. By
Proposition~\ref{prop:doubling_drift} the drift of $V_1$ at a ladder state $j$ is
$c(j+1)^{1-s}-1$, which tends to $-1$ as $j\rightarrow\infty$ because $s>1$; fix $j_0$ with
$c(j+1)^{1-s}-1\leq-\tfrac12$ for $j>j_0$. Then $V_1\geq0$ has finite sublevel sets,
$\mathbb E(V_1(X_1)\mid X_0=x)<+\infty$ for every $x$, and the drift is at most $-\tfrac12$ off the
finite set $\{s_0,s_f,1,\dots,j_0\}$: Foster's criterion for positive recurrence applies, and
$X$ is positive recurrent.

\emph{Step 4 (transience for $s<1$).} Fix $t:=1$ in Step 2. Since $s<1$, the bracket in
\eqref{eq:doubling_Wdrift} is $1-c\left(1-2^{-1}\right)(j+1)^{1-s}\rightarrow-\infty$, so the
displayed term is of order $-(j+1)^{-1-s}$, while
$\Upsilon_1(j)=O\bigl((j+1)^{-3}\bigr)+O\bigl((j+1)^{-2-s}\bigr)$ is smaller by a factor $(j+1)^{-1}$
at least; fix $j_0$ with $\mathbb E(W_1(X_1)\mid X_0=j)\leq W_1(j)$ for $j>j_0$. The function $W_1$ is
non-negative and bounded by $1$, and $W_1(j)\rightarrow0$, so some $x_0>j_0$ has
$W_1(x_0)<(j_0+1)^{-1}$, the infimum of $W_1$ on the finite set $\{s_0,s_f,1,\dots,j_0\}$ off
which $W_1$ is a supermartingale: the standard transience criterion --- a bounded
non-negative supermartingale off a finite set, taking somewhere off that set a value below
its infimum on it --- applies, and $X$ is transient.

\emph{Step 5 ($s=1$ and $c<1$: positive recurrence).} By
Proposition~\ref{prop:doubling_length}, $\mathbb E(\sigma\mid X_0=j)=j/(1-c)$ at every ladder state
$j$ and $\bar\sigma=\bar\jmath/(1-c)$. The return time of $X$ to $s_0$ is one step to $s_f$,
one step into $\{1,\dots,d\}$ and then $\sigma$, so its expectation is
$2+\bar\sigma<+\infty$; with Step 1, $X$ is positive recurrent.

\emph{Step 6 ($s=1$ and $c<1/\ln2$: recurrence).} Define $V_2$ on the whole loop closure by
$V_2(j):=\log(j+1)$ at a ladder state $j$ and $V_2(s_0)=V_2(s_f):=0$. From
$\log(2j+1)-\log(j+1)=\ln2-\tfrac1{2(j+1)}+O(j^{-2})$,
$\log j-\log(j+1)=-\tfrac1{j+1}+O(j^{-2})$ and $\varepsilon(j)=c/(j+1)$, the drift of $V_2$ at
$j$ is $\frac{c\ln2-1}{j+1}+O(j^{-2})$, which is $\leq0$ for $j\geq j_0$ when $c\ln2<1$; fix
such a $j_0$. Since $V_2\geq0$ has finite sublevel sets and its drift is $\leq0$ off the finite
set $\{s_0,s_f,1,\dots,j_0\}$, Foster's criterion for recurrence applies, and $X$ is
recurrent.

\emph{Step 7 ($s=1$ and $1\leq c<2$: $\mathbb E(\sigma\mid X_0=j)=+\infty$, and null recurrence below
$1/\ln2$).} By Proposition~\ref{prop:doubling_drift} the drift is $c-1\geq0$, so
$(X_{n\wedge\sigma})_n$ is a submartingale. Fix an integer $r>1$ and let $\varsigma_r$ be the
exit time of $[1,r)$; the stopped chain takes values in $\{0,1,\dots,2r-2\}$, and from any
state of $[1,r)$ at most $r-1$ consecutive decrements reach $0$, an event of probability at
least $\prod_{i=1}^{r-1}(1-\varepsilon(i))>0$ since $c<2$ makes every $\varepsilon(i)<1$, so
$\varsigma_r<+\infty$ almost surely. Optional stopping for the bounded stopped submartingale
gives $\mathbb E(X_{\varsigma_r}\mid X_0=1)\geq1$. Exiting below means $X_{\varsigma_r}=0$ and exiting
above means $r\leq X_{\varsigma_r}\leq2r-2$, so
$\mathbb E(X_{\varsigma_r}\mid X_0=1)\leq(2r-2)\,\mathbb P(X_{\varsigma_r}\geq r\mid X_0=1)$ and
$\mathbb P(X_{\varsigma_r}\geq r\mid X_0=1)\geq1/(2r)$. Exiting below is $\sigma$ itself, so
$\varsigma_r\leq\sigma$ and $\max_{0\leq n\leq\sigma}X_n\geq X_{\varsigma_r}$, whence
$\mathbb P\bigl(\max_{0\leq n\leq\sigma}X_n\geq r\mid X_0=1\bigr)\geq1/(2r)$; summing over $r\geq2$,
$\mathbb E\bigl(\max_{0\leq n\leq\sigma}X_n\mid X_0=1\bigr)\geq\sum_{r\geq2}\frac1{2r}=+\infty$, and
$\max_{0\leq n\leq\sigma}X_n\leq\sigma$ by Lemma~\ref{lem:doubling_excursion} gives
$\mathbb E(\sigma\mid X_0=1)=+\infty$. For $j\geq2$ the chain reaches $1$ before $0$ --- decrements
are by one --- with probability at least $\prod_{i=2}^{j}(1-\varepsilon(i))>0$, so the strong
Markov property at that hitting time gives $\mathbb E(\sigma\mid X_0=j)=+\infty$ as well. The return
time of $X$ to $s_0$ has expectation
$2+\bar\sigma=2+\sum_{j=1}^{d}\widehat\BackwardPolicy(s_f\rightarrow j)\,\mathbb E(\sigma\mid X_0=j)=+\infty$,
the target row being a probability supported in $\{1,\dots,d\}$ and every
$\mathbb E(\sigma\mid X_0=j)$ being infinite. An irreducible chain whose expected return time
at one state is infinite is not positive recurrent \citep{levin2017markov}: $X$ is not
positive recurrent at any $1\leq c<2$. For $1\leq c<1/\ln2$ the chain is recurrent by Step 6,
and an irreducible recurrent chain whose expected return time at one state is infinite is
null recurrent \citep{levin2017markov}: $X$ is null recurrent there.

\emph{Step 8 ($s=1$ and $c>1/\ln2$: transience).} Put $\phi(t):=c(1-2^{-t})-t$, so that
$\phi(0)=0$ and $\phi'(0)=c\ln2-1>0$; fix $t>0$ with $\phi(t)>0$. At $s=1$ the bracket in
\eqref{eq:doubling_Wdrift} is the constant $-\phi(t)$ and
$\Upsilon_t(j)=O\bigl((j+1)^{-t-2}\bigr)$ by $\varepsilon(j)=c/(j+1)$, so
$\mathbb E(W_t(X_1)\mid X_0=j)-W_t(j)<0$ for $j$ large, and the transience criterion applies as in
Step 4: $X$ is transient.
\end{proof}

At $s=1$ and $c=1/\ln2$ the chain is not positive recurrent, and null recurrence and
transience are left open.

\subsubsection{The cut balance and the Cram\'er root}

This block states the balance that stationarity imposes across a cut of the ladder, and names
the exponent that balance forces on the invariant measure. Proposition~\ref{prop:doubling_cut}
equates the two fluxes across such a cut, on the loop closure and on the truncation alike, and
adds two one-sided ratios. Proposition~\ref{prop:doubling_exponent} identifies the exponent of
any polynomially decaying solution of that balance as a non-zero root of the Cram\'er equation
of Definition~\ref{def:doubling_setting} --- a Cram\'er equation in the exponential variable,
which is what names $p_*$. Lemma~\ref{lem:doubling_cramer_root} collects what is used of that
root: its existence and its uniqueness away from a single value of $c$, and its two crossings,
at $1$ and at $0$, which are the two thresholds of the phase diagram.

\begin{proposition}[cut balance]\label{prop:doubling_cut}
In the setting of Definition~\ref{def:doubling_setting}, let $\varepsilon(j)\in(0,1)$ for
every $j\geq1$ and let $\lambda$ be an invariant probability of the loop-closed chain. Then
for every $m>d$
\begin{equation}\label{eq:doubling_cut}
    \lambda_m\bigl(1-\varepsilon(m)\bigr)\;=\;\sum_{j=\lceil m/2\rceil}^{m-1}\lambda_j\,\varepsilon(j),
\end{equation}
and, for every $j\geq1$ with $2j$ a state of the chain and every $m\geq2$,
\begin{equation}\label{eq:doubling_ratios}
    \lambda_{2j}\;\geq\;\lambda_j\,\varepsilon(j),
    \qquad
    \lambda_m\;\leq\;\frac{\lambda_{m-1}}{1-\varepsilon(m)} .
\end{equation}
On the truncation at an even integer $K\geq d$, with $\lambda^K$ in place of $\lambda$,
\eqref{eq:doubling_cut} holds for every integer $m$ with $d<m\leq K/2$, and
\eqref{eq:doubling_ratios} for every $j\geq1$ with $2j\leq K$ and every integer $m$ with
$2\leq m\leq K$.
\end{proposition}

\begin{proof}
\emph{Step 1 (the identities on the loop closure).} Let $m>d$ and put $A:=[m,+\infty)$.
Stationarity equates the fluxes across $A$.
A state $x\in A$ leaves $A$ only through $x\rightarrow x-1$ with $x-1<m$, that is $x=m$,
with probability $1-\varepsilon(m)$. A state $x\notin A$ enters $A$ only through
$x\rightarrow2x$ with $2x\geq m$, that is $\lceil m/2\rceil\leq x\leq m-1$, with probability
$\varepsilon(x)$. The remaining
transitions into $A$ would come from $s_0$, whose only image is $s_f$, and from $s_f$, whose
image lies in $\{1,\dots,d\}$ with $d<m$. This is \eqref{eq:doubling_cut}. For
\eqref{eq:doubling_ratios}, stationarity at the state $2j$ writes $\lambda_{2j}$ as a sum of
non-negative terms one of which is $\lambda_j\varepsilon(j)$, the doubling edge $j\rightarrow2j$
being present by hypothesis, and stationarity at $m-1$ writes $\lambda_{m-1}$ as a sum of
non-negative terms one of which is $\lambda_m(1-\varepsilon(m))$; dropping the others gives
the two inequalities.

\emph{Step 2 (the identities on the truncation).} Let $K\geq d$ be even, let $d<m\leq K/2$ and
put $A:=[m,K]$. A state $x\in A$ with $x>m$ decrements to $x-1\geq m$ and, when its doubling
edge is present, doubles to $2x\leq K$ with $2x\geq x\geq m$: both images lie in $A$, the top
state $K$ decrementing to $K-1\geq2m-1\geq m$. The doubling edge out of $m$ is present because
$2m\leq K$, so $A$ is left from $m$ alone and with probability $1-\varepsilon(m)$. A state
$x\notin A$ enters $A$ only through $x\rightarrow2x$ with $2x\geq m$, that is
$\lceil m/2\rceil\leq x\leq m-1$, and every such edge is present because $2x\leq2m-2<K$; the
images of $s_0$ and $s_f$ are those of the loop closure and lie outside $A$. Equating the
fluxes across $A$ is \eqref{eq:doubling_cut} for $\lambda^K$. For \eqref{eq:doubling_ratios},
let $j\geq1$ with $2j\leq K$ and let $2\leq m\leq K$: the doubling edge $j\rightarrow2j$ is
present, so stationarity at $2j$ writes $\lambda^K_{2j}$ as a sum of non-negative terms one of
which is $\lambda^K_j\varepsilon(j)$; and stationarity at $m-1$ writes $\lambda^K_{m-1}$ as a
sum of non-negative terms one of which is the flux from $m$ by decrement, equal to
$\lambda^K_m\bigl(1-\varepsilon(m)\bigr)$ when $2m\leq K$ and to $\lambda^K_m$ when $2m>K$, in
either case at least $\lambda^K_m(1-\varepsilon(m))$; dropping the others gives the two
inequalities.
\end{proof}

\begin{proposition}[the decay exponent is a Cram\'er root]\label{prop:doubling_exponent}
In the setting of Definition~\ref{def:doubling_setting}, let $s=1$ and let $c>0$ --- inside the
standing range or not --- so that $\varepsilon(j)=c/(j+1)$, and let $(\lambda_j)_{j\geq1}$ be a
positive sequence satisfying \eqref{eq:doubling_cut} for every $m>d$. If $\lambda_j=Aj^{-p}(1+o(1))$ as
$j\rightarrow\infty$ for some $A>0$ and some $p\neq0$, then $p$ satisfies the Cram\'er
equation \eqref{eq:doubling_cramer}.
\end{proposition}

\begin{proof}
The left side of \eqref{eq:doubling_cut} is $\lambda_m(1-\varepsilon(m))=Am^{-p}(1+o(1))$,
since $\varepsilon(m)\rightarrow0$. On the right,
$\lambda_j\varepsilon(j)=cAj^{-p-1}(1+o(1))$ uniformly for $\lceil m/2\rceil\leq j\leq m-1$,
and the summand $u\mapsto u^{-p-1}$ is monotone with consecutive ratio $1+O(1/m)$ on that
range, so the sum equals $\int_{m/2}^{m}u^{-p-1}du$ up to a relative $O(1/m)$; as $p\neq0$
that integral is $\frac{m^{-p}}{p}\left(2^p-1\right)$. Equating the two sides and dividing by
$Am^{-p}>0$ gives \eqref{eq:doubling_cramer}.
\end{proof}

\begin{lemma}[the Cram\'er root]\label{lem:doubling_cramer_root}
Let $c>0$ and put $\psi(t):=c\bigl(2^{t}-1\bigr)-t$ for $t\in\mathbb R$, so that
\eqref{eq:doubling_cramer} reads $\psi(p)=0$. Then $\psi$ is strictly convex and $\psi(0)=0$;
$\psi$ has exactly one root other than $0$ when $c\neq1/\ln2$, and none when $c=1/\ln2$.
Write $p_*$ for that root and $\tau:=2^{p_*}$ when $c\neq1/\ln2$. Then:
\begin{enumerate}
    \item $p_*>0$ if and only if $c<1/\ln2$, and $p_*>1$ if and only if $c<1$;
    \item the Cram\'er equation reads
    \begin{equation}\label{eq:doubling_cramer_form}
        c\,\frac{\tau-1}{p_*}\;=\;1 ,
    \end{equation}
    and the function $k(v):=c\ln2\cdot2^{p_*v}$ on $[0,1]$ has $\int_0^1k(v)\,dv=1$;
    \item if $0<c<1$, then
    \begin{equation}\label{eq:doubling_cramer_ineq}
    \begin{aligned}
        \tau\;&>\;2,\qquad \tau\;>\;p_*+1,\\
        2\tau\;&>\;p_*+2,\qquad 4\tau\;>\;p_*+3 .
    \end{aligned}
    \end{equation}
\end{enumerate}
\end{lemma}

\begin{proof}
The function $t\mapsto2^{t}$ is strictly convex and $t\mapsto-t$ is affine, so $\psi$ is
strictly convex; $\psi(0)=c(1-1)-0=0$; and a strictly convex function has at most two roots,
so $\psi$ has at most one root other than $0$. Its derivative is
$\psi'(t)=c\ln2\cdot2^{t}-1$, so $\psi'(0)=c\ln2-1$. Since $c>0$,
$\psi(t)\rightarrow+\infty$ as $t\rightarrow+\infty$, and $\psi(t)=-t-c+c2^{t}$ gives
$\psi(t)\rightarrow+\infty$ as $t\rightarrow-\infty$.

If $c<1/\ln2$ then $\psi'(0)<0$, so $\psi<0$ immediately to the right of $0$, and continuity
with $\psi(t)\rightarrow+\infty$ as $t\rightarrow+\infty$ gives a root $p_*>0$. If
$c>1/\ln2$ then $\psi'(0)>0$, so $\psi<0$ immediately to the left of $0$, and
$\psi(t)\rightarrow+\infty$ as $t\rightarrow-\infty$ gives a root $p_*<0$. If $c=1/\ln2$ then
$\psi'(0)=0$, and strict convexity makes $0$ the strict minimum of $\psi$, so $\psi(t)>0$ for
$t\neq0$ and $0$ is the only root. This proves the existence and uniqueness assertions and the
first equivalence of \emph{(1)}.

In the remainder $c\neq1/\ln2$, so that $p_*$ and $\tau$ are defined.

\emph{(1), second equivalence.} Strict convexity with the two roots $0$ and $p_*$ makes $\psi<0$ strictly
between them and $\psi>0$ outside their closed interval, and $\psi(1)=c-1$. If $c<1$ then
$\psi(1)<0$, so $1$ lies strictly between $0$ and $p_*$, and $p_*>0$ by the first
equivalence: hence $p_*>1$. If $c=1$ then $\psi(1)=0$ and $p_*=1$. If $c>1$ then
$\psi(1)>0$, so $1$ lies outside the closed interval with endpoints $0$ and $p_*$, and
$p_*<1$.

\emph{(2)} $\psi(p_*)=0$ reads $c(\tau-1)=p_*$, and $p_*\neq0$: dividing by $p_*$ gives
\eqref{eq:doubling_cramer_form}. Again by $p_*\neq0$,
$\int_0^1c\ln2\cdot2^{p_*v}\,dv=c\ln2\left[\frac{2^{p_*v}}{p_*\ln2}\right]_0^1=c\,\frac{\tau-1}{p_*}$,
which is $1$ by \eqref{eq:doubling_cramer_form}.

\emph{(3)} By \emph{(1)}, $0<c<1$ gives $p_*>1$, hence $\tau=2^{p_*}>2$. The function
$t\mapsto2^{t}-t-1$ is strictly convex and vanishes at $t=0$ and at $t=1$, so it is positive
at every $t>1$; applying this at $t=p_*$, at $t=p_*+1$ and at $t=p_*+2$, each greater than
$1$, gives $\tau>p_*+1$, $2\tau=2^{p_*+1}>p_*+2$ and $4\tau=2^{p_*+2}>p_*+3$.
\end{proof}

\subsubsection{Two-sided decay of the invariant measure}

This block proves that the invariant measure decays at the Cram\'er exponent with no asymptotic
hypothesis, which is the two-sided form of item \emph{(2)} of
Theorem~\ref{theo:doubling_main}. Definition~\ref{def:doubling_decay_notation} fixes the
notation the estimates run in. Lemma~\ref{lem:doubling_averaging} reads the balance across a
cut as an average of the measure over a window of indices below that cut, and
Lemma~\ref{lem:doubling_expansion} expands the total weight of that average to first order in
the reciprocal of the cut, Remark~\ref{rem:doubling_parity} recording what the parity of the
window costs a pure power. Lemma~\ref{lem:doubling_descent} builds a Markov chain descending
the ladder and gives the measure at one index as an expectation along its paths,
Lemma~\ref{lem:doubling_escape} bounds the time that chain spends in a dyadic block, and
Lemma~\ref{lem:doubling_product} bounds the weight a descent accumulates, in expectation.
Theorem~\ref{theo:doubling_decay} carries the two-sided bound: its first item controls the
measure at every index above a block of indices by the measure on that block alone, its second
bounds the measure above and below by the pure power at the Cram\'er exponent.
Corollary~\ref{cor:doubling_tail} reads that bound on the tail and on the ratio of the tail to
the mass at the cut.

\begin{definition}[notation for the decay estimates]\label{def:doubling_decay_notation}
In the setting of Definition~\ref{def:doubling_setting}, let $s=1$ and $0<c<1$, so that
$\varepsilon(j)=c/(j+1)$, and, with $p_*$ the root of Lemma~\ref{lem:doubling_cramer_root},
write $p:=p_*$ and $\tau:=2^{p}$, so that $p>1$ and $\tau>2$ by items \emph{(1)} and
\emph{(3)} of that lemma. For an
integer $m>d$, the \emph{window}, the \emph{parity} and the \emph{cut weights} at $m$ are
\begin{equation}\label{eq:doubling_window}
\begin{aligned}
    W(m)\;&:=\;\{\lceil m/2\rceil,\dots,m-1\},\\
    \delta(m)\;&:=\;m-2\lfloor m/2\rfloor,\\
    q_m(j)\;&:=\;\frac{c}{(j+1)\bigl(1-\varepsilon(m)\bigr)}\quad\bigl(j\in W(m)\bigr).
\end{aligned}
\end{equation}
For $\alpha\in\mathbb R$ and $j\geq1$ put $\Phi_\alpha(j):=j^{-p-1}(j+\alpha)=j^{-p}(1+\alpha/j)$,
and for an integer $m>\max(d,2|\alpha|)$
\begin{equation}\label{eq:doubling_R}
    R_\alpha(m)\;:=\;\frac{1}{\Phi_\alpha(m)}\sum_{j\in W(m)}q_m(j)\,\Phi_\alpha(j),
\end{equation}
and $w_m(j):=q_m(j)\,(m/j)^{p}$ for $j\in W(m)$. For a positive sequence
$(\lambda_j)_{j\geq1}$ --- in particular for the invariant probability of
Definition~\ref{def:doubling_setting} at these parameters --- the
\emph{rescaled profile} is $u_m:=\lambda_m\,m^{p}$. For integers $\ell\geq1$ and $i\geq0$,
$I_i(\ell):=[2^{i}\ell,2^{i+1}\ell)$.
\end{definition}

\begin{lemma}[the averaging form of the cut balance]\label{lem:doubling_averaging}
In the setting of Definitions~\ref{def:doubling_setting} and~\ref{def:doubling_decay_notation},
let $(\lambda_j)_{j\geq1}$ be a positive sequence and let $m>d$ be an integer. Then $q_m>0$ on
$W(m)$, and:
\begin{enumerate}
    \item if $(\lambda_j)_{j\geq1}$ satisfies \eqref{eq:doubling_cut} at $m$, then
    \begin{equation}\label{eq:doubling_avg}
        \lambda_m\;=\;\sum_{j\in W(m)}q_m(j)\,\lambda_j ;
    \end{equation}
    \item if $(\lambda_j)_{j\geq1}$ satisfies \eqref{eq:doubling_avg} at $m$, then $w_m>0$ on
    $W(m)$ and
    \begin{equation}\label{eq:doubling_uavg}
    \begin{aligned}
        u_m\;&=\;\sum_{j\in W(m)}w_m(j)\,u_j,\\
        \sum_{j\in W(m)}w_m(j)\;&=\;R_0(m) .
    \end{aligned}
    \end{equation}
\end{enumerate}
\end{lemma}

\begin{proof}
Since $m\geq1$, $\varepsilon(m)=c/(m+1)\leq c/2<1$, so $1-\varepsilon(m)>0$ and $q_m(j)>0$ on
$W(m)$, $c$ being positive.

\emph{(1)} Substituting $\varepsilon(j)=c/(j+1)$ in \eqref{eq:doubling_cut} and dividing by
$1-\varepsilon(m)>0$ gives \eqref{eq:doubling_avg}, the indices of $W(m)$ running from
$\lceil m/2\rceil$ to $m-1$.

\emph{(2)} Every $j\in W(m)$ satisfies $j\geq1$, so $w_m(j)=q_m(j)(m/j)^{p}>0$. Multiplying
\eqref{eq:doubling_avg} by $m^{p}>0$ and writing $\lambda_j=u_j\,j^{-p}$ gives
$u_m=\sum_{j\in W(m)}q_m(j)(m/j)^{p}u_j$, which is the first line of \eqref{eq:doubling_uavg}.
The second line is \eqref{eq:doubling_R} at $\alpha=0$, where $\Phi_0(j)=j^{-p}$.
\end{proof}

\begin{lemma}[the expansion of $R_\alpha$]\label{lem:doubling_expansion}
In the setting of Definitions~\ref{def:doubling_setting} and~\ref{def:doubling_decay_notation},
put
\begin{equation}\label{eq:doubling_AB}
\begin{aligned}
    A_0\;&:=\;\frac{p\bigl(p+3+2\tau(p-1)\bigr)}{2(p+1)(\tau-1)},\qquad
    A_1\;:=\;A_0-c\,\tau,\\
    \Gamma\;&:=\;\frac{\tau(p-1)+1}{(\tau-1)(p+1)} .
\end{aligned}
\end{equation}
Then $A_1=p\bigl(p+3-4\tau\bigr)\big/\bigl(2(p+1)(\tau-1)\bigr)$,
\begin{equation}\label{eq:doubling_signs}
    0<\Gamma<1,\qquad A_0>0>A_1 ,
\end{equation}
and for every $\alpha\in\mathbb R$, as $m\rightarrow\infty$,
\begin{equation}\label{eq:doubling_Rexp}
    R_\alpha(m)\;=\;1+\frac{A_{\delta(m)}+\Gamma\alpha}{m}+O\bigl(m^{-2}\bigr),
\end{equation}
with an implied constant depending on $c$ and $\alpha$ alone.
\end{lemma}

\begin{proof}
\emph{Step 1 (partial fractions).} For $j\geq1$ and $\alpha\in\mathbb R$,
$(j+\alpha)/\bigl(j(j+1)\bigr)=\alpha/j+(1-\alpha)/(j+1)$ and
$j^{-p}/(j+1)=j^{-p-1}-j^{-p-2}+j^{-p-2}/(j+1)$, so
\begin{equation}\label{eq:doubling_pf}
\begin{aligned}
    \frac{\Phi_\alpha(j)}{j+1}\;&=\;j^{-p-1}+(\alpha-1)\,j^{-p-2}+e_\alpha(j),\\
    e_\alpha(j)\;&:=\;(1-\alpha)\frac{j^{-p-2}}{j+1},\\
    |e_\alpha(j)|\;&\leq\;|1-\alpha|\,j^{-p-3} .
\end{aligned}
\end{equation}

\emph{Step 2 (a window sum).} Let $r>1$ and let $1\leq a<b$ be integers. Euler--Maclaurin for
$f(t)=t^{-r}$ on $[a,b]$, truncated before its first Bernoulli term, gives
\begin{equation}\label{eq:doubling_em}
\begin{aligned}
    \sum_{j=a}^{b-1}j^{-r}\;&=\;\frac{a^{1-r}-b^{1-r}}{r-1}+\frac{a^{-r}-b^{-r}}{2}+E,\\
    |E|\;&\leq\;\frac{r}{12}\bigl(a^{-r-1}+b^{-r-1}\bigr),
\end{aligned}
\end{equation}
the bound on $E$ because $t\mapsto t^{-r}$ is completely monotone on $(0,+\infty)$, so that the
remainder of the expansion after any number of terms is bounded in modulus by the first
omitted term, here $\tfrac1{12}|f'(b)-f'(a)|=\tfrac{r}{12}\bigl(a^{-r-1}-b^{-r-1}\bigr)$.

\emph{Step 3 (the ratio, and its parity).} Fix $\alpha\in\mathbb R$, let
$m>\max(d,2|\alpha|)$, and put $a:=\lceil m/2\rceil=(m+\delta)/2$ with $\delta:=\delta(m)$.
Summing \eqref{eq:doubling_pf} over $W(m)=\{a,\dots,m-1\}$, applying \eqref{eq:doubling_em}
with the endpoints $a$ and $b:=m$ at $r=p+1$ and at $r=p+2$ and bounding
$\sum_{j=a}^{m-1}|e_\alpha(j)|\leq|1-\alpha|\,m\,a^{-p-3}$,
\begin{equation}
\begin{aligned}
    \sum_{j\in W(m)}\frac{\Phi_\alpha(j)}{j+1}
    \;=\;&\frac{a^{-p}-m^{-p}}{p}+\frac{a^{-p-1}-m^{-p-1}}{2}\\
    &+(\alpha-1)\,\frac{a^{-p-1}-m^{-p-1}}{p+1}\\
    &+O\bigl(m^{-p-2}\bigr),
\end{aligned}
\end{equation}
the implied constants depending on $c$ and $\alpha$ alone in this proof. The ceiling enters
only through $a^{-p}=\tau m^{-p}\bigl(1-p\delta/m+O(m^{-2})\bigr)$ and
$a^{-p-1}=2\tau m^{-p-1}\bigl(1+O(m^{-1})\bigr)$, which turn the display into
$\frac{\tau-1}{p}m^{-p}+\Theta_\alpha(\delta)\,m^{-p-1}+O(m^{-p-2})$ with
$\Theta_\alpha(\delta):=\frac{2\tau-1}{2}-\tau\delta+(\alpha-1)\frac{2\tau-1}{p+1}$.
Multiplying by $c/(1-\varepsilon(m))=c\bigl(1+c/m+O(m^{-2})\bigr)$ and dividing by
$\Phi_\alpha(m)=m^{-p}(1+\alpha/m)$, whose reciprocal is
$m^{p}\bigl(1-\alpha/m+O(m^{-2})\bigr)$, gives by \eqref{eq:doubling_R}
\begin{equation}\label{eq:doubling_Rraw}
\begin{aligned}
    R_\alpha(m)\;=\;&c\,\frac{\tau-1}{p}\\
    &+\frac{1}{m}\left[c\,\Theta_\alpha(\delta)+c\,\frac{\tau-1}{p}\,(c-\alpha)\right]\\
    &+O\bigl(m^{-2}\bigr) .
\end{aligned}
\end{equation}
The leading coefficient of \eqref{eq:doubling_Rraw} is $c(\tau-1)/p$, which equals $1$ by
\eqref{eq:doubling_cramer_form}; the same relation applied to the second term leaves
$c\,\Theta_\alpha(\delta)+c-\alpha$ for the coefficient of $1/m$. Its part in $\alpha$ is
$\bigl(c\frac{2\tau-1}{p+1}-1\bigr)\alpha$, and substituting $c=p/(\tau-1)$,
\begin{equation}
\begin{aligned}
    c\,\frac{2\tau-1}{p+1}-1
    \;&=\;\frac{p(2\tau-1)-(\tau-1)(p+1)}{(\tau-1)(p+1)}\\
    \;&=\;\frac{\tau(p-1)+1}{(\tau-1)(p+1)}\;=\;\Gamma .
\end{aligned}
\end{equation}
Its part free of $\alpha$ is
$c\bigl(\frac{2\tau-1}{2}-\frac{2\tau-1}{p+1}+1\bigr)-c\tau\delta$, and the same substitution,
with $(2\tau-1)(p-1)+2(p+1)=2\tau(p-1)+p+3$, gives
\begin{equation}
\begin{aligned}
    c\left(\frac{2\tau-1}{2}-\frac{2\tau-1}{p+1}+1\right)
    \;&=\;\frac{p}{\tau-1}\cdot\frac{2\tau(p-1)+p+3}{2(p+1)}\\
    \;&=\;A_0 ,
\end{aligned}
\end{equation}
so that part is $A_0$ at $\delta=0$ and $A_0-c\tau=A_1$ at $\delta=1$: this is
\eqref{eq:doubling_Rexp} with the coefficients \eqref{eq:doubling_AB}.

\emph{Step 4 (the closed form of $A_1$, and the signs).} By \eqref{eq:doubling_cramer_form},
$c\tau=p\tau/(\tau-1)=2p\tau(p+1)\big/\bigl(2(p+1)(\tau-1)\bigr)$, and
$p\bigl(p+3+2\tau(p-1)\bigr)-2p\tau(p+1)=p\bigl(p+3-4\tau\bigr)$, so
$A_1=A_0-c\tau=p(p+3-4\tau)\big/\bigl(2(p+1)(\tau-1)\bigr)$. By
Lemma~\ref{lem:doubling_cramer_root}\emph{(3)}, $\tau>2$, $\tau>p+1$, $2\tau>p+2$ and
$4\tau>p+3$, and $p>1$. Hence $(\tau-1)(p+1)>0$ and $\tau(p-1)+1>0$, so $\Gamma>0$;
$(\tau-1)(p+1)-\bigl(\tau(p-1)+1\bigr)=2\tau-p-2>0$, so $\Gamma<1$; $p+3+2\tau(p-1)>0$, so
$A_0>0$; and $p+3-4\tau<0$, so $A_1<0$. This is \eqref{eq:doubling_signs}.
\end{proof}

\begin{remark}[the parity of the window]\label{rem:doubling_parity}
In the notation of Definition~\ref{def:doubling_decay_notation}, and with the coefficients
$A_0$, $A_1$ of Lemma~\ref{lem:doubling_expansion}, the two parities of
\eqref{eq:doubling_Rexp} differ at order $1/m$ by $(A_0-A_1)/m=c\tau/m$, which is half of the
heaviest summand $2c\tau/m+O(m^{-2})$ of \eqref{eq:doubling_R} relative to $\Phi_0(m)$.
Heuristically, the window at odd $m$ starts half a step later than $m/2$, at $(m+1)/2$, and in
a continuum reading of the window sum that half step removes half of that summand. Call a
positive sequence $(x_j)_{j\geq1}$ a \emph{supersolution} of \eqref{eq:doubling_avg} if
$x_m\geq\sum_{j\in W(m)}q_m(j)\,x_j$ for every large enough integer $m$, and a
\emph{subsolution} of it if the reverse inequality holds for every large enough integer $m$.
By \eqref{eq:doubling_Rexp} at $\alpha=0$ and \eqref{eq:doubling_signs}, the pure power
$\Phi_0(j)=j^{-p}$ has $R_0(m)>1$ at every large even $m$ and $R_0(m)<1$ at every large odd
$m$, so it is neither a supersolution nor a subsolution of \eqref{eq:doubling_avg}.
\end{remark}

\begin{lemma}[the descent chain]\label{lem:doubling_descent}
Let $s=1$ and $0<c<1$, in the notation of Definitions~\ref{def:doubling_setting}
and~\ref{def:doubling_decay_notation}. Then there are
$c_4>0$ and an integer $\ell_1>d$, depending on $c$ and $d$ alone, with the following
property. The bounds
\begin{equation}\label{eq:doubling_R0}
    \tfrac12\;\leq\;R_0(m)\;\leq\;2
    \qquad\text{and}\qquad
    \bigl|R_0(m)-1\bigr|\;\leq\;\frac{c_4}{m}
\end{equation}
hold for every $m\geq\ell_1$; and for every integer $\ell\geq\ell_1$, with $(Y_n)_{n\geq0}$ the
Markov chain on the integers $\geq\ell$ with, for every state $y\geq\ell$,
$\mathbb P(Y_{n+1}=j\mid Y_n=y)=w_y(j)/R_0(y)$ for every $j\in W(y)$ when $y\geq2\ell$ --- a
probability on $W(y)$, since $w_y>0$ there and $\sum_{j\in W(y)}w_y(j)=R_0(y)>0$ by
\eqref{eq:doubling_R} at $\alpha=0$ and \eqref{eq:doubling_R0} --- and $Y_{n+1}=Y_n$ when
$y<2\ell$, and with $N_\ell:=\inf\{n\geq0:Y_n<2\ell\}$, the following hold for every starting
state $m\geq2\ell$:
\begin{enumerate}
    \item $\lceil Y_n/2\rceil\leq Y_{n+1}\leq Y_n-1$ for every $n<N_\ell$, and $N_\ell\leq m$;
    \item $Y_{N_\ell}\in[\ell,2\ell)$ almost surely;
    \item if $M_1\in\{2\ell,2\ell+1,\dots\}\cup\{+\infty\}$, if $(\lambda_j)_{j\geq1}$ is a
    positive sequence satisfying \eqref{eq:doubling_cut} at every integer $y$ with
    $d<y\leq M_1$, and if $2\ell\leq m\leq M_1$, then, with
    $Z_\ell:=\prod_{n<N_\ell}R_0(Y_n)$, which is positive,
    \begin{equation}\label{eq:doubling_pathid}
        u_m\;=\;\mathbb E\bigl(\,Z_\ell\,u_{Y_{N_\ell}}\,\mid Y_0=m\bigr) .
    \end{equation}
\end{enumerate}
\end{lemma}

\begin{proof}
\eqref{eq:doubling_Rexp} at $\alpha=0$, with the coefficients \eqref{eq:doubling_AB}, reads
$R_0(m)=1+A_{\delta(m)}/m+O(m^{-2})$ with an implied constant depending on $c$ alone, which
gives the second half of \eqref{eq:doubling_R0} with $c_4$ and $\ell_1$ depending only on $c$
and $d$; enlarging $\ell_1$ past $2c_4$ gives the first half. Every state $y\geq2\ell$ that the
chain leaves satisfies $y\geq2\ell_1>d$ and $\varepsilon(y)=c/(y+1)\leq c/2<1$, so $q_y>0$ on
$W(y)$, $c$ being positive, and $w_y(j)=q_y(j)(y/j)^{p}>0$ there; and
$\sum_{j\in W(y)}w_y(j)=R_0(y)\geq\tfrac12$ by \eqref{eq:doubling_R} at $\alpha=0$ and
\eqref{eq:doubling_R0}: the transition law of the statement is a probability on $W(y)$.

\emph{(1)} $j\in W(y)$ means $\lceil y/2\rceil\leq j\leq y-1$. Hence $Y_n\leq m-n$ for
$n\leq N_\ell$, and $Y_{N_\ell-1}\geq2\ell$ forces $N_\ell\leq m-2\ell+1\leq m$.

\emph{(2)} $N_\ell\geq1$ because $Y_0=m\geq2\ell$, and $Y_{N_\ell}<2\ell$ by definition while
$Y_{N_\ell}\geq\lceil Y_{N_\ell-1}/2\rceil\geq\ell$ by \emph{(1)}.

\emph{(3)} By \emph{(1)} every path is finite and \eqref{eq:doubling_pathid} is a finite sum.
We prove $\mathbb E(Z_\ell u_{Y_{N_\ell}}\mid Y_0=y)=u_y$ for every $y$ with $\ell\leq y\leq m$
by induction on $y$. For $y<2\ell$, $N_\ell=0$, the product is empty and $Y_{N_\ell}=y$. Let
$2\ell\leq y\leq m$. Then $y\geq2\ell_1>d$ and $y\leq m\leq M_1$, so $(\lambda_j)_{j\geq1}$
satisfies \eqref{eq:doubling_cut} at $y$ and Lemma~\ref{lem:doubling_averaging}\emph{(1)} then
\emph{(2)} give \eqref{eq:doubling_uavg} at $y$. Every $j\in W(y)$ satisfies
$\ell\leq\lceil y/2\rceil\leq j<y$, so the induction hypothesis and the Markov property give
\begin{equation}
    \mathbb E\bigl(Z_\ell\,u_{Y_{N_\ell}}\mid Y_0=y\bigr)
    \;=\;\sum_{j\in W(y)}\frac{w_y(j)}{R_0(y)}\,R_0(y)\,u_j
    \;=\;u_y
\end{equation}
by \eqref{eq:doubling_uavg}. Positivity of $Z_\ell$ is the first half of \eqref{eq:doubling_R0}.
\end{proof}

\begin{lemma}[the descent halves at a uniform rate]\label{lem:doubling_escape}
Let $s=1$ and $0<c<1$, in the notation of Definitions~\ref{def:doubling_setting}
and~\ref{def:doubling_decay_notation}, with the integer
$\ell_1$ of Lemma~\ref{lem:doubling_descent}. Put
$\gamma:=\sqrt\tau\big/\bigl(2(\sqrt\tau+1)\bigr)$, so that $0<\gamma<\tfrac12$, and
$c_5:=6/\gamma^2$. Then there is an integer $\ell_2\geq2\ell_1$,
depending on $c$ and $d$ alone, such that for every $\ell\geq\ell_2$ the descent chain
$(Y_n)_{n\geq0}$ of Lemma~\ref{lem:doubling_descent} at the level $\ell$ satisfies
\begin{equation}\label{eq:doubling_escape}
    \mathbb P\bigl(Y_1\leq y/\sqrt2\mid Y_0=y\bigr)\;\geq\;\gamma\qquad\text{for every }y\geq2\ell,
\end{equation}
and, for every $i\geq1$ and every $y\in I_i(\ell)$, the exit time
$\varsigma:=\inf\{n\geq0:Y_n\notin I_i(\ell)\}$ obeys
\begin{equation}\label{eq:doubling_sojourn}
\begin{aligned}
    \mathbb P(\varsigma>2r\mid Y_0=y)\;&\leq\;\bigl(1-\gamma^2\bigr)^{r}\quad(r\geq0),\\
    \mathbb E(\varsigma\mid Y_0=y)\;&\leq\;2/\gamma^2,\\
    \mathbb E\bigl(\zeta^{\varsigma}\mid Y_0=y\bigr)\;&\leq\;1+c_5(\zeta-1)
    \quad\bigl(1\leq\zeta\leq1+\gamma^2/6\bigr).
\end{aligned}
\end{equation}
\end{lemma}

\begin{proof}
\emph{Step 1 (one step halves with probability $\gamma$).} Put $a:=\lceil y/2\rceil$ and
$b:=\lfloor y/\sqrt2\rfloor$. Since $w_y(j)\geq c\,y^{p}/\bigl((j+1)j^{p}\bigr)\geq
c\,y^{p}(j+1)^{-p-1}$ and $t\mapsto t^{-p-1}$ is decreasing,
\begin{equation}
\begin{aligned}
    \sum_{j=a}^{b}w_y(j)\;&\geq\;c\,y^{p}\int_{a+1}^{b+2}t^{-p-1}\,dt\\
    &=\;\frac{c\,y^{p}}{p}\Bigl[(a+1)^{-p}-(b+2)^{-p}\Bigr].
\end{aligned}
\end{equation}
Here $a+1\leq(y/2)(1+4/y)$ and $(1+x)^{-p}\geq1-px$ for $x\geq0$, so
$(a+1)^{-p}\geq\tau y^{-p}(1-4p/y)$, while $b+2\geq y/\sqrt2$ gives
$(b+2)^{-p}\leq\sqrt\tau\,y^{-p}$. With $c(\tau-1)/p=1$, that is
\eqref{eq:doubling_cramer_form},
\begin{equation}
    \sum_{j=a}^{b}w_y(j)\;\geq\;\frac{c(\tau-\sqrt\tau)}{p}-\frac{4c\tau}{y}
    \;=\;2\gamma-\frac{4c\tau}{y} .
\end{equation}
With the constant $c_4$ of Lemma~\ref{lem:doubling_descent}, choose $\ell_2$ so that
$y\geq2\ell_2$ forces $a\leq b\leq y-1$, $4c\tau/y\leq\gamma/2$ and $c_4/y\leq\tfrac12$;
dividing by $R_0(y)\leq1+c_4/y$, which is \eqref{eq:doubling_R0}, leaves at least
$(2\gamma-\gamma/2)/(3/2)=\gamma$, which is \eqref{eq:doubling_escape}.

\emph{Step 2 (the three sojourn bounds).} Fix $i\geq1$, $y\in I_i(\ell)$ and $n\geq0$ with
$Y_n\in I_i(\ell)$, and put $G_n:=\{Y_{n+1}\leq Y_n/\sqrt2\}$. On $G_n\cap G_{n+1}$ one has
$Y_{n+2}\leq Y_n/2<2^i\ell$, and on $G_n\setminus\{Y_{n+1}\in I_i(\ell)\}$ one has
$Y_{n+2}\leq Y_{n+1}<2^i\ell$; since $(Y_n)_n$ is non-increasing, $Y_{n+2}\notin I_i(\ell)$ in both
cases, so
\begin{equation}
    \{Y_{n+2}\in I_i(\ell)\}\;\subseteq\;\complement G_n\cup\bigl(G_n\cap\{Y_{n+1}\in I_i(\ell)\}\cap\complement G_{n+1}\bigr).
\end{equation}
Every state of $I_i(\ell)$ is $\geq2^i\ell\geq2\ell$, so \eqref{eq:doubling_escape} applies at
$Y_{n+1}$ on $\{Y_{n+1}\in I_i(\ell)\}$ and at $Y_n$, and conditioning first on $Y_0,\dots,Y_{n+1}$
and then on $Y_0,\dots,Y_n$,
\begin{equation}
\begin{aligned}
    &\mathbb P\bigl(Y_{n+2}\in I_i(\ell)\mid Y_0,\dots,Y_n\bigr)\\
    &\qquad\leq\;1-\gamma\,\mathbb P\bigl(G_n\mid Y_0,\dots,Y_n\bigr)\\
    &\qquad\leq\;1-\gamma^2 .
\end{aligned}
\end{equation}
The times at which a non-increasing chain lies in $I_i(\ell)$ form an interval, so
$\{\varsigma>2r\}=\{Y_{2r}\in I_i(\ell)\}$ under $\mathbb P(\,\cdot\mid Y_0=y)$, and iterating the last display at
$n=0,2,\dots,2r-2$ gives the first line of \eqref{eq:doubling_sojourn}. The second follows from
$\mathbb E(\varsigma\mid Y_0=y)=\sum_{n\geq0}\mathbb P(\varsigma>n\mid Y_0=y)\leq2\sum_{r\geq0}(1-\gamma^2)^{r}$.
For the third, put $\bar\zeta:=\zeta^2$ and $\widetilde\varsigma:=\lceil\varsigma/2\rceil$, so
that $\zeta^{\varsigma}\leq\bar\zeta^{\,\widetilde\varsigma}$ and
$\mathbb P(\widetilde\varsigma>r\mid Y_0=y)\leq(1-\gamma^2)^{r}$. From
$\bar\zeta^{\,\widetilde\varsigma}=1+(\bar\zeta-1)\sum_{r<\widetilde\varsigma}\bar\zeta^{\,r}$,
\begin{equation}
\begin{aligned}
    \mathbb E\bigl(\bar\zeta^{\,\widetilde\varsigma}\mid Y_0=y\bigr)
    \;&=\;1+(\bar\zeta-1)\sum_{r\geq0}\bar\zeta^{\,r}\,\mathbb P(\widetilde\varsigma>r\mid Y_0=y)\\
    \;&\leq\;1+\frac{\bar\zeta-1}{1-\bar\zeta(1-\gamma^2)} ,
\end{aligned}
\end{equation}
and $\zeta-1\leq\gamma^2/6$ with $\gamma<\tfrac12$ gives
$\bar\zeta\leq1+3(\zeta-1)\leq1+\gamma^2/2$, hence $\bar\zeta(1-\gamma^2)\leq1-\gamma^2/2$ and
$\bar\zeta-1\leq3(\zeta-1)$; the ratio is at most $6(\zeta-1)/\gamma^2$.
\end{proof}

\begin{lemma}[the weight of a descent is $1+O(1/\ell)$ in expectation]\label{lem:doubling_product}
Let $s=1$ and $0<c<1$, in the notation of Definitions~\ref{def:doubling_setting}
and~\ref{def:doubling_decay_notation}, with the
constants $\gamma$, $c_5$ and $\ell_2$ of Lemma~\ref{lem:doubling_escape}. Let
$\ell\geq\ell_2$, let $(Y_n)_{n\geq0}$ and $N_\ell$ be the descent chain of
Lemma~\ref{lem:doubling_descent} at the level $\ell$ and its exit time,
let $m\geq2\ell$, let $H$ be a number with $0\leq H\leq\gamma^2\ell/6$, and let $b$ be a
function on the integers $\geq\ell$ with $|b(y)|\leq H/y$ for every integer $y\geq\ell$. Then
\begin{equation}\label{eq:doubling_product}
    e^{-c_5H/\ell}\;\leq\;
    \mathbb E\Bigl(\prod_{n<N_\ell}\bigl(1+b(Y_n)\bigr)\mid Y_0=m\Bigr)\;\leq\;e^{c_5H/\ell} .
\end{equation}
\end{lemma}

\begin{proof}
Every $Y_n$ with $n<N_\ell$ satisfies $Y_n\geq2\ell$, so $|b(Y_n)|\leq H/(2\ell)\leq\gamma^2/12$
and every factor lies in $[\tfrac12,\tfrac32]$. Write
$\Psi(z):=\mathbb E\bigl(\prod_{n<N_\ell}(1+b(Y_n))\mid Y_0=z\bigr)$, $\varphi_i:=\sup_{z\in I_i(\ell)}\Psi(z)$
and $\underline\varphi_i:=\inf_{z\in I_i(\ell)}\Psi(z)$, so that
$\varphi_0=\underline\varphi_0=1$, the product being empty on $I_0(\ell)=[\ell,2\ell)$.

Fix $i\geq1$ and $z\in I_i(\ell)$, and let $\varsigma$ be the exit time from $I_i(\ell)$ of
Lemma~\ref{lem:doubling_escape}. By Lemma~\ref{lem:doubling_descent}\emph{(1)},
$Y_\varsigma\geq\lceil Y_{\varsigma-1}/2\rceil\geq2^{i-1}\ell$ and $Y_\varsigma<2^{i}\ell$, so
$Y_\varsigma\in I_{i-1}(\ell)$; and $\varsigma\leq N_\ell$, every state of $I_i(\ell)$ being $\geq2\ell$. The
factors being positive, the strong Markov property at $\varsigma$ gives
\begin{equation}
    \Psi(z)\;=\;\mathbb E\Bigl(\prod_{n<\varsigma}\bigl(1+b(Y_n)\bigr)\;\Psi(Y_\varsigma)\mid Y_0=z\Bigr),
\end{equation}
and $|b|\leq H2^{-i}/\ell$ on $I_i(\ell)$, so with $\zeta_i:=1+H2^{-i}/\ell\leq1+\gamma^2/6$,
\begin{equation}
    \varphi_i\;\leq\;\varphi_{i-1}\,\mathbb E\bigl(\zeta_i^{\,\varsigma}\mid Y_0=z\bigr)
    \;\leq\;\varphi_{i-1}\Bigl(1+\frac{c_5H2^{-i}}{\ell}\Bigr)
\end{equation}
by \eqref{eq:doubling_sojourn}. Hence $\varphi_i\leq\exp\bigl(c_5H\ell^{-1}\sum_{r\geq1}2^{-r}\bigr)
=e^{c_5H/\ell}$ for every $i$. In the other direction, $(1-x)^{\varsigma}\geq1-\varsigma x$
for $x\in[0,1]$ and \eqref{eq:doubling_sojourn} give
$\mathbb E\bigl((1-H2^{-i}/\ell)^{\varsigma}\mid Y_0=z\bigr)\geq1-2H2^{-i}/(\gamma^2\ell)\geq\tfrac56$,
so $\underline\varphi_i\geq\underline\varphi_{i-1}\bigl(1-2H2^{-i}/(\gamma^2\ell)\bigr)$ and,
by $1-x\geq e^{-2x}$ on $[0,\tfrac12]$, $\underline\varphi_i\geq e^{-4H/(\gamma^2\ell)}\geq
e^{-c_5H/\ell}$. The starting state $m\geq2\ell$ lies in $I_i(\ell)$ for some $i\geq1$, so
$\mathbb E\bigl(\prod_{n<N_\ell}(1+b(Y_n))\mid Y_0=m\bigr)=\Psi(m)$ lies between $\underline\varphi_i$
and $\varphi_i$: this is \eqref{eq:doubling_product}.
\end{proof}

\begin{theorem}[the invariant measure decays at the Cram\'er exponent]\label{theo:doubling_decay}
In the setting of Definitions~\ref{def:doubling_setting} and~\ref{def:doubling_decay_notation},
with the constants $c_4$ and $\ell_1$ of Lemma~\ref{lem:doubling_descent} and $\gamma$, $c_5$
and $\ell_2$ of Lemma~\ref{lem:doubling_escape}, put
$m_0:=\max\bigl(\ell_2,\lceil6c_4/\gamma^2\rceil\bigr)$, an integer depending on $c$ and $d$
alone and larger than $d$, since $\ell_2\geq2\ell_1$ and $\ell_1>d$. Then:
\begin{enumerate}
    \item for every integer $\ell\geq m_0$, every
    $M_1\in\{2\ell,2\ell+1,\dots\}\cup\{+\infty\}$, every positive sequence
    $(\lambda_j)_{j\geq1}$ satisfying \eqref{eq:doubling_cut} at every integer $y$ with
    $d<y\leq M_1$, and every integer $m$ with $2\ell\leq m\leq M_1$,
    \begin{equation}\label{eq:doubling_block}
        e^{-c_5c_4/\ell}\,\min_{\ell\leq j<2\ell}u_j\;\leq\;u_m\;\leq\;
        e^{c_5c_4/\ell}\,\max_{\ell\leq j<2\ell}u_j ,
    \end{equation}
    where $u_j=\lambda_j\,j^{p}$ is the rescaled profile of
    Definition~\ref{def:doubling_decay_notation};
    \item for every positive sequence $(\lambda_j)_{j\geq1}$ satisfying \eqref{eq:doubling_cut}
    at every integer $m>d$ there are constants $0<c_1\leq c_2<+\infty$, depending on $c$, $d$
    and $\lambda_1,\dots,\lambda_{2m_0}$ alone, with
    \begin{equation}\label{eq:doubling_decay}
        c_1\,j^{-p_*}\;\leq\;\lambda_j\;\leq\;c_2\,j^{-p_*}\qquad\text{for every }j\geq1 .
    \end{equation}
\end{enumerate}
\end{theorem}

\begin{proof}
\begin{itemize}
    \item \emph{One block of the profile bounds every index above it, within
    $e^{\pm c_5c_4/\ell}$.} Let $\ell$, $M_1$, $(\lambda_j)_{j\geq1}$ and $m$ be as in
    \emph{(1)}. Since $\ell\geq m_0\geq\ell_2\geq\ell_1$, the descent chain $(Y_n)_{n\geq0}$ of
    Lemma~\ref{lem:doubling_descent} at the level $\ell$ and its exit time $N_\ell$ are
    defined, and item \emph{(3)} of that lemma applies at the range $M_1$ and the starting
    state $m$: it gives \eqref{eq:doubling_pathid},
    $u_m=\mathbb E\bigl(Z_\ell\,u_{Y_{N_\ell}}\mid Y_0=m\bigr)$ with
    $Z_\ell=\prod_{n<N_\ell}R_0(Y_n)$ positive, while item \emph{(2)} of that lemma gives
    $Y_{N_\ell}\in[\ell,2\ell)$ almost surely. The extrema of $u$ over the finite set
    $\{\ell,\dots,2\ell-1\}$ are attained and strictly positive, every $\lambda_j$ being
    positive, so
    \begin{multline}
        \Bigl(\min_{\ell\leq j<2\ell}u_j\Bigr)\,\mathbb E\bigl(Z_\ell\mid Y_0=m\bigr)
        \;\leq\;u_m\\
        \;\leq\;\Bigl(\max_{\ell\leq j<2\ell}u_j\Bigr)\,
        \mathbb E\bigl(Z_\ell\mid Y_0=m\bigr).
    \end{multline}
    Apply Lemma~\ref{lem:doubling_product} at $b:=R_0-1$ and $H:=c_4$. Its hypotheses hold:
    $\ell\geq\ell_2$; $\bigl|R_0(y)-1\bigr|\leq c_4/y$ at every integer $y\geq\ell\geq\ell_1$
    by \eqref{eq:doubling_R0}; $0\leq c_4\leq\gamma^2\ell/6$ by
    $\ell\geq m_0\geq6c_4/\gamma^2$; and $m\geq2\ell$. Its product
    $\prod_{n<N_\ell}\bigl(1+b(Y_n)\bigr)$ is $Z_\ell$, so \eqref{eq:doubling_product} reads
    $e^{-c_5c_4/\ell}\leq\mathbb E\bigl(Z_\ell\mid Y_0=m\bigr)\leq e^{c_5c_4/\ell}$. The
    extrema being positive, the last display then gives \eqref{eq:doubling_block}.
    \item \emph{The rescaled profile is bounded above and below by positive constants.} Let
    $(\lambda_j)_{j\geq1}$ be as in \emph{(2)} and apply \emph{(1)} at $\ell=m_0$ and
    $M_1=+\infty$: every integer $m\geq2m_0$ satisfies
    \begin{equation}
    \begin{aligned}
        e^{-c_5c_4/m_0}\min_{m_0\leq j<2m_0}u_j\;&\leq\;u_m\\
        \;&\leq\;e^{c_5c_4/m_0}\max_{m_0\leq j<2m_0}u_j .
    \end{aligned}
    \end{equation}
    Put $c_1:=e^{-c_5c_4/m_0}\min_{1\leq j<2m_0}u_j$ and
    $c_2:=e^{c_5c_4/m_0}\max_{1\leq j<2m_0}u_j$, attained over a finite set of indices and
    finite and strictly positive, every $\lambda_j$ being positive. The extrema of the last
    display run over a subset of $\{1,\dots,2m_0-1\}$, so $c_1\leq u_m\leq c_2$ at every
    integer $m\geq2m_0$; and $e^{-c_5c_4/m_0}\leq1\leq e^{c_5c_4/m_0}$ gives
    $c_1\leq u_j\leq c_2$ at every $j$ with $1\leq j<2m_0$. Since $u_j=\lambda_j\,j^{p}$ and
    $p=p_*$, this is \eqref{eq:doubling_decay}; and $c_1\leq c_2$, both bounding $u_1$. The
    integer $m_0$ depends on $c$ and $d$ alone, $\ell_2$, $c_4$ and $\gamma$ depending on $c$
    and $d$ alone by Lemmas~\ref{lem:doubling_descent} and~\ref{lem:doubling_escape}, and
    $c_1$ and $c_2$ depend on $c$, $d$ and $\lambda_1,\dots,\lambda_{2m_0-1}$, hence on $c$,
    $d$ and $\lambda_1,\dots,\lambda_{2m_0}$.
\end{itemize}
\end{proof}

\begin{corollary}[the tail, and its ratio to the mass at the cut]\label{cor:doubling_tail}
Let $s=1$ and $0<c<1$, let $(\lambda_j)_{j\geq1}$ be a positive sequence satisfying
\eqref{eq:doubling_cut} for every $m>d$, in the notation of
Definitions~\ref{def:doubling_setting} and~\ref{def:doubling_decay_notation}, let
$0<c_1\leq c_2$ be constants with \eqref{eq:doubling_decay} and let $L$ be the tail of
Definition~\ref{def:doubling_setting}. Then there are numbers $\xi_m$, $m\geq1$, with
\begin{equation}\label{eq:doubling_powertail}
    \sum_{j\geq m}j^{-p_*}\;=\;\frac{m^{1-p_*}}{p_*-1}\,\bigl(1+\xi_m\bigr),
    \qquad 0\;\leq\;\xi_m\;\leq\;\frac{p_*-1}{m} ,
\end{equation}
and for every $m\geq1$
\begin{equation}\label{eq:doubling_tailratio}
\begin{aligned}
    \frac{c_1}{p_*-1}\,m^{1-p_*}\;&\leq\;L(m)\;\leq\;\frac{c_2\,p_*}{p_*-1}\,m^{1-p_*},\\
    \frac{c_1}{c_2(p_*-1)}\,m\;&\leq\;\frac{L(m)}{\lambda_m}\;\leq\;\frac{c_2\,p_*}{c_1(p_*-1)}\,m ;
\end{aligned}
\end{equation}
in particular $L(m)\rightarrow0$ and $L(m)/\lambda_m\rightarrow+\infty$.
\end{corollary}

\begin{proof}
Since $p=p_*>1$ (Definition~\ref{def:doubling_decay_notation}), $t\mapsto t^{-p}$
is decreasing and integrable at $+\infty$ with $\int_m^{+\infty}t^{-p}\,dt=m^{1-p}/(p-1)$, and
comparing the sum with that integral,
\begin{equation}
\begin{aligned}
    \frac{m^{1-p}}{p-1}\;\leq\;\sum_{j\geq m}j^{-p}
    \;&\leq\;m^{-p}+\frac{m^{1-p}}{p-1}\\
    \;&=\;\frac{m^{1-p}}{p-1}\left(1+\frac{p-1}{m}\right) ,
\end{aligned}
\end{equation}
which is \eqref{eq:doubling_powertail} with
$\xi_m:=(p-1)\,m^{p-1}\sum_{j\geq m}j^{-p}-1$; and $m\geq1$ gives $(p-1)/m\leq p-1$, hence
$1\leq1+\xi_m\leq p$. Summing \eqref{eq:doubling_decay} over $j\geq m$ and inserting
\eqref{eq:doubling_powertail} with those two bounds on $1+\xi_m$ gives the first pair of
\eqref{eq:doubling_tailratio}, and dividing it by \eqref{eq:doubling_decay} at $j=m$ gives the
second. Finally $p>1$ sends the upper bound $c_2\,p\,m^{1-p}/(p_*-1)$ of $L(m)$ to $0$ and the
lower bound $c_1m/\bigl(c_2(p_*-1)\bigr)$ of $L(m)/\lambda_m$ to $+\infty$.
\end{proof}

\subsubsection{The sharp asymptotic and its constant}

This block proves that the constant Theorem~\ref{theo:doubling_decay} leaves anywhere between
its two bounds exists, that the convergence to it carries a rate, and that it is a linear
functional of the boundary data; its existence and its determination by that data are item
\emph{(2)} of Theorem~\ref{theo:doubling_main}. Theorem~\ref{theo:doubling_sharp} is the sharp
asymptotic, with the rate. Lemma~\ref{lem:doubling_weight} bounds in mean the deviation from
$1$ of the weight a descent accumulates. Lemma~\ref{lem:doubling_doeblin} bounds below the law
in which a descent leaves a dyadic block, by a fixed fraction of one reference law on the block
beneath it, and Lemma~\ref{lem:doubling_coupling} bounds the total variation between the laws
in which two descents started in one block reach the base block, geometrically in the index of
that block. Remark~\ref{rem:doubling_no_logperiodic} records that the decay carries no
log-periodic term, Corollary~\ref{cor:doubling_tail_sharp} restates the tail and the ratio of
the tail to the mass at the cut with that constant, in place of the two-sided bounds of
Corollary~\ref{cor:doubling_tail}, Proposition~\ref{prop:doubling_constant} identifies the
constant as a non-negative linear functional of the boundary data, and
Remark~\ref{rem:doubling_second_order} records a second-order expansion that is formal and is
not proved here.

\begin{theorem}[the sharp asymptotic]\label{theo:doubling_sharp}
Let $s=1$ and $0<c<1$, in the notation of Definitions~\ref{def:doubling_setting}
and~\ref{def:doubling_decay_notation}, with the integer $m_0$ of
Theorem~\ref{theo:doubling_decay}, and let $(\lambda_j)_{j\geq1}$ be any positive sequence
satisfying \eqref{eq:doubling_cut} for every $m>d$. Then
$\lim_{j\rightarrow\infty}\lambda_j\,j^{p_*}$ exists and lies in $[c_1',c_2']$ for every pair
of constants $0<c_1'\leq c_2'$ for which \eqref{eq:doubling_decay} holds with $c_1'$ and $c_2'$
in place of $c_1$ and $c_2$; writing $C$ for that limit,
\begin{equation}\label{eq:doubling_sharp}
    \lambda_j\;=\;C\,j^{-p_*}\bigl(1+o(1)\bigr)\qquad(j\rightarrow\infty),
\end{equation}
$C$ is determined by $c$, $d$ and $\lambda_1,\dots,\lambda_d$, and there are
$\vartheta\in(0,1)$, depending on $c$ alone, and $c_6>0$, depending on $c$, $d$ and
$\lambda_1,\dots,\lambda_{2m_0}$ alone, with
\begin{equation}\label{eq:doubling_rate}
    \bigl|\lambda_j\,j^{p_*}-C\bigr|\;\leq\;c_6\,j^{-\vartheta}
    \qquad\text{for every }j\geq1 .
\end{equation}
\end{theorem}

\begin{lemma}[the weight of a descent deviates from $1$ by $O(1/\ell)$ in mean]\label{lem:doubling_weight}
In the setting of Definitions~\ref{def:doubling_setting} and~\ref{def:doubling_decay_notation},
with the constant $c_4$ of Lemma~\ref{lem:doubling_descent} and $\gamma$, $c_5$ and $\ell_2$
of Lemma~\ref{lem:doubling_escape}, put $c_3:=8c_5c_4$ and
$\ell_3:=\max\bigl(\ell_2,\lceil2c_4\rceil,\lceil24c_4/\gamma^2\rceil,\lceil4c_4c_5\rceil\bigr)$,
which depend on $c$ and $d$ alone. Let $\ell\geq\ell_3$ be an integer, let $(Y_n)_{n\geq0}$ and
$N_\ell$ be the descent chain of Lemma~\ref{lem:doubling_descent} at the level $\ell$ and its
exit time, and let $Z_\ell:=\prod_{n<N_\ell}R_0(Y_n)$ be the weight of
Lemma~\ref{lem:doubling_descent}\emph{(3)}. Then, for every integer $m\geq2\ell$,
\begin{equation}\label{eq:doubling_weight}
    \mathbb E\bigl(|Z_\ell-1|\mid Y_0=m\bigr)\;\leq\;\frac{c_3}{\ell} .
\end{equation}
\end{lemma}

\begin{proof}
Put $\widetilde b(y):=e^{2c_4/y}-1$ for an integer $y\geq\ell$. By
Lemma~\ref{lem:doubling_descent}\emph{(1)}, $N_\ell\leq m$, so the sums and products below are
finite.

Every state $Y_n$ with $n<N_\ell$ satisfies $Y_n\geq2\ell$ by the definition of $N_\ell$, and
$2\ell\geq\ell\geq\ell_2\geq2\ell_1$, with $\ell_1$ of Lemma~\ref{lem:doubling_descent}, so
\eqref{eq:doubling_R0} gives $1-c_4/Y_n\leq R_0(Y_n)\leq1+c_4/Y_n$, with
$c_4/Y_n\leq c_4/(2\ell)\leq\tfrac14$ because
$\ell\geq2c_4$. Since $1-x\geq e^{-2x}$ on $[0,\tfrac12]$ and $1+x\leq e^{x}$,
\begin{equation}
    \exp\Bigl(-2\sum_{n<N_\ell}\frac{c_4}{Y_n}\Bigr)\;\leq\;Z_\ell\;\leq\;
    \exp\Bigl(\sum_{n<N_\ell}\frac{c_4}{Y_n}\Bigr),
\end{equation}
hence $|Z_\ell-1|\leq\exp\bigl(2\sum_{n<N_\ell}c_4/Y_n\bigr)-1$, both $e^{x}-1$ and $1-e^{-2x}$
being at most $e^{2x}-1$ for $x\geq0$.

That last bound is $\prod_{n<N_\ell}\bigl(1+\widetilde b(Y_n)\bigr)-1$, and
$0\leq\widetilde b(y)\leq4c_4/y$ at every integer $y\geq\ell$: there $2c_4/y\leq2c_4/\ell\leq1$
because $\ell\geq2c_4$, and $e^{x}-1\leq2x$ on $[0,1]$. Lemma~\ref{lem:doubling_product}
applies at $b:=\widetilde b$ and $H:=4c_4$ --- $\ell\geq\ell_2$ and $m\geq2\ell$,
$|\widetilde b(y)|\leq H/y$ at every integer $y\geq\ell$, and $H\leq\gamma^2\ell/6$ because
$\ell\geq24c_4/\gamma^2$ --- and \eqref{eq:doubling_product} bounds
$\mathbb E\bigl(\prod_{n<N_\ell}(1+\widetilde b(Y_n))\mid Y_0=m\bigr)$ by $e^{4c_5c_4/\ell}$.
Finally $4c_5c_4/\ell\leq1$ because $\ell\geq4c_4c_5$, so $e^{4c_5c_4/\ell}-1\leq8c_5c_4/\ell$
by $e^{x}-1\leq2x$ on $[0,1]$ once more: this is \eqref{eq:doubling_weight}.
\end{proof}

\begin{lemma}[the exit from a block charges a fixed fraction of one reference law]\label{lem:doubling_doeblin}
In the setting of Definitions~\ref{def:doubling_setting} and~\ref{def:doubling_decay_notation},
with the integer $\ell_3$ of Lemma~\ref{lem:doubling_weight}, put
$\omega:=\tfrac{c^2}{16}\ln(5/4)\ln(8/5)$, which lies in $(0,1)$ and depends on $c$ alone. Then
there is an integer $\ell_4\geq\ell_3$, depending on $c$ and $d$ alone, with the following
property. Let $\ell\geq\ell_4$ and $i\geq1$ be integers, write $a:=2^{i-1}\ell$, so that
$I_{i-1}(\ell)=[a,2a)$ and $I_i(\ell)=[2a,4a)$, and let $\varpi_i$ be the probability on the
integers of $[\lceil5a/4\rceil,2a)$ proportional to $j\mapsto1/(j+1)$. Let $(Y_n)_{n\geq0}$ be
the descent chain of Lemma~\ref{lem:doubling_descent} at the level $\ell$ and let
$\varsigma:=\inf\{n\geq0:Y_n\notin I_i(\ell)\}$ be its exit time from $I_i(\ell)$, the exit time
of Lemma~\ref{lem:doubling_escape}. Then, for every state $y\in I_i(\ell)$, $\varsigma\leq y$
and $Y_\varsigma\in I_{i-1}(\ell)$ almost surely under $\mathbb P(\,\cdot\mid Y_0=y)$, and
\begin{equation}\label{eq:doubling_doeblin}
    \mathbb P\bigl(Y_\varsigma=j\mid Y_0=y\bigr)\;\geq\;\omega\,\varpi_i(j)
    \qquad\text{for every integer }j .
\end{equation}
\end{lemma}

\begin{proof}
Put $\ell_4:=\max\bigl(\ell_3,\lceil4/\ln(5/4)\rceil\bigr)$, an integer at least $\ell_3$ and,
$\ell_3$ being one, depending on $c$ and $d$ alone. Fix integers $\ell\geq\ell_4$ and $i\geq1$,
put $a:=2^{i-1}\ell$, so that $a\geq\ell\geq\ell_4\geq4/\ln(5/4)>4$, $\ln(5/4)$ being
less than $1$, and hence $a\geq3$ and
$2/a\leq\tfrac12\ln(5/4)\leq\tfrac12\ln(8/5)$, and write
$J:=\{\lceil5a/4\rceil,\dots,2a-1\}$ for the range carrying $\varpi_i$ and
$J':=\{2a,\dots,\lceil5a/2\rceil-1\}$. Both ranges are non-empty, since
$\lceil5a/4\rceil\leq5a/4+1\leq2a-1$ and $2a+1\leq5a/2\leq\lceil5a/2\rceil$ at $a\geq3$;
$J\subseteq[a,2a)=I_{i-1}(\ell)$, and $J'\subseteq[2a,4a)=I_i(\ell)$ because
$\lceil5a/2\rceil-1\leq(5a-1)/2<4a$. Comparing each sum with an integral of $dt/t$ over the
matching interval, and using $\lceil5a/4\rceil+1\leq5a/4+2$ and $\lceil5a/2\rceil\geq5a/2$,
\begin{equation}\label{eq:doubling_windowsums}
\begin{aligned}
    \sum_{j\in J}\frac1{j+1}\;&\geq\;\int_{\lceil5a/4\rceil+1}^{2a+1}\frac{dt}t
    \;\geq\;\ln\frac{2a}{5a/4+2}\\
    \;&=\;\ln\tfrac85-\ln\Bigl(1+\frac8{5a}\Bigr)\;\geq\;\ln\tfrac85-\frac2a,\\
    \sum_{j\in J'}\frac1{j+1}\;&\geq\;\int_{2a+1}^{\lceil5a/2\rceil+1}\frac{dt}t
    \;\geq\;\ln\frac{5a/2}{2a+1}\\
    \;&=\;\ln\tfrac54-\ln\Bigl(1+\frac1{2a}\Bigr)\;\geq\;\ln\tfrac54-\frac2a,
\end{aligned}
\end{equation}
so the two sums are at least $\tfrac12\ln\tfrac85$ and $\tfrac12\ln\tfrac54$.
\begin{itemize}
    \item \emph{The exit is into $I_{i-1}(\ell)$, and $\varsigma\leq y$.} Every state of
    $I_i(\ell)$ is $\geq2^{i}\ell\geq2\ell$, so $n<N_\ell$ for every $n<\varsigma$, whence
    $\varsigma\leq N_\ell\leq y$ by Lemma~\ref{lem:doubling_descent}\emph{(1)}; and
    $\varsigma\geq1$ because $Y_0=y\in I_i(\ell)$. Lemma~\ref{lem:doubling_descent}\emph{(1)}
    at $n=\varsigma-1$ gives $Y_\varsigma\geq\lceil Y_{\varsigma-1}/2\rceil\geq a$, the state
    $Y_{\varsigma-1}$ lying in $I_i(\ell)$ and being therefore $\geq2a$, while
    $Y_\varsigma\leq Y_0<4a$ and $Y_\varsigma\notin[2a,4a)$ force $Y_\varsigma<2a$.
    \item \emph{One step from a state $z\geq2\ell$ gives every $j\in W(z)$ mass at least
    $c/\bigl(2(j+1)\bigr)$.} By Lemma~\ref{lem:doubling_descent} that mass is $w_z(j)/R_0(z)$,
    where $w_z(j)=q_z(j)(z/j)^{p}$ with $(z/j)^{p}\geq1$ because $j\leq z-1$,
    $q_z(j)=c/\bigl((j+1)(1-\varepsilon(z))\bigr)\geq c/(j+1)$ because $0<\varepsilon(z)<1$,
    and $R_0(z)\leq2$ by \eqref{eq:doubling_R0}, $z$ being $\geq2\ell\geq\ell_2\geq\ell_1$,
    with $\ell_1$ of Lemma~\ref{lem:doubling_descent} and $\ell_2$ of
    Lemma~\ref{lem:doubling_escape}.
    \item \emph{Every $y\in I_i(\ell)$ with $y<\lceil5a/2\rceil$ satisfies
    $\mathbb P(Y_\varsigma=j\mid Y_0=y)\geq\tfrac c4\ln\tfrac85\,\varpi_i(j)$ for every integer
    $j$.} Such a $y$ has $J\subseteq W(y)$: $y\leq\lceil5a/2\rceil-1\leq(5a-1)/2$ gives
    $y/2\leq(5a-1)/4$ and hence $\lceil y/2\rceil\leq\lceil5a/4\rceil$, while $2a-1\leq y-1$
    because $y\geq2a$. Since $J\subseteq I_{i-1}(\ell)$, a first step into $J$ is the exit, so
    $\varsigma=1$ and $Y_\varsigma=Y_1$ on $\{Y_1\in J\}$, and the one-step bound gives
    $\mathbb P(Y_\varsigma=j\mid Y_0=y)\geq c/\bigl(2(j+1)\bigr)$ for $j\in J$. On $J$,
    $1/(j+1)=\varpi_i(j)\sum_{j'\in J}1/(j'+1)\geq\tfrac12\ln\tfrac85\,\varpi_i(j)$ by
    \eqref{eq:doubling_windowsums}; off $J$, $\varpi_i(j)=0$.
    \item \emph{Every $y\in I_i(\ell)$ with $y\geq\lceil5a/2\rceil$ satisfies
    $\mathbb P(Y_\varsigma=j\mid Y_0=y)\geq\omega\,\varpi_i(j)$ for every integer $j$.} Such a
    $y$ has $J'\subseteq W(y)$: the largest element $\lceil5a/2\rceil-1$ of $J'$ is at most
    $y-1$, and $\lceil y/2\rceil\leq2a$ because $y\leq4a-1$. The one-step bound and
    \eqref{eq:doubling_windowsums} give
    $\mathbb P(Y_1\in J'\mid Y_0=y)\geq\tfrac c2\sum_{j\in J'}1/(j+1)\geq\tfrac c4\ln\tfrac54$.
    Every state $z$ of $J'$ lies in $I_i(\ell)$ and is $<\lceil5a/2\rceil$, so $\varsigma>1$ on
    $\{Y_1\in J'\}$, the exit time from $I_i(\ell)$ of the chain restarted at time $1$ is
    $\varsigma-1$, and the case of a starting state below $\lceil5a/2\rceil$ applies at $z$; by
    the Markov property at time $1$,
    \begin{equation}
    \begin{aligned}
        \mathbb P\bigl(Y_\varsigma=j\mid Y_0=y\bigr)
        \;&\geq\;\sum_{z\in J'}\mathbb P\bigl(Y_1=z\mid Y_0=y\bigr)\\
        &\qquad\times\;\mathbb P\bigl(Y_\varsigma=j\mid Y_0=z\bigr)\\
        \;&\geq\;\tfrac c4\ln\tfrac54\cdot\tfrac c4\ln\tfrac85\,\varpi_i(j)\\
        \;&=\;\omega\,\varpi_i(j).
    \end{aligned}
    \end{equation}
\end{itemize}
The two cases cover $I_i(\ell)$, and $\tfrac c4\ln\tfrac85\geq\omega$ because
$\tfrac c4\ln\tfrac54\leq1$: this is \eqref{eq:doubling_doeblin}. Finally $\omega>0$, and
$\omega<1$ because $c<1$ and $\ln\tfrac54\ln\tfrac85<1$; $\omega$ depends on $c$ alone.
\end{proof}

\begin{lemma}[two descents from one block reach the base block with nearly the same law]\label{lem:doubling_coupling}
In the setting of Definitions~\ref{def:doubling_setting} and~\ref{def:doubling_decay_notation},
with $\omega$ and the integer $\ell_4$ of Lemma~\ref{lem:doubling_doeblin}, let
$\ell\geq\ell_4$ be an integer and let $(Y_n)_{n\geq0}$ and $N_\ell$ be the descent chain of
Lemma~\ref{lem:doubling_descent} at the level $\ell$ and its exit time. Then, for every integer
$i\geq0$ and every pair of states $y,y'\in I_i(\ell)$,
\begin{equation}\label{eq:doubling_coupling}
\begin{split}
    \TV\Bigl(&\mathbb P\bigl(Y_{N_\ell}\in\cdot\mid Y_0=y\bigr)\,\Big\|\,
    \mathbb P\bigl(Y_{N_\ell}\in\cdot\mid Y_0=y'\bigr)\Bigr)\\
    &\qquad\leq\;(1-\omega)^{i},
\end{split}
\end{equation}
the total variation carrying the $\tfrac12$-convention of \S\ref{app:notation}.
\end{lemma}

\begin{proof}
For a state $z\geq\ell$ put $\pi_z:=\mathbb P(Y_{N_\ell}\in\cdot\mid Y_0=z)$; for an integer
$i\geq1$ let $\varpi_i$ be the reference law of Lemma~\ref{lem:doubling_doeblin} at the level
$\ell$ and $\varsigma$ the exit time from $I_i(\ell)$ of that lemma, and for a state
$y\in I_i(\ell)$ put $\mu_y:=\mathbb P(Y_\varsigma\in\cdot\mid Y_0=y)$ and
$\widetilde\mu_y:=(\mu_y-\omega\varpi_i)/(1-\omega)$, which is legitimate since $\omega<1$. The
proof is by induction on $i$, the claim at $i$ being that
$\TV(\pi_y\|\pi_{y'})\leq(1-\omega)^{i}$ for every pair $y,y'\in I_i(\ell)$.
\begin{itemize}
    \item \emph{The claim holds at $i=0$.} A state $y\in I_0(\ell)=[\ell,2\ell)$ has
    $N_\ell=0$ and $Y_{N_\ell}=y$, so $\pi_y$ and $\pi_{y'}$ are the Dirac masses at $y$ and
    at $y'$, whose total variation is $0$ when $y=y'$ and $1$ when $y\neq y'$: in both cases
    at most $1=(1-\omega)^{0}$.
    \item \emph{At $i\geq1$, $\widetilde\mu_y$ is a probability on $I_{i-1}(\ell)$ and
    $\mu_y=\omega\varpi_i+(1-\omega)\widetilde\mu_y$.} By
    Lemma~\ref{lem:doubling_doeblin}, $\mu_y$ is a probability carried by $I_{i-1}(\ell)$ with
    $\mu_y\geq\omega\varpi_i$, so $\widetilde\mu_y$ is non-negative by
    \eqref{eq:doubling_doeblin}, is carried by $I_{i-1}(\ell)$, $\varpi_i$ being carried by
    that block too, and has total mass $(1-\omega)/(1-\omega)=1$; the identity is the
    definition of $\widetilde\mu_y$.
    \item \emph{At $i\geq1$, $\pi_y=\sum_{z\in I_{i-1}(\ell)}\mu_y(z)\,\pi_z$.} The exit time
    $\varsigma$ is a stopping time with $\varsigma\leq N_\ell$, every state of $I_i(\ell)$
    being $\geq2\ell$, and $Y_\varsigma\in I_{i-1}(\ell)$ by
    Lemma~\ref{lem:doubling_doeblin}; the descent chain at the level $\ell$ is
    time-homogeneous, so the strong Markov property at $\varsigma$ gives the identity.
    \item \emph{The claim at $i-1$ implies the claim at $i$, for $i\geq1$.} Subtracting the
    two mixture identities $\pi_y=\sum_{z\in I_{i-1}(\ell)}\mu_y(z)\,\pi_z$ and
    $\pi_{y'}=\sum_{z\in I_{i-1}(\ell)}\mu_{y'}(z)\,\pi_z$ and cancelling their common part
    $\omega\sum_z\varpi_i(z)\pi_z$ leaves
    $\pi_y-\pi_{y'}=(1-\omega)\bigl(\sum_z\widetilde\mu_y(z)\pi_z
    -\sum_z\widetilde\mu_{y'}(z)\pi_z\bigr)$, and writing both sums against the product
    probability $\widetilde\mu_y\otimes\widetilde\mu_{y'}$ on $I_{i-1}(\ell)^2$,
    \begin{equation}
        \pi_y-\pi_{y'}\;=\;(1-\omega)\!\!\sum_{z,z'\in I_{i-1}(\ell)}\!\!
        \widetilde\mu_y(z)\,\widetilde\mu_{y'}(z')\,\bigl(\pi_z-\pi_{z'}\bigr) .
    \end{equation}
    The block $I_{i-1}(\ell)$ is finite and the coefficients sum to $1$, so the triangle
    inequality for the total-variation norm bounds $\TV(\pi_y\|\pi_{y'})$ by
    $(1-\omega)\max_{z,z'\in I_{i-1}(\ell)}\TV(\pi_z\|\pi_{z'})$, which the claim at $i-1$
    bounds by $(1-\omega)^{i}$.
\end{itemize}
This is \eqref{eq:doubling_coupling}.
\end{proof}

\begin{proof}[Proof of Theorem \ref{theo:doubling_sharp}]
Let $0<c_1\leq c_2$ be the constants of Theorem~\ref{theo:doubling_decay}\emph{(2)}, so that
$c_1\leq u_m\leq c_2$ at every integer $m\geq1$, $u_m=\lambda_m\,m^{p}$ being the rescaled
profile of Definition~\ref{def:doubling_decay_notation} and $p$ being $p_*$; let $c_3$ and
$\ell_3$ be the constants of Lemma~\ref{lem:doubling_weight}, and $\omega$ and $\ell_4$ those
of Lemma~\ref{lem:doubling_doeblin}. For an integer level $\ell$ at least the $\ell_1$ of
Lemma~\ref{lem:doubling_descent}, write $(Y_n)_{n\geq0}$, $N_\ell$ and
$Z_\ell:=\prod_{n<N_\ell}R_0(Y_n)$ for the descent chain of that lemma at the level $\ell$, its
exit time and its weight. Every use of Lemma~\ref{lem:doubling_descent}\emph{(3)} below is at
the range $M_1=+\infty$, whose hypotheses hold: $(\lambda_j)_{j\geq1}$ satisfies
\eqref{eq:doubling_cut} at every integer $y>d$, and every level used is at least
$\ell_4\geq\ell_3\geq\ell_2\geq2\ell_1$, the three inequalities being those of
Lemmas~\ref{lem:doubling_doeblin}, \ref{lem:doubling_weight} and~\ref{lem:doubling_escape}.
\begin{itemize}
    \item \emph{Two states of one dyadic block have profiles differing by at most
    $c_2(1-\omega)^{i}+2c_2c_3/\ell$.} Let $\ell\geq\ell_4$ and $i\geq0$ be integers and let
    $y,y'\in I_i(\ell)$. At $i=0$ both $u_y$ and $u_{y'}$ lie in $(0,c_2]$, so
    $|u_y-u_{y'}|\leq c_2=c_2(1-\omega)^{0}$, which is at most the announced bound,
    $2c_2c_3/\ell$ being non-negative. Let $i\geq1$, so that $y,y'\geq2^{i}\ell\geq2\ell$,
    and run the descent chain at the level $\ell$: splitting
    $Z_\ell\,u_{Y_{N_\ell}}=u_{Y_{N_\ell}}+(Z_\ell-1)\,u_{Y_{N_\ell}}$ in
    \eqref{eq:doubling_pathid} at the starting state $y$ gives
    \begin{equation}
        u_y\;=\;\mathbb E\bigl(u_{Y_{N_\ell}}\mid Y_0=y\bigr)
        \;+\;\mathbb E\bigl((Z_\ell-1)\,u_{Y_{N_\ell}}\mid Y_0=y\bigr),
    \end{equation}
    and the same at $y'$. Since $0<u_j\leq c_2$ at every $j\geq1$, the second term has modulus
    at most $c_2\,\mathbb E(|Z_\ell-1|\mid Y_0=y)\leq c_2c_3/\ell$ by
    \eqref{eq:doubling_weight}, which applies at the level $\ell\geq\ell_3$ and the starting
    state $y\geq2\ell$, and the same at $y'$. Write
    $\pi_z:=\mathbb P(Y_{N_\ell}\in\cdot\mid Y_0=z)$ for the law of the exit state from a
    starting state $z$, so that the two first terms are $\int u\,d\pi_y$ and
    $\int u\,d\pi_{y'}$, both bearing on the restriction of $u$ to $[\ell,2\ell)$, where
    $Y_{N_\ell}$ lives by Lemma~\ref{lem:doubling_descent}\emph{(2)}. Two probabilities $\pi$,
    $\pi'$ on the integers and a bounded function $f$ satisfy
    $\bigl|\int f\,d(\pi-\pi')\bigr|\leq(\sup f-\inf f)\,\TV(\pi\|\pi')$: the signed measure
    $\pi-\pi'$ has total mass $0$, so
    $\int f\,d(\pi-\pi')=\int\bigl(f-\tfrac12(\sup f+\inf f)\bigr)\,d(\pi-\pi')$, the
    integrand of the right side has modulus at most $\tfrac12(\sup f-\inf f)$, and
    $\sum_j|\pi(j)-\pi'(j)|=2\,\TV(\pi\|\pi')$ in the $\tfrac12$-convention in force
    (Definition~\ref{def:doubling_setting}). Apply this at $f=u$ on $[\ell,2\ell)$, whose
    supremum minus infimum there is at most $c_2-c_1\leq c_2$, and at $\pi=\pi_y$,
    $\pi'=\pi_{y'}$, whose total variation is at most $(1-\omega)^{i}$ by
    \eqref{eq:doubling_coupling}, whose hypotheses
    $\ell\geq\ell_4$ and $y,y'\in I_i(\ell)$ hold: the two first terms differ by at most
    $c_2(1-\omega)^{i}$.
    \item \emph{Two states above $2^{i}\ell$ have profiles differing by at most
    $c_2(1-\omega)^{i}+3c_2c_3/\ell$.} Let $\ell\geq\ell_4$ and $i\geq0$ be integers and let
    $m,m'$ be integers with $m,m'\geq2^{i}\ell$. The blocks $I_j(\ell)$, $j\geq0$, partition
    the integers $\geq\ell$, so $m\in I_j(\ell)$ and $m'\in I_{j'}(\ell)$ for unique integers
    $j$ and $j'$, both $\geq i$ because $m,m'\geq2^{i}\ell$, and $j\leq j'$ after exchanging
    $m$ and $m'$ if need be, the claim being symmetric in the two. If $j'=j$, then $m$ and
    $m'$ are two states of the block $I_j(\ell)$, so
    $|u_m-u_{m'}|\leq c_2(1-\omega)^{j}+2c_2c_3/\ell$, which is at most the announced bound
    because $0<1-\omega<1$ and $j\geq i$. Let $j'>j$ and put $\ell':=2^{j}\ell$, so that
    $I_j(\ell)=[\ell',2\ell')$, $\ell'\geq\ell\geq\ell_4$ and $m'\geq2^{j'}\ell\geq2\ell'$.
    Run the descent chain at the level $\ell'$ from the starting state $m'$:
    \eqref{eq:doubling_pathid} gives
    $u_{m'}=\mathbb E\bigl(Z_{\ell'}u_{Y_{N_{\ell'}}}\mid Y_0=m'\bigr)$ with
    $Y_{N_{\ell'}}\in[\ell',2\ell')=I_j(\ell)$ by Lemma~\ref{lem:doubling_descent}\emph{(2)},
    so splitting $Z_{\ell'}u_{Y_{N_{\ell'}}}-u_m=\bigl(u_{Y_{N_{\ell'}}}-u_m\bigr)
    +(Z_{\ell'}-1)\,u_{Y_{N_{\ell'}}}$ and using $0<u_j\leq c_2$ at every $j\geq1$,
    \begin{equation}
    \begin{aligned}
        \bigl|u_{m'}-u_m\bigr|\;\leq\;&
        \mathbb E\bigl(\bigl|u_{Y_{N_{\ell'}}}-u_m\bigr|\bigm| Y_0=m'\bigr)\\
        &+\;c_2\,\mathbb E\bigl(|Z_{\ell'}-1|\mid Y_0=m'\bigr) .
    \end{aligned}
    \end{equation}
    Both $Y_{N_{\ell'}}$ and $m$ are states of the block $I_j(\ell)$, so the first term is at
    most $c_2(1-\omega)^{j}+2c_2c_3/\ell$; and \eqref{eq:doubling_weight} at the level
    $\ell'\geq\ell_3$ and the starting state $m'\geq2\ell'$ bounds the second by
    $c_2c_3/\ell'$. With $(1-\omega)^{j}\leq(1-\omega)^{i}$ and $\ell'\geq\ell$, this is the
    announced bound.
    \item \emph{The rescaled profile converges, and its limit lies in $[c_1',c_2']$ for every
    pair $0<c_1'\leq c_2'$ satisfying \eqref{eq:doubling_decay}.} Let $t>0$, choose an integer
    $\ell\geq\ell_4$ with $3c_2c_3/\ell\leq t/2$ and then an integer $i\geq0$ with
    $c_2(1-\omega)^{i}\leq t/2$, which $0<1-\omega<1$ permits. Two integers
    $m,m'\geq2^{i}\ell$ then have $|u_m-u_{m'}|\leq c_2(1-\omega)^{i}+3c_2c_3/\ell\leq t$, so
    $(u_m)_{m\geq1}$ is a Cauchy sequence and converges to a limit $C$. If
    $0<c_1'\leq c_2'$ satisfy \eqref{eq:doubling_decay}, then $c_1'\leq u_m\leq c_2'$ at every
    $m\geq1$, so $C\in[c_1',c_2']$; and $u_m=\lambda_m\,m^{p_*}$, so $u_m\rightarrow C$ is
    \eqref{eq:doubling_sharp}.
    \item \emph{The convergence carries the rate \eqref{eq:doubling_rate}.} Put
    $\alpha:=\log_2\bigl(1/(1-\omega)\bigr)$ and $\vartheta:=\alpha/(1+\alpha)$, so that
    $\alpha>0$ and $\vartheta\in(0,1)$, $\omega$ lying in $(0,1)$; both depend on $c$ alone,
    $\omega$ doing so by Lemma~\ref{lem:doubling_doeblin}. Put
    $m_3:=\max\bigl(\lceil\ell_4^{1/\vartheta}\rceil,\lceil2^{1/(1-\vartheta)}\rceil\bigr)$,
    an integer depending on $c$ and $d$ alone, and, for an integer $m\geq1$,
    $\ell:=\lceil m^{\vartheta}\rceil$ and $i:=\lfloor\log_2(m/\ell)\rfloor$. At $m\geq m_3$
    these satisfy $\ell\geq\ell_4$ and $i\geq0$: $m\geq\ell_4^{1/\vartheta}$ gives
    $m^{\vartheta}\geq\ell_4$ and hence $\ell\geq\ell_4$; and $m^{\vartheta}\geq1$ gives
    $\ell\leq m^{\vartheta}+1\leq2m^{\vartheta}$, which $m^{1-\vartheta}\geq2$ makes at most
    $m$, so that $i\geq0$. There $2^{i}\ell\leq m$, so $m$ and every integer
    $m'\geq2^{i}\ell$ are two states above $2^{i}\ell$ and satisfy
    $|u_m-u_{m'}|\leq c_2(1-\omega)^{i}+3c_2c_3/\ell$; letting $m'\rightarrow\infty$ gives
    $|u_m-C|\leq c_2(1-\omega)^{i}+3c_2c_3/\ell$. Here $\ell\geq m^{\vartheta}$ gives
    $1/\ell\leq m^{-\vartheta}$, while $2^{i+1}>m/\ell$ and $1-\omega=2^{-\alpha}$ give
    \begin{equation}
    \begin{aligned}
        (1-\omega)^{i}\;&=\;2^{-\alpha i}\;<\;\Bigl(\frac{2\ell}{m}\Bigr)^{\alpha}
        \;\leq\;\Bigl(\frac{4m^{\vartheta}}{m}\Bigr)^{\alpha}\\
        \;&=\;4^{\alpha}\,m^{-\alpha(1-\vartheta)}\;=\;4^{\alpha}\,m^{-\vartheta},
    \end{aligned}
    \end{equation}
    by $\ell\leq2m^{\vartheta}$ and $\alpha(1-\vartheta)=\alpha/(1+\alpha)=\vartheta$. Hence
    $|u_m-C|\leq c_2\bigl(4^{\alpha}+3c_3\bigr)m^{-\vartheta}$ at every integer $m\geq m_3$;
    and at $1\leq m<m_3$, both $u_m$ and $C$ lying in $[c_1,c_2]$,
    $|u_m-C|\leq c_2\leq c_2\,m_3^{\vartheta}\,m^{-\vartheta}$. So
    $c_6:=c_2\max\bigl(4^{\alpha}+3c_3,\,m_3^{\vartheta}\bigr)$ satisfies
    \eqref{eq:doubling_rate}, $u_j$ being $\lambda_j\,j^{p_*}$. It depends on $c$, $d$ and
    $\lambda_1,\dots,\lambda_{2m_0}$ alone: $\alpha$ and $\vartheta$ on $c$; $c_3$, $\ell_4$
    and with them $m_3$ on $c$ and $d$, by Lemmas~\ref{lem:doubling_weight}
    and~\ref{lem:doubling_doeblin}; and $c_2$ on $c$, $d$ and $\lambda_1,\dots,\lambda_{2m_0}$
    by Theorem~\ref{theo:doubling_decay}\emph{(2)}.
    \item \emph{The limit is determined by $c$, $d$ and $\lambda_1,\dots,\lambda_d$.} By
    Lemma~\ref{lem:doubling_averaging}\emph{(1)}, \eqref{eq:doubling_avg} holds at every
    integer $m>d$, and it determines $\lambda_m$ from $c$ and
    $\lambda_{\lceil m/2\rceil},\dots,\lambda_{m-1}$, whose indices lie in $\{1,\dots,m-1\}$;
    by induction it determines the whole sequence, and with it $C$, from $c$, $d$ and
    $\lambda_1,\dots,\lambda_d$.
\end{itemize}
\end{proof}

\begin{remark}[no log-periodic term]\label{rem:doubling_no_logperiodic}
The limit in \eqref{eq:doubling_sharp} exists, so the decay carries no log-periodic term, the
usual failure mode of dyadic recursions. What excludes it is the window of
Definition~\ref{def:doubling_decay_notation}. A descent that halved deterministically would
leave each block $I_i(\ell)$ at a single state, fixed by the dyadic position of the state it
started from; two starts of one block would then share no reference law, the best constant in
\eqref{eq:doubling_doeblin} would be $0$, and every $1$-periodic function of $\log_2m$ would be
compatible with the recursion such a descent would represent, to leading order. The window spreads
the exit law over the integers of $[\lceil5a/4\rceil,2a)$, where $a=2^{i-1}\ell$ is the foot of
the block below, with mass bounded below at every one of them, and
Lemma~\ref{lem:doubling_coupling} leaves only the constants.
\end{remark}

\begin{corollary}[the tail, with the constant]\label{cor:doubling_tail_sharp}
Let $s=1$ and $0<c<1$, let $(\lambda_j)_{j\geq1}$ be a positive sequence satisfying
\eqref{eq:doubling_cut} for every $m>d$, in the notation of
Definitions~\ref{def:doubling_setting} and~\ref{def:doubling_decay_notation}, let $C$ be the
constant of \eqref{eq:doubling_sharp} and let $L$ be the tail of
Definition~\ref{def:doubling_setting}.
Then
\begin{equation}\label{eq:doubling_tail_sharp}
\begin{aligned}
    L(m)\;&=\;\frac{C}{p_*-1}\,m^{1-p_*}\bigl(1+o(1)\bigr),\\
    \frac{L(m)}{\lambda_m}\;&=\;\frac{m}{p_*-1}\bigl(1+o(1)\bigr).
\end{aligned}
\end{equation}
\end{corollary}

\begin{proof}
By \eqref{eq:doubling_powertail},
$\sum_{j\geq m}j^{-p}=\frac{m^{1-p}}{p-1}(1+\xi_m)$ with
$0\leq\xi_m\leq(p-1)/m$. Let $t>0$ and let $m_t$ be such that
$|u_j-C|\leq t$ for $j\geq m_t$, which
Theorem~\ref{theo:doubling_sharp} provides. For $m\geq m_t$,
\begin{equation}
    (C-t)\sum_{j\geq m}j^{-p}\;\leq\;L(m)\;\leq\;(C+t)\sum_{j\geq m}j^{-p},
\end{equation}
so $(p-1)m^{p-1}L(m)$ has all its limit points in $[C-t,C+t]$; as
$t$ is arbitrary, $(p-1)m^{p-1}L(m)\rightarrow C$, which is the first half of
\eqref{eq:doubling_tail_sharp}. Dividing it by $u_m=\lambda_mm^{p}\rightarrow C>0$ gives the
second.
\end{proof}

\begin{proposition}[the constant is a linear functional of the boundary data]\label{prop:doubling_constant}
Let $s=1$ and $0<c<1$, in the notation of Definitions~\ref{def:doubling_setting}
and~\ref{def:doubling_decay_notation}, and let $C$ be the limit of \eqref{eq:doubling_sharp},
which Theorem~\ref{theo:doubling_sharp}
attaches to every positive sequence satisfying \eqref{eq:doubling_cut} at every integer $m>d$.
Then the real sequences $(\lambda_j)_{j\geq1}$ satisfying \eqref{eq:doubling_cut} at every
integer $m>d$ form a real vector space, the map
$(\lambda_j)_{j\geq1}\mapsto(\lambda_1,\dots,\lambda_d)$ is a linear isomorphism of it onto
$\mathbb R^{d}$, and such a sequence is positive if and only if $\lambda_1,\dots,\lambda_d$
are. There are reals $\nu_1,\dots,\nu_d$ with
\begin{equation}\label{eq:doubling_constant}
    C\;=\;\sum_{j=1}^{d}\nu_j\,\lambda_j
\end{equation}
at every positive such sequence, and they satisfy $\nu_j\geq0$ for every $j$,
$\sum_{j=1}^{d}\nu_j>0$, and $\nu_j=0$ whenever $2j\leq d$.
\end{proposition}

\begin{proof}
\emph{Step 1 (the space, its coordinates and its positive cone).} Let $m>d$ be an integer. Then
$\varepsilon(m)=c/(m+1)\leq c/2<1$, so dividing \eqref{eq:doubling_cut} at $m$ by
$1-\varepsilon(m)>0$ and substituting $\varepsilon(j)=c/(j+1)$ turns it into
$\lambda_m=\sum_{j\in W(m)}q_m(j)\lambda_j$, with $q_m>0$ on $W(m)$ by
\eqref{eq:doubling_window}, and every $j\in W(m)$ satisfies $1\leq\lceil m/2\rceil\leq j\leq
m-1$. That identity is linear in $(\lambda_j)_{j\geq1}$, so the sequences satisfying it at
every $m>d$ form a real vector space; and by induction on $m$ it determines $\lambda_m$ for
$m>d$ from $\lambda_1,\dots,\lambda_d$, and produces exactly one sequence of the space for each
prescribed $(\lambda_1,\dots,\lambda_d)\in\mathbb R^{d}$: the restriction to the first $d$
coordinates is a linear isomorphism onto $\mathbb R^{d}$. If $\lambda_1,\dots,\lambda_d>0$, the
same induction with $q_m>0$ gives $\lambda_m>0$ at every $m>d$; conversely a positive sequence
has $\lambda_1,\dots,\lambda_d>0$.

\emph{Step 2 ($C$ is linear on the positive cone).} By Step 1, $\lambda\mapsto\lambda_m$ is a
linear functional of $(\lambda_1,\dots,\lambda_d)$ at every $m\geq1$, hence so is
$\lambda\mapsto\lambda_m\,m^{p_*}$; and by Step 1 the positive sequences are the preimage of
the open positive orthant of $\mathbb R^{d}$ under the isomorphism of that step, hence a
non-empty convex cone, open for the topology that isomorphism transports. On that cone
Theorem~\ref{theo:doubling_sharp} makes $C=\lim_m\lambda_m\,m^{p_*}$ exist, so $C$ is additive
and positively homogeneous there, both properties passing from the functionals
$\lambda\mapsto\lambda_m\,m^{p_*}$ to their limit wherever it exists. Every element
$\lambda''$ of the space is a difference of two elements of the cone: taking $\lambda$ in the
cone and $x>0$ small enough that $\lambda+x\lambda''$ lies in it, which openness allows,
$\lambda''=x^{-1}(\lambda+x\lambda'')-x^{-1}\lambda$ and the cone is stable under
multiplication by $x^{-1}>0$. Setting $C(\lambda-\lambda'):=C(\lambda)-C(\lambda')$ is
unambiguous by additivity and extends $C$ to an additive and $\mathbb R$-homogeneous, that is
linear, functional on the space. Hence
$C=\sum_{j=1}^{d}\nu_j\lambda_j$ with $\nu_j$ the value of that extension at the sequence whose
first $d$ coordinates are the $j$-th vector of the canonical basis of $\mathbb R^{d}$: this is
\eqref{eq:doubling_constant}.

\emph{Step 3 (the signs, and the indices carrying no weight).} At
$\lambda_1=\dots=\lambda_d=1$, \eqref{eq:doubling_constant} reads $\sum_{j=1}^{d}\nu_j=C$,
which is positive: Theorem~\ref{theo:doubling_decay} supplies a pair $0<c_1\leq c_2$ with
\eqref{eq:doubling_decay} and Theorem~\ref{theo:doubling_sharp} puts $C$ in $[c_1,c_2]$. Fix
$j\leq d$, take $\lambda_i=1$ for $i\neq j$ and $\lambda_j=t$: then
$C=\nu_jt+\sum_{i\neq j}\nu_i$ is an affine function of $t$, positive at every $t>0$ by those
same two theorems, whereas an affine function of negative slope is
negative for $t$ large; hence $\nu_j\geq0$. Finally let $2j\leq d$ and let $m>d$. Every
$i\in W(m)$ satisfies $i\geq\lceil m/2\rceil\geq(d+1)/2>d/2\geq j$, so the induction of Step 1
expresses $\lambda_m$ as a linear combination of the $\lambda_i$ with $d/2<i\leq d$ alone;
hence neither $\lambda_m$ for $m>d$ nor $C=\lim_m\lambda_m\,m^{p_*}$ depends on $\lambda_j$,
and $\nu_j=0$.
\end{proof}

The constant is therefore not explicit, and not a function of $c$ alone: it is a linear
functional of the boundary data, and \eqref{eq:doubling_constant} leaves its coefficients
unevaluated.

\begin{remark}[the second order, and the parity wobble --- a formal expansion]\label{rem:doubling_second_order}
In the notation of Definition~\ref{def:doubling_decay_notation}, with $C$ the constant of
Theorem~\ref{theo:doubling_sharp} and with the coefficients $A_0$, $A_1$ and $\Gamma$ of
Lemma~\ref{lem:doubling_expansion}, substituting the ansatz
$u_m=C\bigl(1+\alpha_{\delta(m)}/m+o(1/m)\bigr)$, with coefficients $\alpha_0$ and $\alpha_1$
to be determined and $\alpha_\delta$ standing for $\alpha_{\delta(m)}$, into
\eqref{eq:doubling_uavg}, and using $\sum_{j\in W(m)}q_m(j)\Phi_{\bar\alpha}(j)=R_{\bar\alpha}(m)\Phi_{\bar\alpha}(m)$ from
\eqref{eq:doubling_R} at the constant $\bar\alpha:=(\alpha_0+\alpha_1)/2$ --- the alternating
part of $j\mapsto\alpha_{\delta(j)}$ contributing only at order $m^{-p_*-2}$ --- expansion
\eqref{eq:doubling_Rexp} matches the coefficient of $1/m$ at each parity to $\alpha_\delta=A_\delta+(\Gamma+1)\bar\alpha$. Averaging over $\delta$ gives
$\bar\alpha=-(A_0+A_1)/(2\Gamma)$, and subtracting gives $\alpha_0-\alpha_1=A_0-A_1=c\tau$: the
amplitude of the parity wobble is fixed by \eqref{eq:doubling_AB} with no free constant. This
is a formal matching computation and is not proved here. Proving it needs $|u_m-C|=O(1/m)$,
whereas Theorem~\ref{theo:doubling_sharp} gives \eqref{eq:doubling_rate} with an exponent
$\vartheta<1$ depending on $c$ alone, and the $\lceil m/2\rceil$ of \eqref{eq:doubling_cut}
closes the Mellin route to the exponent $1$.
\end{remark}

\subsubsection{Unboundedness of the diffusion operator}

This block proves items \emph{(3)} and \emph{(4)} of Theorem~\ref{theo:doubling_main}: across
the whole standing range of the family, and against the condition of
\citet{morozov2025revisiting}, the diffusion operator is unbounded on every $L^p(\lambda)$ at
every positive recurrent member, and some pairs of $L^2$ densities cannot be balanced there at
all. Lemma~\ref{lem:doubling_truncation_irreducible} makes the truncation at an even cap a
finite irreducible chain, and Lemma~\ref{lem:doubling_percut} bounds, at one cut of the ladder
and on the loop closure and the truncation alike, the Rayleigh ratio of the centred tail
indicator at that cut. Theorem~\ref{theo:doubling_unbounded} is the general statement at the
finite exponents --- it asks of the doubling probability only a growth condition ---
Lemma~\ref{lem:doubling_ramp} reaches the exponent $+\infty$ along a second family of trial
functions, and Corollary~\ref{cor:doubling_family} makes the finite exponents explicit at
$s=1$ and $c<1$, with a rate and along a family named in advance.
Proposition~\ref{prop:doubling_unsolvable} turns the unboundedness into a pair of $L^2$
densities admitting no exact outflow. Remark~\ref{rem:doubling_geometric} shows that growth
condition is not vacuous: a geometrically decaying doubling probability fails it.

\begin{lemma}[the truncation is a finite irreducible chain]\label{lem:doubling_truncation_irreducible}
In the setting of Definition~\ref{def:doubling_setting}, let $\varepsilon(j)\in(0,1)$ for every
$j\geq1$ and let $K\geq d$ be an even integer. Then the truncation at $K$ is an irreducible
Markov chain on the $K+2$ states $\{s_0,1,\dots,K,s_f\}$; it carries a unique invariant
probability $\lambda^K$, which is positive at every state; and the operator $S$ of
\eqref{eq:doubling_Bhat} formed on it satisfies \eqref{eq:doubling_resolvent}, with
$\widehat B_K<+\infty$.
\end{lemma}

\begin{proof}
\emph{Step 1 (the state set is closed and every retained transition is positive).} A ladder
state $j\leq K$ decrements to $j-1$, read as $s_0$ at $j=1$, and doubles to $2j$ only when that
edge is present, that is when $2j\leq K$; $s_0$ maps to $s_f$, and $s_f$ into $\{1,\dots,d\}$,
which is contained in $\{1,\dots,K\}$ because $d\leq K$. So $\{s_0,1,\dots,K,s_f\}$ is closed
and carries the truncation as a Markov chain on $K+2$ states. The decrement out of a ladder
state $j$ has probability $1-\varepsilon(j)>0$ when $2j\leq K$ and probability $1$ when
$2j>K$; the doubling edge out of a ladder state $j$ with $2j\leq K$ has probability
$\varepsilon(j)>0$; the wrap $s_0\rightarrow s_f$ has probability $1$; and the target row, a
probability supported in $\{1,\dots,d\}$, is positive at some $j_0\in\{1,\dots,d\}$.

\emph{Step 2 (every state leads to the state $1$).} From a ladder state $j$ the $j$ consecutive
decrements reach $s_0$, from $s_0$ the wrap reaches $s_f$, from $s_f$ the target row reaches
$j_0$, and from $j_0$ the $j_0-1$ consecutive decrements reach $1$; each of these transitions
is positive by Step 1.

\emph{Step 3 (the state $1$ leads to every state).} Let $2^r$ be the largest power of two with
$2^r\leq K$, so that $K<2^{r+1}$ and $K/2<2^r$. The doubling edges
$2^i\rightarrow2^{i+1}$ are present for $0\leq i<r$, so $1$ leads to $2^r$, and the decrements
out of $2^r$ lead to every ladder state $n<2^r$: the state $1$ leads to every ladder state
$n\leq2^r$. Let $2^r<n\leq K$. If $n$ is even, then $n/2\leq K/2<2^r$, so $1$ leads to $n/2$,
and the doubling edge $n/2\rightarrow n$ is present because $n\leq K$. If $n$ is odd, then
$n\neq K$, $K$ being even, so $n+1$ is even with $2^r<n+1\leq K$ and $1$ leads to $n+1$ by the
previous case, and the decrement $n+1\rightarrow n$ is present. Hence $1$ leads to every ladder
state $n\leq K$, and to $s_0$ and $s_f$ by Step 2.

\emph{Step 4 (the invariant probability and the diffusion operator).} By Steps 2 and 3 every
state leads to every other, so the truncation is an irreducible chain on a finite state space
and carries a unique invariant probability $\lambda^K$, positive at every state
\citep{levin2017markov}. Lemma~\ref{lem:doubling_operator}\emph{(2)}, applied to that chain,
gives $S$ and \eqref{eq:doubling_resolvent}; $S$ acts on the finite-dimensional space
$L^2(\lambda^K)$, so $\widehat B_K=\|S\|_{L^2(\lambda^K)}<+\infty$.
\end{proof}

\begin{lemma}[the per-cut inequality]\label{lem:doubling_percut}
In the setting of Definition~\ref{def:doubling_setting}, let the backward policy be
\eqref{eq:doubling_policy} with $\varepsilon(j)>0$ for every $j\geq1$ and
$\varepsilon_{\max}=\sup_{j\geq1}\varepsilon(j)<1$, and let $m$ be an integer. Assume that one
of the following holds:
\begin{enumerate}
    \item the loop-closed chain is positive recurrent with invariant probability $\lambda$ and
    $m>d$; put $A:=[m,+\infty)$ and $L:=L(m)$;
    \item $K\geq d$ is an even integer, $\lambda$ is the invariant probability $\lambda^K$ of
    the truncation at $K$, and $d<m\leq K/2$; put $A:=[m,K]$ and $L:=L_K(m)$.
\end{enumerate}
Let $P_\star$ be the function action of the backward kernel of that chain,
let $\Pi$ be the $\lambda$-mean projection and put $f_m:=\mathbf 1_A-L$. Then
\begin{equation}\label{eq:doubling_step1}
    \left(\mathrm{Id}-P_\star\right)\mathbf 1_A
    \;=\;\bigl(1-\varepsilon(m)\bigr)\,\mathbf 1_{\{m\}}
    \;-\;\sum_{x=\lceil m/2\rceil}^{m-1}\varepsilon(x)\,\mathbf 1_{\{x\}} ,
\end{equation}
\begin{equation}\label{eq:doubling_step2}
    \left\|\left(\mathrm{Id}-P_\star\right)\mathbf 1_A\right\|^2_{L^2(\lambda)}
    \;\leq\;\lambda_m\bigl(1-\varepsilon(m)\bigr)\bigl(1-\varepsilon(m)+\varepsilon_{\max}\bigr)
    \;\leq\;2\lambda_m ,
\end{equation}
and $\Pi f_m=0$, $\|f_m\|^2_{L^2(\lambda)}=L\,(1-L)>0$ and
\begin{equation}\label{eq:doubling_step3}
    \frac{\left\|\left(\mathrm{Id}-P_\star\right)f_m\right\|_{L^2(\lambda)}}{\|f_m\|_{L^2(\lambda)}}
    \;\leq\;\sqrt{\frac{2\lambda_m}{L\bigl(1-L\bigr)}} .
\end{equation}
Moreover, for every $p\in[1,+\infty)$,
\begin{equation}\label{eq:doubling_stepp}
\begin{gathered}
    \left\|\left(\mathrm{Id}-P_\star\right)f_m\right\|^p_{L^p(\lambda)}
    \;\leq\;\lambda_m\left(1+\varepsilon_{\max}^{\,p-1}\right)\;\leq\;2\lambda_m,\\
    \|f_m\|^p_{L^p(\lambda)}\;=\;L\,(1-L)^p+(1-L)\,L^p\;\geq\;L\,(1-L)^p .
\end{gathered}
\end{equation}
\end{lemma}

\begin{proof}
\emph{Step 1 (the invariant measure is positive and the tail is not degenerate).} In case
\emph{(1)} $\varepsilon$ satisfies $0<\varepsilon(j)\leq\varepsilon_{\max}<1$ at every
$j\geq1$, so the loop closure is irreducible and $\lambda$ is positive at every state
(Lemma~\ref{lem:doubling_irreducible}); in case
\emph{(2)} $\lambda$ is positive at every state by
Lemma~\ref{lem:doubling_truncation_irreducible}. Since $m>d\geq1$, the state $1$ lies outside
$A$ and the state $m$ lies in $A$, so $0<\lambda_m\leq L\leq1-\lambda_1<1$.

\emph{Step 2 (an exact identity).} $P_\star$ is the function action of the backward kernel
(Lemma~\ref{lem:adjoint}\emph{(3)}), so
$\bigl((\mathrm{Id}-P_\star)\mathbf 1_A\bigr)(x)=\mathbf 1_A(x)-\mathbb P(X_1\in A\mid X_0=x)$, and the
transitions \eqref{eq:doubling_policy} give six cases. For $x\geq m+1$ in $A$ both images lie
in $A$: the decrement because $x-1\geq m$, and the doubling, when its edge is present, because
$2x\geq x\geq m$ and, in case \emph{(2)}, $2x\leq K$. A state $x>K/2$ of case \emph{(2)} falls
here with the decrement as its only image, which lies in $A$; in either case the value is
$1-1=0$. For $x=m$ the doubling edge is present, in case \emph{(2)} because $2m\leq K$, and
the value is
$1-\varepsilon(m)$. For $\lceil m/2\rceil\leq x\leq m-1$ only the doubling edge enters $A$, and
it is present, in case \emph{(2)} because $2x\leq2m-2<K$, so the value is $-\varepsilon(x)$.
For $1\leq x<\lceil m/2\rceil$ both $2x<m$ and $x-1<m$, so the value is $0$. At $s_0$ the image
is $s_f\notin A$, and at $s_f$ the image lies in $\{1,\dots,d\}$ with $d<m$, so both values are
$0$. This is \eqref{eq:doubling_step1}.

\emph{Step 3 (its norm).} The indicators in \eqref{eq:doubling_step1} have disjoint supports,
so for every $p\in[1,+\infty)$
\begin{equation}
\begin{aligned}
    \left\|\left(\mathrm{Id}-P_\star\right)\mathbf 1_A\right\|^p_{L^p(\lambda)}
    \;&=\;\lambda_m\bigl(1-\varepsilon(m)\bigr)^p\\
    &\quad+\sum_{x=\lceil m/2\rceil}^{m-1}\lambda_x\,\varepsilon(x)^p .
\end{aligned}
\end{equation}
The cut balance \eqref{eq:doubling_cut} holds at $m$ by Proposition~\ref{prop:doubling_cut} ---
in case \emph{(1)} because $m>d$, in case \emph{(2)} because $K$ is even and $d<m\leq K/2$ ---
so bounding $\varepsilon(x)^p\leq\varepsilon_{\max}^{\,p-1}\varepsilon(x)$ on the sum gives
$\sum_x\lambda_x\varepsilon(x)^p\leq\varepsilon_{\max}^{\,p-1}\sum_x\lambda_x\varepsilon(x)=\varepsilon_{\max}^{\,p-1}\lambda_m(1-\varepsilon(m))$
and
\begin{equation}
\begin{gathered}
    \left\|\left(\mathrm{Id}-P_\star\right)\mathbf 1_A\right\|^p_{L^p(\lambda)}\\
    \leq\;\lambda_m\bigl(1-\varepsilon(m)\bigr)
    \Bigl(\bigl(1-\varepsilon(m)\bigr)^{p-1}+\varepsilon_{\max}^{\,p-1}\Bigr).
\end{gathered}
\end{equation}
At $p=2$, with $1-\varepsilon(m)\leq1$ and $1-\varepsilon(m)+\varepsilon_{\max}\leq2$, this is
\eqref{eq:doubling_step2}; at a general $p$, with $1-\varepsilon(m)\leq1$ and
$\varepsilon_{\max}<1$, it is the first pair of \eqref{eq:doubling_stepp}.

\emph{Step 4 (the ratio at the cut).} $P_\star\mathbf 1=\mathbf 1$, so
$(\mathrm{Id}-P_\star)f_m=(\mathrm{Id}-P_\star)\mathbf 1_A$. Since $\lambda(A)=L$, the integral
$\int f_m\,d\lambda$ vanishes, so $\Pi f_m=0$; and $f_m$ takes the value $1-L$ on $A$, of
$\lambda$-mass $L$, and $-L$ off $A$, so
$\|f_m\|^p_{L^p(\lambda)}=L(1-L)^p+(1-L)L^p$ for every $p\in[1,+\infty)$, which is the second
pair of \eqref{eq:doubling_stepp} and, at $p=2$, is $L(1-L)$, positive by Step 1.
Dividing \eqref{eq:doubling_step2} by it and taking square roots gives
\eqref{eq:doubling_step3}.
\end{proof}

\begin{theorem}[the diffusion operator is unbounded]\label{theo:doubling_unbounded}
In the setting of Definition~\ref{def:doubling_setting}, let the backward policy be
\eqref{eq:doubling_policy} with $\varepsilon(j)>0$ for every $j\geq1$ and
$\varepsilon_{\max}=\sup_{j\geq1}\varepsilon(j)<1$, and consider the growth condition
\begin{equation}\label{eq:doubling_star}
    \sum_{i=0}^{D-1}\log\frac{1}{\varepsilon(2^i)}\;=\;o\left(2^D\right)\qquad(D\rightarrow\infty).
\end{equation}
Then:
\begin{enumerate}
    \item at $\varepsilon=\varepsilon_{c,s}$ with $s\geq0$ and $0<c<2^s$, the left side of
    \eqref{eq:doubling_star} is $O\left(D^2\right)$ and \eqref{eq:doubling_star} holds;
    \item if the loop-closed chain is positive recurrent with invariant probability $\lambda$
    and \eqref{eq:doubling_star} holds, then for every $p\in[1,+\infty)$
    \begin{equation}\label{eq:doubling_inf}
        \inf\bigl\{\|(\mathrm{Id}-P_\star)f\|_{L^p(\lambda)}:\Pi f=0,\
        \|f\|_{L^p(\lambda)}=1\bigr\}=0
    \end{equation}
    and no bounded operator $S$ on $L^p(\lambda)$ satisfies
    $S(\mathrm{Id}-P_\star)=\mathrm{Id}-\Pi$; at $p=2$ the series
    $\sum_{n\geq0}(P_\star^n-\Pi)$ does not converge in operator norm and
    $\sum_{n\geq0}\widehat\beta_n=+\infty$.
\end{enumerate}
\end{theorem}

\begin{proof}
\emph{(1)} At $\varepsilon=\varepsilon_{c,s}$ with $s\geq0$ and $0<c<2^s$,
$\sum_{i<D}\log(1/\varepsilon(2^i))=\sum_{i<D}\bigl(s\log(2^i+1)-\log c\bigr)=O(D^2)=o(2^D)$.

\emph{(2)} Let the loop-closed chain be positive recurrent with invariant probability
$\lambda$ and let \eqref{eq:doubling_star} hold. Since
$0<\varepsilon(j)\leq\varepsilon_{\max}<1$ at every $j\geq1$, the loop closure is
irreducible and $\lambda_x>0$ for every state $x$ (Lemma~\ref{lem:doubling_irreducible}); in
particular $\lambda_1>0$ and
$0<L(m)<1$ for every $m>1$, and $L(m)\rightarrow0$, $\lambda$ being a probability.

\emph{Step 1 (the ratio $L/\lambda$ is unbounded).} Suppose $\sup_{m>d}L(m)/\lambda_m<+\infty$
and put $Q:=\sup_{m>d}L(m)/\lambda_m$, so that $Q\geq1$ because $L(m)\geq\lambda_m$. Since
$\lambda_m=L(m)-L(m+1)$, the bound $L(m)\leq Q\lambda_m$ gives
$L(m+1)\leq L(m)(1-1/Q)\leq L(m)e^{-1/Q}$ at every $m>d$, hence
$L(m)\leq L(d+1)\,e^{-(m-d-1)/Q}$ for every $m\geq d+1$. Iterating the first inequality of
\eqref{eq:doubling_ratios}, which Proposition~\ref{prop:doubling_cut} gives at every ladder
state, along the powers of two yields
$L(2^D)\geq\lambda_{2^D}\geq\lambda_1\prod_{i=0}^{D-1}\varepsilon(2^i)$. Taking logarithms of
the two bounds at $m=2^D$, for every integer $D$ with $2^D\geq d+1$,
\begin{equation}
    \frac{2^D-d-1}{Q}\;\leq\;\log\frac{L(d+1)}{\lambda_1}+\sum_{i=0}^{D-1}\log\frac1{\varepsilon(2^i)},
\end{equation}
whose right side is $o(2^D)$ by \eqref{eq:doubling_star} while the left side is of order
$2^D$: a contradiction for $D$ large. Hence $\sup_{m>d}L(m)/\lambda_m=+\infty$.

\emph{Step 2 (the infimum, along the centred tail indicators).} Since $\lambda_m>0$ at every
ladder state, dropping finitely many cuts from the supremum of Step 1 leaves it infinite:
$\sup_{m\geq m_1}L(m)/\lambda_m=+\infty$ for every integer $m_1>d$. Fix an integer $m_1>d$
with $1-L(m)\geq\tfrac12$ for every $m\geq m_1$, which is possible because
$L(m)\rightarrow0$. Let $p\in[1,+\infty)$, let $t>0$ and choose $m\geq m_1$ with
$\lambda_m/L(m)\leq\tfrac12(t/2)^p$. Case \emph{(1)} of Lemma~\ref{lem:doubling_percut} holds
at that cut, so $\Pi f_m=0$ and, by \eqref{eq:doubling_stepp} with $1-L(m)\geq\tfrac12$,
$\|(\mathrm{Id}-P_\star)f_m\|_{L^p(\lambda)}\leq(2\lambda_m)^{1/p}$ and
$\|f_m\|_{L^p(\lambda)}\geq\tfrac12L(m)^{1/p}>0$, whence
\begin{equation}
    \frac{\left\|\left(\mathrm{Id}-P_\star\right)f_m\right\|_{L^p(\lambda)}}{\|f_m\|_{L^p(\lambda)}}
    \;\leq\;2\left(\frac{2\lambda_m}{L(m)}\right)^{1/p}\;\leq\;t .
\end{equation}
The function $f_m/\|f_m\|_{L^p(\lambda)}$ is mean-zero of unit $L^p(\lambda)$-norm and $t>0$
was arbitrary, which is \eqref{eq:doubling_inf}.

\emph{Step 3 (the consequences).} Let $p\in[1,+\infty)$. If a bounded $S$ on $L^p(\lambda)$
satisfied $S(\mathrm{Id}-P_\star)=\mathrm{Id}-\Pi$, then
$\|f\|_{L^p(\lambda)}=\|(\mathrm{Id}-\Pi)f\|_{L^p(\lambda)}\leq\|S\|\,\|(\mathrm{Id}-P_\star)f\|_{L^p(\lambda)}$
for every mean-zero $f$, contradicting \eqref{eq:doubling_inf}. At $p=2$, norm convergence of
$\sum_{n\geq0}(P_\star^n-\Pi)$ would produce such an $S$ by
Lemma~\ref{lem:doubling_operator}\emph{(3)}, and $\sum_n\widehat\beta_n<+\infty$ would make
that series absolutely convergent, $\|P_\star^n-\Pi\|_{L^2(\lambda)}=\widehat\beta_n$ by
Lemma~\ref{lem:doubling_operator}\emph{(1)}; both are therefore excluded.
\end{proof}

\begin{lemma}[the clipped ramp]\label{lem:doubling_ramp}
In the setting of Definition~\ref{def:doubling_setting}, let the backward policy be
\eqref{eq:doubling_policy} with $\varepsilon(j)>0$ for every $j\geq1$ and
$\varepsilon_{\max}=\sup_{j\geq1}\varepsilon(j)<1$, let
$\gamma:=\sup_{j\geq1}(j+1)\,\varepsilon(j)$ be finite, and let the loop-closed chain be
positive recurrent with invariant probability $\lambda$. Write $V$ for the ladder height,
$V(j):=j$ at a ladder state $j\geq1$ and $V(s_0)=V(s_f):=0$, and, for an integer $N\geq d$,
put $r_N:=\min(V,N)$ and $\tilde r_N:=r_N-\Pi r_N$. Then $\tilde r_N\in L^\infty(\lambda)$
with $\Pi\tilde r_N=0$,
\begin{equation}\label{eq:doubling_ramp}
\begin{aligned}
    \|\tilde r_N\|_{L^\infty(\lambda)}\;&\geq\;\frac N2,\\
    \left\|\left(\mathrm{Id}-P_\star\right)\tilde r_N\right\|_{L^\infty(\lambda)}
    \;&\leq\;\max\left(1,\gamma,\bar\jmath\right),
\end{aligned}
\end{equation}
and consequently
$\inf\bigl\{\|(\mathrm{Id}-P_\star)f\|_{L^\infty(\lambda)}:\Pi f=0,\ \|f\|_{L^\infty(\lambda)}=1\bigr\}=0$
and no bounded operator $S$ on $L^\infty(\lambda)$ satisfies
$S(\mathrm{Id}-P_\star)=\mathrm{Id}-\Pi$. For the family \eqref{eq:doubling_family} in the
standing range, $\gamma=c\,2^{1-s}$ is finite exactly when $s\geq1$.
\end{lemma}

\begin{proof}
\begin{itemize}
    \item \emph{The image of the clipped ramp is bounded uniformly in $N$.} Indeed, $r_N$ is
    bounded, so $(P_\star r_N)(x)=\mathbb E(r_N(X_1)\mid X_0=x)$ at every state, and the
    transitions \eqref{eq:doubling_policy} give five cases. At a ladder state $j$ with
    $2j\leq N$ the two images $2j$ and $j-1$ --- the latter read as $s_0$ at $j=1$, where
    $r_N(s_0)=0=j-1$ --- are at most $N$, so
    $(P_\star r_N)(j)=\varepsilon(j)\,2j+\bigl(1-\varepsilon(j)\bigr)(j-1)=j-1+\varepsilon(j)(j+1)$
    and $\bigl((\mathrm{Id}-P_\star)r_N\bigr)(j)=1-\varepsilon(j)(j+1)$, which lies in
    $[1-\gamma,1]$. At a ladder state $j$ with $2j>N\geq j$ the images carry $r_N(2j)=N$ and
    $r_N(j-1)=j-1$, so $(P_\star r_N)(j)=j-1+\varepsilon(j)(N-j+1)$ and the value is
    $1-\varepsilon(j)(N-j+1)$, again in $[1-\gamma,1]$ because $2j>N$ gives
    $N-j+1\leq j+1$. At a ladder state $j>N$ every image carries the value $N=r_N(j)$, so the
    value is $0$. At $s_0$ the only image is $s_f$ and $r_N(s_0)=r_N(s_f)=0$, so the value is
    $0$. At $s_f$ the target row is a probability supported in $\{1,\dots,d\}$ of mean
    $\bar\jmath$, on which $r_N=V$ because $N\geq d$, so the value is $-\bar\jmath$. Hence
    $(\mathrm{Id}-P_\star)r_N$ takes its values in $\bigl[-\max(\gamma-1,\bar\jmath),1\bigr]$
    and $\|(\mathrm{Id}-P_\star)r_N\|_{L^\infty(\lambda)}\leq\max(1,\gamma,\bar\jmath)$.
    \item \emph{Centring leaves the image unchanged and half of the clip.} Indeed,
    $0\leq r_N\leq N$ and $\lambda$ is a probability, so $\Pi r_N$ is the constant
    $\lambda(r_N)\in[0,N]$; $P_\star\mathbf 1=\mathbf 1$ gives
    $(\mathrm{Id}-P_\star)\tilde r_N=(\mathrm{Id}-P_\star)r_N$, and $\Pi^2=\Pi$ gives
    $\Pi\tilde r_N=0$. Since $0<\varepsilon(j)\leq\varepsilon_{\max}<1$ at every $j\geq1$, the
    chain is irreducible and $\lambda$ is positive at every state
    (Lemma~\ref{lem:doubling_irreducible}): an $L^\infty(\lambda)$-norm is a
    supremum over states. The values $r_N(s_0)=0$ and $r_N(N)=N$ are both attained, so
    $\tilde r_N$ takes the values $-\lambda(r_N)$ and $N-\lambda(r_N)$ and
    $\|\tilde r_N\|_{L^\infty(\lambda)}\geq\max\bigl(\lambda(r_N),N-\lambda(r_N)\bigr)\geq N/2$,
    which with the uniform bound on the image of the clipped ramp is \eqref{eq:doubling_ramp}.
    \item \emph{The infimum vanishes, and no bounded inverse exists.} Indeed,
    $\|\tilde r_N\|_{L^\infty(\lambda)}\geq N/2>0$, so
    $h_N:=\tilde r_N/\|\tilde r_N\|_{L^\infty(\lambda)}$ satisfies $\Pi h_N=0$,
    $\|h_N\|_{L^\infty(\lambda)}=1$ and, by \eqref{eq:doubling_ramp},
    $\|(\mathrm{Id}-P_\star)h_N\|_{L^\infty(\lambda)}\leq2\max(1,\gamma,\bar\jmath)/N$, which
    tends to $0$ as $N\rightarrow\infty$. If a bounded $S$ on $L^\infty(\lambda)$ satisfied
    $S(\mathrm{Id}-P_\star)=\mathrm{Id}-\Pi$, then
    $1=\|h_N\|_{L^\infty(\lambda)}\leq\|S\|\,\|(\mathrm{Id}-P_\star)h_N\|_{L^\infty(\lambda)}$
    for every $N$, which that limit forbids.
\end{itemize}
Finally $(j+1)\,\varepsilon_{c,s}(j)=c\,(j+1)^{1-s}$ is non-increasing in $j$ for $s\geq1$,
with supremum $c\,2^{1-s}$ at $j=1$, and is unbounded for $s<1$, $c$ being positive.
\end{proof}

\begin{corollary}[an explicit divergent family]\label{cor:doubling_family}
In the setting of Definition~\ref{def:doubling_setting}, let $s=1$ and $0<c<1$, so that the
loop-closed chain is positive recurrent (Proposition~\ref{prop:doubling_phase}\emph{(3)}) and
its invariant probability $\lambda$, its tail $L$ and its centred tail indicators $f_m$,
$m>d$, are those of that definition. Then there are $c_7>0$, depending only on $c$, $d$ and
$\lambda_1,\dots,\lambda_{2m_0}$ with $m_0$ the integer of
Theorem~\ref{theo:doubling_decay}, and an integer $m_2>d$ with
\begin{equation}\label{eq:doubling_explicit}
    \frac{\|f_m\|_{L^2(\lambda)}}
         {\bigl\|\left(\mathrm{Id}-P_\star\right)f_m\bigr\|_{L^2(\lambda)}}
    \;\geq\;c_7\sqrt m\qquad\text{for every }m\geq m_2 .
\end{equation}
\end{corollary}

\begin{proof}
The loop closure is irreducible and positive recurrent by
Proposition~\ref{prop:doubling_phase}\emph{(3)}, so $\lambda$ exists and is positive at every
state. Since $\varepsilon_{c,1}(j)=c/(j+1)\in(0,c/2]$ for every $j\geq1$ and $c/2<1$,
Proposition~\ref{prop:doubling_cut} gives \eqref{eq:doubling_cut} at every $m>d$, so
Theorem~\ref{theo:doubling_decay} supplies constants $0<c_1\leq c_2$ with
\eqref{eq:doubling_decay} and Corollary~\ref{cor:doubling_tail} the bounds
\eqref{eq:doubling_tailratio} and the limit $L(m)\rightarrow0$; and
Lemma~\ref{lem:doubling_percut} applies in its case \emph{(1)} at every cut $m>d$, its
$\varepsilon_{\max}$ being $\varepsilon_{c,1}(1)=c/2$. The denominator of
\eqref{eq:doubling_explicit} is positive by \eqref{eq:doubling_step1}, whose right side has
$L^2(\lambda)$-norm at least $\sqrt{\lambda_m}\,\bigl(1-\varepsilon(m)\bigr)>0$. Fix an
integer $m_2>d$ with $1-L(m)\geq\tfrac12$ for every $m\geq m_2$, possible because
$L(m)\rightarrow0$. For $m\geq m_2$, \eqref{eq:doubling_step3} and the second pair of
\eqref{eq:doubling_tailratio} give
\begin{equation}
    \frac{\left\|\left(\mathrm{Id}-P_\star\right)f_m\right\|_{L^2(\lambda)}}{\|f_m\|_{L^2(\lambda)}}
    \;\leq\;\sqrt{\frac{4\lambda_m}{L(m)}}
    \;\leq\;\frac{2}{\sqrt m}\sqrt{\frac{c_2(p_*-1)}{c_1}} ,
\end{equation}
which is \eqref{eq:doubling_explicit} with $c_7:=\tfrac12\sqrt{c_1/\bigl(c_2(p_*-1)\bigr)}$:
$c_1$ and $c_2$ depend on $c$, $d$ and $\lambda_1,\dots,\lambda_{2m_0}$ alone by
Theorem~\ref{theo:doubling_decay}, and $p_*$ on $c$ alone.
\end{proof}

\begin{proposition}[an $L^2$ pair with no exact $L^2$ outflow]\label{prop:doubling_unsolvable}
In the setting of Definition~\ref{def:doubling_setting}, let the backward policy be
\eqref{eq:doubling_policy} with $\varepsilon(j)>0$ for every $j\geq1$ and
$\varepsilon_{\max}=\sup_{j\geq1}\varepsilon(j)<1$, let the loop-closed chain be positive
recurrent with invariant probability $\lambda$, and assume \eqref{eq:doubling_star}. Then:
\begin{enumerate}
    \item $\left(\mathrm{Id}-P_\star\right)L^2(\lambda)$ is a dense proper linear subspace of
    $\ker\Pi$;
    \item there are probability densities $\finit,\fterm\in L^2(\lambda)$ such that no
    $f\in L^2(\lambda)$ satisfies $\left(\mathrm{Id}-P_\star\right)f=\finit-\fterm$.
\end{enumerate}
\end{proposition}

\begin{proof}
Write $A$ for the restriction of $\mathrm{Id}-P_\star$ to $\ker\Pi$. It is a bounded operator
of the Hilbert space $\ker\Pi$ into itself: $P_\star$ is a contraction of $L^2(\lambda)$ and
$P_\star=P^*$ by Lemma~\ref{lem:doubling_operator}\emph{(1)}, and $\Pi P_\star=\Pi$ gives
$\Pi(\mathrm{Id}-P_\star)=0$. Since $P_\star\mathbf 1=\mathbf 1$ gives
$(\mathrm{Id}-P_\star)\Pi=0$, the range of $A$ is
$\left(\mathrm{Id}-P_\star\right)L^2(\lambda)$.
\begin{itemize}
    \item \emph{The only fixed points of $P_\star$ and of $P$ in $L^2(\lambda)$ are the
    constants.} Indeed, $0<\varepsilon(j)\leq\varepsilon_{\max}<1$ at every $j\geq1$, so the
    loop closure is irreducible and $\lambda$ is positive at every state
    (Lemma~\ref{lem:doubling_irreducible}), and Lemma~\ref{lem:doubling_fixed_points} gives
    \eqref{eq:doubling_fixed}.
    \item \emph{$A$ is injective with dense range.} Indeed, $\ker A=\{f\in\ker\Pi:P_\star f=f\}$
    is $\{0\}$ since the only fixed points of $P_\star$ in $L^2(\lambda)$ are the constants
    (Lemma~\ref{lem:doubling_fixed_points}), a constant of zero mean being $0$. For
    $f,g\in\ker\Pi$,
    $\langle Af,g\rangle_\lambda=\langle f,(\mathrm{Id}-P)g\rangle_\lambda$ by $P_\star=P^*$,
    and $\int Pg\,d\lambda=\langle g,P_\star\mathbf 1\rangle_\lambda=\int g\,d\lambda=0$ puts
    $(\mathrm{Id}-P)g$ in $\ker\Pi$: the adjoint of $A$ in $\ker\Pi$ is the restriction of
    $\mathrm{Id}-P$, so $\ker A^*=\{g\in\ker\Pi:Pg=g\}=\{0\}$ by
    Lemma~\ref{lem:doubling_fixed_points} and the closure of the range of
    $A$ is $(\ker A^*)^\perp=\ker\Pi$, the orthogonal complement taken in $\ker\Pi$.
    \item \emph{$A$ is not surjective.} Indeed, \eqref{eq:doubling_inf} at $p=2$
    (Theorem~\ref{theo:doubling_unbounded}\emph{(2)}, whose hypotheses are those of this
    statement) gives
    $\inf\{\|Af\|_{L^2(\lambda)}:f\in\ker\Pi,\ \|f\|_{L^2(\lambda)}=1\}=0$. Were $A$
    surjective it would be a bounded bijection of the Hilbert space $\ker\Pi$ onto itself, hence
    with a bounded inverse by the open mapping theorem \citep{brezis2011functional}, and
    $\|f\|_{L^2(\lambda)}\leq\|A^{-1}\|\,\|Af\|_{L^2(\lambda)}$ on $\ker\Pi$ would contradict
    that infimum. With the density of the range of $A$ this is \emph{(1)}.
    \item \emph{A pair of probability densities has its difference outside the range of $A$.}
    Indeed, fix $\psi\in\ker\Pi$ outside the range of $A$, which exists since $A$ is not
    surjective, and put $Z:=1+\lambda(\psi^+)$, which equals $1+\lambda(\psi^-)$
    because $\lambda(\psi)=0$. The functions $\finit:=(\mathbf 1+\psi^-)/Z$ and
    $\fterm:=(\mathbf 1+\psi^+)/Z$ are non-negative, lie in $L^2(\lambda)$ --- $\psi^\pm$ being
    dominated by $|\psi|$ and $\lambda$ a probability --- and have integral $1$: they are
    probability densities, and $\finit-\fterm=-\psi/Z$. If some $f\in L^2(\lambda)$ satisfied
    $(\mathrm{Id}-P_\star)f=\finit-\fterm$, then $(\mathrm{Id}-P_\star)\Pi f=0$ would give
    $A(f-\Pi f)=-\psi/Z$ with $f-\Pi f\in\ker\Pi$, putting $\psi=-Z\,A(f-\Pi f)$ in the range of
    $A$; no such $f$ exists, which is \emph{(2)}.
\end{itemize}
\end{proof}

\begin{remark}[the condition \eqref{eq:doubling_star} indicts the family, not the graph]\label{rem:doubling_geometric}
Condition \eqref{eq:doubling_star} is not vacuous. For a geometrically decaying doubling
probability $\varepsilon(j)=a\rho^j$ with $\rho\in(0,1)$, the left side is
$D\log(1/a)+(2^D-1)\log(1/\rho)$, of order $2^D$, and \eqref{eq:doubling_star} fails. The
contrast is measured in \S\ref{sec:doubling_measured}, with a stated ceiling on it: no proof
of boundedness is claimed for any policy on this graph.
\end{remark}

\subsubsection{The truncation, and the proof of the main theorem}

This block proves item \emph{(5)} of Theorem~\ref{theo:doubling_main} and then assembles the
five items. Corollary~\ref{cor:doubling_truncation} restores a finite constant on the
truncation at an even cap, bounds the invariant measure there by two constants whose ratio does
not move with the cap, and turns the per-cut inequality into a growth of that constant in the
cap.

\begin{corollary}[the truncation, and a proved $\sqrt K$]\label{cor:doubling_truncation}
In the setting of Definition~\ref{def:doubling_setting}, let $s=1$ and $0<c<1$, and for an even
integer $K\geq d$ let $\lambda^K$, $L_K$ and $\widehat B_K$ be the invariant probability, the
tail and the constant \eqref{eq:doubling_Bhat} of the truncation at $K$. Then:
\begin{enumerate}
    \item for every even integer $K\geq d$ and every integer $m$ with $d<m\leq K/2$,
    \begin{equation}\label{eq:doubling_percut_K}
        \widehat B_K\;\geq\;\sqrt{\frac{L_K(m)\bigl(1-L_K(m)\bigr)}{2\,\lambda^K_m}} ;
    \end{equation}
    \item there are integers $m_0$ and $K_0\geq d$ and a constant $c_9\geq1$, depending only
    on $c$ and $d$, such that for every even integer $K\geq K_0$ there are constants
    $0<C^K_-\leq C^K_+$ with $C^K_+/C^K_-\leq c_9$ and, for every integer $m$ with
    $m_0\leq m\leq K/2$,
    \begin{equation}\label{eq:doubling_decayK}
        C^K_-\,m^{-p_*}\;\leq\;\lambda^K_m\;\leq\;C^K_+\,m^{-p_*} ;
    \end{equation}
    \item there is $c_8>0$, depending only on $c$, $d$ and $\bar\jmath$, with
    \begin{equation}\label{eq:doubling_sqrtK}
        \widehat B_K\;\geq\;c_8\sqrt K\qquad\text{for every even integer }K\geq K_0 ,
    \end{equation}
    $K_0$ being the integer of \emph{(2)}.
\end{enumerate}
\end{corollary}

\begin{proof}
\emph{Step 1 (the truncation and its diffusion operator).} Let $K\geq d$ be even. Since
$\varepsilon_{c,1}(j)=c/(j+1)\in(0,c/2]$ for every $j\geq1$ and $c/2<1$,
Lemma~\ref{lem:doubling_truncation_irreducible} makes the truncation at $K$ a finite
irreducible chain with a unique invariant probability $\lambda^K$, positive at every state,
and gives $S$ with \eqref{eq:doubling_resolvent} and $\widehat B_K<+\infty$.

\emph{Step 2 (the identities below the cut).} By the truncation clause of
Proposition~\ref{prop:doubling_cut}, $\lambda^K$ satisfies \eqref{eq:doubling_cut} at every
integer $m$ with $d<m\leq K/2$ and \eqref{eq:doubling_ratios} at every $j\geq1$ with $2j\leq K$
and every integer $m$ with $2\leq m\leq K$.

\emph{Step 3 (the per-cut inequality).} Let $d<m\leq K/2$. Case \emph{(2)} of
Lemma~\ref{lem:doubling_percut} holds at that cut, its $\varepsilon_{\max}$ being
$\varepsilon_{c,1}(1)=c/2$, so the function $f_m=\mathbf 1_{[m,K]}-L_K(m)$ has $\Pi f_m=0$ and
$\|f_m\|^2_{L^2(\lambda^K)}=L_K(m)\bigl(1-L_K(m)\bigr)>0$, and \eqref{eq:doubling_step3}
bounds its Rayleigh ratio. By \eqref{eq:doubling_resolvent},
$\|f_m\|=\|(\mathrm{Id}-\Pi)f_m\|=\|S(\mathrm{Id}-P_\star)f_m\|\leq\widehat B_K\|(\mathrm{Id}-P_\star)f_m\|$,
which with \eqref{eq:doubling_step3} is \eqref{eq:doubling_percut_K} and proves \emph{(1)}.

\emph{Step 4 (the decay of $\lambda^K$, with constants uniform in $K$).} Let $m_0$ be the
integer of Theorem~\ref{theo:doubling_decay}, which depends on $c$ and $d$ alone and exceeds
$d$, and put $J_0:=\{m_0,\dots,2m_0-1\}$ and $K_0:=8m_0+2$. Let $K\geq K_0$ be even, so that
$K/2$ is an integer with $K/2\geq4m_0+1\geq2m_0$, and let $(\lambda_j)_{j\geq1}$ be the
sequence with $\lambda_j:=\lambda^K_j$ for $1\leq j\leq K$ and $\lambda_j:=1$ for $j>K$,
positive because $\lambda^K$ is positive at every state
(Lemma~\ref{lem:doubling_truncation_irreducible}). By the truncation clause of
Proposition~\ref{prop:doubling_cut} it satisfies \eqref{eq:doubling_cut} at every integer $y$
with $d<y\leq K/2$, that identity at such a $y$ carrying the indices
$\lceil y/2\rceil,\dots,y$ alone, all of them at most $K$. Hence
Theorem~\ref{theo:doubling_decay}\emph{(1)} applies at $\ell=m_0$ and $M_1=K/2$ --- its
hypotheses being $\ell\geq m_0$, $M_1$ an integer at least $2\ell$, and the cut balance at
every integer $y$ with $d<y\leq M_1$ --- and gives, for every integer $m$ with
$2m_0\leq m\leq K/2$,
\begin{equation}
\begin{aligned}
    e^{-c_5c_4/m_0}\,\min_{j\in J_0}\lambda^K_j\,j^{p_*}\;&\leq\;\lambda^K_m\,m^{p_*}\\
    \;&\leq\;e^{c_5c_4/m_0}\,\max_{j\in J_0}\lambda^K_j\,j^{p_*},
\end{aligned}
\end{equation}
where $c_4$ is the constant of Lemma~\ref{lem:doubling_descent} and $c_5$ that of
Lemma~\ref{lem:doubling_escape}, the rescaled profile of
Definition~\ref{def:doubling_decay_notation} being $u_j=\lambda_j\,j^{p}$ with $p=p_*$ and
$\lambda_j=\lambda^K_j$ at every $j\leq K$. Put
$C^K_-:=e^{-c_5c_4/m_0}\min_{j\in J_0}\lambda^K_j\,j^{p_*}$ and
$C^K_+:=e^{c_5c_4/m_0}\max_{j\in J_0}\lambda^K_j\,j^{p_*}$, attained over the finite set $J_0$,
finite and strictly positive; $C^K_-\leq C^K_+$, the minimum being at most the maximum and
$c_4$ and $c_5$ positive. Then \eqref{eq:doubling_decayK} holds at every integer $m$ with
$m_0\leq m\leq K/2$: at $m\in J_0$ because $e^{-c_5c_4/m_0}\leq1\leq e^{c_5c_4/m_0}$, and at
$2m_0\leq m\leq K/2$ by the last display. Since $p_*>0$ and $m_0\leq j<2m_0$ on $J_0$, one has
$\max_{j\in J_0}\lambda^K_j\,j^{p_*}\leq(2m_0)^{p_*}\max_{j\in J_0}\lambda^K_j$ and
$\min_{j\in J_0}\lambda^K_j\,j^{p_*}\geq m_0^{p_*}\min_{j\in J_0}\lambda^K_j$, so
\begin{equation}
\begin{aligned}
    \frac{C^K_+}{C^K_-}\;&=\;e^{2c_5c_4/m_0}\,
    \frac{\max_{j\in J_0}\lambda^K_j\,j^{p_*}}{\min_{j\in J_0}\lambda^K_j\,j^{p_*}}\\
    \;&\leq\;e^{2c_5c_4/m_0}\,2^{p_*}\,
    \frac{\max_{j\in J_0}\lambda^K_j}{\min_{j\in J_0}\lambda^K_j} .
\end{aligned}
\end{equation}
Chaining the second inequality of \eqref{eq:doubling_ratios} across $J_0$ bounds
$\lambda^K_{j'}$ by $\lambda^K_j\prod_{i=j+1}^{j'}(1-\varepsilon(i))^{-1}$ for $j<j'$ in $J_0$,
and dropping all but the term $j=i-1$ of \eqref{eq:doubling_cut} at
$i\in\{m_0+1,\dots,2m_0-1\}$, legitimate because $d<m_0<i\leq2m_0-1\leq K/2$, gives
$\lambda^K_i\geq\lambda^K_i(1-\varepsilon(i))\geq\lambda^K_{i-1}\varepsilon(i-1)$, hence
$\lambda^K_{j'}\geq\lambda^K_j\prod_{i=j}^{j'-1}\varepsilon(i)$: the ratio
$\max_{j\in J_0}\lambda^K_j/\min_{j\in J_0}\lambda^K_j$ is bounded by a constant depending only
on $c$ and $m_0$. Since $m_0$ depends on $c$ and $d$ alone, $p_*$ on $c$ alone, and $c_4$ and
$c_5$ on $c$ and $d$ alone by Lemmas~\ref{lem:doubling_descent}
and~\ref{lem:doubling_escape}, there is $c_9\geq1$ depending only on $c$ and $d$ with
$C^K_+/C^K_-\leq c_9$ for every even $K\geq K_0$; and $K_0=8m_0+2$ exceeds $d$ and depends on
$c$ and $d$ alone, $m_0$ doing so. This proves \emph{(2)}.

\emph{Step 5 (the ratio at a cut of order $K$).} Let $K\geq K_0$ be even and
$m:=\lfloor K/4\rfloor$, so that $m\geq2m_0>d$ and $2m\leq K/2$. By \eqref{eq:doubling_decayK}
on $\{m,\dots,2m\}$,
\begin{equation}
\begin{aligned}
    L_K(m)\;&\geq\;\sum_{j=m}^{2m}\lambda^K_j\;\geq\;C^K_-\,(m+1)\,(2m)^{-p_*}\\
    &\geq\;C^K_-\,2^{-p_*}m^{1-p_*},
\end{aligned}
\end{equation}
while $\lambda^K_m\leq C^K_+m^{-p_*}$, so $L_K(m)/\lambda^K_m\geq2^{-p_*}c_9^{-1}m$.

\emph{Step 6 (the mass below the cut).} By Step 1 the truncation is a finite irreducible
chain, so $\lambda^K(s_0)$ is the reciprocal of the expected return time of that chain to
$s_0$ \citep{levin2017markov}. A return from $s_0$ is the wrap $s_0\rightarrow s_f$, the step
from $s_f$ to a state of the target row, and the run from that state to $s_0$, so its expected
length is $2+\bar\sigma_K$ with $\bar\sigma_K$ the functional of
Lemma~\ref{lem:doubling_supersolution}. Stationarity at $s_0$ writes $\lambda^K(s_0)$ as the
flux into $s_0$, which comes from the state $1$ alone and equals
$\lambda^K_1\bigl(1-\varepsilon(1)\bigr)$, the doubling edge out of $1$ being present because
$2\leq K$. Hence $\lambda^K_1\geq\lambda^K(s_0)=1/(2+\bar\sigma_K)$, and
$\bar\sigma_K\leq\bar\jmath/(1-c)$ by Lemma~\ref{lem:doubling_supersolution}, so for every
$m\geq2$
\begin{equation}
    1-L_K(m)\;\geq\;\lambda^K_1\;\geq\;\Bigl(2+\frac{\bar\jmath}{1-c}\Bigr)^{-1},
\end{equation}
a bound free of $K$.

\emph{Step 7 (assembly).} Let $K\geq K_0$ be even and $m:=\lfloor K/4\rfloor$, so that
$m\geq K/8$, $K$ being at least $8$. Then $d<m\leq K/2$, and \eqref{eq:doubling_percut_K} with
Steps 5 and 6 gives
\begin{equation}
\begin{aligned}
    \widehat B_K\;&\geq\;\sqrt{\frac{2^{-p_*}c_9^{-1}m}{2\bigl(2+\bar\jmath/(1-c)\bigr)}}
    \;\geq\;c_8\sqrt K,\\
    c_8\;&:=\;\Bigl(2^{\,p_*+4}\,c_9\,\bigl(2+\bar\jmath/(1-c)\bigr)\Bigr)^{-1/2} ,
\end{aligned}
\end{equation}
which is \eqref{eq:doubling_sqrtK} and proves \emph{(3)}: $p_*$ depends on $c$ alone and
$m_0$ and $c_9$ on $c$ and $d$ alone, so $c_8$ depends only on $c$, $d$ and $\bar\jmath$, with
$K_0$ the integer of \emph{(2)}.
\end{proof}

\begin{proof}[Proof of Theorem \ref{theo:doubling_main}]
\emph{(1)} Irreducibility of the loop closure on the whole standing range is
Proposition~\ref{prop:doubling_phase}. Row \emph{(a)} of the table is
Proposition~\ref{prop:doubling_phase}\emph{(2)}; rows \emph{(b)}, \emph{(c)}, \emph{(d)} and
\emph{(e)} are Proposition~\ref{prop:doubling_phase}\emph{(3)}, which classifies every
$c\in(0,2)$ but $1/\ln2$ and denies positive recurrence at every $1\leq c<2$; and row
\emph{(f)} is Proposition~\ref{prop:doubling_phase}\emph{(1)}. Rows \emph{(b)} and \emph{(f)}
are therefore exactly the positive recurrent members of the standing range.

\emph{(2)} The chain is positive recurrent by
Proposition~\ref{prop:doubling_phase}\emph{(3)}, so $\lambda$ exists, and $\lambda>0$ at every
state by irreducibility; $p_*>1$ because $c<1$ (Lemma~\ref{lem:doubling_cramer_root}). Since
$\varepsilon_{c,1}(j)\in(0,1)$ for every $j\geq1$, Proposition~\ref{prop:doubling_cut} gives
\eqref{eq:doubling_cut} for every $m>d$, so the positive sequence $(\lambda_j)_{j\geq1}$
satisfies the hypotheses of Theorem~\ref{theo:doubling_sharp}, which supplies the limit $C$ of
\eqref{eq:doubling_sharp} and its dependence on $c$, $d$ and $\lambda_1,\dots,\lambda_d$, and
of Corollary~\ref{cor:doubling_tail_sharp}, which gives the first line of
\eqref{eq:doubling_tail_sharp}. That limit is positive: Theorem~\ref{theo:doubling_decay}
supplies a pair $0<c_1\leq c_2$ with \eqref{eq:doubling_decay} and
Theorem~\ref{theo:doubling_sharp} puts $C$ in $[c_1,c_2]$;
Proposition~\ref{prop:doubling_constant} identifies it as
\eqref{eq:doubling_constant}.

\emph{(3)} On the standing range $\varepsilon=\varepsilon_{c,s}$ takes its values in
$(0,\varepsilon_{\max}]$ with $\varepsilon_{\max}<1$ (Lemma~\ref{lem:doubling_range}),
and \eqref{eq:doubling_star} holds by
Theorem~\ref{theo:doubling_unbounded}\emph{(1)}; at a positive recurrent member the remaining
hypothesis of Theorem~\ref{theo:doubling_unbounded}\emph{(2)} holds as well, so
\eqref{eq:doubling_inf} and the exclusion of a bounded $S$
hold at every $p\in[1,+\infty)$, the operator $\mathrm{Id}-P_\star$ acting on each
$L^p(\lambda)$ by Lemma~\ref{lem:doubling_operator}\emph{(1)}. Rows \emph{(b)} and \emph{(f)}
have $s\geq1$, so $\gamma=\sup_{j\geq1}(j+1)\varepsilon_{c,s}(j)=c\,2^{1-s}$ is
finite and Lemma~\ref{lem:doubling_ramp} gives the same two conclusions at $p=+\infty$. At
$p=2$ the excluded operator is the $S$ of \eqref{eq:doubling_Bhat}, so $\widehat B=+\infty$ by
the convention of Definition~\ref{def:doubling_setting}, and the two remaining consequences are
those of Theorem~\ref{theo:doubling_unbounded}\emph{(2)}. Finiteness of $\bar\sigma$ at $s=1$,
$c<1$ is
Proposition~\ref{prop:doubling_length}, and the inequality with $c_7$ and $m_2$ is
\eqref{eq:doubling_explicit} of Corollary~\ref{cor:doubling_family}.

\emph{(4)} At a positive recurrent member of the standing range the hypotheses of
Proposition~\ref{prop:doubling_unsolvable} hold: $\varepsilon=\varepsilon_{c,s}$ takes its
values in $(0,\varepsilon_{\max}]$ with $\varepsilon_{\max}<1$
(Lemma~\ref{lem:doubling_range}), and \eqref{eq:doubling_star} holds by
Theorem~\ref{theo:doubling_unbounded}\emph{(1)}. Item \emph{(2)} of
Proposition~\ref{prop:doubling_unsolvable} is the claim.

\emph{(5)} Since $\varepsilon_{c,1}(j)=c/(j+1)\in(0,1)$ for every $j\geq1$,
Lemma~\ref{lem:doubling_truncation_irreducible} makes the truncation at an even $K\geq d$ a
finite irreducible chain on $K+2$ states with a unique invariant probability $\lambda^K$ and
$\widehat B_K<+\infty$. The constants $c_8$ and $K_0$ are those of \eqref{eq:doubling_sqrtK} in
Corollary~\ref{cor:doubling_truncation}\emph{(3)}.
\end{proof}

Corollary~\ref{cor:doubling_truncation} is the bridge to the measurements, and it is loose in
its exponent: what \eqref{eq:doubling_sqrtK} proves is a growth $\sqrt K$, and what
\S\ref{sec:doubling_cost} measures is $K$.

\subsection{The transition, and the two bounds}\label{sec:doubling_measured}

This subsection measures on the doubling graph what is predicted for it and what is claimed
of every finite chain: the closed form for the expected backward length of
Proposition~\ref{prop:doubling_length}, the decay exponent of
Proposition~\ref{prop:doubling_exponent}, the stable bound of
Theorem~\ref{theo:RL_CV_bound_full}, and the universality certificate of
Theorem~\ref{theo:universality_L2_full} with the self-entropy of
Theorem~\ref{theo:IL_CV_bound} that prices it. An infinite matrix has no singular values, so
every constant below is computed on the truncation at $K$ of
Definition~\ref{def:doubling_setting}, which
Lemma~\ref{lem:doubling_truncation_irreducible} makes a finite irreducible chain on $K+2$
states carrying the invariant probability $\lambda^K$, the tail $L_K$ and the constant
$\widehat B_K$. Every measurement below uses $d=8$, a target row of mean $\bar\jmath=4.5$,
and a cap $K$ that is a power of two, hence even as that definition requires. A quantity that
settles as $K$ grows is
finite on the infinite graph; one that grows with $K$ is a divergence the truncation was
hiding. Everything below is exact linear algebra on that chain --- linear solves, a singular
value, and the exact terminal law of the inference chain --- with no training, no sampling and
no estimator.

Two scope limits attach to the numbers. The evidence code caps $\varepsilon$ at $0.9$; the cap
binds at $c=2.5$, $s=1$, and there only at $j=1$, where it replaces the value $1.25$ --- not a
probability --- by $0.9$ and the drift $c-1=1.5$ of Proposition~\ref{prop:doubling_drift} by
$0.800$, so that column is a measurement on the capped chain, outside the standing range
$0<c<2^s$ of Definition~\ref{def:doubling_setting} and hence outside every statement of this
appendix that assumes that range --- Proposition~\ref{prop:doubling_drift} and
Theorem~\ref{theo:doubling_main} among them. It is legitimate evidence for the decay exponent
all the same. Proposition~\ref{prop:doubling_exponent} reads \eqref{eq:doubling_cut} as an
identity between numbers and puts no upper bound on $c$, so it reaches the capped chain; and
the cap does not touch that identity, since at $d=8$ every window of \eqref{eq:doubling_cut}
with $m>d$ starts at $\lceil9/2\rceil=5$ and $\varepsilon(1)$ --- the only value the cap moves
--- enters it nowhere, the capped chain satisfying it with $\varepsilon_{2.5,1}$ unchanged.
What the cap moves is the boundary data $\lambda_1,\dots,\lambda_d$, hence the constant and
not the exponent. And the two constants written $\widehat B$, the operator norm
\eqref{eq:doubling_Bhat} and the mixing sum of Lemma~\ref{lem:sigma_mixing}, are numerically
distinct --- $17.23$ against $18.39$ at $c=1/4$, $K=64$, in the order
Lemma~\ref{lem:doubling_operator}\emph{(3)} fixes --- the tables below reporting the first.

\paragraph{The closed form is exact, and what is left is truncation.}
Proposition~\ref{prop:doubling_length} predicts $\bar\sigma=\bar\jmath/(1-c)=4.5/(1-c)$, and
Table~\ref{tab:doubling_length} measures it at $K=4096$. The residual at $c=0.75$ and $c=0.90$
is truncation, not error: the deficit decays like $K^{-(p_*-1)}$, and at $c=0.9$ it is $7.735$
at $K=2048$ and $6.420$ at $K=4096$, a ratio of $1.205$ against $2^{p_*-1}=2^{0.2691}=1.205$ ---
and past $1/\ln2$, where $p_*$ turns negative, the same exponent is a growth, $\bar\sigma$ rising
over $K=64$ to $4096$ with fitted exponent $1.018$ against $1-p_*=1.0146$ at $c=1.45$ and $2.76$
against $2.7638$ at $c=2.5$.

\begin{table}[tb]
\centering
\small
\begin{tabular}{@{}lccc@{}}
\hline
$c$ & closed form & measured & relative error \\
\hline
$0.25$ & $6$  & $6.000000$ & $5.8\cdot10^{-10}$ \\
$0.50$ & $9$  & $8.99991$  & $9.7\cdot10^{-6}$ \\
$0.75$ & $18$ & $17.8927$  & $6.0\cdot10^{-3}$ \\
$0.90$ & $45$ & $38.5808$  & $1.4\cdot10^{-1}$ \\
\hline
\end{tabular}
\caption{Expected backward length on the truncation at $K=4096$, against the closed form
$\bar\jmath/(1-c)$ of Proposition~\ref{prop:doubling_length}.}
\label{tab:doubling_length}
\end{table}

\paragraph{The decay exponent is the Cram\'er root, and the growth of $\widehat B_K$ ignores
the transition.}
The exponent the theory identifies is the non-zero root $p_*$ of \eqref{eq:doubling_cramer},
and the fit below is empirical: Proposition~\ref{prop:doubling_exponent} is a statement about
an infinite sequence with a polynomial asymptotic, and $\lambda^K$ is a finite vector with no
$j\rightarrow\infty$ behaviour to have. What transfers to $\lambda^K$ is
Corollary~\ref{cor:doubling_truncation}\emph{(2)}: at $c<1$ and past its threshold $K_0$ it
gives the two-sided bound \eqref{eq:doubling_decayK}, whose two constants have a ratio bounded
uniformly in $K$ --- enough to pin the exponent, not to read it to three decimals. Fitted on the dyadic
band $[K/16,K/2]$ at $K=4096$, the decay exponent of $\lambda^K$ matches $p_*$ to three
decimals at all eight values of $c$, negative roots included
(Table~\ref{tab:doubling_exponents}). Over the same
range $\widehat B_K$ doubles with every doubling of the truncation, at every $c$, on both
sides of both thresholds of
Proposition~\ref{prop:doubling_phase}\emph{(3)}. The $c=0.25$ row is the one that separates
the two conditions: there the expected backward length is $6.000$ and the condition of
\citet{morozov2025revisiting} holds comfortably, and $\widehat B_K$ grows all the same. The
contrast Remark~\ref{rem:doubling_geometric} predicts appears where \eqref{eq:doubling_star}
fails: for the geometric $\varepsilon(j)=a\rho^j$ at $a=0.5$ and $\rho=0.9$, $\widehat B_K$ is
$12.74$, $20.74$, $27.12$, $27.75$ at $K=32$, $64$, $128$, $256$ --- saturating, and not
resolved past $K=256$, where the invariant mass of the high states falls below what the
stationary solve separates from zero.

\begin{table}[tb]
\centering
\small
\begin{tabular}{@{}lccccccc@{}}
\hline
 & \multicolumn{2}{c}{decay exponent} & \multicolumn{4}{c}{$\widehat B_K$} & \\
\cline{2-3}\cline{4-7}
$c$ & predicted $p_*$ & fitted & $K=64$ & $K=256$ & $K=1024$ & $K=4096$ & fitted exponent \\
\hline
$0.25$ & $4.1314$  & $4.130$  & $17.23$ & $66.4$  & $262.7$ & $1048$ & $0.989$ \\
$0.50$ & $2.6599$  & $2.659$  & $26.52$ & $110.8$ & $444.0$ & $1774$ & $1.008$ \\
$0.75$ & $1.7189$  & $1.719$  & $35.71$ & $170.9$ & $735.3$ & $3026$ & $1.064$ \\
$0.90$ & $1.2691$  & $1.269$  & $39.33$ & $199.5$ & $910.5$ & $3936$ & $1.104$ \\
$1.00$ & $1.0000$  & $1.000$  & $40.32$ & $205.8$ & $951.1$ & $4172$ & $1.113$ \\
$1.20$ & $0.5161$  & $0.516$  & $38.76$ & $186.2$ & $815.2$ & $3412$ & $1.074$ \\
$1.45$ & $-0.0146$ & $-0.015$ & $32.75$ & $141.1$ & $576.9$ & $2321$ & $1.022$ \\
$2.50$ & $-1.7638$ & $-1.765$ & $14.78$ & $59.27$ & $237.2$ & $949.2$ & $1.001$ \\
\hline
\end{tabular}
\caption{The decay exponent of $\lambda^K$, fitted on the dyadic band $[K/16,K/2]$ at
$K=4096$, against the non-zero root $p_*$ of \eqref{eq:doubling_cramer}; and $\widehat B_K$
over six octaves, with its fitted exponent in $K$.}
\label{tab:doubling_exponents}
\end{table}

\paragraph{The stable bound holds in every row, and its tightness does not see the truncation.}
Theorem~\ref{theo:RL_CV_bound_full} was checked on $9{,}216$ rows: eight values of $c$, four
truncations ($K=64$ to $512$), four targets crossed with three initial flows, eight
perturbation scales and three exponents $q$. The bound holds in every one of them, and it is
non-vacuous --- right-hand side below $1$, the only way it says anything, since $\TV\leq1$
always --- in $2{,}115$, at the smaller perturbations. Its measured tightness at $q=1$ --- the
right-hand side of \eqref{equ:stable_bound}, whose constant is \eqref{equ:stable_constants},
divided by $\TV$ --- moves by less than $10\%$ over the truncation:
$6.33$, $6.00$, $6.59$, $6.47$ at $K=64$, $128$, $256$, $512$, with a best ratio over the
whole sweep of $6.00$. That is the band measured in
Section~\ref{sec:experiments} across six orders of magnitude of state space, here on a graph
whose every mixing constant is infinite. At $q=2$ and $q=+\infty$ the tightness degrades like
$N^{1-1/q}$ in the number of states $N=K+2$, from $22.5$ to $64.2$ and from $120$ to $1024$
over the same range.

\paragraph{The certificate closes exactly; its bill is set by how far up the ladder the
discrepancy reaches.}
For every pair $(\Finit,\target)$ of the sweep, the certificate
$\foutstareta=(\eta-S\theta)^+$ of Theorem~\ref{theo:universality_L2_full} at
$\eta:=\|S\theta\|_{L^\infty(\lambda)}$ closes flow matching to machine precision --- largest
residual over the sweep $6.0\cdot10^{-8}$, on the cell whose outflow is $1.3\cdot10^9$, a
relative $5\cdot10^{-17}$ --- and its $\LossWKLFM$ loss equals the self-entropy $\mathcal H$
of Theorem~\ref{theo:IL_CV_bound} computed against $\nuB=\lambda$. Its outflow bill is
$\int\foutstareta\,d\lambda=\eta$ exactly in every row, as $\Pi S=0$ requires, and
Table~\ref{tab:doubling_bill} gives its growth in the cap at $c=0.5$
(Figure~\ref{fig:doubling}, right). The unbounded operator charges nothing at all when
$\theta$ is supported at the foot of the ladder, and $K^{2.9}$ as soon as it reaches into it
--- faster than $\widehat B_K$ itself, because $\eta$ is an $L^\infty$ norm and the states
attaining it carry $\lambda\sim K^{-p_*}$. These natural pairs do not saturate the operator:
their amplification $\|S\theta\|/\|\theta\|$ reaches $0.53$ of $\widehat B_K$ at $K=64$ and
$0.38$ at $K=512$. Saturating it is what the critical direction of \S\ref{sec:doubling_cost}
does.

\begin{table}[tb]
\centering
\small
\begin{tabular}{@{}p{0.30\linewidth}ccccc@{}}
\hline
initial flow $\rightarrow$ target & $K=64$ & $K=128$ & $K=256$ & $K=512$ & slope \\
\hline
$\delta_{s_0}\rightarrow$ target row (foot of the ladder) & $11.35$ & $11.35$ & $11.35$ & $11.35$ & $\mathbf{0.00}$ \\
$\delta_{s_0}\rightarrow\lambda$ & $78.7$ & $167$ & $343$ & $697$ & $1.05$ \\
$\delta_{s_0}\rightarrow$ uniform on the ladder & $7641$ & $5.76\cdot10^4$ & $4.25\cdot10^5$ & $3.09\cdot10^6$ & $2.89$ \\
$\delta_{s_0}\rightarrow$ top half of the ladder & $1.44\cdot10^4$ & $1.10\cdot10^5$ & $8.24\cdot10^5$ & $6.03\cdot10^6$ & $2.90$ \\
\hline
\end{tabular}
\caption{The certificate's outflow bill $\eta=\|S\theta\|_{L^\infty(\lambda)}$ at $c=0.5$, by
$(\Finit,\target)$ pair and truncation, $\delta_{s_0}$ being the Dirac mass at the source,
with the slope fitted over $K=64$ to $512$.}
\label{tab:doubling_bill}
\end{table}

\subsection{An unbounded diffusion operator}\label{sec:doubling_cost}

The extremal direction of $S$, the lower bound \eqref{eq:doubling_percut_K} at its best cut and
the price of both are measured on $40$ cells: eight values of $c$ crossed with five
truncations, $K=64$ to $K=1024$. The extremal direction is a unit vector $\theta_*$ of
$L^2(\lambda^K)$ with $\|S\theta_*\|_{L^2(\lambda^K)}=\widehat B_K$, attained because the
truncation is finite (Lemma~\ref{lem:doubling_truncation_irreducible}).

\paragraph{The proved bound holds everywhere, and is loose by exactly the factor it drops.}
Inequality \eqref{eq:doubling_percut_K} holds in every one of the $40$ cells. Its fitted
exponent in $K$ is $0.473$ to $0.510$ --- $0.498$, $0.509$, $0.500$, $0.500$ at $c=0.25$,
$1.0$, $1.45$, $2.5$ --- against a measured exponent of $\widehat B_K$ between $0.983$ and
$1.139$. The best cut therefore realizes the exponent \eqref{eq:doubling_sqrtK} proves and no
more: Corollary~\ref{cor:doubling_truncation} is tight in its own exponent and loose against
the measured $\widehat B_K\sim K$ by a factor $\sqrt K$, the ratio of $\widehat B_K$ to its
certified lower bound running from $7.79$ at $K=64$ to $87.38$ at $K=1024$ and growing like
$K^{0.51}$ (Figure~\ref{fig:doubling}, left).

\paragraph{The critical direction is a centred tail indicator.}
The cut $m_*$ maximizing \eqref{eq:doubling_percut_K} sits at $m_*/K=0.32$ to $0.49$, stable
across six octaves, and the overlap between the measured $\theta_*$ and the centred
indicator $\mathbf 1_{[m_*,K]}-\lambda^K([m_*,K])$ --- the trial function $f_{m_*}$ of
Lemma~\ref{lem:doubling_percut} --- is
$0.71$ to $0.84$, stable in $K$ ($0.719$ and $0.719$ at $c=0.25$ for $K=512$ and $1024$;
$0.826$ at $c=0.9$), and clears the registered $0.8$ at four of the eight values of $c$. The
sharp indicator captures three quarters of the critical direction, and the true maximizer is
that indicator smoothed over a band.

\paragraph{Nothing diverges at finite $K$; the cost does.}
Table~\ref{tab:doubling_cost} fits, in $K$ over $K=64$ to $1024$, the exponents of
$\widehat B_K$, of its lower bound \eqref{eq:doubling_percut_K}, of the certificate level
$\eta_*:=\|S\theta_*\|_{L^\infty(\lambda)}$ along the critical direction, and of the training
rate $\varrho=g''(1)w_{\min}/\widehat B_K^{\,2}$ of
Theorem~\ref{theo:db_stable_frozen_full} in its flow-matching instance, at its divergence
generator $g=(\log x)^2$ and its training weight $w\geq w_{\min}$ taken to be $1$. That
theorem holds at every constant $\widehat B\geq1$ satisfying \eqref{eq:coercivity}, and the
operator norm \eqref{eq:doubling_Bhat} of the truncation is one. Taking adjoints in
\eqref{eq:doubling_resolvent} (Lemma~\ref{lem:doubling_operator}\emph{(2)}), with
$P_\star=P^*$ by Lemma~\ref{lem:adjoint}\emph{(3)} and $\Pi$ self-adjoint, gives
$S^*(\mathrm{Id}-P)=\mathrm{Id}-\Pi$ with $\|S^*\|=\|S\|=\widehat B_K$, so
$\|h-\Pi h\|_{L^2(\lambda^K)}\leq\widehat B_K\,\|(\mathrm{Id}-P)h\|_{L^2(\lambda^K)}$ for every
$h\in L^2(\lambda^K)$ --- \eqref{eq:coercivity} on the truncation --- and $\widehat B_K\geq14.78$
in all $40$ cells, the smallest value at $c=2.5$, $K=64$. For the detailed-balance loss,
Lemma~\ref{lem:lift_coercivity} supplies the constant $1+\widehat B_K$ instead. The operator
norm never exceeds the mixing sum, $\|S\|\leq\sum_n\widehat\beta_n$
(Lemma~\ref{lem:doubling_operator}\emph{(3)}), so the rate read at $\widehat B_K$ is at least
the one summable mixing supplies through Lemma~\ref{lem:sigma_mixing}.

\begin{table}[tb]
\centering
\small
\begin{tabular}{@{}lcccc@{}}
\hline
$c$ & $\widehat B_K$ & lower bound & $\eta_*$ & $\varrho$ \\
\hline
$0.25$ & $0.983$ & $0.473$ & $2.545$ & $-1.966$ \\
$0.50$ & $1.015$ & $0.481$ & $1.821$ & $-2.030$ \\
$0.75$ & $1.089$ & $0.495$ & $1.413$ & $-2.179$ \\
$0.90$ & $1.132$ & $0.507$ & $1.265$ & $-2.264$ \\
$1.00$ & $1.139$ & $0.510$ & $1.190$ & $-2.278$ \\
$1.20$ & $1.097$ & $0.504$ & $1.077$ & $-2.195$ \\
$1.45$ & $1.034$ & $0.497$ & $1.076$ & $-2.067$ \\
$2.50$ & $1.001$ & $0.503$ & $1.008$ & $-2.002$ \\
\hline
\end{tabular}
\caption{Exponents in $K$, fitted over $K=64$ to $1024$: the constant $\widehat B_K$, its
certified lower bound \eqref{eq:doubling_percut_K} at the best cut, the certificate level
$\eta_*$ along the critical direction, and the training rate $\varrho$ of
Theorem~\ref{theo:db_stable_frozen_full} in its flow-matching instance, read at $\widehat B_K$.}
\label{tab:doubling_cost}
\end{table}

Let $G_K$ be the finite marked graph with source $s_0$, ladder $\{1,\dots,K\}$, sink $s_f$,
the forward edges $s_0\rightarrow1$, $j\rightarrow j+1$ for $j<K$ and $2j\rightarrow j$ for
$2j\leq K$, and a terminating edge $j\rightarrow s_f$ from each state $j\leq d$ at which the
target row is positive. It is path-connected in the sense of
Definition~\ref{def:path_connected}: the increments carry $s_0$ to every ladder state, every
ladder state halves down to $1$ --- after one increment when it is odd, the top state $K$
carrying the edge $K\rightarrow K/2$ because $K$ is even --- and the increments carry $1$ to a
state of the target row's support, which terminates. The backward policy of the truncation at
$K$ --- \eqref{eq:doubling_policy} at the states $j$ with $2j\leq K$, the decrement with
probability $1$ above them --- is positive on every edge of the loop closure of $G_K$, whose
backward chain is that truncation, so Theorem~\ref{theo:universality_graphs} supplies an
exactly balanced flow whatever the mixing, and the measurement finds it: the certificate's
flow-matching residual and the total variation of its terminal law are at machine precision
at every cap --- $\TV\leq1.2\cdot10^{-8}$ at worst, on the cell whose outflow is $10^9$, a
relative $10^{-16}$, and $5\cdot10^{-15}$ at $K=1024$ for the foot-loaded target. What
degrades is the bill and the clock: $\eta_*$ grows like $K^{1.0}$ to $K^{2.5}$ and $\varrho$
collapses like $K^{-2}$. The hitting-time envelope degrades with it and for a different
reason. At every ladder state $\sigma\geq\max_{0\leq n\leq\sigma}X_n$
(Lemma~\ref{lem:doubling_excursion}), and $\sigma_*$ maximizes $\mathbb E(\sigma\mid X_0=x)$ over the
states, the top state $K$ among them, so $\sigma_*\geq K$ on the truncation at $K$. The
envelope $\widehat B_\sigma=\sigma_*\lambda_{\min}^{-1/2}$ of
Proposition~\ref{prop:morozov_rate}, read on $G_K$ with
$\lambda_{\min}=\min_x\lambda^K(x)$ the smallest invariant mass \eqref{eq:occupation}, is
then at least $K$ as well --- $\lambda_{\min}\leq1$ --- and diverges with the cap even at
$c<1$, where $\bar\sigma$ converges to $\bar\jmath/(1-c)$. Theorem~\ref{theo:RL_CV_bound_full},
which carries no mixing hypothesis, is untouched throughout: it holds in all $9{,}216$ rows
of \S\ref{sec:doubling_measured} at every cap.

Divergence is therefore not observable on this graph. Universality holds at every finite cap,
and in the limit what fails is the hypothesis of Theorem~\ref{theo:universality_L2};
convergence does not stop but slows quadratically in the cap; and the certificate
of Theorem~\ref{theo:IL_CV_bound} stays exact while its outflow bill runs away. What the study
does not reach is the infinite chain itself --- every number above is a truncation, and the
statements about the limit are Theorem~\ref{theo:doubling_unbounded} and
Proposition~\ref{prop:doubling_unsolvable}. What the second of them settles is that
\emph{exact} $L^2$ flow matching fails there, for a pair of $L^2$ densities the loop closure
cannot balance at all; whether the \emph{weak} form fails --- a defect infimum of zero, which
the range of $\mathrm{Id}-P_\star$ being dense in $\ker\Pi$ leaves open --- this appendix does
not decide.

\begin{figure}[tb]
\centering
\includegraphics[width=\linewidth]{figures/fig8_diffusion.pdf}
\caption{\textbf{Left:} $\widehat B_K$ against the truncation $K$, log-log, one curve per
$c$, with the proved lower bound \eqref{eq:doubling_percut_K} at its best cut as a dashed line.
The bound holds in all $40$ cells and its measured slope is half that of $\widehat B_K$
itself --- $K^{0.47}$ to $K^{0.51}$ against $K^{0.98}$ to $K^{1.14}$ ---
so the gap between the proved and the measured growth is itself a $\sqrt K$.
\textbf{Right:} the certificate bill
$\eta=\|S\theta\|_{L^\infty(\lambda)}$ against $K$ for four $(\Finit,\target)$ pairs at
$c=0.5$: flat at $11.35$ for the target row at the foot of the ladder, and $K^{2.9}$ for the
two ladder-loaded targets. The unbounded operator charges nothing where $\theta$ does not
reach up the ladder.}
\label{fig:doubling}
\end{figure}

%% Appendix I -- the doubling graph, certified: a machine-checked theorem and the ideas behind it.
\section{\appendixItitle}\label{app:lean}

This appendix carries the mathematics of Appendix~\ref{app:doubling} --- the doubling graph
\eqref{eq:counterexample_graph} and the family of backward policies it carries --- with the
proofs delegated to a machine-checked certificate rather than written out in natural language,
and it is meant to replace that appendix rather than to stand beside it.
Theorem~\ref{theo:lean_main} is the certified statement; it is narrower than
Theorem~\ref{theo:doubling_main} in six places, which \S\ref{sec:lean_scope} names, with
where each is proved instead when it is proved at all. \S\ref{sec:lean_model} quotes the formal model the
theorem quantifies over, \S\ref{sec:lean_certificate} the fifteen Lean theorems that carry it
and the axioms they rest on, and \S\ref{sec:lean_ideas} says, in prose and without proving
anything, how each is proved. The numbers measured on truncations are not repeated here: they
stay in \S\ref{sec:doubling_measured} and \S\ref{sec:doubling_cost}, which this appendix does
not replace.

\begin{definition}[standing notation and hypotheses of this appendix]\label{def:lean_setting}
The graph is \eqref{eq:counterexample_graph}, carrying the source $s_0$, the ladder
$1,2,3,\dots$, the sink $s_f$ and the forward edges $s_0\rightarrow1$, $j\rightarrow j+1$,
$2j\rightarrow j$ and $j\rightarrow s_f$; the predecessors of a ladder state $j$ are $j-1$ and
$2j$, so the backward chain at $j$ either \emph{decrements} to $j-1$ or \emph{doubles} to $2j$,
and the \emph{loop closure} adds the wrap $s_f\rightarrow s_0$. The backward policy is
\begin{equation}\label{eq:lean_policy}
    \BackwardPolicy(j\rightarrow2j)=\varepsilon(j),\qquad
    \BackwardPolicy(j\rightarrow j-1)=1-\varepsilon(j)\quad(j\geq1),\qquad
    \BackwardPolicy(s_0\rightarrow s_f)=1,
\end{equation}
with $j-1=0$ read as $s_0$, and the \emph{target row}
$\widehat\BackwardPolicy(s_f\rightarrow\cdot)$ is a probability supported in $\{1,\dots,d\}$ of
mean $\bar\jmath$, the integer $d\geq1$ being a parameter of the setting and $\bar\jmath$
the mean of its row. The family studied throughout is
\begin{equation}\label{eq:lean_family}
    \varepsilon_{c,s}(j):=\frac{c}{(j+1)^s},
\end{equation}
and its \emph{standing range} is $s\geq0$, $0<c<2^s$, on which $\varepsilon_{c,s}$ takes its
values in $(0,1)$ with supremum $\varepsilon_{\max}=c\,2^{-s}<1$.

Write $(X_n)_{n\geq0}$ for the backward chain of the loop closure, $\sigma$ for its hitting
time of $s_0$, and $\bar\sigma$ for the expected backward-trajectory length of
\citet{morozov2025revisiting}, the average of $\mathbb E(\sigma\mid X_0=j)$ over the target
row. When the loop closure carries an invariant probability, $\lambda$ denotes one,
$\lambda_j:=\lambda(\{j\})$ the mass of a ladder state, $L(m):=\lambda([m,+\infty))$ the tail,
and $f_m:=\mathbf 1_{[m,+\infty)}-L(m)$, for $m>d$, the \emph{centred tail indicator} at the
cut $m$. The operator $P_\star$ is the \emph{function action} of the backward kernel,
$(P_\star f)(x)=\mathbb E(f(X_1)\mid X_0=x)$, defined at every function on the states and a
contraction of $L^p(\lambda)$ for every $p\in[1,+\infty]$; $P$ is the density action of that
kernel on $L^2(\lambda)$; $\Pi$ is the $\lambda$-mean projection
$f\mapsto\bigl(\int f\,d\lambda\bigr)\mathbf 1$ and $\ker\Pi$ the functions of zero
$\lambda$-mean in $L^2(\lambda)$; $\widehat\beta_n:=\|P^n-\Pi\|_{L^2(\lambda)}$; and
\begin{equation}\label{eq:lean_Bhat}
    S:=\left(\mathrm{Id}-P_\star+\Pi\right)^{-1}-\Pi,\qquad
    \widehat B:=\|S\|_{L^2(\lambda)},
\end{equation}
with the convention $\widehat B:=+\infty$ when no bounded operator $R$ on $L^2(\lambda)$
satisfies $R(\mathrm{Id}-P_\star)=\mathrm{Id}-\Pi$. The \emph{Cram\'er equation} of the family
is
\begin{equation}\label{eq:lean_cramer}
    c\left(2^p-1\right)\;=\;p ,
\end{equation}
and $p_*=p_*(c)$ denotes its non-zero real root. For an even integer $K\geq d$, the
\emph{truncation at $K$} is the chain obtained from the loop closure by deleting every doubling
edge $j\rightarrow2j$ with $2j>K$, the decrement out of such a state carrying probability $1$;
$\lambda^K$ denotes an invariant probability of it and $\widehat B_K$ its constant
\eqref{eq:lean_Bhat}. Definition~\ref{def:doubling_setting} carries the same objects
at length.
\end{definition}

\begin{theorem}[the doubling graph, in the form the certificate carries]\label{theo:lean_main}
In the setting of Definition~\ref{def:lean_setting}, with $(c,s)$ in the standing range:
\begin{enumerate}
    \item \emph{(irreducibility and the phase diagram)} the loop closure is irreducible and
    every invariant probability of it is positive at every state; and, for $c>0$, it carries no
    invariant probability at $s<1$, at least one at $s=1$ and $0<c<1$, none at $s=1$ and
    $1\leq c$, and at least one at $s>1$;
    \item \emph{(the invariant measure at $s=1$)} at $s=1$ and $0<c<1$, with $\lambda$ an
    invariant probability and $p_*>1$ the root of \eqref{eq:lean_cramer}, there is $C>0$ with
    $\lambda_m m^{p_*}\rightarrow C$ and $(p_*-1)m^{p_*-1}L(m)\rightarrow C$ --- that is
    $L(m)=\frac C{p_*-1}m^{1-p_*}\bigl(1+o(1)\bigr)$ --- and there are $\vartheta\in(0,1)$ and
    $c_6>0$ with $\left|\lambda_m m^{p_*}-C\right|\leq c_6m^{-\vartheta}$ for every $m\geq1$;
    the constant is determined by the boundary data in two senses: two invariant probabilities
    of the loop closure that agree at the ladder states $1,\dots,d$ agree along the whole
    ladder; and there are $\nu_1,\dots,\nu_d\geq0$ with $\sum_j\nu_j>0$ and $\nu_j=0$
    whenever $2j\leq d$ such that every invariant probability $\lambda'$ of the loop closure
    has $\lambda'_m m^{p_*}\rightarrow\sum_{j\leq d}\nu_j\lambda'_j$;
    \item \emph{(the diffusion operator)} at every $(c,s)$ of the standing range and every
    invariant probability $\lambda$: for every $p\in[1,+\infty)$ the centred tail indicators
    drive the $L^p(\lambda)$ Rayleigh quotient to zero --- for every $t>0$ some $m>d$ has
    $\|(\mathrm{Id}-P_\star)f_m\|^p_{L^p(\lambda)}\leq t\|f_m\|^p_{L^p(\lambda)}$ --- and no
    constant $B$ satisfies $\|f\|^p_{L^p(\lambda)}\leq B\|(\mathrm{Id}-P_\star)f\|^p_{L^p(\lambda)}$
    for every bounded $f$ of zero $\lambda$-mean; at $p=+\infty$, no constant $B$ satisfies
    $M\leq B\,D$ whenever a bounded $f$ of zero mean attains $|f|\geq M$ somewhere and has
    $|(\mathrm{Id}-P_\star)f|\leq D$ everywhere; and at $p=2$, no bounded operator $R$ on
    $L^2(\lambda)$ satisfies $R(\mathrm{Id}-P_\star)=\mathrm{Id}-\Pi$, so $\widehat B=+\infty$,
    the partial sums of $\sum_{n\geq0}(P_\star^n-\Pi)$ do not converge in operator norm, and
    $\sum_{n\geq0}\widehat\beta_n=+\infty$. At $s=1$ and $0<c<1$ the expected
    backward-trajectory length is finite, $\bar\sigma=\bar\jmath/(1-c)$, and there are
    $0<c_1\leq c_2$ and an integer $m_2>d$ with
    \begin{equation}\label{eq:lean_rate}
        c_1\,m\,\left\|\left(\mathrm{Id}-P_\star\right)f_m\right\|^2_{L^2(\lambda)}
        \;\leq\;4\,c_2\left(p_*-1\right)\left\|f_m\right\|^2_{L^2(\lambda)}
    \end{equation}
    for every $m\geq m_2$ --- a $\sqrt m$ rate along the trial functions, and not a bound on
    $\widehat B$, which is $+\infty$ there;
    \item \emph{(an unsolvable flow-matching problem)} at every $(c,s)$ of the standing range
    and every invariant probability $\lambda$, the subspace
    $\left(\mathrm{Id}-P_\star\right)L^2(\lambda)$ is a dense proper subspace of $\ker\Pi$, and
    there are probability densities $\finit,\fterm\in L^2(\lambda)$ such that no
    $f\in L^2(\lambda)$ satisfies $\left(\mathrm{Id}-P_\star\right)f=\finit-\fterm$;
    \item \emph{(the truncation)} for every even $K\geq d$ the truncation at $K$ is irreducible
    on its states and every invariant probability of it is positive at every one of them; it
    carries exactly one such probability $\lambda^K$, and admits a bounded operator $R$ on
    $L^2(\lambda^K)$ with
    $R(\mathrm{Id}-P_\star)=(\mathrm{Id}-P_\star)R=\mathrm{Id}-\Pi$ and $\Pi R=R\Pi=0$, so
    that $\widehat B_K<+\infty$; and at $s=1$ and $0<c<1$ there are $c_8>0$ and an integer $K_0$
    such that $c_8K\leq\|R\|^2$ for every even $K\geq K_0$ and every bounded $R$ on
    $L^2(\lambda^K)$ with $R(\mathrm{Id}-P_\star)=\mathrm{Id}-\Pi$ --- that is
    $\widehat B_K\geq\sqrt{c_8}\sqrt K$.
\end{enumerate}
\end{theorem}

The two conditions come apart on this graph. That of \citet{morozov2025revisiting} is repaired
by every member with $s=1$ and $c<1$, which is what $\bar\sigma=\bar\jmath/(1-c)$ says; the
summable-mixing hypothesis of Theorem~\ref{theo:universality_L2} is repaired by no positive
recurrent member, which is item \emph{(3)} read at $p=2$. What item \emph{(4)} refutes is
\emph{exact} $L^2$ flow matching --- a pair of $L^2$ densities the loop closure cannot balance
at all --- and not the weak universality of Theorem~\ref{theo:universality_L2_full}, whose
conclusion is a defect infimum of zero: the range of $\mathrm{Id}-P_\star$ in $\ker\Pi$ being
dense is exactly why the weak form is untouched.

\subsection{The model}\label{sec:lean_model}

A machine-checked proof certifies that each proof proves its statement; what it cannot certify
is that the statement is the paper's. Four modelling decisions carry that burden here, and each
is legible off the source below.

Three names collide with the paper's, and a reader checking a quotation needs the dictionary
before the first of them. In the Lean, $\texttt{S}$ is a $\texttt{Setting}$ --- the data of the
graph and the policy --- where the paper's $S$ is the diffusion operator \eqref{eq:lean_Bhat},
which appears there under the name $\texttt{Sop}$; $\texttt{L}$ is an invariant probability and
$\texttt{L.lam}$ is the paper's $\lambda$, where the paper's $L(m)$ is the tail, which is
$\texttt{L.tailMass}$; and $\texttt{p}$ is the exponent the paper writes $p_*$.

The states are a ladder indexed by $\mathbb N$ and a sink, and \emph{the source $s_0$ is the
ladder index $0$}, which is Definition~\ref{def:lean_setting}'s own convention --- so the
decrement is uniformly $\texttt{lad}\ (j+1)\rightarrow\texttt{lad}\ j$ with no case at the foot
of the ladder. A single parameter $\texttt{cap}$ of type $\texttt{Option}\ \mathbb N$ carries
both chains: $\texttt{none}$ is the loop closure, where every doubling edge is present, and
$\texttt{some}\ K$ the truncation at $K$, where $\texttt{HasDouble}$ deletes the edge out of $j$
as soon as $2j>K$. The predicate $\texttt{OnChain}$ names the states the truncation lives on.
\begin{lstlisting}[style=lean]
inductive St where
  | lad : ℕ → St
  | sink : St
  deriving DecidableEq, Repr

def HasDouble : Option ℕ → ℕ → Prop
  | none, _ => True
  | some K, j => 2 * j ≤ K

def OnChain : Option ℕ → St → Prop
  | none, _ => True
  | some K, .lad j => j ≤ K
  | some _, .sink => True
\end{lstlisting}

The standing hypotheses of Definition~\ref{def:lean_setting}, minus the family, are one
structure. The target row is finitely supported by hypothesis, so every expression involving it
is a finite sum; $\varepsilon_{\max}$ is a field rather than a supremum, because only
$\varepsilon(j)\leq\varepsilon_{\max}<1$ is ever used; and $\bar\jmath$ is not a field at all
but is derived from the row, $\texttt{Setting.jbar}$. The family \eqref{eq:lean_family} and the
Cram\'er equation \eqref{eq:lean_cramer} are two functions. The structure is quoted with its
field docstrings elided, which is the one place a quotation below is not the whole block.
\begin{lstlisting}[style=lean]
structure Setting where
  d : ℕ
  d_pos : 1 ≤ d
  eps : ℕ → ℝ
  eps_pos : ∀ ⦃j⦄, 1 ≤ j → 0 < eps j
  epsMax : ℝ
  eps_le : ∀ ⦃j⦄, 1 ≤ j → eps j ≤ epsMax
  epsMax_lt_one : epsMax < 1
  row : ℕ → ℝ
  row_nonneg : ∀ j, 0 ≤ row j
  row_supp : ∀ ⦃j⦄, row j ≠ 0 → 1 ≤ j ∧ j ≤ d
  row_sum : ∑ j ∈ Finset.Icc 1 d, row j = 1

noncomputable def epsCS (c s : ℝ) (j : ℕ) : ℝ := c / ((j : ℝ) + 1) ^ s

noncomputable def psi (c t : ℝ) : ℝ := c * ((2 : ℝ) ^ t - 1) - t
\end{lstlisting}

The kernel enters chiefly through $P_\star$, and $P_\star$ is given \emph{in closed form as a
total function} on all of $\texttt{St}\rightarrow\mathbb R$ --- three cases for the three kinds
of state, and no integrability side condition anywhere. That is the second decision, and it is
what lets every statement at the finite exponents, and at $p=+\infty$, be an inequality between
real numbers rather than a statement about a measure. It does not carry everything: the density
action enters as well, in $\widehat\beta_n$, and the support of the kernel is encoded separately
by $\texttt{Edge}$, which is what the irreducibility statements read.
\begin{lstlisting}[style=lean]
noncomputable def pstar (S : Setting) (cap : Option ℕ) (f : St → ℝ) : St → ℝ
  | .lad 0 => f .sink
  | .lad (j + 1) =>
      if HasDouble cap (j + 1) then
        S.eps (j + 1) * f (.lad (2 * (j + 1))) + (1 - S.eps (j + 1)) * f (.lad j)
      else
        f (.lad j)
  | .sink => ∑ k ∈ Finset.Icc 1 S.d, S.row k * f (.lad k)
\end{lstlisting}

The third decision is that ``an invariant probability'' is a structure and not a theorem.
Invariance is carried in the integral form $\int P_\star f\,d\lambda=\int f\,d\lambda$ against
bounded test functions, which is the form every use consumes; positivity is carried on chain,
because one state type serves two chains and the truncation genuinely does not live on all of
it. Every statement of Theorem~\ref{theo:lean_main} that names $\lambda$ therefore quantifies
over such a structure rather than asserting that one exists --- existence is
$\texttt{main\_phase}$, and it is the hard half of item \emph{(1)}. Irreducibility is the
reflexive-transitive closure of the positive transitions of \eqref{eq:lean_policy}, and
$\texttt{PreStat}$ is $\texttt{Stat}$ without positivity, which is what makes ``irreducible
implies positive everywhere'' a theorem rather than an assumption.
\begin{lstlisting}[style=lean]
inductive Edge (S : Setting) (cap : Option ℕ) : St → St → Prop
  | wrap : Edge S cap (.lad 0) .sink
  | dec (j : ℕ) : Edge S cap (.lad (j + 1)) (.lad j)
  | dbl (j : ℕ) (h : HasDouble cap (j + 1)) : Edge S cap (.lad (j + 1)) (.lad (2 * (j + 1)))
  | row (k : ℕ) (hk : k ∈ Finset.Icc 1 S.d) (hpos : 0 < S.row k) : Edge S cap .sink (.lad k)

abbrev Reach (S : Setting) (cap : Option ℕ) : St → St → Prop :=
  Relation.ReflTransGen (Edge S cap)

structure Stat (S : Setting) (cap : Option ℕ) where
  lam : St → ℝ
  nonneg : ∀ x, 0 ≤ lam x
  pos : ∀ ⦃x⦄, OnChain cap x → 0 < lam x
  vanish : ∀ ⦃x⦄, ¬ OnChain cap x → lam x = 0
  summable : Summable lam
  total : ∑' x, lam x = 1
  inv : ∀ f : St → ℝ, (∃ C, ∀ x, |f x| ≤ C) →
          ∑' x, lam x * pstar S cap f x = ∑' x, lam x * f x

structure PreStat (S : Setting) (cap : Option ℕ) where
  lam : St → ℝ
  nonneg : ∀ x, 0 ≤ lam x
  vanish : ∀ ⦃x⦄, ¬ OnChain cap x → lam x = 0
  summable : Summable lam
  total : ∑' x, lam x = 1
  inv : ∀ f : St → ℝ, (∃ C, ∀ x, |f x| ≤ C) →
          ∑' x, lam x * pstar S cap f x = ∑' x, lam x * f x
\end{lstlisting}

The fourth decision is where the norms live. At finite $p$ the statements are in a real
\emph{mass} layer, $\texttt{L.mass}\ p\ f=\|f\|^p_{L^p(\lambda)}$, a plain sum of non-negative
reals; the tail $L(m)$ and the centred tail indicator $f_m$ are functions of the same kind, and
so are the growth condition on $\varepsilon$ along the powers of two, the recursion for
$\mathbb E(\sigma\wedge n\mid X_0=\cdot)$ and the truncated averages whose supremum is
$\bar\sigma$.
\begin{lstlisting}[style=lean]
noncomputable def Stat.mass (L : Stat S cap) (p : ℝ) (f : St → ℝ) : ℝ :=
  ∑' x, |f x| ^ p * L.lam x

noncomputable def Stat.tailMass (L : Stat S cap) (m : ℕ) : ℝ := ∑' x, L.lam x * tailInd cap m x

noncomputable def Stat.centredTail (L : Stat S cap) (m : ℕ) : St → ℝ :=
  fun x => tailInd cap m x - L.tailMass m

def GrowthCond (S : Setting) : Prop :=
  Tendsto (fun D : ℕ => (∑ i ∈ Finset.range D, Real.log (1 / S.eps (2 ^ i))) / 2 ^ D)
    atTop (𝓝 0)

noncomputable def hitExp (S : Setting) (cap : Option ℕ) : ℕ → St → ℝ
  | 0 => fun _ => 0
  | n + 1 => fun x =>
      match x with
      | .lad 0 => 0
      | .sink => 0
      | .lad (j + 1) => 1 + pstar S cap (hitExp S cap n) (.lad (j + 1))

noncomputable def sbar (S : Setting) (cap : Option ℕ) (n : ℕ) : ℝ :=
  ∑ k ∈ Finset.Icc 1 S.d, S.row k * hitExp S cap n (.lad k)
\end{lstlisting}

Only the $L^2$ operator statements --- items \emph{(3)} at $p=2$, \emph{(4)} and \emph{(5)} ---
need a genuine Hilbert space, and there they run on Mathlib's $L^p$ spaces over the measure
$\lambda$ read as a counting measure with density, with $P_\star$ and $\Pi$ continuous linear
maps on it, $\widehat\beta_n$ the norm of $P^n-\Pi$ for the density action, and $\ker\Pi$ the
orthogonal complement of the constants. That $P_\star$ is the $L^2(\lambda)$-adjoint of $P$ is a
theorem of the library, $\texttt{Stat.adjoint\_pstarL2}$, and not a convention; it is not among
the fifteen quoted below, and item \emph{(3)}'s clause on $\sum_n\widehat\beta_n$ is where the
theorem is consumed.
\begin{lstlisting}[style=lean]
noncomputable def mu : Measure St :=
  Measure.count.withDensity (fun x => ENNReal.ofReal (L.lam x))

noncomputable def pstarL2 (hrow : RowOnChain S cap) : Lp ℝ 2 L.mu →L[ℝ] Lp ℝ 2 L.mu :=
  LinearMap.mkContinuous
    { toFun := L.pstarLp
      map_add' := L.pstarLp_add hrow
      map_smul' := fun c F => L.pstarLp_smul hrow c F } 1
    fun F => by simpa using L.norm_pstarLp_le F

noncomputable def piL2 : Lp ℝ 2 L.mu →L[ℝ] Lp ℝ 2 L.mu :=
  (innerSL ℝ L.oneLp).smulRight L.oneLp

noncomputable def betaHat (hrow : RowOnChain S cap) (n : ℕ) : ℝ :=
  ‖(L.densL2 hrow) ^ n - L.piL2‖

noncomputable def partialSum (P Pi : E →L[ℝ] E) (N : ℕ) : E →L[ℝ] E := ∑ n ∈ Finset.range N, (P ^ n - Pi)

noncomputable def defect : Lp ℝ 2 L.mu →L[ℝ] Lp ℝ 2 L.mu := 1 - L.pstarL2 (rowOnChain_none S)

noncomputable def kerPi : Submodule ℝ (Lp ℝ 2 L.mu) := (ℝ ∙ L.oneLp)ᗮ

noncomputable def defectRange : Submodule ℝ (Lp ℝ 2 L.mu) :=
  LinearMap.range (L.defect : Lp ℝ 2 L.mu →ₗ[ℝ] Lp ℝ 2 L.mu)
\end{lstlisting}

\subsection{The certificate, and its axioms}\label{sec:lean_certificate}

Fifteen theorems carry Theorem~\ref{theo:lean_main}, all of them in the module
$\texttt{GFNBounds/Doubling/Main.lean}$ of the supplementary library, and each is quoted below
in full, statement only. Item \emph{(1)} is two: irreducibility with positivity, and the phase
diagram, each of whose four rows is the existence or the non-existence of an invariant
probability.
\begin{lstlisting}[style=lean]
theorem main_irreducible (P : PreStat S none) :
    (∀ x y : St, Reach S none x y) ∧ (∀ y : St, 0 < P.lam y) :=

theorem main_phase {c s : ℝ} (hc : 0 < c) (heps : ∀ j, S.eps j = epsCS c s j) :
    (s < 1 → IsEmpty (Stat S none)) ∧
      (s = 1 → c < 1 → Nonempty (Stat S none)) ∧
      (s = 1 → 1 ≤ c → IsEmpty (Stat S none)) ∧
      (1 < s → Nonempty (Stat S none)) := by
\end{lstlisting}

Item \emph{(2)} is three: the limit $C$, the tail asymptotic and the polynomial rate in one
statement; then the two senses in which the boundary data determine $C$ --- that they determine
the whole ladder, and that $C$ is a non-negative linear combination of them, with coefficients
the statement asserts to exist rather than exhibits.
\begin{lstlisting}[style=lean]
theorem main_invariant_measure {c : ℝ} (hc0 : 0 < c) (hc1 : c < 1)
    (heps : ∀ j, S.eps j = epsCS c 1 j) (L : Stat S none) :
    ∃ p : ℝ, psi c p = 0 ∧ 1 < p ∧ ∃ C : ℝ, 0 < C ∧
      Tendsto (fun m : ℕ => L.lam (.lad m) * (m : ℝ) ^ p) atTop (𝓝 C) ∧
      Tendsto (fun m : ℕ => (p - 1) * (m : ℝ) ^ (p - 1) * L.tailMass m) atTop (𝓝 C) ∧
      ∃ ϑ c₆ : ℝ, 0 < ϑ ∧ ϑ < 1 ∧ 0 < c₆ ∧
        ∀ m : ℕ, 1 ≤ m → |L.lam (.lad m) * (m : ℝ) ^ p - C| ≤ c₆ * (m : ℝ) ^ (-ϑ) := by
\end{lstlisting}
\begin{lstlisting}[style=lean]
theorem main_constant_determined {c : ℝ} (hc0 : 0 < c) (hc1 : c < 1)
    (heps : ∀ j, S.eps j = epsCS c 1 j) (L L' : Stat S none)
    (hagree : ∀ j : ℕ, 1 ≤ j → j ≤ S.d → L.lam (.lad j) = L'.lam (.lad j)) :
    ∀ m : ℕ, 1 ≤ m → L.lam (.lad m) = L'.lam (.lad m) := by

theorem main_constant_functional {c : ℝ} (hc0 : 0 < c) (hc1 : c < 1)
    (heps : ∀ j, S.eps j = epsCS c 1 j) :
    ∃ p : ℝ, psi c p = 0 ∧ 1 < p ∧ ∃ ν : ℕ → ℝ, (∀ j, 0 ≤ ν j) ∧
      0 < ∑ j ∈ Finset.Icc 1 S.d, ν j ∧ (∀ j, 1 ≤ j → 2 * j ≤ S.d → ν j = 0) ∧
      ∀ L : Stat S none,
        Tendsto (fun m : ℕ => L.lam (.lad m) * (m : ℝ) ^ p) atTop
          (𝓝 (∑ j ∈ Finset.Icc 1 S.d, ν j * L.lam (.lad j))) := by
\end{lstlisting}

Item \emph{(3)} is six. The first three are the unboundedness, at the finite exponents, at
$p=+\infty$, and at $p=2$ where it becomes $\widehat B=+\infty$, a divergent operator series
and a divergent mixing sum.
\begin{lstlisting}[style=lean]
theorem main_unbounded {c s : ℝ} (hc : 0 < c) (hs : 0 ≤ s)
    (heps : ∀ j, S.eps j = epsCS c s j) (L : Stat S none) {p : ℝ} (hp : 1 ≤ p) :
    (∀ t : ℝ, 0 < t → ∃ m, S.d < m ∧
        L.mass p (fun x => L.centredTail m x - pstar S none (L.centredTail m) x)
          ≤ t * L.mass p (L.centredTail m)) ∧
      ¬ ∃ B : ℝ, ∀ f : St → ℝ, (∃ C, ∀ x, |f x| ≤ C) → (∑' x, L.lam x * f x = 0) →
          L.mass p f ≤ B * L.mass p (fun x => f x - pstar S none f x) := by

theorem main_unbounded_infty {c s : ℝ} (hc : 0 < c)
    (heps : ∀ j, S.eps j = epsCS c s j) (L : Stat S none) :
    ¬ ∃ B : ℝ, ∀ (f : St → ℝ) (M Dd : ℝ),
        (∃ C, ∀ x, |f x| ≤ C) → (∑' x, L.lam x * f x = 0) → (∃ x, M ≤ |f x|) →
        (∀ x, |f x - pstar S none f x| ≤ Dd) → M ≤ B * Dd := by
\end{lstlisting}
\begin{lstlisting}[style=lean]
theorem main_unbounded_two {c s : ℝ} (hc : 0 < c) (hs : 0 ≤ s)
    (heps : ∀ j, S.eps j = epsCS c s j) (L : Stat S none) :
    (¬ ∃ Sop : Lp ℝ 2 L.mu →L[ℝ] Lp ℝ 2 L.mu,
        Sop * (1 - L.pstarL2 (rowOnChain_none S)) = 1 - L.piL2) ∧
      (¬ ∃ U : Lp ℝ 2 L.mu →L[ℝ] Lp ℝ 2 L.mu,
        Tendsto (partialSum (L.pstarL2 (rowOnChain_none S)) L.piL2) atTop (𝓝 U)) ∧
      ¬ Summable (L.betaHat (rowOnChain_none S)) := by
\end{lstlisting}

The other three are the backward length and the rate. $\texttt{main\_sigmaBar\_eq}$ gives the
closed form of item \emph{(3)}; $\texttt{main\_sigmaBar\_le}$ read at
$\texttt{cap}=\texttt{none}$ is what makes that length finite in the first place, and its other
half --- the same bound uniformly over every truncation --- is more than the theorem states; and
$\texttt{main\_rate}$ is the $\sqrt m$ rate \eqref{eq:lean_rate} along the trial functions.
\begin{lstlisting}[style=lean]
theorem main_sigmaBar_eq {c : ℝ} (hc0 : 0 < c) (hc1 : c < 1)
    (heps : ∀ j, S.eps j = epsCS c 1 j) :
    ⨆ n, sbar S none n = S.jbar / (1 - c) :=

theorem main_sigmaBar_le {c : ℝ} (hc0 : 0 < c) (hc1 : c < 1) {cap : Option ℕ}
    (heps : ∀ j, S.eps j = epsCS c 1 j) (n : ℕ) :
    sbar S cap n ≤ S.jbar / (1 - c) :=

theorem main_rate {c : ℝ} (hc0 : 0 < c) (hc1 : c < 1)
    (heps : ∀ j, S.eps j = epsCS c 1 j) (L : Stat S none) :
    ∃ p : ℝ, psi c p = 0 ∧ 1 < p ∧ ∃ c₁ c₂ : ℝ, 0 < c₁ ∧ c₁ ≤ c₂ ∧
      ∃ m₂ : ℕ, S.d < m₂ ∧ ∀ m : ℕ, m₂ ≤ m →
        c₁ * (m : ℝ) * L.mass 2 (fun x => L.centredTail m x - pstar S none (L.centredTail m) x)
          ≤ 4 * c₂ * (p - 1) * L.mass 2 (L.centredTail m) := by
\end{lstlisting}

Item \emph{(4)} is one statement, carrying both halves: the range is dense and proper, and a
pair of probability densities lies outside it.
\begin{lstlisting}[style=lean]
theorem main_unsolvable {c s : ℝ} (hc : 0 < c) (hs : 0 ≤ s)
    (heps : ∀ j, S.eps j = epsCS c s j) (L : Stat S none) :
    (L.defectRange ≤ L.kerPi ∧ L.defectRange.topologicalClosure = L.kerPi ∧
        L.defectRange ≠ L.kerPi) ∧
      ∃ fi fe : St → ℝ,
        (∀ x, 0 ≤ fi x) ∧ (∀ x, 0 ≤ fe x) ∧
        MemLp fi 2 L.mu ∧ MemLp fe 2 L.mu ∧
        ∑' x, L.lam x * fi x = 1 ∧ ∑' x, L.lam x * fe x = 1 ∧
        ¬ ∃ f : St → ℝ, MemLp f 2 L.mu ∧
            ∀ᵐ x ∂L.mu, f x - pstar S none f x = fi x - fe x :=
\end{lstlisting}

Item \emph{(5)} is three: the truncation is irreducible, it carries exactly one invariant
probability and a bounded diffusion operator, and that operator's norm grows at least like
$\sqrt K$.
\begin{lstlisting}[style=lean]
theorem main_truncation_irreducible {K : ℕ} (hK : Even K) (hdK : S.d ≤ K)
    (P : PreStat S (some K)) :
    (∀ x y : St, OnChain (some K) y → Reach S (some K) x y) ∧
      (∀ ⦃y : St⦄, OnChain (some K) y → 0 < P.lam y) :=

theorem main_truncation_bhat {K : ℕ} (hK : Even K) (hdK : S.d ≤ K) :
    Nonempty (Stat S (some K)) ∧
      (∀ L L' : Stat S (some K), L.lam = L'.lam) ∧
      ∀ L : Stat S (some K), ∃ Sop : Lp ℝ 2 L.mu →L[ℝ] Lp ℝ 2 L.mu,
        (1 - L.pstarL2 (rowOnChain_some hdK)) * Sop = 1 - L.piL2 ∧
        Sop * (1 - L.pstarL2 (rowOnChain_some hdK)) = 1 - L.piL2 ∧
        L.piL2 * Sop = 0 ∧ Sop * L.piL2 = 0 :=
\end{lstlisting}
\begin{lstlisting}[style=lean]
theorem main_truncation_sqrtK {c : ℝ} (hc0 : 0 < c) (hc1 : c < 1)
    (heps : ∀ j, S.eps j = epsCS c 1 j) :
    ∃ c₈ : ℝ, 0 < c₈ ∧ ∃ K₀ : ℕ, ∀ K : ℕ, K₀ ≤ K →
      ∀ (L : Stat S (some K)) (hrow : RowOnChain S (some K))
        (Sop : Lp ℝ 2 L.mu →L[ℝ] Lp ℝ 2 L.mu),
        Sop * (1 - L.pstarL2 hrow) = 1 - L.piL2 → c₈ * (K : ℝ) ≤ ‖Sop‖ ^ 2 := by
\end{lstlisting}

\begin{table*}[t]
\centering
\caption{Each item of Theorem~\ref{theo:lean_main}, the declarations of
$\texttt{GFNBounds/Doubling/Main.lean}$ that carry it, and the modules of
$\texttt{GFNBounds/Doubling/}$ in which the mathematics is proved.}
\label{tab:lean_map}
\small
\begin{tabular}{@{}c p{0.40\linewidth} p{0.36\linewidth}@{}}
\hline
item & declarations & modules \\
\hline
\emph{(1)} & \texttt{main\_irreducible}, \texttt{main\_phase} &
\texttt{Irreducible}, \texttt{PhaseEmpty}, \texttt{PhaseExists}, \texttt{StatExists} \\
\emph{(2)} & \texttt{main\_invariant\_measure}, \texttt{main\_constant\_determined},
\texttt{main\_constant\_functional} &
\texttt{DecayNotation}, \texttt{Product}, \texttt{Sharp}, \texttt{SharpFull},
\texttt{ConstantFunctional} \\
\emph{(3)} & \texttt{main\_unbounded}, \texttt{main\_unbounded\_infty},
\texttt{main\_unbounded\_two}, \texttt{main\_sigmaBar\_eq}, \texttt{main\_sigmaBar\_le},
\texttt{main\_rate} &
\texttt{Unbounded}, \texttt{Ramp}, \texttt{UnboundedL2}, \texttt{Family}, \texttt{Product},
\texttt{Supersolution}, \texttt{Length} \\
\emph{(4)} & \texttt{main\_unsolvable} & \texttt{Unsolvable} \\
\emph{(5)} & \texttt{main\_truncation\_irreducible}, \texttt{main\_truncation\_bhat},
\texttt{main\_truncation\_sqrtK} &
\texttt{Irreducible}, \texttt{TruncationStat}, \texttt{OperatorL2}, \texttt{Truncation},
\texttt{Product}, \texttt{RayleighBridge} \\
\hline
\end{tabular}
\end{table*}

The library is Lean 4 with Mathlib pinned at the tag \texttt{v4.31.0}, revision
\texttt{fabf563a7c95}. It compiles with $\texttt{warningAsError}:=\texttt{true}$, which turns
Lean's own ``declaration uses \texttt{sorry}'' warning into a compile error, so
$\texttt{sorry}$-freeness is enforced by the compiler rather than by discipline; a permissive
scaffold library exists for work in progress and is never imported by the strict one, so
graduating a result is a file move. An axiom sweep over the strict library reports
$4{,}006$ declarations in $99$ modules, each depending only on $\texttt{propext}$,
$\texttt{Classical.choice}$ and $\texttt{Quot.sound}$ --- Lean's three standard axioms, and
nothing else --- and a second audit parses $\texttt{\#print axioms}$ for the $245$ declarations
of an explicit audit file and fails on $\texttt{sorryAx}$ or on any axiom outside that set. A
kernel replay then re-checks every declaration of the strict library from the compiled
\texttt{.olean} files, with Mathlib imported and trusted, so that the certificate rests on the
kernel and not on the elaborator that produced the environment. A traceability check maps each
statement of the three appendices the library covers --- the $41$ of
Appendix~\ref{app:doubling}, the $37$ of the appendix of proofs and the $6$ of
Appendix~\ref{sec:silva_comparison} --- to its Lean declarations and to a \texttt{sha256} of
that statement's own \LaTeX{} block, so that editing a statement marks it stale until the map
is re-affirmed; it reports $84$ statements mapped and every digest current. Coverage of the
long-form appendix this one replaces is $29$ statements closed, $9$ partial and $3$ open.
\emph{Partial} means that some items of a multi-part statement are proved and others are not:
Lemma~\ref{lem:doubling_operator} is carried in its doubling-graph instance alone,
Lemma~\ref{lem:doubling_expansion} and Remark~\ref{rem:doubling_parity} are blocked on an
Euler--Maclaurin estimate Mathlib does not have, and Lemma~\ref{lem:doubling_weight} is carried
in a consequence form; none of the nine touches a clause of Theorem~\ref{theo:lean_main}. The
open three are remarks that appendix does not prove either, and
\texttt{docs/COVERAGE.md} gives the per-statement account. The whole gate is one command,
\texttt{make check}, last run green on 2026-09-13.

What this buys is one half of the reading and not both. The compiler certifies that each proof
proves its statement; what is left to the reader is whether the statement is the paper's, and
that is a question about \S\ref{sec:lean_model} and about \S\ref{sec:lean_scope} --- which is
why the model is quoted in full above, and why the scope below is part of the claim rather than
an apology attached to it.

\subsection{The ideas}\label{sec:lean_ideas}

This subsection says what the proofs do. It proves nothing, and nothing below is used anywhere
else in the paper.

\paragraph{One identity carries the invariant measure.}
Stationarity across the cut $A=[m,+\infty)$, at any $m>d$, has exactly two terms. The only way
out of $A$ is the decrement from $m$, which carries the probability $1-\varepsilon(m)$; the
only way in is a doubling from the window $W(m):=\{\lceil m/2\rceil,\dots,m-1\}$, since $s_0$
goes only to $s_f$ and $s_f$ lands in $\{1,\dots,d\}$ with $d<m$. Equating the two fluxes,
\begin{equation}\label{eq:lean_cut}
    \lambda_m\bigl(1-\varepsilon(m)\bigr)\;=\;\sum_{j\in W(m)}\lambda_j\,\varepsilon(j) ,
\end{equation}
and everything the certificate proves about the invariant measure is this one identity. The
certificate takes \eqref{eq:lean_cut}, and not the chain, as its hypothesis: the decay results
are stated for an arbitrary positive sequence satisfying the balance, with no invariance, no
normalization and no recurrence. That is why they apply verbatim to $\lambda^K$ below the cap,
where the same two-term count holds for every $d<m\leq K/2$, and it is why the library needs no
Markov-chain theory at all to prove them.

\paragraph{Why the exponent solves $c(2^p-1)=p$.}
Doubling multiplies the index, so the natural variable on the ladder is $\log_2$ of the state
and \eqref{eq:lean_cut} is a Laplace transform of one dyadic step. Feed the power ansatz
$\lambda_j\approx A j^{-p}$ into the balance at $s=1$, where $\varepsilon(j)=c/(j+1)$: the left
side is $Am^{-p}(1+o(1))$, the window sum is a Riemann sum for
$cA\int_{m/2}^mu^{-p-1}\,du=cAm^{-p}(2^p-1)/p$, and dividing out $Am^{-p}$ leaves the
Cram\'er equation \eqref{eq:lean_cramer} and nothing else. Its non-zero root is $>1$ exactly
when $c<1$ and $>0$ exactly when $c<1/\ln2$, so the two thresholds of the phase diagram are the
two crossings of a single root.

\paragraph{From ansatz to theorem: a chain that descends the ladder.}
An ansatz is not a theorem, and what turns it into one is a change of variable. Rescale to
$u_j:=\lambda_j j^{p_*}$; then \eqref{eq:lean_cut} says that $u$ at a level is a weighted average
of its values on the window below, the weights summing not to $1$ but to $R_0(m)=1+O(1/m)$.
Read the normalized weights as the transition law of a Markov chain descending the ladder. One
of its steps falls by at least one and by at most half --- from $y$ it lands in
$\{\lceil y/2\rceil,\dots,y-1\}$ --- so the descent from $m$ reaches a fixed base block
$[\ell,2\ell)$ without overshooting it, in at most $m$ steps. That count is linear and carries
nothing; what makes the descent fast is that the halving holds \emph{in probability}. From $y$
the chain lands at or below $y/\sqrt2$ with probability at least a fixed $\gamma$, so the time
it spends in any one dyadic block has expectation at most $2/\gamma^2$, and there are only
$\log_2(m/\ell)$ blocks between the start and the base block. The profile at an index is the
expectation of the accumulated weight times the value of $u$ where the descent lands, and it is
that bounded expected sojourn per block --- not a bounded number of steps --- which lets the
$1+O(1/m)$ deviations telescope over a whole descent into a single bounded factor
$e^{\pm O(1/\ell)}$, the constant depending on $c$ alone. The profile at every index above the
base block is then squeezed between the extremes of the profile on that one finite block, which
is the two-sided decay $c_1 j^{-p_*}\leq\lambda_j\leq c_2 j^{-p_*}$ --- with the same pair
$0<c_1\leq c_2$ that \eqref{eq:lean_rate} binds --- and no asymptotic hypothesis anywhere.

\paragraph{From two constants to one: a coupling.}
The sandwich leaves the constant anywhere between $c_1$ and $c_2$, and closing that gap is a
mixing statement about the descent. Two descents started in one dyadic block leave it with laws
that overlap: the exit law charges at least a fixed fraction $\omega$ of one reference law on
the block beneath --- a Doeblin minorization --- so the two descents can be coupled, and the
total variation between their landing laws falls geometrically in the number of blocks crossed.
The sandwich becomes a limit: $\lambda_j j^{p_*}\rightarrow C$ for some $C>0$, with a
polynomial rate $\left|\lambda_j j^{p_*}-C\right|\leq c_6 j^{-\vartheta}$ for some $c_6>0$ and
the exponent $\vartheta=\ln\bigl(1/(1-\omega)\bigr)/\ln\bigl(2/(1-\omega)\bigr)$ that the
minorization produces.

Two flags belong in the text beside it. The exponent is explicit and \textbf{not effective}: at
$c=1/2$ the proved $\omega$ is $0.0016$ against a block-exit minorization measured at $0.18$ on
exact exit laws, so the proved $\vartheta$ is $0.0024$ against $0.22$, while the observed rate
is $O(1/m)$. And $C$ is a non-negative linear functional of the boundary data and no more:
\eqref{eq:lean_cut} is homogeneous, so no equation in $c$ alone can determine it.

\paragraph{One trial function makes the operator unbounded.}
Apply $\mathrm{Id}-P_\star$ to the indicator of a tail and almost everything cancels: what
survives is supported on $\{m\}\cup W(m)$ --- the mass leaving at $m$, the mass entering over
the window --- so $\|(\mathrm{Id}-P_\star)f_m\|^2_{L^2(\lambda)}\leq2\lambda_m$, while
$\|f_m\|^2_{L^2(\lambda)}=L(m)\bigl(1-L(m)\bigr)$ because $f_m$ is a centred indicator. The
Rayleigh ratio at the cut $m$ is therefore at most
$\sqrt{2\lambda_m/\bigl(L(m)(1-L(m))\bigr)}$, and the operator is unbounded as soon as
$L(m)/\lambda_m$ is unbounded --- as soon as the tail fails to be comparable to its own last
term.

It cannot be comparable. Were $L(m)/\lambda_m$ bounded by $Q$, the tail would fall geometrically
at the rate $e^{-1/Q}$; but iterating the doubling edges along the powers of two gives
$L(2^D)\geq\lambda_1\prod_{i<D}\varepsilon(2^i)$ at every integer $D\geq1$, and the growth
condition --- that $\sum_{i<D}\log\bigl(1/\varepsilon(2^i)\bigr)$ is $o(2^D)$ as
$D\rightarrow\infty$ --- makes that far too large. The mechanism is
geometric and not spectral: the doubling edges plant mass exponentially far out along the
ladder, and no policy whose doubling probability decays subexponentially can keep the tail
comparable to its last term. The condition indicts the family and not the graph --- a
geometrically decaying $\varepsilon(j)=a\rho^j$ fails it, and no boundedness is claimed for any
policy here. At $p=+\infty$ the sharp indicator gives nothing and clipped ramps take its place:
their height grows with the clipping level while their defect stays bounded by a constant.

\paragraph{From unbounded to unsolvable.}
Restricted to $\ker\Pi$, the operator $\mathrm{Id}-P_\star$ is injective with dense range ---
only the constants are fixed by $P_\star$, only the constants are fixed by $P$, and the two
facts are adjoint to each other --- but its Rayleigh infimum is zero, so it cannot also be
surjective: the open mapping theorem would hand back the bounded inverse that the infimum
forbids. Any $\psi$ of zero mean outside the range then splits into two probability densities
$\finit,\fterm$ whose difference is a multiple of $\psi$, and that pair cannot be balanced in
$L^2(\lambda)$ at all. Two flags again. The pair comes from non-surjectivity and is
\textbf{not explicit}; and because the range is \emph{dense}, weak universality --- a defect
infimum of zero --- is untouched, and remains open on the infinite chain.

\paragraph{The truncation restores a finite constant, and prices it.}
Deleting the doubling edges above an even cap leaves a finite irreducible chain, so $\lambda^K$
exists, is unique and is positive, and $\widehat B_K$ is finite. Finiteness becomes growth by
reading the same per-cut inequality at a cut near $K/4$ and feeding it the two-sided decay,
which \eqref{eq:lean_cut} transfers to $\lambda^K$ unchanged, and $\widehat B_K$ grows at
least like $\sqrt K$.
This is loose, and by exactly one factor of $\sqrt K$ --- what \S\ref{sec:doubling_cost}
measures is $K$.

\paragraph{And the backward length stays finite.}
At $s=1$ the drift of the backward chain is $c-1$ at every ladder state, independent of the
state --- $s=1$ is the only exponent at which it is --- so the linear function $j\mapsto j/(1-c)$
solves the hitting-time recursion exactly, which bounds $\mathbb E(\sigma\mid X_0=j)$ above. The
matching lower bound is the one place where formalizing changed the proof.
Appendix~\ref{app:doubling} reaches it by optional stopping on a martingale, dominated by the
integrable excursion maximum; the certificate builds no martingale and never consumes the
excursion bound. It iterates the ladder operator on the identity $x\mapsto x$ instead --- that
operator sends the identity to $x\mapsto x-(1-c)$, so the martingale identity becomes an
induction on the iterate --- and it replaces dominated convergence by one explicit
supersolution, $W(m)=\sum_{i\leq m}\bigl(\lfloor\log_2i\rfloor+1\bigr)$, which the ladder
operator raises by at most $2$ and which grows like $m\log_2m$, so that the identity is
dominated by an arbitrarily small multiple of it plus a constant. The iterates of the identity therefore vanish, which is
$\mathbb E(\sigma\mid X_0=j)=j/(1-c)$; averaging over the target row gives
$\bar\sigma=\bar\jmath/(1-c)$. This is the divorce the appendix exists for: the condition of
\citet{morozov2025revisiting} is repaired at every $c<1$, and $\widehat B=+\infty$ at every one
of them.

\subsection{What the certificate does not carry}\label{sec:lean_scope}

Six clauses and conventions of Theorem~\ref{theo:doubling_main} are absent from
Theorem~\ref{theo:lean_main}, and they are the price of the rest being certified.

The first is uniqueness of the invariant probability on the loop closure. Positive recurrence
in Theorem~\ref{theo:doubling_main}\emph{(1)} carries it; $\texttt{main\_phase}$ delivers
existence alone, and the library discloses that it does not prove uniqueness there. At a finite
cap uniqueness is proved --- it is half of item \emph{(5)} --- by an argument that does not
transfer to the infinite chain.

The words \emph{transient} and \emph{null recurrent} of the phase diagram, and the clause
$\mathbb E(\sigma\mid X_0=j)=+\infty$ at $1\leq c<2$, are not stated. Mathlib \texttt{v4.31.0}
carries no discrete-time Markov chain theory and the library builds none, so each row of item
\emph{(1)} reads instead as the existence or the non-existence of an invariant probability,
which is what the rest of the appendix uses it for; Proposition~\ref{prop:doubling_phase}
carries the labels. The count of $K+2$ states of the truncation is not stated either: the Lean
state type is infinite and the truncation is carried by a predicate on it. The functional
$\bar\sigma$ is the supremum of the truncated hitting expectations, defined by the paper's own
recursion, rather than an expectation over a chain. And Appendix~\ref{app:doubling}'s three
unproved remarks --- the renewal reading of the balance, the absence of a log-periodic term,
and the second-order expansion --- are not formalized, that appendix not proving them either.

Four dependence clauses of constants go the same way. Theorem~\ref{theo:doubling_main}\emph{(3)}
says that its $c_7$ depends only on $c$, $d$ and an initial block of the invariant measure, and
Theorem~\ref{theo:doubling_sharp} says as much of its $c_6$; item \emph{(2)} above binds the
$\nu_j$ existentially, where they are in fact determined by $c$ and $d$ alone; and item
\emph{(5)} binds $c_8$ and $K_0$ the same way, where
Theorem~\ref{theo:doubling_main}\emph{(5)} names what they depend on. None of the four is
carried by the quoted existentials. They are carried instead, as explicit
formulas in $c$, $d$ and the initial block of the profile, by the library's explicit-constant
lemmas $\texttt{Decay.decay\_two\_sided\_explicit}$ and $\texttt{Decay.sharp\_explicit}$, and
for the $\nu_j$ by $\texttt{Decay.constant\_functional}$, whose coefficients depend on $c$ and
$d$ alone.

Two constants also change shape between the two statements, which is worth naming because a
reader comparing them will trip on it: Theorem~\ref{theo:doubling_main}\emph{(3)} writes its
rate as $\|f_m\|\geq c_7\sqrt m\,\|(\mathrm{Id}-P_\star)f_m\|$ where \eqref{eq:lean_rate}
carries the pair $0<c_1\leq c_2$ from which that $c_7$ is assembled, and it writes
$\widehat B_K\geq c_8\sqrt K$ where item \emph{(5)} states the same bound squared, so that the
$c_8$ of item \emph{(5)} is the square of the $c_8$ of
Theorem~\ref{theo:doubling_main}\emph{(5)}.

One clause of the certificate drops a hypothesis: the $\sqrt K$ lower bound is stated for every
$K\geq K_0$, even or not. The drop is empty. At an odd $K\geq K_0$ the truncation carries no
invariant probability at all, so the statement quantifies over nothing --- the state $K$ is on
chain and would have to carry positive mass, while its only sources are a decrement from $K+1$,
which is off chain, a doubling from $K/2$, which is not an integer, and the target row, which is
supported in $\{1,\dots,d\}$ while $K\geq K_0>d$. So this is a vacuous generalization and not a
strengthening, and item \emph{(5)} states the even case, which is the one the setting defines.

Where Appendix~\ref{app:doubling} says a thing is open, formal or non-explicit, the certificate
says so too, rather than laundering it.

The library is included as anonymized supplementary material together with its audit scripts;
\texttt{make check} builds it, runs the axiom sweep and the axiom audit, and re-derives the
per-statement map that ties each of the $84$ statements of the three appendices the library
covers to the declarations certifying it.




\end{document}
